% !TEX root = Anleitung.tex
% Teil "Die Mathematik hinter den Features" -- wird von Anleitung.tex mit % !TEX root = Anleitung.tex
% Teil "Die Mathematik hinter den Features" -- wird von Anleitung.tex mit % !TEX root = Anleitung.tex
% Teil "Die Mathematik hinter den Features" -- wird von Anleitung.tex mit % !TEX root = Anleitung.tex
% Teil "Die Mathematik hinter den Features" -- wird von Anleitung.tex mit \input{Mathe} eingebunden.

\part{Die Mathematik hinter den Features}\label{sec:Mathe}

Vorab will ich erwähnen, dass Teile dieses Abschnitts nicht meine eigene Arbeit sind. Gerade wenn
es darum geht, das Programm effizient und robust zu schreiben, oder wenn sich das Skript Methoden der
weiterführenden Mathematik bzw.\ Numerik bedient, stammen Ideen von KI, z.\,B.\ die
$LDL^T$-Zerlegung (ich bin selbst erst im 2.~Semester, deswegen halten sich meine mathematischen
Kenntnisse noch in Grenzen). Das gilt übrigens auch für große Teile der Umsetzung. Aus diesem Grund
halte ich die Erklärungen zu diesen Aspekten hier etwas kürzer.

Nichtsdestotrotz ist dieser Teil für alle, die wissen wollen, \emph{warum} das Programm so rechnet, wie es rechnet.
Zum Benutzen braucht ihr ihn nicht. Grundlegende Sachen wie Hebelarme, Steiner und Gleichgewicht seien hier nicht weiter erläutert. Hier
werden die Stellen behandelt, an denen das Programm mehr macht als die Vorlesung oder es anders macht,
weil ein Rechner eben nicht wie ein Mensch rechnet. Ich habe versucht, zu jedem Punkt dazuzuschreiben,
was er für euch in der Klausur bedeutet.

Ein paar Konventionen, die überall gelten:

\begin{itemize}
  \item Im \textbf{Tragwerksskript} zeigt $x$ nach \emph{rechts} und $y$ nach \emph{unten}. Drehwinkel,
        Momente und Lastwinkel zählen trotzdem so, wie ihr sie seht: gegen den Uhrzeigersinn
        positiv.
  \item Jeder Stab hat seine eigene Koordinate $x \in [0, L]$ vom Anfangs- zum Endknoten, die
        Tangente $\vec t = (\Delta x, \Delta y)/L$ und die Normale $\vec n = (-\Delta y, \Delta x)/L$.
        Die Schnittgrößen folgen der TM1-Konvention mit gestrichelter Faser auf der Seite von
        $\vec n$.
  \item Im \textbf{Querschnittsskript} gilt die Konvention aus dem Gross (TM2):
        $I_y = \int z^2\,dA$, $I_z = \int y^2\,dA$ und $I_{yz} = -\int yz\,dA$, also mit Minus.
\end{itemize}

%---------------------------------------------------------------------
\section{Das Weggrößenverfahren im Programm}\label{sec:mathe-wgv}

Alles, was ihr im Tragwerksskript seht (Verläufe, Auflagerkräfte, Verformungen,
Explosionszeichnung), kommt aus einer einzigen Rechnung: dem Weggrößenverfahren, also genau dem
allgemeinen WGV aus \textbf{Statik}. Unbekannt sind die Verschiebungen der Knoten, nicht die
Kräfte.

\textbf{Unbekannte:} Jeder Knoten hat drei Freiheitsgrade, $u_\xi$, $u_\eta$ und $\varphi$. Bei
einem ungedrehten Lager ist $\xi = x$ und $\eta = y$, bei einem gedrehten Lager zeigen sie in die
Lagerrichtungen. Ein Lager \emph{streicht} einfach den gehaltenen Freiheitsgrad.

\textbf{Stabsteifigkeit:} Für jeden Stab steht die bekannte $6 \times 6$-Matrix aus $EA/L$,
$12EI/L^3$, $6EI/L^2$, $4EI/L$ und $2EI/L$. Sie verknüpft die sechs Endverschiebungen
$(u_a, w_a, \varphi_a, u_b, w_b, \varphi_b)$ im lokalen System des Stabes mit den sechs
Endkräften. Zwei Drehungen bringen sie ins globale System: die Stabrichtung und (bei einem
gedrehten Lager) die Lagerrichtung,
\[
  K_e = T_K\,T^T\,k\,T\,T_K^T .
\]
So ist ein gedrehtes Lager kein Sonderfall, sondern nur eine andere Basis am Knoten. Matrix,
Transformation und Drehrichtungen sind dieselben wie in der WGV-Formelsammlung aus Statik, die
FEM-Formelsammlung hat dieselben Elemente, zählt die Querrichtung aber nach oben. Die
Unterschiede stehen in Abschnitt~\ref{sec:wgv-vorzeichen}.

\textbf{Lasten:} Knotenlasten gehen direkt in den Lastvektor (vorher um den Lastwinkel gedreht).
Streckenlasten gehen als Festeinspannkräfte ein. Das sind die Endkräfte, die ein
beidseitig eingespannter Stab unter der Last hätte, z.\,B.\ $\mp qL/2$ und $\pm qL^2/12$ bei
Gleichlast (die Volleinspannmomente aus Tafel~A9 der Formelsammlung, mit denselben Vorzeichen).
Im Gesamtsystem stehen sie mit umgekehrtem Vorzeichen, also als Knotenersatzlasten. Für Trapezlasten stehen die entsprechenden Formeln im Code, Lastfunktionen wie
\code{x\^{}2} integriert das Programm mit den Hermite-Formfunktionen
$N_1 = 1 - 3\xi^2 + 2\xi^3$, $N_2 = L(\xi - 2\xi^2 + \xi^3)$ usw.\ (arbeitsäquivalente
Knotenlasten, Abschnitt~\ref{sec:mathe-festeinspann}). Temperatur wirkt wie eine Last $EA\,\varepsilon_T$ bzw.\ $EI\,\kappa_T$ an den
Stabenden. Globale Linienlasten $p_X(x)$, $p_Y(x)$ werden vorher in Quer- und Längsanteil zerlegt,
bei einer projizierten Last ist $x$ die globale Koordinate ab dem Stabanfang, $x_{global} = f\,x$
mit dem Projektionsfaktor $f = |\Delta y|/L$ bzw.\ $|\Delta x|/L$, und je Stablänge wirkt
$f\,p(f\,x)$.

\textbf{Lösen:} Das Gesamtsystem $K\,u = F$ wird mit einer $LDL^T$-Zerlegung gelöst. Wird dabei
ein Diagonalelement $D_j$ (fast) null, fehlt dem System in dieser Richtung jede Steifigkeit. Es
ist ein Mechanismus, und ihr bekommt \emph{System ist kinematisch!}.

\textbf{Zurückrechnen:} Aus den Knotenverschiebungen folgen die Endkräfte jedes Stabes,
$s = k\,v + p$ (Steifigkeit mal Endverschiebung plus Festeinspannkraft), und aus dem Anfangswert
durch Integrieren der Lasten die Verläufe, z.\,B.\
$M(x) = M_a + Q_a x - m\,x - \int_0^x q(t)\,(x-t)\,dt$. Das ist exakt, nicht genähert. Deshalb
braucht jeder Bereich zwischen Lasten, Lagern und Gelenken nur einen einzigen Stab.

\textbf{Auflagerkräfte:} Die Summe der Stabendkräfte an einem Knoten minus der Knotenlast ist die
Kraft, die Lager und Federn auf den Knoten ausüben. Davon zeigt das Programm
$R_\xi = \vec R\cdot\vec\xi$ und $R_\eta = -\vec R\cdot\vec\eta$, deshalb zeigt $R_\eta$ nach oben,
während die Knotenverschiebung \code{ü\_y} auf derselben Menüseite nach unten positiv ist.

\subsection{Festeinspannkräfte beliebiger Lasten}\label{sec:mathe-festeinspann}

Für das Gesamtsystem braucht das Programm bzw. das WGV von einer Streckenlast nur, was sie an den Stabenden
bewirkt, die Festeinspannkräfte. Für Gleich- und Trapezlasten stehen diese In Tafeln oder in der FOS, für eine
Lastfunktion wie $q(x) = x^2$ oder $q_0 \sin(\pi x/L)$ nicht. Das Programm rechnet sie deshalb für
jede Lastfunktion mit den Hermite-Funktionen aus, und zwar exakt, nicht genähert.

\textbf{Die Hermite-Funktionen.} Im lokalen System läuft $x$ von $0$ bis $L$, $\xi = x/L$, $w$ zeigt
zur gestrichelten Faser und $\varphi = -w'$ gegen den Uhrzeigersinn. Die vier Funktionen sind
\begin{align*}
  N_1 &= 1 - 3\xi^2 + 2\xi^3       && w(0) = 1, \\
  N_2 &= L\,(\xi - 2\xi^2 + \xi^3) && w'(0) = 1, \text{ also } \varphi_a = -1, \\
  N_3 &= 3\xi^2 - 2\xi^3           && w(L) = 1, \\
  N_4 &= L\,(-\xi^2 + \xi^3)       && w'(L) = 1, \text{ also } \varphi_b = -1,
\end{align*}
und die jeweils drei anderen Endwerte sind null. Jede ist eine kubische Parabel, löst also
$EI\,w'''' = 0$ exakt und stellt somit eine homogene Lösung der DGL dar.
Gleichzeitig erfüllt sie die vier von ihr geforderten Randbedingungen. Damit ist sie genau die Biegelinie des
beidseitig eingespannten Stabes, dem man diese eine Einheitsverformung aufzwingt.

Das WGV zählt die vier Endverschiebungen eines Stabes aber als $v = (w_a, \varphi_a, w_b, \varphi_b)$,
mit $\varphi$ gegen den Uhrzeigersinn positiv. Die Einheitsverformung zum Freiheitsgrad $j$ in
dieser Zählung heißt im Folgenden $S_j$ (S für Formfunktion, engl.\ \emph{shape}):
\[
  S_1 = N_1, \qquad S_2 = -N_2, \qquad S_3 = N_3, \qquad S_4 = -N_4 .
\]
Bei den Drehungen kehrt sich das Vorzeichen um, weil $N_2$ und $N_4$ zu $\varphi = -1$ gehören.

\textbf{Beweisidee mit dem Satz von Betti.} Die folgende Herleitung soll die Idee zeigen und
ist kein formal vollständiger Beweis. So nennen wir alle $F_i$ Kräfte und alle Verschiebungen Wege,
obwohl bei den Drehungen Momente und Winkel gemeint sind, und die Einheitsverschiebung \enquote{$1$}
hat eigentlich eine Einheit (m bzw.\ rad). Das formal sauber mitzuschleppen ändert am Ergebnis
nichts, macht die Rechnung aber deutlich schwerer zu lesen.

Die Festeinspannkräfte $F_i$ sind die Kräfte und
Momente, die die Einspannungen auf den belasteten Stab ausüben, jeweils in positiver Richtung des
Freiheitsgrades (das $p$ in $s = k\,v + p$). Man schneidet den Stab frei. Die Einspannungen sind
dann nur noch Kräfte an den Stabenden, und man hat \emph{ein} System mit zwei Lastgruppen:

\begin{itemize}
  \item Lastgruppe~1: die Last $q(x)$ und die Endkräfte $F_j$, die die Enden festhalten (gesucht).
        Die Biegelinie ist $w_1$ mit $w_1 = w_1' = 0$ an beiden Enden, alle Endverschiebungen sind
        also $v_j^{(1)} = 0$.
  \item Lastgruppe~2: nur die Endkräfte $k_{ji}$, die den Freiheitsgrad $i$ um $1$ verschieben und
        alle anderen festhalten. Die Biegelinie ist $S_i(x)$, die Endverschiebungen sind
        $v_j^{(2)} = \delta_{ij}$ (Kronecker-Delta: $1$ für $j = i$, sonst $0$).
\end{itemize}

Beide Gruppen sind für sich im Gleichgewicht, deshalb stört es nicht, dass der freigeschnittene
Stab nicht gelagert ist. Nach Betti leistet Lastgruppe~1 auf den Wegen von Lastgruppe~2 dieselbe
Arbeit wie Lastgruppe~2 auf den Wegen von Lastgruppe~1:
\[
  \int_0^L q(x)\,S_i(x)\,dx + F_j\,v_j^{(2)} = k_{ji}\,v_j^{(1)} .
\]
Über das doppelte $j$ wird auf beiden Seiten summiert (Einsteinsche Summenkonvention). Setzt man
die Verschiebungen ein, bleibt links $F_j\,\delta_{ij} = F_i$ und rechts null. Die innere Arbeit
fehlt dabei nicht, sie ist der Grund für die Gleichheit:
\[
  W_a^{1 \to 2} = \int_0^L M_1\,\kappa_2\,dx = \int_0^L \frac{M_1 M_2}{EI}\,dx
  = \int_0^L M_2\,\kappa_1\,dx = W_a^{2 \to 1} .
\]
Das erste Gleichheitszeichen ist der Arbeitssatz (äußere gleich innere Arbeit) mit zwei
verschiedenen Zuständen: die Kräfte und Momente $M_1$ aus Lastgruppe~1, die im Gleichgewicht sind,
und die Verformung $\kappa_2$ aus Lastgruppe~2, die verträglich ist. Einer der beiden Zustände
spielt dabei die Rolle des virtuellen. Im PvV aus \textbf{TM1} ist es die Verschiebung (hier die
von Lastgruppe~2), beim Prinzip der virtuellen Kräfte aus \textbf{Statik} wäre es die Kraftgruppe.
Der Satz selbst gilt in beiden Lesarten, solange der eine Zustand im Gleichgewicht und der andere
verträglich ist. Das letzte Gleichheitszeichen ist derselbe Satz mit vertauschten Rollen: Kräfte
aus Lastgruppe~2, Verformung aus Lastgruppe~1. In der Mitte steht nur $\kappa = M/EI$. Betti
schreibt die beiden Enden dieser Kette hin. Die Kette enthält nur die Biegearbeit. Das Programm rechnet nach Euler-Bernoulli, der Stab ist
also schubstarr ($\gamma = 0$), und die Querkraft leistet keine innere Arbeit. Nach Timoschenko käme
in der Mitte noch $\int_0^L Q_1 Q_2/(GA_s)\,dx$ dazu. Betti gilt dann weiterhin, nur wären die $S_i$
andere Kurven, die auch von $GA_s$ abhängen.

Dass die gemischte innere Arbeit hier null ist, obwohl die Schnittgrößen aus Lastgruppe~2 auf
der Biegung aus Lastgruppe~1 lokal sehr wohl Arbeit leisten, sieht man auch direkt: $M_2$ ist ohne
Streckenlast linear, also $M_2'' = 0$, und zweimal partiell integriert bleibt
\[
  \int_0^L M_2\,\kappa_1\,dx = -\int_0^L M_2\,w_1''\,dx
  = \Bigl[-M_2\,w_1' + M_2'\,w_1\Bigr]_0^L - \int_0^L M_2''\,w_1\,dx = 0 ,
\]
weil $w_1$ und $w_1'$ an beiden Enden verschwinden. Also gilt für jede Lastfunktion
\[
  F_i = -\int_0^L q(x)\,S_i(x)\,dx .
\]
Mit $S_1 = N_1$, $S_2 = -N_2$, $S_3 = N_3$ und $S_4 = -N_4$ ergibt das genau die Formeln im
Programm:
\begin{align*}
  Q_a &= -\int_0^L q\,N_1\,dx, & M_a &= \int_0^L q\,N_2\,dx, \\
  Q_b &= -\int_0^L q\,N_3\,dx, & M_b &= \int_0^L q\,N_4\,dx.
\end{align*}

\textbf{Anschaulich} ist $S_i(x)$ die Einflusslinie der Festeinspannkraft $i$ (Satz von
Müller-Breslau): Eine Einzelkraft $P$ an der Stelle $a$ erzeugt $F_i = -P\,S_i(a)$, und das Integral
summiert nur die vielen kleinen Einzellasten $q\,dx$ auf.

\textbf{Kontrolle mit Gleichlast:} $\int_0^L N_1\,dx = L/2$, $\int_0^L N_2\,dx = L^2/12$ und
$\int_0^L N_4\,dx = -L^2/12$. Damit ist $Q_a = Q_b = -qL/2$ (die Last zeigt nach unten, die
Einspannung drückt dagegen) und $M_a = qL^2/12$, $M_b = -qL^2/12$, genau die Tafelwerte. Auch das
Vorzeichen von $M_a$ passt: Ohne Einspannung würde sich das linke Ende im Uhrzeigersinn drehen, die
Einspannung hält gegen den Uhrzeigersinn dagegen.

\textbf{Die anderen Lastarten.}
\begin{itemize}
  \item Ein \textbf{Streckenmoment} $m(x)$ (gegen den Uhrzeigersinn positiv) leistet Arbeit auf der
        Verdrehung, $\int m\,\varphi\,dx = -\int m\,w'\,dx$. Mit Betti wird daraus
        $F_i = \int_0^L m\,S_i'\,dx$, also $Q_a = \int m\,N_1'\,dx$, $M_a = -\int m\,N_2'\,dx$ und
        entsprechend am Ende $b$. Bei konstantem $m$ ist $\int N_1'\,dx = N_1(L) - N_1(0) = -1$ und
        $\int N_2'\,dx = 0$: $Q_a = -m$, $Q_b = +m$ und keine Endmomente. Das ist das Kräftepaar aus
        Abschnitt~\ref{sec:wgv-vorzeichen} und der Grund, warum das Programm vom $m\,l/2$ der
        WGV-Tafel abweicht.
  \item Für eine \textbf{Längslast} $n(x)$ sind die Formfunktionen linear, $1 - \xi$ und $\xi$ (sie
        lösen $EA\,u'' = 0$ exakt), und es gilt $N_a = -\int n\,(1-\xi)\,dx$, $N_b = -\int n\,\xi\,dx$.
  \item \textbf{Temperatur} ist keine Last, sondern eine eingeprägte Verzerrung. Die
        Festeinspannkräfte stehen direkt da: $N_a = EA\,\varepsilon_T$, $N_b = -EA\,\varepsilon_T$,
        $M_a = EI\,\kappa_T$, $M_b = -EI\,\kappa_T$.
\end{itemize}

\textbf{Zwei Wege im Code.} Für konstante Lasten und Trapezlasten
$q = p_a(1-\xi) + p_b\,\xi$ sind die Integrale schon ausgerechnet:
\begin{align*}
  Q_a &= -\tfrac{L}{60}\,(21p_a + 9p_b),   & M_a &= \tfrac{L^2}{60}\,(3p_a + 2p_b), \\
  Q_b &= -\tfrac{L}{60}\,(9p_a + 21p_b),   & M_b &= -\tfrac{L^2}{60}\,(2p_a + 3p_b).
\end{align*}
Hängt die Last von $x$ ab (auch bei $m(x)$, $n(x)$ und den globalen Linienlasten), rechnet das CAS
genau die Integrale oben aus. Ein Beispiel ist $q(x) = q_0 \sin(\pi x/L)$. Für die Querkraft
braucht man $\int_0^L q\,N_1\,dx$, und das geht ohne Rechnung in drei Schritten:
\begin{enumerate}
  \item $N_1 + N_3 = 1$: Schiebt man beide Stabenden um $1$ nach unten, verschiebt sich der ganze
        Stab um $1$. Also ist $\int_0^L q\,N_1\,dx + \int_0^L q\,N_3\,dx = \int_0^L q\,dx$.
  \item $N_3$ ist $N_1$ gespiegelt, $N_3(\xi) = N_1(1-\xi)$. Ist die Last symmetrisch zur
        Stabmitte, so wie der Sinus, sind beide Integrale deshalb gleich groß.
  \item Also ist jedes die Hälfte: $\int_0^L q\,N_1\,dx = \frac12 \int_0^L q\,dx
        = \frac12 \cdot \frac{2q_0L}{\pi} = \frac{q_0L}{\pi}$, und $Q_a = -q_0L/\pi$.
\end{enumerate}
Anschaulich: Bei symmetrischer Last trägt jede Einspannung die halbe Last. Für das Moment
braucht man $\int_0^1 \sin(\pi\xi)\,(\xi - 2\xi^2 + \xi^3)\,d\xi = 2/\pi^3$, also
$M_a = 2q_0L^2/\pi^3 \approx 0{,}0645\,q_0L^2$.

\textbf{Warum das exakt ist.} Betti gilt für jede Last, und die $S_i$ sind exakte Lösungen der
homogenen Gleichung. Also sind die Festeinspannkräfte exakt, damit auch die Knotenverschiebungen aus
$K\,u = F$ und die Stabendkräfte $s = k\,v + p$. Nicht exakt wäre nur die Hermite-Interpolation
$\sum_j v_j S_j(x)$ im Feld, denn die kennt die Last nicht. Die benutzt das Programm aber nie als
Ergebnis (siehe unten).

\textbf{Zwei Zerlegungen derselben Biegelinie.} Die Biegelinie eines Stabes ist immer
Partikulärlösung plus homogene Lösung. Welche Partikulärlösung man nimmt, ist aber nicht
festgelegt: Zwei Partikulärlösungen unterscheiden sich immer um eine kubische Parabel, also um eine
homogene Lösung. Das Programm benutzt zwei davon, und beide beschreiben dieselbe Linie.

\begin{itemize}
  \item \emph{Die Anfangswert-Lösung} ist das Vierfachintegral der Querlast,
        \[
          EI\,w_p(x) = \int_0^x q(t)\,\frac{(x-t)^3}{6}\,dt ,
        \]
        mit $w_p = w_p' = w_p'' = w_p''' = 0$ bei $x = 0$, so wie ihr sie in \textbf{TM2} aufstellt.
        Dazu kommt die homogene Lösung, deren vier Konstanten die Randbedingungen erfüllen. Das
        Rand-LGS schreibt sie als $c_1 x^3/6 + c_2 x^2/2 + c_3 x + c_4$ und holt die Konstanten aus
        allen Rand- und Übergangsbedingungen (Abschnitt~\ref{sec:mathe-randlgs}). Der alte Export
        \code{v\_EI\_stab} schreibt sie mit den Hermite-Funktionen, und die Koeffizienten sind die
        Endverschiebungen aus dem WGV:
        \[
          EI\,w = EI\,w_p + EI w_a\,N_1 + EI w_a'\,N_2 + \bigl(EI w_b - EI w_p(L)\bigr)\,N_3
                  + \bigl(EI w_b' - EI w_p'(L)\bigr)\,N_4 .
        \]
        $w_p$ bringt am Ende $b$ schon eigene Werte mit, die ziehen die letzten beiden Koeffizienten
        wieder ab, damit die Summe dort genau die Endverschiebung hat. Die Koeffizienten sind
        Durchbiegung und Neigung, so wie ihr die Randbedingungen in \textbf{TM2} aufstellt. In der
        Zählung des WGV ist $w_a = v_1$, $w_a' = -\varphi_a = -v_2$, $w_b = v_3$ und $w_b' = -v_4$.
  \item \emph{Die Lösung des eingespannten Stabes} ist
        \[
          w_0(x) = w_p(x) - w_p(L)\,N_3(x) - w_p'(L)\,N_4(x) .
        \]
        Auch sie löst $EI\,w'''' = q$, hat aber an \emph{beiden} Enden $w_0 = w_0' = 0$: Das ist die
        Biegelinie des beidseitig eingespannten Stabes unter der Last. Mit ihr zerfällt die Linie in
        \[
          w(x) = w_0(x) + \sum_j v_j\,S_j(x) ,
        \]
        und die Koeffizienten des homogenen Teils sind genau die Endverschiebungen $v_j$, also die
        Unbekannten des WGV. Diese Zerlegung ist die Idee hinter $s = k\,v + p$: Der homogene Teil
        $\sum_j v_j S_j$ liefert die Endkräfte $k\,v$. $w_0$ verschiebt die Stabenden nicht
        ($w_0 = w_0' = 0$ an beiden Enden), stört also die Verträglichkeit an den Knoten nicht und
        liefert nur Endkräfte, und das sind genau die Festeinspannkräfte $p$. Mehr braucht das WGV
        von der Last nicht.
\end{itemize}

Bei Gleichlast ist $EI\,w_p = qx^4/24$ mit $EI\,w_p(L) = qL^4/24$ und $EI\,w_p'(L) = qL^3/6$, und
daraus wird
\[
  EI\,w_0 = \frac{qL^4}{24}\bigl(\xi^4 - (3\xi^2 - 2\xi^3) - 4(\xi^3 - \xi^2)\bigr)
          = \frac{q\,x^2 (L-x)^2}{24},
\]
die bekannte Biegelinie des eingespannten Balkens. Ihr Moment am linken Ende ist
$-EI\,w_0''(0) = -qL^2/12$ (Zug oben), und die Einspannung hält mit $M_a = +qL^2/12$ gegen den
Uhrzeigersinn dagegen. Der Satz von Betti erspart es, $w_0$ überhaupt auszurechnen: Die
Festeinspannkräfte kommen direkt aus den Integralen über $q\,S_i$.

Die fertige Biegelinie ist in beiden Zerlegungen dieselbe. Das sieht man, wenn man die Formel des
alten Exports umsortiert. Mit $w_a'\,N_2 = -v_2\,N_2 = v_2\,S_2$ und genauso am Ende $b$ wird aus ihr
\[
  EI\,w = \underbrace{EI\,\bigl(w_p - w_p(L)\,N_3 - w_p'(L)\,N_4\bigr)}_{EI\,w_0}
          + EI \sum_j v_j\,S_j ,
\]
also Buchstabe für Buchstabe die zweite Zerlegung, nur anders gruppiert. Ausgerechnet wird $w_0$
dabei nie. Für die Knotengleichungen braucht das WGV nur seine Endkräfte $p$, und die kommen nach
Betti direkt aus den Integralen über $q\,S_i$. Die Verläufe rechnet das Programm danach aus den
exakten Endkräften (Gleichgewicht) bzw.\ mit der Anfangswert-Lösung plus Hermite-Anteil, wie im
ersten Punkt. Die zweite Zerlegung dient hier also nur der Erklärung, warum $s = k\,v + p$ exakt ist.

\subsection{Gelenke und die Knotenverdrehung}\label{sec:mathe-gelenke}

Hier kommt die Warnung aus Abschnitt~\ref{sec:Eingabe_Knoten} her. Ein \textbf{M-Gelenk} wird nicht
herauskondensiert, sondern bekommt einen eigenen Freiheitsgrad: Das Stabende am Gelenk hat seine
eigene Verdrehung. Zu diesem Freiheitsgrad gehört eine eigene Gleichgewichtsgleichung, und an ihm
greift kein äußeres Moment an. Somit ist das Gelenkmoment null. Für die angezeigte Knotenverdrehung
heißt das:

\begin{itemize}
  \item \textbf{Angelenk} (Gelenk am Stabende, andere Stäbe biegesteif am Knoten): Die
        Knotenverdrehung ist die Verdrehung der \emph{biegesteif} angeschlossenen Stäbe. Der Stab
        mit dem Gelenk dreht sich anders, seine Endverdrehung steht im Stabmenü auf Seite~7
        ($\varphi_a$, $\varphi_b$), der Knick im Gelenk ist die Differenz.
  \item \textbf{Vollgelenk} (Knoten gelenkig): Der Knoten hat \emph{gar keine} eigene Verdrehung,
        jedes Stabende dreht sich für sich. Die $0$ im Knotenmenü ist nur ein Platzhalter, die
        echten Verdrehungen stehen wieder auf Seite~7 der Stäbe.
  \item \textbf{Nur Angelenke an einem Knoten}, aber kein Vollgelenk und kein Lager, das die
        Drehung hält: Dann hat der Knoten eine Verdrehung, an der nichts hängt. Das Programm
        meldet \emph{kinematisch}. Deshalb der Stolperstein im Anhang: In diesem Fall braucht ihr
        das Vollgelenk.
\end{itemize}

\textbf{N- und Q-Gelenke} werden dagegen herauskondensiert (siehe Dallmann Baustatik 3): Im Stab wird der entsprechende Anteil
der Steifigkeit gestrichen, die Festeinspannkräfte werden statisch gleichwertig
umverteilt, und nach der Rechnung wird die Klaffung im Gelenk wieder ausgerechnet. Die
Endverschiebungen auf Seite~7 enthalten sie. Was dabei gestrichen wird, betrifft immer nur die
freigegebene Richtung:

\begin{itemize}
  \item \textbf{N-Gelenk:} Die Dehnsteifigkeit fällt weg. Am Gelenk ist $N = 0$, im Stab bleibt
        nur der Anteil aus Längslasten. Die Biegung ist davon nicht betroffen.
  \item \textbf{Q-Gelenk:} Am Gelenk ist $Q = 0$, ohne Querlast also im ganzen Stab, und wegen
        $M' = Q$ bleibt von der Biegesteifigkeit nur ein konstantes Moment übrig. Die Normalkraft
        trägt der Stab ganz normal weiter.
\end{itemize}

%---------------------------------------------------------------------
\section{Bögen: Teilstäbe und Schnittgrößen über den Winkel}\label{sec:mathe-bogen}

\subsection{Geometrie und Teilstäbe}

Ein Bogen ist ein Kreisstück mit Radius $R$ über der Sehne $L$ vom Stabanfang zum Stabende. Sein
Öffnungswinkel ist $\alpha = 2\arcsin\frac{L}{2R}$, und jeder Punkt auf ihm hat einen Winkel
$\varphi \in [0, \alpha]$ ab dem Stabanfang. Für die Rechnung ersetzt das Programm den Bogen durch
$n$ gerade Teilstäbe (Obermenü Seite~3, Standard 20). Das erste und das letzte Teilstück liegen
\emph{tangential} am Bogen und sind so lang wie $R\alpha/n$, die übrigen Teilungspunkte liegen auf
dem Bogen. Die tangentialen Endstücke haben einen Zweck: Die Stabendkräfte des ersten Teilstücks
sind dann direkt Normal- und Querkraft \emph{am Bogen}, und von dort startet der Export.

Die Teilstäbe erben Steifigkeiten, Gelenke an den Bogenenden, Temperatur und die Lasten:

\begin{itemize}
  \item $q$ (radial) und $n$ (tangential) bekommt jeder Teilstab mit den Werten an seinen beiden
        Endpunkten, dazwischen linear. Eine Lastfunktion wertet das Programm an den
        Teilungspunkten aus, mit $x = \varphi$ im Winkelmodus des Rechners.
  \item Lasten je Bogenlänge ($p_X$, $p_Y$ ohne Projektion, $m$) werden je Teilstab mit
        $R\,\Delta\varphi / l$ skaliert ($l$ = Länge des Teilstabs). Damit trägt jeder Teilstab genau
        die Last seines Bogenstücks, und die Summe ist exakt, z.\,B.\ das Eigengewicht $gR\alpha$.
  \item Projizierte Lasten gehen je Grundriss- bzw.\ Höhenlänge des Teilstabs ein.
\end{itemize}

\textbf{Warum konstante radiale Lasten exakt sind.} Ein konstanter Druck $q$ senkrecht auf einen
beliebigen Weg von $A$ nach $B$ hat die Resultierende $q$ mal den um $90^\circ$ gedrehten Vektor
$B - A$. Sein Moment um einen Punkt $P$ ist $\frac q2\left(|B-P|^2 - |A-P|^2\right)$. Beides hängt nur
von den Endpunkten ab, nicht vom Weg dazwischen. Der Linienzug der Teilstäbe hat dieselben Endpunkte
wie der Bogen, also stimmen bei einem statisch bestimmten System die Lagerkräfte und die
Schnittgrößen an den Teilungspunkten auf dem Bogen exakt: Sie folgen dort allein aus dem
Gleichgewicht, und das sieht nur die Resultierenden. Wie biegsam die Stäbe sind, spielt dabei keine
Rolle. Bei einem statisch unbestimmten System hängt die Kraftverteilung dagegen auch von der
Steifigkeit ab, und der Linienzug ist etwas anders steif als der Bogen, dann gilt das nur bis auf
den Teilungsfehler (siehe \emph{Genauigkeit} unten).
Beim Eigengewicht stimmt dagegen nur die Summe exakt. Die Hebel weichen um Größenordnung
$(\alpha/n)^2$ ab, mit 20 Teilstäben bei einem Viertelkreis um etwa 0,04\,\%.

\subsection{Schnittgrößen über den Winkel}

Der Export schreibt die Schnittgrößen als Funktionen des Winkels $p$ ab Stabanfang. Dafür schneidet
er den Bogen bei $p$ und schreibt das Gleichgewicht des Stücks von $0$ bis $p$ auf. Am Anfang wirken
$N_A$, $Q_A$, $M_A$ aus der Rechnung (erstes, tangentiales Teilstück). Mit $\varepsilon = +1$, wenn
die gestrichelte Faser innen liegt ($R > 0$), und $\varepsilon = -1$, wenn sie außen liegt, gilt
ohne Last auf dem Bogen
\begin{align*}
  N(p) &= N_A\cos p + \varepsilon\,Q_A\sin p , \\
  Q(p) &= Q_A\cos p - \varepsilon\,N_A\sin p , \\
  M(p) &= M_A + Q_A\,R\sin p - \varepsilon\,N_A\,R\,(1-\cos p) .
\end{align*}
Die ersten beiden Zeilen drehen nur die Kraft vom Anfang in die Tangente bei $p$. In der dritten
stehen die Hebel: Der Schnittpunkt liegt um $R\sin p$ entlang der Anfangstangente und um
$R(1-\cos p)$ zum Mittelpunkt hin versetzt.
Hier sieht man auch, warum im Bogen Normalkraft und Biegung gekoppelt sind: $N_A$ hat am
gekrümmten Stab den Hebel $R(1-\cos p)$ und erzeugt ein Moment. Beim geraden Stab ($R \to \infty$
mit $R\,p = x$) geht dieser Hebel gegen $x^2/(2R) \to 0$, und $Q_A R\sin p$ wird zu $Q_A\,x$. Das
ist reines Gleichgewicht am unverformten Stab, also Theorie I.~Ordnung. Theorie II.~Ordnung hieße,
das Gleichgewicht am \emph{verformten} System aufzustellen (Knicken, Durchschlagen flacher Bögen),
das rechnet das Programm weder bei Bögen noch bei geraden Stäben.

Last auf dem Bogenstück kommt als Integral über den Lastwinkel $t$ von $0$ bis $p$ dazu (im Bogenmaß,
$ds = R\,dt$):
\begin{align*}
  \Delta N(p) &= -R\int_0^p \Bigl(\varepsilon\,q(t)\sin(p-t) + n(t)\cos(p-t)\Bigr)\,dt , \\
  \Delta Q(p) &= -R\int_0^p \Bigl(q(t)\cos(p-t) - \varepsilon\,n(t)\sin(p-t)\Bigr)\,dt , \\
  \Delta M(p) &= -R\int_0^p \Bigl(R\,q(t)\sin(p-t) - \varepsilon\,R\,n(t)\bigl(1-\cos(p-t)\bigr)
                + m(t)\Bigr)\,dt .
\end{align*}
Auch hier stehen die Hebel direkt da. Eine radiale Kraft beim Winkel $t$ geht durch den
Mittelpunkt, ihr Abstand zum Schnittpunkt bei $p$ ist $R\sin(p-t)$. Eine tangentiale Kraft liegt
auf der Tangente bei $t$, und der Schnittpunkt hat von dieser Tangente den Abstand $R(1-\cos(p-t))$.
Globale Linienlasten $p_X$, $p_Y$ gehen genauso ein, nur mit Hebeln aus den Koordinaten. Ihre
Richtung relativ zur Tangente ändert sich mit $p$.

\textbf{Winkelmodus.} Steht der Rechner auf DEG, laufen $p$ und $t$ in Grad. Dann ist
$ds = R\,\frac{\pi}{180}\,dt$, und genau diesen Faktor setzt das Programm vor die Integrale. Den
Modus erkennt es daran, ob $\cos(180) = -1$ ist. $\sin p$ und $\cos p$ rechnet der Rechner
ohnehin im eingestellten Modus, deshalb stimmen die Funktionen in jedem Modus.

\textbf{Genauigkeit.} Die Formeln sind das exakte Gleichgewicht am Kreisbogen. Ungenau sein können
nur die Startwerte $N_A$, $Q_A$, $M_A$, denn sie stammen aus der Rechnung mit Teilstäben. Bei einem
statisch bestimmten System mit Lasten, die die Teilstäbe exakt abbilden (Knotenlasten, konstante
radiale Last, projizierte Lasten), sind auch sie exakt. Sonst gehen die Abweichungen mit mehr
Teilstäben zurück: bei Eigengewicht, tangentialen Lasten und Lastfunktionen, und bei statisch
unbestimmten Systemen, in denen die Steifigkeit des Linienzugs etwas von der des Bogens abweicht.

%---------------------------------------------------------------------
\section{Starre Stäbe: Zwangsbedingungen und Nullraum}\label{sec:mathe-starr}

Warum nicht einfach $EA = 10^{15}$? Weil der Rechner mit etwa 16 gültigen Stellen rechnet. Steht
in derselben Matrix eine Feder mit $c = 1$ neben $10^{15}$, bleibt von der Feder nach der
Zerlegung praktisch nichts übrig. Das ist der Zahlenmüll aus Abschnitt~\ref{sec:starr}. Das
Programm macht es deshalb exakt.

\textbf{1. Zwangsbedingungen.} Ein starrer Anteil bekommt in der Steifigkeitsmatrix die Steifigkeit
\emph{null}. Stattdessen schreibt das Programm für jeden starren Stab auf, was \enquote{starr}
bedeutet, und zwar in den lokalen Endverschiebungen:
\begin{align*}
  \text{dehnsteif:}\quad & u_b - u_a = \varepsilon_T L, \\
  \text{biegesteif:}\quad & \varphi_a - \varphi_b = -\kappa_T L, \qquad
                            w_b - w_a + L\,\varphi_a = -\kappa_T \frac{L^2}{2} .
\end{align*}
Ohne Temperatur steht rechts null: gleiche Länge, gleiche Endverdrehung, Enden auf einer Geraden.
Alle Bedingungen zusammen, ins globale System gedreht, ergeben $C\,u = g$.

\textbf{2. Kern der Matrix (=Nullraum).} $C$ wird mit Gauß-Jordan auf reduzierte Stufenform gebracht. Jede Zeile hat
dann einen \emph{abhängigen} Freiheitsgrad (die Pivotspalte), der durch die übrigen, \emph{freien}
Freiheitsgrade ausgedrückt wird. Damit lässt sich jede erlaubte Verschiebung schreiben als
\[
  u = T\,u_r + u_0 , \qquad C\,T = 0 .
\]
Die Spalten von $T$ spannen den Nullraum bzw. den Kern von $C$ auf, also genau die Bewegungen, die die starren
Stäbe mitmachen. $u_r$ sind nur noch die freien Freiheitsgrade, $u_0$ ist der Temperaturanteil.

\textbf{3. Projizieren.} Eingesetzt in $K\,u = F$ und von links mit $T^T$ multipliziert:
\begin{equation} \label{eq:LGSred}
  \left(T^T K\,T\right) u_r = T^T\left(F - K\,u_0\right).
\end{equation}
Das ist nichts anderes als das Prinzip der virtuellen Verschiebungen: Als virtuelle
Verschiebungen sind nur noch solche zugelassen, die die Zwangsbedingungen erfüllen, also Spalten
von $T$. Das neue System ist kleiner als das alte und genauso gut konditioniert wie ein System
ohne starre Stäbe, weil nirgends eine riesige Zahl steht.

\textbf{4. Kräfte in den starren Stäben.} Nach Schritt~3 kennt man die Verschiebungen $u$, aber
noch keine Kraft in einem starren Stab, denn ihr Anteil in $K$ ist null. Setzt man $u$ in das
Gleichgewicht der Knoten ein, bleibt ein Rest
\[
  r = F - K\,u .
\]
Das sind die Knotenkräfte, die weder die elastischen Stäbe noch die Federn aufnehmen. Diesen Rest
müssen die starren Stäbe tragen.

Das widerspricht nicht $K\,u = F$. In $K$ fehlen die starren Stäbe, denn ihre Kraft lässt sich
nicht aus $u$ ausrechnen: Ein starrer Stab verformt sich bei jeder Kraft um null. Das vollständige
Gleichgewicht der Knoten lautet
\[
  K\,u + (\text{Kräfte der starren Stäbe}) = F ,
\]
und $r$ ist genau der fehlende Anteil. Das kennt man von den Lagern: Auch an einem gehaltenen
Freiheitsgrad bleibt ein Rest, und der ist die Auflagerkraft. Ein starrer Stab ist so etwas wie ein
inneres Lager, das statt einer absoluten eine relative Verschiebung festhält.

\textbf{Woher die $\lambda$ kommen.} Schritt~3 hat nicht $K\,u = F$ gelöst, sondern nur
\eqref{eq:LGSred}. Alles auf eine Seite gebracht und $u = u_0 + T\,u_r$ eingesetzt:
\begin{align*}
  0 &= T^T F - T^T K\,u_0 - T^T K\,T\,u_r \\
    &= T^T\bigl(F - K\,(u_0 + T\,u_r)\bigr) = T^T\,(F - K\,u) = T^T r .
\end{align*}
Die Zeilen von $T^T$ sind die Spalten von $T$, also eine Basis des Kerns von $C$. $T^T r = 0$ sagt
deshalb: $r$ steht auf jedem dieser Basisvektoren und damit auf dem ganzen Kern von $C$ senkrecht.

Daraus folgt, dass sich $r$ aus den Zeilen von $C$ zusammensetzen lässt. Das ist reine lineare
Algebra aus \textbf{HM1}. $C$ hat $m$ Zeilen (eine je Zwangsbedingung) und $n$ Spalten (eine je
Freiheitsgrad). Im $\mathbb R^n$ liegen damit zwei Unterräume:
\begin{itemize}
  \item $\operatorname{Kern} C$: alle $x$ mit $C\,x = 0$,
  \item $\operatorname{Bild} C^T$: alle Vektoren $C^T\lambda$, also alle Kombinationen der Zeilen
        von $C$.
\end{itemize}
Gezeigt wird $r \in \operatorname{Bild} C^T$, und zwar in vier Schritten.

\emph{(a) $r$ steht senkrecht auf $\operatorname{Kern} C$.} Das haben wir gerade gesehen.

\emph{(b) $\operatorname{Bild} C^T$ steht senkrecht auf $\operatorname{Kern} C$.} Für ein $x$ aus
$\operatorname{Kern} C$ und einen beliebigen Vektor $C^T\lambda$ gilt
\[
  (C^T\lambda)^T x = \lambda^T C\,x = \lambda^T 0 = 0 .
\]
Das erste Gleichheitszeichen ist nur die Rechenregel $(AB)^T = B^T A^T$. Insbesondere liegt außer
der Null kein Vektor in beiden Räumen, denn er müsste auf sich selbst senkrecht stehen.

\emph{(c) Zusammen füllen beide Räume den ganzen $\mathbb R^n$.} Nach der Dimensionsformel ist
$\dim\operatorname{Kern} C + \dim\operatorname{Bild} C = n$, und weil Zeilenrang gleich Spaltenrang
ist, gilt $\dim\operatorname{Bild} C^T = \dim\operatorname{Bild} C$. Also ist
$\dim\operatorname{Kern} C + \dim\operatorname{Bild} C^T = n$. Weil sich die beiden Räume nur in der
Null treffen, lässt sich jeder Vektor des $\mathbb R^n$ eindeutig als Summe aus einem Teil in
$\operatorname{Kern} C$ und einem Teil in $\operatorname{Bild} C^T$ schreiben.

\emph{(d) $r$ liegt in $\operatorname{Bild} C^T$.} Man zerlegt $r = r_K + r_B$ mit $r_K$ in
$\operatorname{Kern} C$ und $r_B$ in $\operatorname{Bild} C^T$. Nach (a) und (b) stehen $r$ und
$r_B$ senkrecht auf $\operatorname{Kern} C$, also auch $r_K = r - r_B$. $r_K$ liegt aber selbst in
$\operatorname{Kern} C$, steht also senkrecht auf sich selbst, und das kann nur der Nullvektor.
Damit ist $r = r_B$.

Ein Beispiel mit Zahlen: $30$ Freiheitsgrade und ein starrer Stab, der dehn- und biegesteif ist
($3$ unabhängige Zeilen). $\operatorname{Kern} C$ ist $27$-dimensional, das sind die $27$
Unbekannten $u_r$ in Schritt~3. $\operatorname{Bild} C^T$ ist $3$-dimensional, das sind die drei
$\lambda$.

Es gibt also Zahlen $\lambda_k$, eine je Zwangsbedingung, mit
\[
  r = C^T \lambda = \sum_k \lambda_k\,C_k^T ,
\]
wobei $C_k$ die $k$-te Zeile von $C$ ist. Weil $r$ nach dem vollständigen Gleichgewicht genau die
Kräfte der starren Stäbe sind, heißt das: Die starren Stäbe wirken mit $C^T\lambda$ auf die Knoten,
und jede Zwangsbedingung~$k$ trägt $\lambda_k\,C_k^T$ dazu bei. Das vollständige Gleichgewicht
lautet damit
\[
  K\,u + C^T\lambda = F .
\]
Mit $T^T$ multipliziert fällt $\lambda$ wegen $T^T C^T = (C\,T)^T = 0$ heraus, das ist Schritt~3.
Die übrigen Gleichungen erfüllt Schritt~4 mit der passenden Wahl von $\lambda$. Die $\lambda$ sind
übrigens die Lagrange-Multiplikatoren, die man bekommt, wenn man das Minimum der potentiellen
Energie $\frac12\,u^T K\,u - F^T u$ unter der Nebenbedingung $C\,u = g$ sucht.

\textbf{Mechanisch gelesen.} Das Skalarprodukt aus Kraft und Verschiebung ist die Arbeit,
\enquote{senkrecht} heißt also \enquote{keine Arbeit}. (a) sagt dann: Die Kräfte der starren Stäbe
leisten auf keiner Bewegung Arbeit, die die starren Stäbe erlauben. Genau das zeichnet eine
Zwangskraft aus. Eine reibungsfreie Führung leistet auf der Bewegung, die sie zulässt, auch keine
Arbeit, sie verhindert nur alle anderen. Physikalisch steckt diese Annahme schon in Schritt~3: Dort
sind als virtuelle Verschiebungen nur solche zugelassen, die die Zwangsbedingungen erfüllen. Alles
danach ist reine Mathematik.

Eindeutig sind die $\lambda$ nur, wenn die Zeilen von $C$ linear unabhängig sind, also
$\operatorname{Kern} C^T = \{0\}$ gilt. Dann gehört zu jeder Kraft aus $\operatorname{Bild} C^T$
genau ein $\lambda$. Ist $\operatorname{Kern} C^T$ größer, gibt es Kräfte in den starren Stäben,
die auf keinen Knoten wirken. Das sind die Eigenspannungszustände weiter unten.

\textbf{Warum die $\lambda$ Schnittgrößen sind.} Das folgt aus dem Gleichgewicht des einzelnen
starren Stabes, und zwar mit demselben Werkzeug wie eben: Kern, Bild und Dimensionsformel, diesmal
für einen Stab allein. Im lokalen System hat ein Stab die sechs Endverschiebungen
$v = (u_a, w_a, \varphi_a, u_b, w_b, \varphi_b)$ und sechs Endkräfte $s$ in denselben Richtungen.
Die drei Zeilen eines biege- und dehnsteifen Stabes bilden die $3 \times 6$-Matrix
\[
  C_\ell = \begin{pmatrix} -1 & 0 & 0 & 1 & 0 & 0 \\ 0 & 0 & 1 & 0 & 0 & -1 \\ 0 & -1 & L & 0 & 1 & 0 \end{pmatrix}
  \qquad
  \begin{matrix} \text{Längszeile: } u_b - u_a \\ \text{Drehzeile: } \varphi_a - \varphi_b \\
                 \text{Querzeile: } w_b - w_a + L\,\varphi_a \end{matrix}
\]

\emph{(a) Starr heißt: nur Starrkörperbewegungen.} Ein Stab kann sich auf drei Arten bewegen, ohne
sich zu verformen: längs verschieben, quer verschieben und um das Ende $a$ drehen. In
Endverschiebungen sind das die Spalten von
\[
  R = \begin{pmatrix} 1 & 0 & 0 & 1 & 0 & 0 \\ 0 & 1 & 0 & 0 & 1 & 0 \\ 0 & 0 & 1 & 0 & -L & 1 \end{pmatrix}^T
\]
(bei der Drehung um $\varphi = 1$ wandert das Ende $b$ wegen $\varphi = -w'$ um $w_b = -L$).
Einsetzen zeigt $C_\ell R = 0$: Starrkörperbewegungen verformen nichts, $\operatorname{Bild} R
\subseteq \operatorname{Kern} C_\ell$. Beide Matrizen haben Rang~$3$, nach der Dimensionsformel ist
also $\dim \operatorname{Kern} C_\ell = 6 - 3 = 3 = \dim \operatorname{Bild} R$, und damit
$\operatorname{Kern} C_\ell = \operatorname{Bild} R$. Die Zwangsbedingungen lassen genau die
Starrkörperbewegungen zu, nicht mehr und nicht weniger.

\emph{(b) Gleichgewicht heißt: keine Arbeit auf Starrkörperbewegungen.} Zuerst ein Stab ohne
Stablasten, die kommen in (d) dazu. Endkräfte $s$ sind genau dann im Gleichgewicht, wenn
\[
  R^T s = \begin{pmatrix} s_1 + s_4 \\ s_2 + s_5 \\ s_3 - L\,s_5 + s_6 \end{pmatrix} = 0 ,
\]
das sind Summe der Kräfte längs, Summe der Kräfte quer und Summe der Momente um $a$. Anders gelesen
ist das das PvV für einen starren Körper aus \textbf{TM1}: Auf jeder Starrkörperbewegung leisten die
Kräfte keine Arbeit. Die Endkräfte, die für sich im Gleichgewicht sind, bilden also den Raum
$\operatorname{Kern} R^T$.

\emph{(c) Dieser Raum wird genau von den Zeilen aufgespannt.} Wie beim Nullraum oben in zwei
Schritten: $R^T C_\ell^T = (C_\ell R)^T = 0$, also $\operatorname{Bild} C_\ell^T
\subseteq \operatorname{Kern} R^T$. Dimension: $\dim \operatorname{Bild} C_\ell^T = \operatorname{Rang}
C_\ell = 3$ und $\dim \operatorname{Kern} R^T = 6 - \operatorname{Rang} R = 3$. Also
\[
  \operatorname{Kern} R^T = \operatorname{Bild} C_\ell^T .
\]
Jede Gruppe von Endkräften im Gleichgewicht ist eine Kombination der drei Zeilen, und weil die
Zeilen linear unabhängig sind, auf genau eine Weise.

\emph{(d) Anwendung auf den starren Stab.} Auf den starren Stab wirken zwei Dinge: die Lasten auf
dem Stab ($q$, $n$, $m$) und die Endkräfte $s$, mit denen die Knoten auf seine Enden wirken. $s$
sind also nur die Endkräfte, die Lasten stehen getrennt. Zusammen sind beide im Gleichgewicht.
Die Festeinspannkräfte $p$ sind ebenfalls Endkräfte, die dieselben Lasten im Gleichgewicht
halten. Dass sie zu einem beidseitig eingespannten Stab gehören, spielt hier keine Rolle. Man
braucht nur \emph{irgendeinen} Satz Endkräfte, der die Lasten allein ausgleicht, und $p$ hat das
Programm aus dem WGV ohnehin. Zieht man die beiden Gleichgewichtsgruppen
\enquote{Lasten $+ s$} und \enquote{Lasten $+ p$} voneinander ab, fallen die Lasten heraus: $s - p$
ist eine Gruppe reiner Endkräfte im Gleichgewicht, also genau der Fall aus (b) ohne Stablasten.
Ausgeschrieben: Nach (b) gilt $R^T (s - p) = 0$, also liegt
$s - p$ in $\operatorname{Kern} R^T$, und das ist nach (c) $\operatorname{Bild} C_\ell^T$. Es gibt
also genau ein $\lambda_\ell \in \mathbb R^3$ mit $s - p = C_\ell^T \lambda_\ell$, umgestellt
\[
  s = p + C_\ell^T \lambda_\ell .
\]
Ins globale System gedreht und über alle starren Stäbe summiert sind das genau die Knotenkräfte
$C^T\lambda$ von oben, die $\lambda_\ell$ eines Stabes sind seine Einträge in $\lambda$. Ist ein Stab
nur dehnsteif, fehlen Dreh- und Querzeile, die Biegung liefert dann wie bei jedem Stab $k\,v$.
Ein N-Gelenk streicht die Längszeile, ein Q-Gelenk die Querzeile. Das hat zwei Gründe, die
zusammenpassen. Kinematisch kann sich das Stabende in der freigegebenen Richtung gegenüber dem
Knoten verschieben. Die Zeile würde aber den Knoten selbst festhalten, also mehr verlangen, als der
Stab tut. Und weil N- und Q-Gelenke herauskondensiert werden, hat das Stabende in dieser Richtung
keinen eigenen Freiheitsgrad, mit dem man die Zeile richtig schreiben könnte. Statisch ist die
zugehörige Schnittgröße am Gelenk null. Die Lasten trägt schon $p$, das für Stäbe mit N- oder
Q-Gelenk passend umverteilt ist, ein $\lambda$ für diese Zeile wäre also ohnehin null. Ein
M-Gelenk streicht dagegen nichts: Das Stabende hat dort einen eigenen Drehfreiheitsgrad
(Abschnitt~\ref{sec:mathe-gelenke}). Die Drehzeile verbindet die Verdrehung des Stabes mit diesem
Freiheitsgrad statt mit dem Knoten, und $M = 0$ am Gelenk folgt aus dessen Gleichung. $p$ darf man dabei unverändert nehmen, weil die
Festeinspannkräfte nach Abschnitt~\ref{sec:mathe-festeinspann} gar nicht von $EI$ abhängen.

\emph{(e) Was die drei Zahlen bedeuten.} Für jede Zeile liest man den Schnittgrößenverlauf des
unbelasteten Stabes direkt ab, mit $N(0) = -s_1$, $Q(0) = -s_2$, $M(0) = -s_3$ wie bei jedem Stab:
\begin{align*}
  \text{Längszeile:}\quad & s = \lambda_N\,(-1, 0, 0, 1, 0, 0)^T: && N(x) = \lambda_N, \\
  \text{Drehzeile:}\quad  & s = \lambda_M\,(0, 0, 1, 0, 0, -1)^T: && M(x) = -\lambda_M, \\
  \text{Querzeile:}\quad  & s = \lambda_Q\,(0, -1, L, 0, 1, 0)^T: && Q(x) = \lambda_Q, \quad
                                                                     M(x) = -\lambda_Q\,(L - x) .
\end{align*}
$\lambda_N$ ist also die Normalkraft, $\lambda_Q$ die Querkraft und $-\lambda_M$ das Moment am Ende
$b$. Zusammen $M(x) = -\lambda_M - \lambda_Q\,(L - x)$, dazu kommt der Anteil aus $p$. Mechanisch
passt das: Die Arbeit von $\lambda_k\,C_k^T$ auf den Endverschiebungen ist $\lambda_k\,(C_k\,v)$,
die Normalkraft arbeitet also auf der Längenänderung, das Moment auf der gegenseitigen Verdrehung
und die Querkraft darauf, wie weit das Ende $b$ von der Tangente am Ende $a$ abweicht. Das ist mit
dem Satz in Abschnitt~\ref{sec:starr} gemeint, die Schnittgrößen starrer Stäbe würden \enquote{aus
dem Gleichgewicht bestimmt}: Sie sind die $\lambda$, und die kommen allein aus dem Rest $r$, also aus
dem Gleichgewicht der Knoten, nicht aus einer Steifigkeit.

\textbf{Wie das Programm sie ausrechnet.} $r = C^T\lambda$ hat eine Gleichung je Freiheitsgrad, aber
nur eine Unbekannte je Zwangsbedingung, also meist mehr Gleichungen als Unbekannte. Weil $r$ nach
Schritt~3 eine Kombination der Zeilen von $C$ ist, geht es trotzdem exakt auf. Von links mit $C$
multipliziert wird daraus ein kleines quadratisches System
\[
  (C\,C^T)\,\lambda = C\,r ,
\]
das eindeutig lösbar ist, solange die Zeilen von $C$ linear unabhängig sind.

\textbf{Starre Teile, die sich gegenseitig halten.} Beispiele sind zwei dehnstarre Stäbe
nebeneinander zwischen denselben zwei Knoten oder ein biegestarrer Zweifeldträger (zwei starre
Stäbe über drei Lager). Dann sind einige Zeilen von $C$ linear abhängig. Das merkt schon die
Elimination in Schritt~2: Eine solche Zeile wird null, nachdem die früheren Pivotspalten abgezogen
sind. Das Programm protokolliert dabei, welche Kombination der ursprünglichen Zeilen null ergeben
hat:
\[
  \sum_k \mu_k\,C_k = 0 \qquad\Longleftrightarrow\qquad C^T\mu = 0 .
\]
Physikalisch ist $\mu$ ein Satz von Kräften in den starren Stäben, der auf keinen Knoten eine Kraft
ausübt, ein \emph{Eigenspannungszustand}. Bei den zwei Stäben nebeneinander ist das Zug~$1$ im
einen und Druck~$1$ im anderen.

Für die Kinematik ist eine abhängige Zeile harmlos, ohne Temperatur ist sie automatisch erfüllt. Für
das Gleichgewicht heißt sie: $r = C^T\lambda$ hat unendlich viele Lösungen
\[
  \lambda = \lambda_p + \sum_i t_i\,\mu_i .
\]
Das ist die Aufteilung, die man von jedem LGS kennt: eine partikuläre Lösung plus beliebige
Lösungen der homogenen Gleichung $C^T\mu = 0$. Und es ist genau die Aufteilung aus dem KGV in
\textbf{Statik}:
\begin{itemize}
  \item $\lambda_p$ rechnet das Programm mit $(C\,C^T)\,\lambda = C\,r$ nur über die unabhängigen
        Zeilen, die abhängigen bekommen $\lambda = 0$. Das ist, als würde man die überzähligen
        Zwangsbedingungen aufschneiden: ein statisch bestimmtes Hauptsystem mit seinem
        $0$-Zustand.
  \item Die $\mu_i$ sind die $1$-Zustände. Jedes $\mu_i$ hat an seiner abhängigen Zeile eine $1$,
        dazu kommen die Kräfte in den übrigen starren Stäben, die diese $1$ im Gleichgewicht halten.
  \item Die $t_i$ sind die statisch Überzähligen $X_i$.
\end{itemize}
Wie im KGV legt das Gleichgewicht die $X_i$ nicht fest. Dort kommen sie aus der Verträglichkeit,
und die braucht Nachgiebigkeiten. Genau die haben starre Stäbe nicht: Jede Wahl der $t_i$ ergibt
dieselben Verschiebungen, und die Starrkörpermechanik kann nicht sagen, welche stimmt. Das hängt
davon ab, wie starr die Stäbe \emph{im Verhältnis zueinander} wirklich sind.

\textbf{Der Grenzwert und die Ergänzungsenergie.} Das Programm trifft dieselbe Annahme wie der
Grenzwert im Rand-LGS: Alle starren Stäbe bekommen $EA = EI = 1/\varepsilon$ mit
$\varepsilon \to 0$. Dann verformen sie sich ein kleines bisschen. In den Größen, die die Zeilen
von $C$ messen, ist ihre elastische Verformung
\[
  C\,u - g = \varepsilon\,W\,\lambda ,
\]
wobei $W$ die Nachgiebigkeit der starren Stäbe mit $EA = EI = 1$ ist (blockweise je Stab, siehe
unten). Das ist das Elastizitätsgesetz der fast starren Stäbe.

Für das Folgende muss man zwei Dinge auseinanderhalten: $u$ sind die Knotenverschiebungen, $C\,u$
sind die Verformungen der starren Stäbe, gemessen in den Größen der Zwangsbedingungen
(Längenänderung, gegenseitige Verdrehung, Abweichung von der Tangente, siehe (e)). Die Stabkräfte
$\lambda$ arbeiten auf $C\,u$, ihre Knotenkräfte $C^T\lambda$ auf $u$, und beides ist dieselbe
Arbeit: $\lambda^T (C\,u) = (C^T\lambda)^T u$, das ist wieder nur $(AB)^T = B^T A^T$.

Jetzt kommt die Verträglichkeit, mit dem Arbeitssatz wie im KGV. Ein Eigenspannungszustand $\mu_i$
übt auf keinen Knoten eine Kraft aus, $C^T\mu_i = 0$. Seine Arbeit auf den wirklichen Verformungen
der starren Stäbe ist deshalb null:
\[
  \mu_i^T (C\,u) = (C^T\mu_i)^T u = 0^T u = 0 .
\]
Multipliziert man das Elastizitätsgesetz $C\,u - g = \varepsilon\,W\,\lambda$ von links mit
$\mu_i^T$, steht links genau diese Arbeit, und es bleibt
\[
  \underbrace{\mu_i^T C\,u}_{=\,0} - \mu_i^T g = \varepsilon\,\mu_i^T W\,\lambda
  \qquad\Longrightarrow\qquad
  \varepsilon\,\mu_i^T W\,\lambda + \mu_i^T g = 0 .
\]
Die Temperatur $g$ bleibt dabei stehen. Ohne Temperatur fällt der zweite Term weg, das $\varepsilon$ kürzt sich heraus, und mit
$\lambda = \lambda_p + \sum_j t_j\,\mu_j$ bleibt
\[
  \sum_j \underbrace{\mu_i^T W\,\mu_j}_{\delta_{ij}}\;t_j
  = -\underbrace{\mu_i^T W\,\lambda_p}_{\delta_{i0}} .
\]
Das sind die Elastizitätsgleichungen des KGV (hier ist $\delta_{ij}$ wieder eine $\delta$-Zahl und
kein Kronecker-Delta). Weil $\varepsilon$ herausfällt, hängt das Ergebnis nicht davon ab,
\emph{wie} starr die Stäbe sind, nur davon, dass sie alle gleich starr sind. Deshalb existiert der
Grenzwert. Passt die Temperatur dagegen nicht zu den abhängigen Zeilen ($\mu_i^T g \ne 0$), muss
$\varepsilon\,\mu_i^T W \lambda$ einen festen Wert annehmen, und für $\varepsilon \to 0$ wird
$\lambda$ unendlich groß. Das ist die unendliche Zwangskraft, zu der das Programm einen Hinweis
zeigt.

\textbf{Wo die Ergänzungsenergie hineinspielt.} Die Gleichungen oben sind genau die Bedingung dafür,
dass
\[
  \Pi^*(t) = \tfrac12\,\lambda^T W \lambda , \qquad \lambda = \lambda_p + \sum_i t_i\,\mu_i ,
\]
minimal wird, denn $\partial\Pi^*/\partial t_i = \mu_i^T W \lambda$. Das ist der Satz von Menabrea:
Unter allen Kraftzuständen im Gleichgewicht stellt sich der mit der kleinsten Ergänzungsenergie
ein. \enquote{Kleinste Ergänzungsenergie} ist dabei nur eine andere Schreibweise für
\enquote{verträglich}, das KGV in \textbf{Statik} rechnet mit dem Prinzip der virtuellen Kräfte
genau dasselbe. Die Ergänzungsenergie ist also kein zusätzliches Prinzip, sondern die
Verträglichkeit in Energieform. Genau genommen ist die Ergänzungsenergie der fast starren Stäbe
$\varepsilon\,\Pi^*$, denn ihre Nachgiebigkeit ist $\varepsilon\,W$ und nicht $W$. Ein fester
Faktor ändert aber nichts daran, wo das Minimum liegt, deshalb steht $\varepsilon$ oben gar nicht
erst in $\Pi^*$.

\textbf{$W$ ausrechnen} geht wie eine $\delta$-Zahl mit der Koppeltafel, mit $EA = EI = 1$ und für
jeden Stab für sich (verschiedene Stäbe koppeln nicht):
\begin{itemize}
  \item Längszeile: $N = 1$ konstant, $\int N^2\,dx = L$.
  \item Drehzeile: $M = 1$ konstant, $\int M^2\,dx = L$.
  \item Querzeile: $M$ linear von $L$ auf $0$, $\int M^2\,dx = L^3/3$.
  \item Dreh- mal Querzeile: $\int 1 \cdot (L - x)\,dx = L^2/2$. Längs- mal Biegezeile: $0$.
\end{itemize}
Die Querkraftverformung fehlt, wie überall im Programm (Euler-Bernoulli). Dass $EA$ und $EI$ beide
$1$ sind, ist eine Annahme und keine Physik. Sie spielt keine Rolle, solange in einem
Eigenspannungszustand nur dehnstarre oder nur biegestarre Zeilen vorkommen, denn dann kürzt sich der
gemeinsame Faktor. Mischen sich beide, hängt die Aufteilung vom angenommenen Verhältnis $EA : EI$
ab.

So teilen sich zwei dehnstarre Stäbe nebeneinander eine Last je zur Hälfte, und ein starrer
Zweifeldträger unter Gleichlast hat die bekannten Lagerkräfte $\frac38\,ql$, $\frac{10}8\,ql$,
$\frac38\,ql$, genau wie mit dem KGV und dem Stützmoment als $X_1$. Eine gelbe Box weist darauf
hin, dass diese Aufteilung angenommen ist, aber nur, wenn sie das Ergebnis auch ändert. Ist
$\lambda_p$ auf allen Zeilen null, die über $W$ mit einem $\mu_i$ gekoppelt sind, dann ist
$\delta_{i0} = \mu_i^T W \lambda_p = 0$ für \emph{jede} Wahl der Nachgiebigkeiten, also $t_i = 0$, egal wie
starr die Stäbe im Verhältnis zueinander sind. Ein Beispiel ist eine starre Platte, die oben und
unten auf Loslagern steht: Die Längskraft zwischen den beiden Lagern ist unbestimmt, aber ohne Last
in Plattenrichtung immer null.

\textbf{Das Gesamtbild.} Beide Aufteilungen in partikulär und homogen gehören zusammen. Die
Kinematik $C\,u = g$ und das Gleichgewicht $C^T\lambda = r$ benutzen dieselbe Matrix, einmal so und
einmal transponiert:

\begin{center}
\begin{tabular}{@{}p{2.9cm}p{5.9cm}p{5.9cm}@{}}
  \toprule
   & \textbf{Kinematik} $C\,u = g$ & \textbf{Gleichgewicht} $C^T\lambda = r$ \\
  \midrule
  partikulär       & $u_0$ (Temperatur) & $\lambda_p$ (Hauptsystem) \\
  homogen          & $T\,u_r$: Bewegungen, die die starren Stäbe erlauben
                   & $\sum_i t_i\,\mu_i$: Eigenspannungszustände \\
  festgelegt durch & Gleichgewicht (Schritt~3) & Verträglichkeit \\
  als Energie      & kleinste potentielle Energie & kleinste Ergänzungsenergie \\
  wie im           & WGV & KGV \\
  \bottomrule
\end{tabular}
\end{center}

Der Kern von $C$ sind die Bewegungen, die die starren Teile frei lassen. Der Kern von $C^T$ besteht
aus den Eigenspannungszuständen $\mu_i$, seine Dimension ist der Grad der statischen
Unbestimmtheit der starren Teile untereinander, also die Anzahl der $t_i$ (im KGV der $X_i$).

Die rote Box \emph{Starrmodus \& statisch unbestimmt} kommt übrigens schon dann, wenn das System
statisch unbestimmt ist und alle Stäbe $EA, EI \ge 10^7$ haben, also auch mit den
Standardwerten.

%---------------------------------------------------------------------
\section{Das Rand-LGS: die Biegelinie wie von Hand}\label{sec:mathe-randlgs}

Der Export \emph{Neue Biegelinie} rechnet unabhängig vom Weggrößenverfahren genau das, was ihr in
\textbf{TM2} von Hand macht, nur für jedes System und mit allen Bedingungen auf einmal.

\subsection{Ansatz je Stab}

Mit $\varphi = -w'$, $Q' = -q$, $M' = Q - m$ und $N' = -n$ setzt das Programm für jeden Stab an

\[
  EI\,w(x) = w_p(x) + c_1\frac{x^3}{6} + c_2\frac{x^2}{2} + c_3 x + c_4, \qquad
  EA\,u(x) = u_p(x) + C_1 x + C_2 .
\]

Das sind sechs Unbekannte je Stab. Die Partikulärlösungen baut es nicht aus einer Tabelle,
sondern als Integral, das der CAS-Kern exakt auswertet:

\begin{align*}
  w_p(x) &= \int_0^x q(t)\,\frac{(x-t)^3}{6}\,dt + \int_0^x m(t)\,\frac{(x-t)^2}{2}\,dt
            - EI\,\kappa_T\,\frac{x^2}{2}, \\
  u_p(x) &= -\int_0^x n(t)\,(x-t)\,dt + EA\,\varepsilon_T\,x .
\end{align*}

Die Idee dahinter: Man denkt sich die Streckenlast in viele kleine Einzellasten $q(t)\,dt$ an den
Stellen $t$ zerlegt. Bei einer Einzellast $1$ an der Stelle $t$ springt die Querkraft dort um $-1$
(wegen $Q' = -q$), und ab $t$ kommt der Anteil $EI\,w = (x-t)^3/6$ dazu. Das ist $c_1 x^3/6$ mit
$c_1 = 1$, nur um $t$ verschoben, davor ist der Anteil null. Weil alles linear ist, addieren sich
die Beiträge aller kleinen Einzellasten, und aus der Summe wird das Integral. Genauso gibt ein
Einzelmoment $1$ bei $t$ den Anteil $(x-t)^2/2$ in $EI\,w$ und eine Längskraft $1$ den Anteil
$-(x-t)$ in $EA\,u$. In der Mathematik heißt so ein
Integral \emph{Faltung}. Deshalb funktioniert es für jede Lastfunktion, für Trapezlasten und für
die globalen Linienlasten (die vorher auf Quer- und Längsrichtung des Stabes zerlegt werden). Und deshalb heißt die Integrationsvariable \code{t}, und
ihr könnt \code{t} nicht als Symbol verwenden.

Aus dem Ansatz folgen die Größen an beiden Stabenden als \emph{lineare} Ausdrücke in den
Konstanten, am Anfang z.\,B.\ $w = c_4/EI$, $\varphi = -c_3/EI$, $M = -c_2$, $Q = -c_1$,
$u = C_2/EA$, $N = C_1$.

\subsection{Bedingungen an den Knoten}

An jedem Knoten schreibt das Programm die Bedingungen, die ihr auch aufschreiben würdet, nur
systematisch:

\begin{itemize}
  \item \textbf{Kinematik:} Alle Stabenden ohne N- bzw.\ Q-Gelenk haben dieselbe
        Knotenverschiebung $\vec U$, also $u_e = \vec U\cdot\vec t_e$ und
        $w_e = \vec U\cdot\vec n_e$, alle Enden ohne M-Gelenk dieselbe Verdrehung. $\vec U$ selbst
        wird durch das erste passende Stabende ausgedrückt, damit keine zusätzlichen Unbekannten
        entstehen.
  \item \textbf{Gelenke:} Ein M-Gelenk gibt $M_e = 0$, ein N-Gelenk $N_e = 0$, ein Q-Gelenk
        $Q_e = 0$.
  \item \textbf{Lager:} In einer gehaltenen Richtung $\vec U\cdot\vec e = 0$, in einer freien
        Richtung das Gleichgewicht der Stabendkräfte, Federn und Lasten in dieser Richtung. Die
        Richtungen sind $\vec\xi = (\cos\alpha, -\sin\alpha)$ und $\vec\eta = (\sin\alpha, \cos\alpha)$,
        so gehen gedrehte Lager ohne Sonderfall.
  \item \textbf{Drehung:} Einspannung $\varphi = 0$, sonst $\sum M_e - c_m\,\varphi + M = 0$.
\end{itemize}

Am Ende muss die Anzahl der Gleichungen genau $6 \cdot (\text{Stäbe})$ sein (plus zwei je Knoten,
an dem alle Stäbe über N- oder Q-Gelenke hängen, siehe unten). Sonst kommt \emph{nicht quadratisch}.

\textbf{Warum solche Knoten eigene Verschiebungen brauchen.} Normalerweise wird $\vec U$ durch ein
Stabende ohne N- und Q-Gelenk ausgedrückt: Dessen $u_e$ und $w_e$ sind die beiden Komponenten von
$\vec U$ in Stab- und Querrichtung, damit ist $\vec U$ ganz bekannt. Ein Stabende mit N-Gelenk darf
sich aber längs gegen den Knoten verschieben, es liefert nur noch $w_e = \vec U\cdot\vec n_e$. Eines
mit Q-Gelenk liefert nur $u_e = \vec U\cdot\vec t_e$. Hat jedes Stabende am Knoten so ein Gelenk,
gibt es kein Ende, das beide Komponenten festlegt. Dann bekommt $\vec U$ die eigenen Unbekannten
$u_x(K)$, $u_y(K)$, und jedes Ende schreibt seine Bedingung gegen $\vec U$. Normalerweise
\enquote{verbraucht} das Bezugsende seine zwei Komponenten, um $\vec U$ festzulegen, und liefert
selbst keine Gleichung. Hier verbraucht keines etwas, daher die zwei Gleichungen mehr. Beispiel:
zwei Stäbe, die sich an einem Knoten mit je einem Q-Gelenk treffen. Jeder sagt nur, wie weit sich
der Knoten in seiner Längsrichtung bewegt, erst beide zusammen bestimmen $\vec U$.

\subsection{Lösen: erst wie von Hand, dann das CAS}

Ein LGS mit 30 symbolischen Unbekannten direkt in \code{simult} zu stecken, ist auf dem Rechner
langsam und liefert unübersichtliche Ausdrücke. Das Programm macht deshalb vorher, was ihr auch
machen würdet:

\begin{enumerate}
  \item Jede Gleichung mit nur \emph{einer} Unbekannten wird sofort aufgelöst und der Wert überall
        eingesetzt (z.\,B.\ $M = 0$ am Gelenk liefert direkt $c_2$). Das wiederholt sich, solange
        etwas passiert.
  \item Eine Unbekannte, die nur noch in einer einzigen Gleichung vorkommt, wird mit dieser
        Gleichung beiseitegelegt und am Schluss aus ihr berechnet.
  \item Gleichungen der Form $a\,x_j + b\,x_m = 0$ mit Zahlen $a$, $b$ werden zu $x_j = f\,x_m$
        und eingesetzt.
\end{enumerate}

Was übrig bleibt, zerfällt oft in unabhängige Blöcke (Unbekannte, die nie zusammen in einer
Gleichung stehen). Jeder Block wird einzeln mit \code{exact(simult(\dots))} gelöst, danach wird
rückwärts eingesetzt. Aus demselben Protokoll entsteht auch \code{randbed}: Eine Größe steht dort
nur dann als Wert, wenn ihre Konstanten in einem \emph{früheren} Schritt schon bekannt waren,
also genau in der Reihenfolge, in der ihr es auf dem Papier bekommen würdet.

\textbf{Die Aufteilung in \enquote{nur Biegelinie} und \enquote{nur Längslinie}} kommt aus
demselben Protokoll: Das Programm verfolgt für jede Konstante der Biegelinie zurück, aus welchen
Gleichungen sie bestimmt wurde. Tauchen darin keine Längslinien-Konstanten auf, die nicht schon
vorher feststanden, reichen die Bedingungen der Biegelinie allein.

Bei Rahmen mit dehnsteifen Stäben reicht das nicht. Dort bestimmt das LGS die Biegelinie zusammen
mit $C_1$ und $C_2$ der dehnsteifen Stäbe, weil z.\,B.\ die Seitenverschiebung $u$ des Riegels in
die Biegelinie der Stiele eingeht. Bei der direkten Methode ist $u = C_2$ und $N = C_1$ aber an
beiden Stabenden dieselbe Unbekannte (ohne Längslast und Temperatur). Das Programm nimmt deshalb
eine Bedingung, in der nur eine dieser Unbekannten steht, löst sie danach auf und setzt sie in die
anderen Bedingungen ein, genau wie von Hand. Mathematisch ist das ein Gauß-Schritt mit dieser
Unbekannten als Pivot, das LGS bleibt dabei regulär. Bleiben danach genau $4\cdot(\text{Stäbe})$
Bedingungen übrig, in denen nur noch Konstanten der Biegelinie stehen, genügen sie für $w$. Bei
dehnbaren Stäben ist $u(L) = (C_1 L + C_2)/EA$, dort hängt die Längsverschiebung am Stabende von
beiden Konstanten ab, und das Durchreichen geht nicht.

\subsection{Starre Stäbe im Rand-LGS: direkt oder Grenzwert}

Für einen starren Stab gibt es kein $EI\,w'' = -M$ mehr, trotzdem braucht das Programm sechs
Gleichungen. Zwei Wege:

\begin{itemize}
  \item \textbf{direkt} (Standard): Der biegesteife Stab bewegt sich als Starrkörper,
        $w(x) = c_4 + c_3 x - \kappa_T x^2/2$ und beim dehnsteifen $u(x) = C_2 + \varepsilon_T x$.
        $c_1$, $c_2$ und $C_1$ bleiben als reine Kraftgrößen stehen und werden aus dem
        Gleichgewicht bestimmt.

        Warum das geht: Auch beim elastischen Stab sind $c_1$, $c_2$ und $C_1$ eigentlich
        Kraftgrößen, nämlich die Schnittgrößen am Stabanfang, $Q(0) = -c_1$, $M(0) = -c_2$,
        $N = C_1$. In der Biegelinie stehen sie nur deshalb, weil die Krümmung aus dem Moment kommt:
        Sie gehen als $c_1 x^3/(6EI)$ und $c_2 x^2/(2EI)$ in $w$ ein, also immer geteilt durch
        $EI$. $c_3$ und $c_4$ dagegen sind echte Weggrößen ($c_3/EI = w'(0)$, $c_4/EI = w(0)$).
        Wird der Stab starr, $EI \to \infty$, verschwindet bei endlichen Kräften der Beitrag von
        $c_1$ und $c_2$ aus der Biegelinie, und $c_3$, $c_4$ schreibt man gleich ohne $EI$
        ($c_4 = w(0)$, $c_3 = w'(0)$). Die Schnittgrößen bleiben aber da: $Q(x)$, $M(x)$ und $N(x)$
        folgen aus $c_1$, $c_2$, $C_1$ und den Lasten allein über das Gleichgewicht am Stabstück
        ($Q' = -q$, $M' = Q$, $N' = -n$), ganz ohne Steifigkeit.

        Die Bedingungen an den Knoten sind dieselben wie beim elastischen Stab. Jede enthält aber
        entweder nur Kraftgrößen ($c_1$, $c_2$, $C_1$) oder nur Weggrößen ($c_3$, $c_4$, $C_2$) des
        starren Stabes: Kinematik und Lager mit gehaltener Richtung enthalten Wege, Gelenk-,
        Gleichgewichts- und Drehbedingungen Kräfte. Beispiel: starrer Kragarm, eingespannt am
        Anfang, Kraft $F$ (in $+w$) am freien Ende. Die Einspannung gibt $w(0) = c_4 = 0$ und
        $w'(0) = c_3 = 0$, also bewegt sich nichts. Das freie Ende gibt $Q(L) = -c_1 = F$ und
        $M(L) = -c_2 - c_1 L = 0$, also $c_1 = -F$, $c_2 = F L$ und $M(0) = -F L$. Das ist genau die
        Rechnung des statisch bestimmten Kragarms von Hand. In \code{randbed} heißen die
        Kraftgrößen eines starren Stabes deshalb $Q_i$, $M_i$, $N_i$ statt $-EI\,w_i'''$ usw.
  \item \textbf{Grenzwert}: Jeder starre Stab bekommt $EI = 1/\varepsilon$ bzw.\
        $EA = 1/\varepsilon$ (im CAS heißt $\varepsilon$ \code{rleps}), das System wird symbolisch in
        $\varepsilon$ gelöst und danach \code{limit(\dots, rleps, 0)} gebildet.
\end{itemize}

Warum überhaupt der Grenzwert? Das direkte Verfahren wird singulär, wenn sich starre Teile
\emph{gegenseitig} statisch unbestimmt halten. Beispiel: zwei dehnstarre Stäbe hintereinander
zwischen zwei Festlagern, Horizontallast in der Mitte. Beide Stäbe ändern ihre Länge nicht, also
ist $u = 0$ überall. Aber wie sich die Last auf die beiden Stäbe verteilt, sagt euch das
Gleichgewicht nicht, nur $N_1 - N_2 = F$. Mit $EA = 1/\varepsilon$ ist das System regulär, und
der Grenzwert $\varepsilon \to 0$ liefert die Aufteilung, die zwei \enquote{gleich starren} Stäben
entspricht. Genau in diesem Fall wechselt das Programm von selbst auf den Grenzwert. In
\code{randinfo} steht dann \emph{direkt singulär, starre Teile gegenseitig statisch unbestimmt}.

\subsection{Knotenverschiebungen aus den Randwerten}\label{sec:mathe-knotenverschiebung}

Für die Knotenverschiebung (Knotenmenü Seite~4) muss nichts Neues gelöst werden. Die
Kinematik-Bedingungen drücken $\vec U$ ja schon durch ein Referenz-Stabende aus, und zwar durch das erste
Stabende am Knoten ohne N- und Q-Gelenk. Dessen Längs- und Querverschiebung stehen nach dem Lösen
fest, $u(0)$ und $w(0)$ am Stabanfang bzw.\ $u(L)$ und $w(L)$ am Stabende. Zurück ins globale
System:
\[
  U_x = u\,t_x + w\,n_x, \qquad U_y = u\,t_y + w\,n_y, \qquad
  \vec t = \frac{(\Delta x, \Delta y)}{L}, \quad \vec n = \frac{(-\Delta y, \Delta x)}{L} .
\]
Das sind dieselben $\vec t$ und $\vec n$ wie in den Bedingungen, deshalb zählt $U_y$ wie die
$y$-Achse des Programms nach unten. Die Verdrehung kommt vom ersten Stabende ohne M-Gelenk,
\[
  \varphi = -w' = -\frac{w_p'(x) + c_1 x^2/2 + c_2 x + c_3}{EI},
\]
beim biegesteifen Stab $\varphi = \kappa_T x - c_3$. Hängen alle Stäbe über N- oder Q-Gelenke am
Knoten, hat das LGS für ihn ohnehin die eigenen Unbekannten $u_x(K)$, $u_y(K)$, die werden direkt
genommen.

Damit das beim zweiten Knoten nicht wieder ein ganzes LGS kostet, hebt das Programm die Linien des
letzten Rand-LGS mit einem Schlüssel auf: dem Fingerabdruck aller Eingaben (derselbe, an dem die
Hinweisbox erkennt, ob sich etwas geändert hat. Er enthält auch die Knoten- und Stabnummern) und
den Schaltern, die das Rand-LGS sonst beeinflussen (Methode für starre Stäbe, symbolischer Modus,
Zahlenformat). Passt der Schlüssel, werden nur die Randwerte eingesetzt, das sind drei kurze
CAS-Aufrufe statt eines LGS. Gleicher Fingerabdruck heißt wirklich gleiches System, auch wenn ein
Knoten zwischendurch gelöscht und genau so wieder eingebaut wurde. Werden dabei Nummern
vertauscht, ändert sich der Fingerabdruck, und das ist auch richtig so: Dann gehören die Namen
\code{wv1}, \code{wv2}, \dots\ zu anderen Stäben.

%---------------------------------------------------------------------
\section{Symbolischer Modus: rechnen mit Platzhaltern}\label{sec:mathe-symbolisch}

Die Finite-Elemente-Rechnung für die Diagramme ist eine reine Zahlenrechnung. Symbolisch
invertiert der Rechner keine Steifigkeitsmatrix in vernünftiger Zeit. Der Trick:

\begin{enumerate}
  \item Jedes Symbol bekommt einen \textbf{Platzhalterwert}, standardmäßig~1. Die Länge $l$
        bekommt, wenn die Geometrie kein eigenes Symbol hat, den Rasterabstand.
  \item \textbf{Superposition:} Alle Lasten werden auf null gesetzt und dann \emph{einzeln} wieder
        eingeschaltet. Für jede Last rechnet das Programm die Stabendkräfte und teilt sie durch den
        Platzhalter der Last und die passende Potenz von $l$. Übrig bleibt ein reiner
        Zahlenfaktor~$c$:
        \[
          N = c\,F, \quad M = c\,F\,l \quad\text{(Einzellast)}, \qquad
          N = c\,q\,l, \quad M = c\,q\,l^2 \quad\text{(Streckenlast)} .
        \]
        Die Steifigkeitsmatrix wird dafür nur einmal zerlegt, jede Last kostet dann nur noch ein
        Vorwärts- und Rückwärtseinsetzen.
  \item Der Zahlenfaktor wird \textbf{zurück in einen exakten Ausdruck} verwandelt (siehe unten)
        und alle Anteile werden addiert.
\end{enumerate}

Das funktioniert, weil die Rechnung linear ist und alle Längen Vielfache von $l$ sind. Genau hier
kommen die Einschränkungen her: Eine Temperatur passt nicht in das Schema \enquote{Last mal
Länge hoch $k$} (deshalb fehlt sie in den Diagrammwerten; die Exporte \emph{Neue Biegelinie} und
\emph{Neue Längslinie}, also \code{wneu}, \code{uneu}, \code{wv}, \code{uv}, haben sie, weil das
Rand-LGS gar nicht mit Superposition rechnet), und ein Zahlen-$EI$
neben einem symbolischen $EI$ macht die Steifigkeitsverhältnisse vom Platzhalterwert abhängig
(deshalb der gelbe Hinweis). Die Stelle des Maximums in einem Feld sucht das Programm numerisch
und rastet sie auf typische Werte wie $1/\sqrt 3$ oder einfache Brüche ein.

\textbf{Aus einer Zahl wird wieder ein Bruch.} Aus $0{,}708333\ldots$ soll wieder
$\frac{17}{24}$ werden. Dafür entwickelt das Programm die Zahl in einen \emph{Kettenbruch}
$a_0 + \cfrac{1}{a_1 + \cfrac{1}{a_2 + \dots}}$ und nimmt den ersten Näherungsbruch, der bis
auf $10^{-12}$ stimmt (Nenner höchstens 10\,000). Klappt das nicht und steht das Zahlenformat auf
\emph{Wurzel/Bruch}, versucht es dasselbe mit dem \emph{Quadrat} der Zahl und zieht aus Zähler und
Nenner die Quadratfaktoren heraus. So wird aus $0{,}7071\ldots$ wieder $\frac{\sqrt 2}{2}$ und
aus $86{,}60\ldots$ ein $50\sqrt 3$. Scheitert beides, sucht es zuletzt noch einen Bruch mit
kleinem Nenner (höchstens 100), der nur bis auf die eingestellte Toleranz $10^{-x}$ stimmen muss
(Obermenü Seite~3, \emph{CAS Bruch-Tol.}). Das ist die einzige Stelle, an der diese Einstellung
wirkt, und gerechnet wird dabei in Lua, nicht im CAS. Erst wenn auch das nicht klappt, bleibt eine
Dezimalzahl stehen. Im Format \emph{Dezimal} fallen alle drei Versuche weg.

%---------------------------------------------------------------------
\section{Kraftgrößenverfahren: Hauptsystem und \texorpdfstring{$\delta$}{delta}-Zahlen}\label{sec:mathe-kgv}

\textbf{Das Hauptsystem} sucht das Programm \enquote{gierig}: Es löst der Reihe nach
Einspannmomente, Lagerkräfte, Momente an biegesteifen Stabanschlüssen und zuletzt unbelastete
Pendelstäbe. Es behält jede Lösung nur, wenn das System danach noch tragfähig ist (die
Steifigkeitsmatrix lässt sich noch zerlegen). Sobald nichts mehr geht, ist das System statisch
bestimmt, und die Anzahl der gelösten Bindungen ist der Grad $n$.

\textbf{Die $\delta$-Zahlen} integriert das Programm nicht über $\int M_i M_j/(EI)\,dx$ (so rechnet
ihr sie in \textbf{Statik} mit dem Prinzip der virtuellen Kräfte), sondern es
rechnet für jeden Zustand ($X_j = 1$ bzw.\ Lastzustand) die Verformung des Hauptsystems mit dem
Weggrößenverfahren und liest die zu $X_i$ gehörende Klaffung ab: Verschiebung am gelösten Lager,
gegenseitige Verdrehung am eingefügten Gelenk, Längenänderung am gelösten Pendelstab. Nach
Maxwell--Betti ist das dasselbe wie das Arbeitsintegral, inklusive der Anteile aus Normalkraft,
Federn und endlichem $EA$, die man von Hand gern vergisst. Ein gelöster Pendelstab ist im
Hauptsystem ganz entfernt. Seine eigene Längenänderung $L/(EA)$ fehlt deshalb in den
Knotenverschiebungen und wird zu $\delta_{jj}$ addiert. Bei gedrehten Lagern liegen die
Knotenfreiheitsgrade schon in Lagerrichtung ($u_\xi$, $u_\eta$), die Klaffung eines gelösten Lagers
ist genau dieser Freiheitsgrad, und die Einheitskräfte eines gelösten Pendelstabes werden in
diese Richtungen zerlegt.

%---------------------------------------------------------------------
\section{Polplan, Verschiebungszustand und Arbeitssatz}\label{sec:mathe-polplan}

\subsection{Scheiben finden}

Zuerst muss das Programm wissen, welche Stäbe zusammen eine starre Scheibe bilden. Das passiert
in zwei Durchgängen:

\begin{enumerate}
  \item \textbf{Ohne Lager:} An jedem Knoten ohne Vollgelenk fasst das Programm alle Stäbe
        zusammen, deren Enden dort kein Gelenk haben (\emph{Union-Find}: jede Gruppe hat einen
        Vertreter, zwei Gruppen werden vereinigt, indem man einen Vertreter unter den anderen
        hängt). Mit diesen Scheiben und nur den inneren Gelenken läuft der Polplan einmal durch.
        Findet er zwei Scheiben, die gegeneinander unverschieblich sind, z.\,B.\ die drei Stäbe
        eines Dreigelenk-Dreiecks, werden sie zu einer Scheibe zusammengefasst.
  \item \textbf{Mit Lagern:} Stäbe, die biegesteif an einer Einspannung hängen (oder an einem
        Festlager mit Drehfeder), gehören zur Scheibe~0, dem Boden. Mit diesen zusammengefassten
        Scheiben und den Lagern entsteht der Polplan, den ihr seht.
\end{enumerate}

\subsection{Pole und Polstrahlen}

An jedem Knoten schaut das Programm, welche Scheiben sich dort treffen und welche Relativbewegung
zwischen ihnen möglich ist:

\begin{longtable}{p{0.34\linewidth}p{0.60\linewidth}}
  \toprule
  \textbf{Verbindung} & \textbf{Folge im Polplan} \\
  \midrule \endhead
  M-Gelenk, Festlager & Pol im Knoten (Drehung möglich, keine Verschiebung) \\
  Rollenlager & Polstrahl durch den Knoten senkrecht zur Rollrichtung, also entlang der
    Lagerkraft \\
  N- oder Q-Gelenk & Pol im Unendlichen (reine Verschiebung längs bzw.\ quer zum Stab) \\
  Einspannung, biegesteifer Anschluss & kein Pol: die Scheiben sind starr verbunden \\
  \bottomrule
\end{longtable}

Danach wiederholt das Programm (höchstens zehnmal) zwei Schritte, bis nichts Neues mehr dazukommt:

\begin{enumerate}
  \item \textbf{Satz von den drei Polen:} Für je drei Scheiben $i$, $j$, $k$ liegen die Pole
        $(i,j)$, $(j,k)$ und $(i,k)$ auf einer Geraden. Sind zwei davon bekannt, entsteht ein
        Polstrahl für den dritten. Liegt einer der beiden im Unendlichen, geht der Strahl durch
        den anderen in dessen Richtung.
  \item \textbf{Strahlen schneiden:} Zwei Polstrahlen derselben Scheibenpaarung schneiden sich im
        Pol. Sind sie parallel, liegt der Pol im Unendlichen in dieser Richtung.
\end{enumerate}

Findet es für eine Paarung zwei verschiedene Pole, oder liegt ein Pol nicht auf einem Strahl, auf
dem er liegen müsste, sind die beiden Scheiben gegeneinander unverschieblich (im Polplan rot). Im
Durchgang ohne Lager werden solche Paare zu einer Scheibe zusammengefasst (siehe oben).

\subsection{Der Verschiebungszustand}

Eine Scheibe, die sich um ihren Hauptpol $(c_x, c_y)$ um den kleinen Winkel $\varphi$ dreht,
verschiebt jeden Punkt um $\vec r_\perp \varphi$, also senkrecht zum Polstrahl vom Pol zu diesem
Punkt. Das Programm gibt einer Startscheibe $\varphi = 1$ (bevorzugt eine, deren Hauptpol auf einem
Knoten liegt) und reicht die Bewegung über die Nebenpole weiter: Der Nebenpol $(A,B)$ ist der Punkt,
der auf beiden Scheiben dieselbe Verschiebung hat. Kennt man die Verschiebung dieses Punktes aus
Scheibe~$A$ und den Hauptpol von Scheibe~$B$, folgt $\varphi_B$. Hat eine Scheibe ihren Hauptpol im
Unendlichen, verschiebt sie sich nur.

\textbf{Verschobene Scheiben unterwegs.} Die Kette reißt an einer Scheibe, die sich nur verschiebt,
nicht ab. Bei ihr bewegt sich jeder Punkt um denselben Vektor $\vec v$, also auch der Nebenpol zur
nächsten Scheibe~$B$. Damit ist die Verschiebung des Nebenpols bekannt, und es geht weiter wie bei
einer drehenden Scheibe: Liegt der Hauptpol von $B$ im Endlichen, folgt $\varphi_B$ aus $\vec v$,
liegt er im Unendlichen, verschiebt sich auch $B$ nur. Liegt der Nebenpol $(A,B)$ selbst im
Unendlichen, drehen beide Scheiben gleich, $\varphi_B = \varphi_A = 0$: $B$ verschiebt sich dann auch
nur, oder sie steht still, wenn ihr Hauptpol im Endlichen liegt. Eine reine Verschiebung ist nur der
Grenzfall einer Drehung mit dem Pol im Unendlichen.

\textbf{Nebenpol gleich Hauptpol.} Der Hauptpol einer Scheibe ist der Punkt, der sich auf ihr nicht
bewegt. Fällt der Nebenpol $(A,B)$ mit dem Hauptpol von $B$ zusammen, bewegt sich dieser Punkt also
auf $B$ nicht, und weil er auf $A$ dieselbe Verschiebung haben muss, auch auf $A$ nicht. Kommt aus
der Startscheibe $A$ trotzdem eine Verschiebung ungleich null an, passt das nicht zusammen: Der Punkt
kann nicht zugleich stehen und sich bewegen. Dann kann sich $A$ gar nicht bewegen, die Startscheibe
war falsch gewählt, und das Programm nimmt die nächste.

Ein Beispiel ist euer Dreigelenktragwerk mit dem Gelenk $G$ senkrecht über dem linken Lager $A$,
bei dem ihr das linke Festlager zum Loslager macht (verschieblich in $x$). Das rechte Festlager $B$
ist der Hauptpol $(\mathrm{II})$ der rechten Scheibe, das Gelenk der Nebenpol $(\mathrm{I,II})$. Der
Hauptpol $(\mathrm{I})$ der linken Scheibe liegt auf der Senkrechten durch $A$ (Loslager) und nach
dem Satz von den drei Polen auf der Geraden durch $(\mathrm{II})$ und $(\mathrm{I,II})$, also auf
$BG$. Beide Geraden schneiden sich in $G$, damit ist $G$ zugleich $(\mathrm{I})$ und
$(\mathrm{I,II})$. Startet das Programm mit der rechten Scheibe, dreht sie um $B$ und schiebt $G$
seitlich weg, aber auf der linken Scheibe steht $G$ als Hauptpol still: Widerspruch, die rechte
Scheibe kann sich nicht bewegen. Mit der linken Scheibe als Start geht es: Sie dreht um $G$, ihr
Lager $A$ wandert waagerecht (senkrecht zum Polstrahl $GA$, und das lässt das Loslager zu). $G$
selbst bewegt sich dabei nicht. Die rechte Scheibe hat damit zwei Punkte, die stehen bleiben ($B$
und $G$), und bewegt sich deshalb gar nicht, $\varphi_\mathrm{II} = 0$. Das ist kein Fehler,
sondern die Kinematik dieses Systems: Nur die linke Scheibe kippt um das Gelenk.

\subsection{Der Arbeitssatz}\label{sec:mathe-arbeitssatz}

Nach dem Lösen einer Bindung ist das System eine kinematische Kette mit genau einem Freiheitsgrad.
Für jede virtuelle Bewegung, die die übrigen Bindungen zulassen, gilt
\[
  \delta W = \sum_k \vec F_k \cdot \delta\vec u_k + \sum_m M_m\,\delta\varphi_m = 0 ,
\]
und weil es nur einen Freiheitsgrad gibt, sind alle Wege Vielfache eines einzigen Parameters. Der
Export (Abschnitt~\ref{sec:arbeitssatz}) macht daraus drei Schritte.

\textbf{Weg einer Kraft (\code{pvv\_kin}, Spalte~2).} Eine Kraft mit Richtung
$\vec e = (e_x, e_y)$ ($y$ nach unten) greift im Punkt $P$ auf Scheibe~$i$ an. Ihr Weg ist die
Verschiebung von $P$ in Richtung $\vec e$:
\begin{itemize}
  \item Scheibe~$i$ dreht um ihren Hauptpol $(c_x, c_y)$ um $\varphi_i$ (gegen den Uhrzeigersinn).
        Jeder Punkt bewegt sich senkrecht zum Polstrahl vom Pol zu ihm, projiziert auf $\vec e$ ist das
        \[
          \delta_e = \varphi_i\,\bigl[e_x\,(P_y - c_y) - e_y\,(P_x - c_x)\bigr] .
        \]
        Die eckige Klammer ist der Hebel der Kraft um den Pol.
  \item Scheibe~$i$ verschiebt sich nur (Hauptpol im Unendlichen) um $\delta_i$ in Richtung
        $\vec t_i$: $\delta_e = \delta_i\,(\vec e\cdot\vec t_i)$.
  \item Ein Moment auf Scheibe~$i$ arbeitet am Drehwinkel: $\delta\varphi = \varphi_i$.
\end{itemize}
Verschieben \emph{und} drehen zugleich gibt es nicht als eigenen Fall: Eine kleine ebene
Starrkörperbewegung ist immer eine Drehung um einen Punkt, eben den Pol. Eine Verschiebung plus
Drehung ist nur eine Drehung um einen anderen Pol, die reine Verschiebung der Grenzfall mit dem Pol
im Unendlichen.
Streckenlasten dürfen exakt durch ihre Resultierende $R$ ersetzt werden. Auf einem geraden Stab
der Länge $L$, der zu einer starren Scheibe gehört, ändert sich die virtuelle Verschiebung linear
mit der Stabkoordinate $x$, denn ein Starrkörper verschiebt jeden Punkt linear mit seinem Ort:
$\delta\vec u(x) = \delta\vec u_A + \frac{x}{L}\,(\delta\vec u_B - \delta\vec u_A)$. Damit ist die
Arbeit der Streckenlast
\[
  \int_0^L q(x)\,\vec e\cdot\delta\vec u(x)\,dx
  = \vec e\cdot\delta\vec u_A \underbrace{\int_0^L q\,dx}_{R}
    + \vec e\cdot\frac{\delta\vec u_B - \delta\vec u_A}{L} \underbrace{\int_0^L q\,x\,dx}_{S = R\,x_R}
  = R\;\vec e\cdot\delta\vec u(x_R) ,
\]
also die Arbeit der Resultierenden $R$ an ihrem Angriffspunkt $x_R = S/R$. Auf einem Bogen ist die
Verschiebung nicht linear in der Bogenlänge, das Ergebnis gilt aber genauso: Dort wird die Last
durch ihre Resultierende (Kraft und Richtung) auf ihrer Wirkungslinie ersetzt. Die freigesetzte
Größe ist eine äußere Last in der gezeichneten Richtung, bei einem gelösten Gelenk arbeitet sie an
der gegenseitigen Verdrehung bzw.\ Verschiebung der beiden Seiten.

\textbf{Verhältnis der Scheiben (\code{pvv\_pol}).} Der Nebenpol $P_{AB}$ ist der Punkt, der auf
Scheibe~$A$ und Scheibe~$B$ dieselbe Verschiebung hat. Schreibt man diese Verschiebung mit der
Formel oben einmal über $A$ und einmal über $B$ hin, entsteht eine lineare Beziehung
$\varphi_B = \kappa_{AB}\,\varphi_A$. Die Verschiebung ist ein Vektor mit $x$- und $y$-Komponente,
das sind eigentlich zwei Gleichungen. Beide liefern aber dasselbe $\kappa_{AB}$: Über $A$ steht die
Verschiebung senkrecht auf dem Polstrahl von $(A)$ zum Nebenpol, über $B$ senkrecht auf dem von
$(B)$, und weil $(A)$, $(A,B)$ und $(B)$ auf einer Geraden liegen, sind beide Richtungen gleich.
Deshalb reicht eine Komponente. Das Programm nimmt die betragsmäßig größere, damit es nicht durch
eine Zahl nahe null teilt. Liegt der Nebenpol zweier drehender Scheiben im Unendlichen, drehen beide
gleich, $\varphi_A = \varphi_B$.

Von der Bezugsscheibe aus (der Startscheibe des Verschiebungszustands, ihr Winkel heißt
$\varphi_0$) laufen diese Beziehungen über die Nebenpole durch alle Scheiben und liefern
$\varphi_i = \kappa_i\,\varphi_0$. Das sind die unteren Zeilen von \code{pvv\_kin}. Gibt es mehr
Nebenpole als nötig, sind die übrigen Proben. Als Bezug nimmt das Programm eine Scheibe, die sich
dreht, wenn es eine gibt. Eine Scheibe, die stehen bleibt, bekommt einfach $\kappa_i = 0$, ihre
Lasten leisten keine Arbeit. Verschieben sich alle bewegten Scheiben nur (alle ihre Hauptpole im
Unendlichen), gibt es keinen Drehwinkel als Bezug. Dann ist die Bezugsgröße die Verschiebung $\delta_0$ der Startscheibe, und
alle anderen Bewegungen werden als Vielfache von $\delta_0$ geschrieben (siehe unten).

\textbf{Einsetzen (\code{w\_eq2}).} Mit $\varphi_i = \kappa_i\,\varphi_0$ wird jeder Weg zu
$h_k\,\kappa_{i(k)}\,\varphi_0$ (Spalte~3 von \code{pvv\_kin}). Dabei ist $i(k)$ die Scheibe, an der
die Kraft $F_k$ angreift, und $h_k$ ihr Hebel um den Pol dieser Scheibe, also die eckige Klammer
von oben, $h_k = e_x\,(P_y - c_y) - e_y\,(P_x - c_x)$. Das ist der senkrechte Abstand der
Wirkungslinie vom Pol mit Vorzeichen, nicht der Abstand des Angriffspunktes. Verschiebt sich die
Scheibe nur, tritt an die Stelle von $h_k\,\varphi_i$ der Weg $\delta_i\,(\vec e\cdot\vec t_i)$ aus
der Liste oben, $h_k$ ist dann der Richtungskosinus $\vec e\cdot\vec t_i$. Der Arbeitssatz zerfällt in
\[
  \delta W = \varphi_0 \Bigl(\sum_k F_k\,h_k\,\kappa_{i(k)} + \sum_m M_m\,\kappa_{i(m)}\Bigr) = 0 .
\]
Da $\varphi_0$ beliebig ist, muss die Klammer verschwinden. Sie ist linear in der freigesetzten
Größe, also eine Gleichung mit einer Unbekannten. Verschiebt sich die Bezugsscheibe nur, steht
$\delta_0$ statt $\varphi_0$ da, und die Faktoren der Drehwinkel haben die Einheit $1/$Länge.
Intern zählt das Programm Drehwinkel im Uhrzeigersinn. Im Export ist das Vorzeichen umgedreht,
damit alles zur Konvention der Vorlesung (gegen den Uhrzeigersinn positiv) passt.

%---------------------------------------------------------------------
\section{Gleichgewicht der Teilsysteme (LGS-Matrix)}\label{sec:mathe-gleichgewicht}

Für den Export \emph{LGS-Matrix} schneidet das Programm das System an allen Gelenken auf, genau
so, wie Seite~2 der Explosionszeichnung es zeigt. Beide benutzen dieselben drei Regeln im Code:
\code{gelenkKnoten} (welche Stäbe hängen mit Gelenk, welche biegesteif am Knoten, wie viele
Teilsysteme treffen sich), \code{gelenkRichtung} (positive Richtung der Gelenkgrößen) und
\code{lastStabKnoten} (an welchem Teilsystem Lasten und Lager eines Gelenkknotens hängen). Jedes
Teilsystem liefert drei Gleichungen: Kräfte nach rechts, nach oben und Momente gegen den
Uhrzeigersinn. Unbekannt sind die Lagerreaktionen und die Größen, die ein Gelenk überträgt: beim
M-Gelenk die beiden Kraftkomponenten $Gx$, $Gy$, beim N-Gelenk $Gq$ und $Gm$, beim Q-Gelenk $Gn$
und $Gm$. Treffen sich mehr als zwei Teilsysteme in einem Gelenk, bekommt jeder Gelenkstab eigene
Gelenkkräfte, gibt es dort keinen biegesteifen Rest, sagen zwei zusätzliche Gleichungen, dass der
Gelenkbolzen selbst im Gleichgewicht ist. Bei nur zwei Teilsystemen braucht der Bolzen keine
eigenen Gleichungen: Die Gelenkkraft steht in beiden mit entgegengesetztem Vorzeichen (actio =
reactio), damit ist sein Gleichgewicht schon erfüllt. Hängt am Gelenk ein biegesteifer Rest, gehört
der Bolzen zu ihm. Ein Beispiel ist eine Stütze, die durch den Knoten durchläuft, und zwei Riegel,
die gelenkig an ihr hängen. Die Gelenkkräfte der Riegel greifen dann direkt an der Stütze an und
stehen in deren drei Gleichungen. Der Bolzen ist kein eigener Körper, also gibt es für ihn keine
eigenen Gleichungen, sein Gleichgewicht steckt im Gleichgewicht der Stütze. Erst wenn alle Stäbe
gelenkig am Knoten hängen, gehört der Bolzen zu keinem Teilsystem und braucht seine zwei
Gleichungen selbst. Im symbolischen Modus stehen Lasten mit ihren Namen und
Hebelarme als Vielfache von $l$ in der Matrix.

Den \textbf{Momentenbezugspunkt} wählt das Programm so, wie ihr es in \textbf{TM1} lernt: Es
sucht den Punkt, durch den möglichst viele unbekannte Kräfte gehen, also die senkrechte Linie mit den
meisten unbekannten Vertikalkräften und die waagerechte mit den meisten unbekannten
Horizontalkräften, am liebsten ein Knoten, der auf beiden liegt. Die Momentengleichung wird dadurch
so kurz wie möglich.

Die Werte auf Seite~1 der Explosionszeichnung kommen aus dem Weggrößenverfahren. Löst ihr das LGS
(mit den oben beschriebenen Befehlen oder von Hand), müssen dieselben Zahlen herauskommen. Das ist
eure Kontrolle.

\textbf{Lösungsweg (\code{lgs\_weg}):} Gesucht ist eine Reihenfolge, in der jede Gleichung möglichst
nur eine neue Unbekannte enthält. Dafür braucht das Programm mehr Gleichungen als die in \code{lgs}.
Alle sind Linearkombinationen der Zeilen aus \code{lgs}, also nichts Neues, nur anders
zusammengefasst.

\textbf{Anderer Momentenpunkt.} Eine Kraft mit den Komponenten $F_\rightarrow$ (nach rechts) und
$F_\uparrow$ (nach oben) im Punkt $(x, y)$ hat um $P$ das Moment
$M_P = (x - x_P)\,F_\uparrow - (y - y_P)\,F_\rightarrow$ (gegen den Uhrzeigersinn positiv, hier $y$
nach oben). Die Pfeile sind dieselben wie in den Namen der Kraftzeilen von \code{lgs}. Um einen
anderen Punkt $Q$ ist es
\[
  M_Q = (x - x_Q)\,F_\uparrow - (y - y_Q)\,F_\rightarrow
      = M_P + (x_P - x_Q)\,F_\uparrow - (y_P - y_Q)\,F_\rightarrow .
\]
Über alle Kräfte eines Teilsystems summiert (reine Momente hängen nicht vom Punkt ab) heißt das
\[
  \textstyle\sum M_Q = \sum M_P + (x_P - x_Q) \sum F_\uparrow - (y_P - y_Q) \sum F_\rightarrow .
\]
Die Momentengleichung um $Q$ ist also die Momentenzeile plus Vielfache der beiden Kraftzeilen
desselben Teilsystems, die Faktoren sind Hebelarme. Im Programm zeigt $y$ nach unten, dort steht
deshalb $+(y_P - y_Q)$.

\textbf{Teilsysteme zusammenfassen.} Für eine Gruppe werden die Kraftzeilen addiert und die
Momentenzeilen, jede vorher auf denselben Punkt $Q$ umgerechnet. Eine Gelenkkraft zwischen zwei
Teilsystemen der Gruppe steht in beiden mit entgegengesetztem Vorzeichen und greift am selben Punkt
an. Sie fällt deshalb aus den Kräften \emph{und} aus den Momenten heraus. Alle Teilsysteme zusammen
sind das Gesamtsystem.

\textbf{Kräfte quer zu einer Richtung.} Zeigt eine Unbekannte in Richtung $(d_x, d_y)$ (wieder mit $y$
nach unten), dann enthält $d_y \sum F_\rightarrow + d_x \sum F_\uparrow$ (normiert) diese Unbekannte
nicht mehr.

\textbf{Warum höchstens drei Gleichungen je Scheibe?} Jede Gleichgewichtsbedingung einer Scheibe ist
entweder eine Kraftsumme in eine Richtung $\vec e$, also $e_x \sum F_\rightarrow + e_y \sum
F_\uparrow$, oder eine Momentensumme um einen Punkt, nach der Formel oben $\sum M_P$ plus Vielfache
der beiden Kraftsummen. Alle liegen also in dem Raum, den $\sum F_\rightarrow$, $\sum F_\uparrow$ und
$\sum M_P$ aufspannen, und der hat höchstens die Dimension~3. Mehr als drei linear unabhängige
Gleichungen bekommt ihr deshalb nie, egal um wie viele Punkte ihr die Momente bildet. Drei
Momentengleichungen um $A$, $B$ und $C$ sind genau dann unabhängig, wenn die drei Punkte nicht auf
einer Geraden liegen. Liegen sie auf einer Geraden, ist die dritte eine Kombination der beiden
anderen. Ein Durchlaufträger auf drei Lagern bleibt deshalb einfach statisch unbestimmt, auch wenn
man um jedes Lager die Momente bildet. Anschaulich (PvV): Jede Gleichung gehört zu einer virtuellen
Starrkörperbewegung, einer Verschiebung in Richtung $\vec e$ bzw.\ einer Drehung um $P$, und eine
Scheibe hat in der Ebene nur drei Freiheitsgrade.

\textbf{Welche Gruppen und Punkte?} Das Programm sucht die Gleichungen nicht gezielt, sondern stellt
alle aus einem festen Vorrat auf und wählt danach aus:

\begin{itemize}
  \item \textbf{Gruppen:} jedes Teilsystem einzeln, das Gesamtsystem und bei Rahmen alle
        zusammenhängenden Gruppen. Zwei Teilsysteme hängen zusammen, wenn sie eine
        Gelenkunbekannte teilen. Ohne gemeinsame Gelenkkraft hebt sich beim Zusammenfassen nichts
        auf. Die Gruppen entstehen, indem an jede Gruppe ein Nachbar angehängt wird, bis die
        Zahl unter \emph{Max. Teilgruppen} erreicht ist (ab Werk 30). Im Fachwerk gibt es nur die Knoten und das
        Gesamtsystem, denn eine Gruppe von Knoten wäre ein Ritterschnitt.
  \item \textbf{Momentenpunkte $Q$:} Knoten und Schnittpunkte zweier Wirkungslinien von
        Unbekannten bzw.\ Pendelachsen (innerhalb der doppelten Systemgröße), aber je Gruppe nur
        die, an denen eine \emph{äußere} Kraft der Gruppe herausfällt. Äußere Kräfte sind die, die
        nach dem Zusammenfassen noch in der Kraftsumme stehen. Gelenkkräfte zwischen zwei
        Teilsystemen der Gruppe heben sich auf und zählen nicht. Ein Knoten zählt, wenn die
        Wirkungslinie einer äußeren Kraft durch ihn geht, ein Schnittpunkt, wenn beide Linien zu
        äußeren Kräften gehören, und er wird nach diesem Paar benannt. Für \code{TS 1+2} ist das
        Gelenk zwischen TS~1 und TS~2 deshalb nur dann ein Kandidat, wenn noch eine andere äußere
        Kraft durch es hindurchgeht. $Q$ wird also nicht vorher ausgesucht: Jeder dieser Punkte
        liefert eine Gleichung, und die Auswahl unten nimmt die, in der nur noch eine Unbekannte
        übrig ist.
  \item \textbf{Kräfte quer} zu jeder schrägen äußeren Unbekannten und jeder schrägen Pendelachse.
\end{itemize}

Gleichungen, die sich nur um einen Faktor unterscheiden, bleiben einmal übrig, und zwar in der
einfachsten Form. Weil dabei nur mit Zahlen gerechnet wird, dauert das auch bei einigen hundert
Kandidaten kaum Zeit.

\textbf{Benennung.} Der Name nennt die Art, den Punkt und die Gruppe: \code{$\Sigma$M P5 TS 1+2}
sind die Momente um Knoten~5 an den Teilsystemen 1 und 2 zusammen, \code{$\Sigma$M S(A\_h,B\_v)
Ges} die Momente des Gesamtsystems (\code{Ges}) um den Schnittpunkt der Wirkungslinien von $A_h$ und
$B_v$, \code{$\Sigma$F quer zu B\_h TS 2} die Kräfte an Teilsystem~2 quer zu $B_h$. Neue Gleichungen
bekommen die nächste freie Nummer nach denen aus \code{lgs}.

\textbf{Pendelstäbe.} Ein Teilsystem eines Rahmens ohne Lasten, an dem nur an zwei Punkten $U$ und
$V$ Anschlusskräfte angreifen (Lager- bzw.\ Gelenkkräfte, keine Momente), nennen wir hier
Pendelstab. Ist es ein gerader Stab, ist das der Pendelstab aus \textbf{TM1} mit reiner
Normalkraft. Ist es geknickt (z.\,B.\ ein unbelasteter Winkel zwischen zwei Gelenken), hat es zwar
Momente und Querkräfte, für das Gleichgewicht verhält es sich aber genauso: Die Anschlusskraft liegt
auf der Verbindungslinie $UV$. Nur darum geht es hier. Sein
Momentengleichgewicht um $U$ enthält nur die beiden Kraftkomponenten bei $V$ und sagt: Die Kraft bei
$V$ liegt auf der Achse $UV$. Liegt der Punkt $Q$ einer anderen Momentengleichung auf dieser Achse,
stehen die beiden Komponenten bei $V$ dort im selben Verhältnis wie im Pendelgleichgewicht, nur mit
einem Faktor $\lambda$. Zieht man $\lambda$-mal das Pendelgleichgewicht ab, fallen beide heraus. Das
Programm prüft das für jede Gleichung und jeden Pendelstab mit den Zahlen und nimmt die Achse als
weitere Wirkungslinie für die Schnittpunkte dazu. So findet sich auch beim Dreigelenkrahmen mit
verschieden hohen Lagern und einer unbelasteten Hälfte ein Punkt, der direkt eine Lagerkraft
liefert. Ist die Achse waagerecht oder senkrecht, braucht es den Umweg nicht, dann hat das
Pendelgleichgewicht selbst schon nur eine Unbekannte. Dieselbe Gleichung entsteht oft auf zwei
Wegen, etwa am Gesamtsystem (Lagerkräfte über das Pendelgleichgewicht um das Gelenkende
herausgerechnet) und an der anderen Scheibe allein (Gelenkkräfte über das Pendelgleichgewicht um
das Lager herausgerechnet). Das Programm behält dann den Namen mit dem Gesamtsystem, weil das der
Vorstellung entspricht, dass die Lagerkraft auf der Pendelachse wirkt.

\textbf{Auswahl.} Dann geht es schrittweise: Gibt es eine Gleichung mit genau einer noch unbekannten
Größe, ist das der nächste Schritt. Unter mehreren gewinnt die einfachste: erst die aus \code{lgs},
dann andere Momentenpunkte am selben Teilsystem, dann das Gesamtsystem, dann Zusammenfassungen und
Schnittpunkte, zuletzt Gleichungen, in denen ein Pendelgleichgewicht steckt. Bei Gleichstand gewinnt
die mit weniger Unbekannten, danach die, die zuerst aufgestellt wurde (Knoten vor Schnittpunkten,
nach Nummern). Lagergrößen (Lager und Federn) haben dabei Vorrang, aber nur innerhalb derselben
Stufe: Ein Einzelschritt für eine Lagergröße geht vor jedem anderen Einzelschritt. Gibt es keinen,
nimmt das Programm einen Einzelschritt, nach dem eine Lagergröße als Einzelschritt folgt (einen
Schritt vorausgeschaut), sonst den einfachsten. Ein Block kommt erst, wenn es gar keinen
Einzelschritt mehr gibt. Gibt es keinen, sucht das Programm zwei Gleichungen mit denselben zwei neuen
Unbekannten, die voneinander unabhängig sind (Rang~2), danach drei mit drei. Unter gleich großen
Blöcken gewinnt der mit mehr Lagergrößen, dann der aus den einfacheren Gleichungen. Dafür hat jede
Gleichung eine Stufe (\emph{Klasse}) nach der Reihenfolge oben: $0$ für eine Zeile aus \code{lgs},
höher für andere Momentenpunkte, das Gesamtsystem, Zusammenfassungen, Schnittpunkte und zuletzt
Gleichungen mit Pendelgleichgewicht. Verglichen wird die Summe der Stufen der Gleichungen im Block,
die kleinere gewinnt. Erst wenn auch das nicht geht, löst es alle übrigen Unbekannten als einen
Block aus Gleichungen von \code{lgs}. Ein solcher Block wird durch Einsetzen gelöst: Die Gleichung
mit den wenigsten Unbekannten wird nach einer davon aufgelöst und in die übrigen eingesetzt, das
wiederholt sich, bis die letzte Gleichung nur noch eine Unbekannte hat. Dann wird von hinten
zurückeingesetzt. Bei gleich vielen Unbekannten nimmt das Programm am liebsten eine Unbekannte, die
in ihrer Gleichung den Faktor $+1$ oder $-1$ hat, denn dann entsteht beim Auflösen kein Bruch.

Die Auswahl läuft mit Zahlen, denn die Faktoren vor den Unbekannten hängen nur von der Geometrie ab,
auch im symbolischen Modus. Die Texte und Ergebnisse in \code{lgs\_weg} rechnet danach das CAS
exakt aus, sie sind also dieselben wie die Lösung des ganzen LGS.

\textbf{Option LGS wie Lösungsweg.} Der Lösungsweg benutzt genau $n$ Gleichungen für $n$
Unbekannte. Ordnet man sie nach den Schritten, hat ihre Matrix Blockdreiecksform: Jeder Schritt
enthält außer schon bekannten Größen nur seine neuen Unbekannten, und deren Block ist regulär
(Einzelschritt: Faktor ungleich null, Block: voller Rang). Damit ist die ganze Matrix regulär. Das
Programm lässt die Gleichungen aus \code{lgs} stehen, die der Weg benutzt, und setzt jede
abgeleitete Gleichung auf den Platz einer Gleichung, die er nicht benutzt, am liebsten auf eine,
aus der sie mit entsteht. Danach prüft es den Rang des neuen Systems noch einmal mit Zahlen. Allgemein
gilt: Ersetzt man Zeile $r$ durch $\sum_i f_i\, r_i$, bleibt die Lösung eindeutig genau dann, wenn
$f_r \neq 0$ ist (die Determinante ändert sich um den Faktor $f_r$).

\textbf{Gegenprüfung mit der FEM.} Für jede Unbekannte bestimmt das Programm aus den FEM-Ergebnissen
den Wert mit demselben Namen und derselben Richtung wie auf Seite~2: Lagerreaktionen direkt,
Gelenkgrößen aus den Stabendkräften des angeschlossenen Stabes (ohne Knotenlast und Lagerkraft, wenn
die zu diesem Stab gehören), Stabkräfte im Fachwerk aus der Normalkraft. Das vergleicht es mit dem
Wert aus \code{lgs\_weg}, im symbolischen Modus mit denselben Ersatzwerten für die Symbole, mit denen
die FEM gerechnet hat. Weicht etwas um mehr als $10^{-4}$ des größten Wertes ab, kommt die Warnung.

\textbf{Resultierende von Streckenlasten:} $R = \int_0^L q\,dx$ und das statische Moment
$S = \int_0^L q\,x\,dx$ um den Stabanfang, die Wirkungslinie liegt bei $x_R = S/R$. Für Trapezlasten
sind das die bekannten Formeln, für Lastfunktionen integriert das Programm mit der Simpsonregel
über elf Stützstellen. Die ist für $q\,x$ bis zu einem quadratischen $q$ exakt.

%---------------------------------------------------------------------
\section{Querschnittswerte}\label{sec:mathe-qswerte}

\textbf{Vielecke} (Rechteck, Dreieck, Trapez, Polygon, dicke Linie) rechnet das Programm alle mit
derselben Formel, der Gaußschen Trapezformel. Mit den Eckpunkten $(y_i, z_i)$ gegen den
Uhrzeigersinn und $c_i = y_i z_{i+1} - y_{i+1} z_i$ ist
\[
  A = \tfrac12 \sum c_i, \qquad
  y_S = \frac{\sum (y_i + y_{i+1})\,c_i}{6A}, \qquad
  \int z^2\,dA = \tfrac1{12} \sum \left(z_i^2 + z_i z_{i+1} + z_{i+1}^2\right) c_i ,
\]
entsprechend für $z_S$, $\int y^2\,dA$ und $\int yz\,dA$. Das ist exakt, es ist der Gaußsche
Integralsatz, angewandt auf die Randkanten. Deshalb ist es auch egal, wie kompliziert das Vieleck
ist, solange es sich nicht selbst schneidet.

\textbf{Kreisausschnitt und -abschnitt} mit halbem Öffnungswinkel $\alpha$: Der Schwerpunkt liegt
auf der Winkelhalbierenden im Abstand $\frac{2R\sin\alpha}{3\alpha}$ (Ausschnitt) bzw.\
$\frac{4R\sin^3\alpha}{3(2\alpha - \sin 2\alpha)}$ (Abschnitt). Die Trägheitsmomente werden erst um
die Winkelhalbierende bestimmt (beim Abschnitt wird das Dreieck zum Mittelpunkt abgezogen) und dann
ins KOS gedreht.

\textbf{Dünnwandige Linien} sind Mittellinie mal Dicke: $A = L\,t$,
$I_y = t L^3 \sin^2\beta/12$ usw. Den Eigenanteil $L t^3/12$ lässt das Programm weg, genau wie
die Vorlesung, weil er von höherer Ordnung in $t$ ist. \textbf{Dünnwandige Kreise und Bögen}
ebenso: $A = R\,\Delta\varphi\,t$, $I_y = R^3 t \int \sin^2\varphi\,d\varphi$ um den Mittelpunkt
(Vollkreis $\pi R^3 t$), Schwerpunkt des Bogens im Abstand $R\,\frac{\sin(\Delta\varphi/2)}{\Delta\varphi/2}$
auf der Winkelhalbierenden.

\textbf{Löcher} gehen mit negativer Fläche in die Summen ein, danach setzt der Satz von Steiner alles
zusammen. \textbf{Numerisch addierte} Elemente legt das Programm auf ein Raster von $40 \times 40$
Zellen über ihr umschreibendes Rechteck und zählt nur die Zellen, die nicht schon in einem anderen
Element liegen. Deshalb die Warnung, das ist eine Näherung.

\textbf{Hauptachsen:} Mit $m = \frac{I_y + I_z}2$ und
$d = \sqrt{\left(\frac{I_y - I_z}2\right)^2 + I_{yz}^2}$ ist $I_{1,2} = m \pm d$ und
\[
  \alpha = \tfrac12 \operatorname{atan2}(2 I_{yz},\ I_y - I_z) .
\]
Das \code{atan2} (der Arkustangens, der das Vorzeichen von Zähler und Nenner kennt) sorgt dafür,
dass $\alpha$ immer zur Achse~1 mit $I_1$ gehört. Mit dem einfachen $\tan 2\alpha$ aus der
Formelsammlung müsstet ihr das von Hand prüfen, dort kann auch die Achse mit $I_2$ herauskommen.

\textbf{Widerstandsmomente} $W = I/e$ nehmen als Randabstand $e$ den größten Abstand der Elemente
vom Schwerpunkt. Bei dünnwandigen Elementen ist das die Mittellinie, $t/2$ kommt nicht dazu.

%---------------------------------------------------------------------
\section{Normalspannung und neutrale Faser}\label{sec:mathe-sigma}

Mit $\Delta = I_y I_z - I_{yz}^2$ und $y$, $z$ vom Schwerpunkt aus rechnet das Programm immer mit
\[
  \sigma_x = \frac NA + \frac{M_y I_z - M_z I_{yz}}{\Delta}\,z
                      - \frac{M_z I_y - M_y I_{yz}}{\Delta}\,y .
\]
Für $I_{yz} = 0$ bleibt das bekannte $\frac NA + \frac{M_y}{I_y} z - \frac{M_z}{I_z} y$ übrig. Die
Formel kommt direkt aus $M_y = \int \sigma z\,dA$ und $M_z = -\int \sigma y\,dA$ mit dem
linearen Ansatz $\sigma = a + b\,y + c\,z$, also zwei Gleichungen für $b$ und $c$.

Weil $\sigma_x$ linear in $y$ und $z$ ist, speichert das Programm nur den Gradienten $(c_y, c_z)$ und
den Wert $c_0$ im Ursprung des gezeigten KOS. Daraus entstehen beide Gleichungen in der Tabelle: die
KOS-Form direkt, die HAS-Form, indem der Gradient um den Hauptachsenwinkel gedreht wird. Die
neutrale Faser ist die Gerade $c_0 + c_y y + c_z z = 0$.

Die \textbf{Extremwerte} sucht das Programm nur an den Eckpunkten des Querschnitts (Bögen alle
$10^\circ$ abgetastet). Eine lineare Funktion hat ihr Maximum immer in einer Ecke der konvexen
Hülle. Bei dünnwandigen Profilen sind das die Endpunkte der Mittellinien, also wieder die
Näherung aus der Vorlesung.

%---------------------------------------------------------------------
\section{Schubfluss in dünnwandigen Profilen}\label{sec:mathe-schub}

\subsection{Offene Profile}

Aus den dünnwandigen Linien baut das Programm einen \emph{Profilgraph}: Die Linien werden an
Kreuzungen, Berührpunkten und an den Schwerachsen geteilt, Punkte, die zusammenfallen, werden ein
Knoten. \textbf{Kreise und Bögen} kommen als Kette kurzer Sehnen dazu, geteilt an ihren Enden, an
Endpunkten anderer Elemente auf dem Bogen, an Kreuzungen mit geraden Linien und an den Schwerachsen,
dazwischen gleichmäßig mit höchstens $15^\circ$ je Sehne. Die Sehnen bestimmen nur die Knoten und
den Zusammenhang (Zellen, Symmetrie). Jede Sehne kennt ihren Bogen (Mittelpunkt, Radius, Winkel
$\varphi_1$ bis $\varphi_2$), und gerechnet wird auf dem Bogen: Länge $R\,|\varphi_2 - \varphi_1|$
statt der Sehne, Stützstellen auf dem Bogen. Dann läuft es los:

\begin{itemize}
  \item An jedem \textbf{freien Ende} ist der Schubfluss null, dort beginnt ein Pfad mit
        $S_y = S_z = 0$.
  \item Entlang eines Segments werden die statischen Momente aufintegriert,
        $S_z \mathrel{+}= (y - y_S)\,t\,\Delta s$ und $S_y \mathrel{+}= (z - z_S)\,t\,\Delta s$. Weil
        der Integrand auf einer geraden Linie linear ist, ist das mit der Mittelpunktsregel exakt.
        $S$ selbst ist quadratisch in $s$. Auf einem Bogen mit Mittelpunkt $(y_C, z_C)$ ist
        $ds = R\,d\varphi$, und das Integral ist geschlossen lösbar:
        \[
          \Delta S_z = t R \left[(y_C - y_S)\,\Delta\varphi + R\,(\sin\varphi_b - \sin\varphi_a)\right],
          \quad
          \Delta S_y = t R \left[(z_C - z_S)\,\Delta\varphi - R\,(\cos\varphi_b - \cos\varphi_a)\right]
        \]
        (für den Lauf gegen den Uhrzeigersinn, sonst mit umgekehrtem Vorzeichen).
  \item An einer \textbf{Verzweigung} gilt das Knotengleichgewicht: Kommen alle Äste bis auf einen an,
        startet der letzte mit der Summe der ankommenden statischen Momente.
\end{itemize}

Aus den statischen Momenten folgt der Schubfluss mit der gekoppelten Formel
\[
  \tau\,t = -\frac{Q_y\,(I_y S_z + I_{yz} S_y) + Q_z\,(I_z S_y + I_{yz} S_z)}{I_y I_z - I_{yz}^2} ,
\]
die für $I_{yz} = 0$ in das bekannte $\tau = -\frac{Q_z S_y}{I_y t}$ bzw.\ $-\frac{Q_y S_z}{I_z t}$
zerfällt. $\tau$ ist positiv in Laufrichtung $s$.

\subsection{Eine geschlossene Zelle: woher kommt \texorpdfstring{$q_0$}{q0}?}\label{sec:mathe-zelle}

Eine Zelle hat kein freies Ende, an dem der Schubfluss null sein müsste. Das Programm schneidet sie
deshalb an einer Stelle auf, rechnet wie beim offenen Profil und bekommt einen Verlauf, der sich
vom richtigen nur um einen \textbf{konstanten Umlaufanteil} $q_0$ unterscheidet. Diesen bestimmt
es auf zwei Wegen:

\textbf{(a) Symmetrie.} Liegt eine Symmetrieachse parallel zur Querkraft, ist der richtige Verlauf
antisymmetrisch und auf der Achse null. Das Programm schneidet die Zelle deshalb möglichst auf der
Achse auf. Klappt das nicht genau, verschiebt es den Umlaufanteil so, dass ein Punkt und sein
Spiegelpunkt entgegengesetzt gleiche Werte haben. Die Symmetrie selbst prüft es, indem es jede
Linie an 13 Stellen abtastet (jede Bogensehne an 3) und schaut, ob der gespiegelte Punkt wieder auf
einer Linie gleicher Dicke liegt. Weil Bögen auch an den Schwerachsen geteilt und die Stücke
gleichmäßig in Sehnen zerlegt werden, sind die Sehnen spiegelbildlicher Bögen selbst
spiegelbildlich.

\textbf{(b) Verträglichkeit.} Ohne passende Symmetrieachse kommt die Bedingung, dass sich die Zelle
unter einer Querkraft im Schubmittelpunkt nicht verdrillt:
\[
  \oint \frac{q}{G\,t}\,ds = 0 .
\]
Der Kniff ist, diese Bedingung nicht für die gerade eingegebene Querkraft zu lösen, sondern für
die statischen Momente selbst. Dazu sortiert man die Formel für den Schubfluss nach $S_z$ und $S_y$:
\[
  q = -\frac{a\,S_z + b\,S_y}{\Delta} , \qquad
  \begin{pmatrix} a \\ b \end{pmatrix}
  = \underbrace{\begin{pmatrix} I_y & I_{yz} \\ I_{yz} & I_z \end{pmatrix}}_{J}
    \begin{pmatrix} Q_y \\ Q_z \end{pmatrix} ,
\]
mit $\Delta = I_y I_z - I_{yz}^2 = \det J$. $a$, $b$, $G$ und $\Delta$ sind entlang der Zelle
konstant, also wird die Verträglichkeit zu
\[
  a \oint \frac{S_z}{t}\,ds + b \oint \frac{S_y}{t}\,ds = 0 .
\]
\emph{Warum $J$ regulär sein muss.} Die statischen Momente hängen nur vom Querschnitt ab, nicht von
der Last. Sollen sie für \emph{jede} Querkraft den richtigen Schubfluss liefern, muss die Gleichung
oben für jedes $(Q_y, Q_z)$ gelten. Weil $J$ regulär ist, gibt es zu jedem gewünschten Paar
$(a, b)$ eine Querkraft, die es erzeugt, nämlich $(Q_y, Q_z)^T = J^{-1}(a, b)^T$. Man darf also
$(a, b) = (1, 0)$ und $(a, b) = (0, 1)$ einsetzen und bekommt die beiden Bedingungen einzeln:
\[
  \oint \frac{S_z}{t}\,ds = 0 \quad\text{und}\quad \oint \frac{S_y}{t}\,ds = 0 .
\]
Wäre $J$ singulär, lägen alle Paare $(a, b)$, die überhaupt vorkommen können, auf einer Geraden durch
den Ursprung. Dann wüsste man nur, dass eine bestimmte Kombination der beiden Umlaufintegrale
verschwindet, und könnte die beiden statischen Momente nicht einzeln festlegen.

Regulär ist $J$ bei jedem echten Querschnitt: Nach der Cauchy-Schwarz-Ungleichung aus \textbf{HM1}
ist $I_{yz}^2 = \bigl(\int y z\,dA\bigr)^2 \le \int y^2\,dA \int z^2\,dA = I_z I_y$, also
$\Delta \ge 0$. Gleichheit gibt es nur, wenn die ganze Fläche auf einer Geraden durch den Schwerpunkt
liegt, und dann versagt schon die Biegeformel, in der $\Delta$ im Nenner steht.

Streng genommen braucht man die Regularität nur für diese Richtung des Schlusses. Dass zwei
einzeln verschobene statische Momente für jede Querkraft den richtigen Schubfluss liefern, folgt
schon daraus, dass $q_0 = -\oint (q_{\text{offen}}/t)\,ds \big/ \oint ds/t$ linear in $a$ und $b$
ist, also in einen Teil zu $S_z$ und einen zu $S_y$ zerfällt. Die Regularität sagt zusätzlich, dass
es keine andere Wahl der beiden Konstanten gibt. Deshalb sind auch die angezeigten $S$-Verläufe
eindeutig.

Das Programm verschiebt also jedes statische Moment in der Zelle um seine eigene Konstante $s_0$.
Aus $\oint (S + s_0)/t\,ds = 0$ folgt
\[
  s_0 = -\,\frac{\oint S/t\;ds}{\oint ds/t}
\]
(beide Integrale im Umlaufsinn gegen den Uhrzeigersinn). Damit stimmen nicht nur $\tau$, sondern
auch die angezeigten $S$-Verläufe. Integriert wird mit der Simpsonregel, die für das quadratische
$S/t$ exakt ist. Auf einem Bogen ist $S$ eine glatte Winkelfunktion, mit 16 Teilen je Sehne
von höchstens $15^\circ$ ist Simpson dort auf etwa $10^{-10}$ genau. Beim symmetrischen Profil kommt
so übrigens genau die antisymmetrische Lösung aus (a) heraus.

\subsection{Schubmittelpunkt}

Der Schubmittelpunkt ist der Punkt, durch den die Resultierende des Schubflusses geht. Das
Programm rechnet den Schubfluss für die Einheitslasten $Q_y = 1$ und $Q_z = 1$ (bei einer Zelle
inklusive $q_0$) und dessen Moment um den Schwerpunkt,
\[
  M_q = \int \left[(y - y_S)\,t_z - (z - z_S)\,t_y\right] q\,ds ,
\]
mit der Tangentenrichtung $(t_y, t_z)$ der Linie. Weil dasselbe Moment aus der Querkraft im
Schubmittelpunkt kommen muss, $M_q = (y_M - y_S)\,Q_z - (z_M - z_S)\,Q_y$, ist
\[
  y_M = y_S + M_q(Q_z = 1), \qquad z_M = z_S - M_q(Q_y = 1) .
\]
Der Integrand ist kubisch in $s$, die Simpsonregel also wieder exakt (auf Bögen, mit der
Bogentangente, wieder auf etwa $10^{-10}$). Die Momententabelle zeigt genau diese Summe, aufgeteilt
nach Elementen. Gegen die Handrechnung: offener Halbkreis $e = 4R/\pi$, D-Profil aus Halbkreis und
Platte und Kreisrohr stimmen auf etwa $10^{-9}$.

\subsection{Massive Querschnitte}

Für $\tau = \frac{Q\,S}{I\,b}$ teilt das Programm den Querschnitt in etwa 160 waagerechte (bzw.\
senkrechte) Streifen. Die Breite $b(z)$ kommt exakt aus allen Elementen: Löcher abgezogen, Polygone
über die Schnittpunkte der Kanten, Kreise über $2\sqrt{R^2 - \Delta^2}$, Kreisausschnitt und
-abschnitt als Kreis geschnitten mit den Halbebenen der Randstrahlen bzw.\ der Sehne. Die
Streifengrenzen liegen auf allen Knickstellen von $b$ (Ecken, Rechteckkanten, Kreisscheitel,
Lochränder) und auf der Schwerachse, dazwischen gleich weit. In jedem Streifen ist $b$ damit glatt,
und $S(z)$ wird mit drei Gauß-Punkten je Streifen von oben aufsummiert (an Knickstellen fein
unterteilt, wegen der Wurzel am Kreisscheitel).

$\tau$ wertet das Programm auf den Streifengrenzen aus, also auch genau am oberen und unteren Rand.
Dort ist $S = 0$, damit $\tau = 0$. Der kleine Rest, den die Quadratur am zweiten Rand übrig lässt,
wird anteilig nach der Fläche abgezogen. An Spitzen und am Kreisscheitel geht $b \to 0$, aber $S$
schneller, der Grenzwert $\tau \to 0$ wird direkt eingesetzt. An einem Breitensprung (Flansch und
Steg, Lochrand) gibt es zwei Stützstellen, eine mit der Breite darunter und eine mit der darüber,
der Verlauf zeigt also den Sprung. Früher wurde $\tau$ in der Streifenmitte ausgewertet, aber mit
$S$ bis zur Streifenunterkante: Der oberste Wert war deshalb nie null (beim Rechteck etwa
$2{,}5\,\%$ von $\tau_{max}$), und der Rand selbst kam gar nicht vor.

Wirken $Q_y$ und $Q_z$ gleichzeitig, ist das angezeigte $\tau_{max}$
$\sqrt{\tau_{y,max}^2 + \tau_{z,max}^2}$, also eine Abschätzung nach oben, denn die beiden Maxima
liegen meist nicht an derselben Stelle.

%---------------------------------------------------------------------
\section{Torsion und Verwölbung}\label{sec:mathe-torsion}

\textbf{Torsion:} Beim offenen Profil ist $I_T = \sum \frac{\xi}{3} L\,t^3$ und
$\tau = M_T\,t/I_T$. Bei einer Zelle kommt die Fläche $A_m$ der Mittellinie aus der Gaußschen
Trapezformel und $\oint ds/t = \sum L_i/t_i$, damit $I_T = 4A_m^2/\oint ds/t$ und
$\tau = M_T/(2A_m t)$. Läuft die Zelle über einen Bogen, nimmt die Trapezformel die Sehnen, und je
Sehne mit dem Winkel $\Delta\varphi$ kommt das Kreissegment
\[
  \Delta A = \frac{R^2}{2}\left(\Delta\varphi - \sin\Delta\varphi\right)
\]
zwischen Sehne und Bogen dazu (mit dem Vorzeichen des Umlaufs). $L_i$ ist dort die Bogenlänge
$R\,\Delta\varphi$. Beides ist exakt, die Sehnen sind also nur eine Hilfe für die Topologie.

\textbf{Zelle mit offenen Ästen} (z.\,B.\ ein Kasten mit überstehenden Gurten): Die Zelle trägt
nach Bredt das ganze $M_T$, und $I_T$ ist ihr Bredt-Wert. Daraus folgt die Verdrillung
$\vartheta = M_T/(G\,I_T)$, und die offenen Äste müssen sich genauso verdrillen wie die Zelle. Ein
Ast der Dicke $t$ hat dann am Rand $\tau = G\,\vartheta\,t = M_T\,t/I_T$. Streng genommen nehmen die
Äste damit zusätzlich das Moment $G\,\vartheta\,I_{T,\text{offen}}$ auf, und alle Schubspannungen
zusammen ergeben $M_T\,(1 + I_{T,\text{offen}}/I_{T,\text{Zelle}})$ statt $M_T$. Das ist eine
Näherung: $I_{T,\text{offen}}$ wächst mit $t^3$, $I_{T,\text{Zelle}}$ mit $A_m^2$, und der
Unterschied liegt bei üblichen Profilen im Promillebereich. Bei einem Kasten $200 \times 100$ mit
$t = 10$ und zwei $50$ langen Gurten derselben Dicke sind es etwa $0{,}1\,\%$. Genauso lässt Bredt
den kleinen offenen Anteil $\frac13 L\,t^3$ der Zellwände selbst weg. Exakt im Gleichgewicht wäre
$I_T = I_{T,\text{Zelle}} + I_{T,\text{offen}}$ mit einer Aufteilung von $M_T$ im Verhältnis der
beiden Steifigkeiten.

Ein
dünnwandiges Kreisrohr ist die Zelle mit $A_m = \pi R^2$ und $\oint ds/t = 2\pi R/t$, also
$I_T = 2\pi R^3 t$. Das Programm erkennt das Rohr aus einem einzelnen Kreis gesondert. Über die
Sehnenkette käme dasselbe heraus, ein Rohr aus zwei Halbkreisen rechnet so. Das Torsionsmoment
aus platzierten Kräften ist das Moment um den \emph{Schubmittelpunkt}, denn um den dreht sich der
Querschnitt.

\textbf{Verwölbung:} Auf einem geraden Segment ist der Fahrstrahlabstand $r_\perp$ linear, damit
ist die Verwölbung nach Formelsammlung geschlossen integrierbar:
\[
  u_x(s) = u_{x,0} + \frac{q_T}{G\,t}\,s - \vartheta\left[r_1 s + (r_2 - r_1)\frac{s^2}{2L}\right],
  \qquad q_T = \frac{M_T}{2A_m}, \quad \vartheta = \frac{M_T}{G\,I_T} ,
\]
wobei der Bredt-Anteil $q_T/(Gt)$ nur in der Zelle auftritt (mit Vorzeichen je nach Umlaufsinn).
Auf einem Bogen ist $r_\perp$ nicht linear, das Integral aber ebenfalls geschlossen lösbar. Mit
dem Pol $(y_M, z_M)$ und dem Winkel $\varphi$ ab dem Sehnenanfang $\varphi_1$ ist
\[
  \int r_\perp\,ds = R\left[(y_C - y_M)(\sin\varphi - \sin\varphi_1)
    - (z_C - z_M)(\cos\varphi - \cos\varphi_1) + R\,(\varphi - \varphi_1)\right] .
\]
Ein Kreisrohr verwölbt sich damit nicht: Mit dem Pol im Mittelpunkt ist $r_\perp = R$, und der
Bredt-Anteil hebt $R\,\vartheta$ genau auf.
Das Programm läuft von einem Knoten auf einer Symmetrieachse aus (dort ist $u_x$ bei symmetrischen
Profilen null) über den Profilgraph und setzt die Werte Segment für Segment fort. Beim Schließen der
Zelle muss es wieder beim Startwert ankommen, weil $\oint r_\perp\,ds = 2A_m$ ist. Das ist die
eingebaute Probe. Zum Schluss wird die Konstante so gewählt, dass $\int u_x\,t\,ds = 0$ ist, also
keine mittlere Längsverschiebung übrig bleibt.

%---------------------------------------------------------------------
\section{Kernfläche}\label{sec:mathe-kern}

Die Kernfläche ist die Menge aller Lastpunkte, deren neutrale Faser den Querschnitt nicht
schneidet. Gerechnet wird rückwärts:

\begin{enumerate}
  \item Das Programm sammelt die Eckpunkte des Querschnitts (Bögen abgetastet) und bildet ihre
        \textbf{konvexe Hülle}. Einbuchtungen spielen keine Rolle, denn eine Gerade, die die Hülle
        berührt, berührt auch den Querschnitt nicht innen.
  \item Jede Kante der Hülle ist eine mögliche \textbf{neutrale Faser}. Mit dem nach außen
        zeigenden Einheitsvektor $\vec n$ der Kante und ihrem Abstand $d$ vom Schwerpunkt gehört
        dazu der Lastpunkt
        \[
          \vec e = -\frac{1}{A\,d}
          \begin{pmatrix} I_z & -I_{yz} \\ -I_{yz} & I_y \end{pmatrix} \vec n .
        \]
        Für ein Rechteck und die obere Kante ist $\vec n = (0,1)$, $d = h/2$, und es kommt
        $e_z = -h/6$ heraus, das bekannte Kerndrittel.
  \item Die Kernpunkte aller Kanten, der Reihe nach verbunden, sind die Kernfläche. An Bögen legt
        das Programm alle $5^\circ$ eine Tangente an, beim Vollkreis 36 Stück.
\end{enumerate}

\textbf{Woher die Matrix kommt.} Die Formel für $\vec e$ ist die Formel für $\sigma_x$ aus
Abschnitt~\ref{sec:mathe-sigma}, rückwärts gelesen. Eine Normalkraft $N$ im Punkt
$\vec e = (e_y, e_z)$, gemessen vom Schwerpunkt, erzeugt $M_y = N e_z$ und $M_z = -N e_y$. In
$\sigma_x$ eingesetzt und nach $y$ und $z$ sortiert:
\[
  \sigma_x = \frac NA + \frac N\Delta\Bigl[\bigl(I_y e_y + I_{yz} e_z\bigr)\,y
                                       + \bigl(I_{yz} e_y + I_z e_z\bigr)\,z\Bigr]
           = \frac NA\Bigl(1 + \frac A\Delta\,\vec x^T J\,\vec e\Bigr) ,
\]
mit $\vec x = (y, z)$ und der Trägheitsmatrix $J = \begin{pmatrix} I_y & I_{yz} \\ I_{yz} & I_z
\end{pmatrix}$, $\Delta = \det J$. Die neutrale Faser ist $\sigma_x = 0$, also die Gerade
\[
  \vec x^T \Bigl(\frac A\Delta\,J\,\vec e\Bigr) = -1 .
\]
Die Kante mit dem nach außen zeigenden Einheitsvektor $\vec n$ und dem Abstand $d$ vom Schwerpunkt
ist die Gerade $\vec x^T \vec n = d$, oder genauso $\vec x^T(-\vec n/d) = -1$. Eine Gerade, die
nicht durch den Schwerpunkt geht, lässt sich nur auf eine Weise in der Form $\vec x^T \vec g = -1$
schreiben. Die neutrale Faser ist also genau dann die Kante, wenn die beiden Vektoren gleich sind:
\[
  \frac A\Delta\,J\,\vec e = -\frac{\vec n}{d}
  \qquad\Longrightarrow\qquad
  \vec e = -\frac{1}{A\,d}\,\bigl(\Delta\,J^{-1}\bigr)\,\vec n ,
  \qquad
  \Delta\,J^{-1} = \begin{pmatrix} I_z & -I_{yz} \\ -I_{yz} & I_y \end{pmatrix} .
\]
Die Matrix in der Formel ist also die Inverse der Trägheitsmatrix ohne den Vorfaktor $1/\det J$
(Diagonale vertauscht, Nebendiagonale mit anderem Vorzeichen). Auch hier braucht man, dass $J$
regulär ist, siehe Abschnitt~\ref{sec:mathe-zelle}.

\textbf{Warum man die Kernpunkte geradlinig verbinden darf.} Dreht man die neutrale Faser um eine
Ecke $\vec P$ der Hülle von einer Kante zur nächsten, gehen alle diese Geraden durch $\vec P$, also
gilt $\vec P^T(\vec n/d) = 1$. Der Vektor $\vec n/d$ wandert damit auf einer Geraden, und weil
$\vec e$ linear von $\vec n/d$ abhängt, wandert auch der Lastpunkt auf einer Geraden. Jeder Ecke der
Hülle entspricht so eine gerade Kante der Kernfläche, jeder Kante der Hülle eine Ecke.

Im Hauptachsensystem fällt die Kopplung weg ($I_{\eta\zeta} = 0$). Schreibt man die Kante mit
ihren Achsenabschnitten als $\eta/a_\eta + \zeta/a_\zeta = 1$, ist $\vec n/d = (1/a_\eta,
1/a_\zeta)$, und aus der Formel wird die bekannte Form $\eta_K = -i_\zeta^2/a_\eta$,
$\zeta_K = -i_\eta^2/a_\zeta$ mit $i^2 = I/A$. Deshalb beschriftet das
Programm die umhüllenden Geraden mit ihren Abschnitten auf den Hauptachsen, gemessen vom
Schwerpunkt, und nicht mit Koordinaten im gezeichneten KOS.

%---------------------------------------------------------------------
\section{Spannungskreis: vom Schnitt zum Zustand}\label{sec:mathe-kreis}

\textbf{Lineare Gleichungen.} Mit $\sigma_m$, $R_c$ und $R_s$ als Unbekannten ist jede Angabe mit
bekanntem Winkel eine Zeile
\[
  \sigma = \sigma_m + R_c\cos 2\alpha + R_s\sin 2\alpha, \qquad
  \tau = -R_c\sin 2\alpha + R_s\cos 2\alpha .
\]
\textbf{Punkte ohne Winkel} erfüllen $(\sigma - \sigma_m)^2 + \tau^2 = R_c^2 + R_s^2$. Zieht man
diese Gleichung für zwei solche Punkte voneinander ab, fallen alle quadratischen Terme weg und es
bleibt eine lineare Gleichung für $\sigma_m$ allein:
$2(\sigma_2 - \sigma_1)\,\sigma_m = \sigma_2^2 + \tau_2^2 - \sigma_1^2 - \tau_1^2$. Bei mehreren
Punkten ohne Winkel zieht das Programm die Gleichung des ersten Punktes von jeder anderen ab. Jede
Differenz ist linear und kommt zu den linearen Zeilen. Quadratisch bleibt damit nur noch eine
Gleichung, die des ersten Punktes selbst: Die Information der anderen Punkte steckt vollständig in
den Differenzen.

\textbf{Rang.} Die linearen Zeilen löst das Programm mit Gauß-Jordan und Spaltenpivotsuche und
zählt dabei, wie viele Pivotelemente es findet (Toleranz relativ zum größten Eintrag). Bleibt in
einer Zeile ohne Pivot rechts etwas übrig, widersprechen sich die Angaben.

\begin{itemize}
  \item \textbf{Rang 3:} Der Zustand ist eindeutig. Ein Punkt ohne Winkel muss dann nur noch auf
        dem Kreis liegen (Probe).
  \item \textbf{Rang 2 und ein Punkt ohne Winkel:} Die Lösungen der linearen Zeilen bilden eine
        Gerade $\vec x = \vec x_0 + t\,\vec d$ im $(\sigma_m, R_c, R_s)$-Raum. In die quadratische
        Gleichung eingesetzt, wird daraus $a t^2 + b t + c = 0$ mit keiner, einer oder zwei Lösungen.
        Genau daher kommt \emph{zwei Zustände möglich}.
  \item \textbf{Weniger:} Es fehlen Angaben.
\end{itemize}

Schnitte ohne Winkel mit nur $\sigma$ oder nur $\tau$ geben keine Gleichung, aber eine
Ungleichung: $|\sigma - \sigma_m| \le R$ bzw.\ $|\tau| \le R$. Damit sortiert das Programm von zwei
möglichen Zuständen den aus, auf dessen Kreis der Wert nicht passt.

\textbf{Die Spannungsscheibe} ist das Vieleck, dessen Kanten alle einen Kreis mit Radius $r$
berühren. Jeder Schnitt $i$ mit dem Normaleneinheitsvektor $\vec n_i$ gibt eine Kante: die Gerade
$\vec n_i \cdot \vec x = r$, senkrecht zu $\vec n_i$ im Abstand $r$ vom Mittelpunkt. Zur Scheibe
gehört alles auf der Seite des Mittelpunktes, also die Halbebene $\vec n_i \cdot \vec x \le r$. Die
Scheibe ist die Menge der Punkte, die in allen diesen Halbebenen zugleich liegen. Liegen zwei
benachbarte Normalen $160^\circ$ oder mehr auseinander, wäre die Scheibe
offen (oder eine sehr spitze Nadel). Dann ergänzt das Programm den $x$- oder $y$-Schnitt, dessen
Normale am nächsten an der Mitte der größten Lücke liegt, und das so oft, bis die Lücke kleiner ist.
 eingebunden.

\part{Die Mathematik hinter den Features}\label{sec:Mathe}

Vorab will ich erwähnen, dass Teile dieses Abschnitts nicht meine eigene Arbeit sind. Gerade wenn
es darum geht, das Programm effizient und robust zu schreiben, oder wenn sich das Skript Methoden der
weiterführenden Mathematik bzw.\ Numerik bedient, stammen Ideen von KI, z.\,B.\ die
$LDL^T$-Zerlegung (ich bin selbst erst im 2.~Semester, deswegen halten sich meine mathematischen
Kenntnisse noch in Grenzen). Das gilt übrigens auch für große Teile der Umsetzung. Aus diesem Grund
halte ich die Erklärungen zu diesen Aspekten hier etwas kürzer.

Nichtsdestotrotz ist dieser Teil für alle, die wissen wollen, \emph{warum} das Programm so rechnet, wie es rechnet.
Zum Benutzen braucht ihr ihn nicht. Grundlegende Sachen wie Hebelarme, Steiner und Gleichgewicht seien hier nicht weiter erläutert. Hier
werden die Stellen behandelt, an denen das Programm mehr macht als die Vorlesung oder es anders macht,
weil ein Rechner eben nicht wie ein Mensch rechnet. Ich habe versucht, zu jedem Punkt dazuzuschreiben,
was er für euch in der Klausur bedeutet.

Ein paar Konventionen, die überall gelten:

\begin{itemize}
  \item Im \textbf{Tragwerksskript} zeigt $x$ nach \emph{rechts} und $y$ nach \emph{unten}. Drehwinkel,
        Momente und Lastwinkel zählen trotzdem so, wie ihr sie seht: gegen den Uhrzeigersinn
        positiv.
  \item Jeder Stab hat seine eigene Koordinate $x \in [0, L]$ vom Anfangs- zum Endknoten, die
        Tangente $\vec t = (\Delta x, \Delta y)/L$ und die Normale $\vec n = (-\Delta y, \Delta x)/L$.
        Die Schnittgrößen folgen der TM1-Konvention mit gestrichelter Faser auf der Seite von
        $\vec n$.
  \item Im \textbf{Querschnittsskript} gilt die Konvention aus dem Gross (TM2):
        $I_y = \int z^2\,dA$, $I_z = \int y^2\,dA$ und $I_{yz} = -\int yz\,dA$, also mit Minus.
\end{itemize}

%---------------------------------------------------------------------
\section{Das Weggrößenverfahren im Programm}\label{sec:mathe-wgv}

Alles, was ihr im Tragwerksskript seht (Verläufe, Auflagerkräfte, Verformungen,
Explosionszeichnung), kommt aus einer einzigen Rechnung: dem Weggrößenverfahren, also genau dem
allgemeinen WGV aus \textbf{Statik}. Unbekannt sind die Verschiebungen der Knoten, nicht die
Kräfte.

\textbf{Unbekannte:} Jeder Knoten hat drei Freiheitsgrade, $u_\xi$, $u_\eta$ und $\varphi$. Bei
einem ungedrehten Lager ist $\xi = x$ und $\eta = y$, bei einem gedrehten Lager zeigen sie in die
Lagerrichtungen. Ein Lager \emph{streicht} einfach den gehaltenen Freiheitsgrad.

\textbf{Stabsteifigkeit:} Für jeden Stab steht die bekannte $6 \times 6$-Matrix aus $EA/L$,
$12EI/L^3$, $6EI/L^2$, $4EI/L$ und $2EI/L$. Sie verknüpft die sechs Endverschiebungen
$(u_a, w_a, \varphi_a, u_b, w_b, \varphi_b)$ im lokalen System des Stabes mit den sechs
Endkräften. Zwei Drehungen bringen sie ins globale System: die Stabrichtung und (bei einem
gedrehten Lager) die Lagerrichtung,
\[
  K_e = T_K\,T^T\,k\,T\,T_K^T .
\]
So ist ein gedrehtes Lager kein Sonderfall, sondern nur eine andere Basis am Knoten. Matrix,
Transformation und Drehrichtungen sind dieselben wie in der WGV-Formelsammlung aus Statik, die
FEM-Formelsammlung hat dieselben Elemente, zählt die Querrichtung aber nach oben. Die
Unterschiede stehen in Abschnitt~\ref{sec:wgv-vorzeichen}.

\textbf{Lasten:} Knotenlasten gehen direkt in den Lastvektor (vorher um den Lastwinkel gedreht).
Streckenlasten gehen als Festeinspannkräfte ein. Das sind die Endkräfte, die ein
beidseitig eingespannter Stab unter der Last hätte, z.\,B.\ $\mp qL/2$ und $\pm qL^2/12$ bei
Gleichlast (die Volleinspannmomente aus Tafel~A9 der Formelsammlung, mit denselben Vorzeichen).
Im Gesamtsystem stehen sie mit umgekehrtem Vorzeichen, also als Knotenersatzlasten. Für Trapezlasten stehen die entsprechenden Formeln im Code, Lastfunktionen wie
\code{x\^{}2} integriert das Programm mit den Hermite-Formfunktionen
$N_1 = 1 - 3\xi^2 + 2\xi^3$, $N_2 = L(\xi - 2\xi^2 + \xi^3)$ usw.\ (arbeitsäquivalente
Knotenlasten, Abschnitt~\ref{sec:mathe-festeinspann}). Temperatur wirkt wie eine Last $EA\,\varepsilon_T$ bzw.\ $EI\,\kappa_T$ an den
Stabenden. Globale Linienlasten $p_X(x)$, $p_Y(x)$ werden vorher in Quer- und Längsanteil zerlegt,
bei einer projizierten Last ist $x$ die globale Koordinate ab dem Stabanfang, $x_{global} = f\,x$
mit dem Projektionsfaktor $f = |\Delta y|/L$ bzw.\ $|\Delta x|/L$, und je Stablänge wirkt
$f\,p(f\,x)$.

\textbf{Lösen:} Das Gesamtsystem $K\,u = F$ wird mit einer $LDL^T$-Zerlegung gelöst. Wird dabei
ein Diagonalelement $D_j$ (fast) null, fehlt dem System in dieser Richtung jede Steifigkeit. Es
ist ein Mechanismus, und ihr bekommt \emph{System ist kinematisch!}.

\textbf{Zurückrechnen:} Aus den Knotenverschiebungen folgen die Endkräfte jedes Stabes,
$s = k\,v + p$ (Steifigkeit mal Endverschiebung plus Festeinspannkraft), und aus dem Anfangswert
durch Integrieren der Lasten die Verläufe, z.\,B.\
$M(x) = M_a + Q_a x - m\,x - \int_0^x q(t)\,(x-t)\,dt$. Das ist exakt, nicht genähert. Deshalb
braucht jeder Bereich zwischen Lasten, Lagern und Gelenken nur einen einzigen Stab.

\textbf{Auflagerkräfte:} Die Summe der Stabendkräfte an einem Knoten minus der Knotenlast ist die
Kraft, die Lager und Federn auf den Knoten ausüben. Davon zeigt das Programm
$R_\xi = \vec R\cdot\vec\xi$ und $R_\eta = -\vec R\cdot\vec\eta$, deshalb zeigt $R_\eta$ nach oben,
während die Knotenverschiebung \code{ü\_y} auf derselben Menüseite nach unten positiv ist.

\subsection{Festeinspannkräfte beliebiger Lasten}\label{sec:mathe-festeinspann}

Für das Gesamtsystem braucht das Programm bzw. das WGV von einer Streckenlast nur, was sie an den Stabenden
bewirkt, die Festeinspannkräfte. Für Gleich- und Trapezlasten stehen diese In Tafeln oder in der FOS, für eine
Lastfunktion wie $q(x) = x^2$ oder $q_0 \sin(\pi x/L)$ nicht. Das Programm rechnet sie deshalb für
jede Lastfunktion mit den Hermite-Funktionen aus, und zwar exakt, nicht genähert.

\textbf{Die Hermite-Funktionen.} Im lokalen System läuft $x$ von $0$ bis $L$, $\xi = x/L$, $w$ zeigt
zur gestrichelten Faser und $\varphi = -w'$ gegen den Uhrzeigersinn. Die vier Funktionen sind
\begin{align*}
  N_1 &= 1 - 3\xi^2 + 2\xi^3       && w(0) = 1, \\
  N_2 &= L\,(\xi - 2\xi^2 + \xi^3) && w'(0) = 1, \text{ also } \varphi_a = -1, \\
  N_3 &= 3\xi^2 - 2\xi^3           && w(L) = 1, \\
  N_4 &= L\,(-\xi^2 + \xi^3)       && w'(L) = 1, \text{ also } \varphi_b = -1,
\end{align*}
und die jeweils drei anderen Endwerte sind null. Jede ist eine kubische Parabel, löst also
$EI\,w'''' = 0$ exakt und stellt somit eine homogene Lösung der DGL dar.
Gleichzeitig erfüllt sie die vier von ihr geforderten Randbedingungen. Damit ist sie genau die Biegelinie des
beidseitig eingespannten Stabes, dem man diese eine Einheitsverformung aufzwingt.

Das WGV zählt die vier Endverschiebungen eines Stabes aber als $v = (w_a, \varphi_a, w_b, \varphi_b)$,
mit $\varphi$ gegen den Uhrzeigersinn positiv. Die Einheitsverformung zum Freiheitsgrad $j$ in
dieser Zählung heißt im Folgenden $S_j$ (S für Formfunktion, engl.\ \emph{shape}):
\[
  S_1 = N_1, \qquad S_2 = -N_2, \qquad S_3 = N_3, \qquad S_4 = -N_4 .
\]
Bei den Drehungen kehrt sich das Vorzeichen um, weil $N_2$ und $N_4$ zu $\varphi = -1$ gehören.

\textbf{Beweisidee mit dem Satz von Betti.} Die folgende Herleitung soll die Idee zeigen und
ist kein formal vollständiger Beweis. So nennen wir alle $F_i$ Kräfte und alle Verschiebungen Wege,
obwohl bei den Drehungen Momente und Winkel gemeint sind, und die Einheitsverschiebung \enquote{$1$}
hat eigentlich eine Einheit (m bzw.\ rad). Das formal sauber mitzuschleppen ändert am Ergebnis
nichts, macht die Rechnung aber deutlich schwerer zu lesen.

Die Festeinspannkräfte $F_i$ sind die Kräfte und
Momente, die die Einspannungen auf den belasteten Stab ausüben, jeweils in positiver Richtung des
Freiheitsgrades (das $p$ in $s = k\,v + p$). Man schneidet den Stab frei. Die Einspannungen sind
dann nur noch Kräfte an den Stabenden, und man hat \emph{ein} System mit zwei Lastgruppen:

\begin{itemize}
  \item Lastgruppe~1: die Last $q(x)$ und die Endkräfte $F_j$, die die Enden festhalten (gesucht).
        Die Biegelinie ist $w_1$ mit $w_1 = w_1' = 0$ an beiden Enden, alle Endverschiebungen sind
        also $v_j^{(1)} = 0$.
  \item Lastgruppe~2: nur die Endkräfte $k_{ji}$, die den Freiheitsgrad $i$ um $1$ verschieben und
        alle anderen festhalten. Die Biegelinie ist $S_i(x)$, die Endverschiebungen sind
        $v_j^{(2)} = \delta_{ij}$ (Kronecker-Delta: $1$ für $j = i$, sonst $0$).
\end{itemize}

Beide Gruppen sind für sich im Gleichgewicht, deshalb stört es nicht, dass der freigeschnittene
Stab nicht gelagert ist. Nach Betti leistet Lastgruppe~1 auf den Wegen von Lastgruppe~2 dieselbe
Arbeit wie Lastgruppe~2 auf den Wegen von Lastgruppe~1:
\[
  \int_0^L q(x)\,S_i(x)\,dx + F_j\,v_j^{(2)} = k_{ji}\,v_j^{(1)} .
\]
Über das doppelte $j$ wird auf beiden Seiten summiert (Einsteinsche Summenkonvention). Setzt man
die Verschiebungen ein, bleibt links $F_j\,\delta_{ij} = F_i$ und rechts null. Die innere Arbeit
fehlt dabei nicht, sie ist der Grund für die Gleichheit:
\[
  W_a^{1 \to 2} = \int_0^L M_1\,\kappa_2\,dx = \int_0^L \frac{M_1 M_2}{EI}\,dx
  = \int_0^L M_2\,\kappa_1\,dx = W_a^{2 \to 1} .
\]
Das erste Gleichheitszeichen ist der Arbeitssatz (äußere gleich innere Arbeit) mit zwei
verschiedenen Zuständen: die Kräfte und Momente $M_1$ aus Lastgruppe~1, die im Gleichgewicht sind,
und die Verformung $\kappa_2$ aus Lastgruppe~2, die verträglich ist. Einer der beiden Zustände
spielt dabei die Rolle des virtuellen. Im PvV aus \textbf{TM1} ist es die Verschiebung (hier die
von Lastgruppe~2), beim Prinzip der virtuellen Kräfte aus \textbf{Statik} wäre es die Kraftgruppe.
Der Satz selbst gilt in beiden Lesarten, solange der eine Zustand im Gleichgewicht und der andere
verträglich ist. Das letzte Gleichheitszeichen ist derselbe Satz mit vertauschten Rollen: Kräfte
aus Lastgruppe~2, Verformung aus Lastgruppe~1. In der Mitte steht nur $\kappa = M/EI$. Betti
schreibt die beiden Enden dieser Kette hin. Die Kette enthält nur die Biegearbeit. Das Programm rechnet nach Euler-Bernoulli, der Stab ist
also schubstarr ($\gamma = 0$), und die Querkraft leistet keine innere Arbeit. Nach Timoschenko käme
in der Mitte noch $\int_0^L Q_1 Q_2/(GA_s)\,dx$ dazu. Betti gilt dann weiterhin, nur wären die $S_i$
andere Kurven, die auch von $GA_s$ abhängen.

Dass die gemischte innere Arbeit hier null ist, obwohl die Schnittgrößen aus Lastgruppe~2 auf
der Biegung aus Lastgruppe~1 lokal sehr wohl Arbeit leisten, sieht man auch direkt: $M_2$ ist ohne
Streckenlast linear, also $M_2'' = 0$, und zweimal partiell integriert bleibt
\[
  \int_0^L M_2\,\kappa_1\,dx = -\int_0^L M_2\,w_1''\,dx
  = \Bigl[-M_2\,w_1' + M_2'\,w_1\Bigr]_0^L - \int_0^L M_2''\,w_1\,dx = 0 ,
\]
weil $w_1$ und $w_1'$ an beiden Enden verschwinden. Also gilt für jede Lastfunktion
\[
  F_i = -\int_0^L q(x)\,S_i(x)\,dx .
\]
Mit $S_1 = N_1$, $S_2 = -N_2$, $S_3 = N_3$ und $S_4 = -N_4$ ergibt das genau die Formeln im
Programm:
\begin{align*}
  Q_a &= -\int_0^L q\,N_1\,dx, & M_a &= \int_0^L q\,N_2\,dx, \\
  Q_b &= -\int_0^L q\,N_3\,dx, & M_b &= \int_0^L q\,N_4\,dx.
\end{align*}

\textbf{Anschaulich} ist $S_i(x)$ die Einflusslinie der Festeinspannkraft $i$ (Satz von
Müller-Breslau): Eine Einzelkraft $P$ an der Stelle $a$ erzeugt $F_i = -P\,S_i(a)$, und das Integral
summiert nur die vielen kleinen Einzellasten $q\,dx$ auf.

\textbf{Kontrolle mit Gleichlast:} $\int_0^L N_1\,dx = L/2$, $\int_0^L N_2\,dx = L^2/12$ und
$\int_0^L N_4\,dx = -L^2/12$. Damit ist $Q_a = Q_b = -qL/2$ (die Last zeigt nach unten, die
Einspannung drückt dagegen) und $M_a = qL^2/12$, $M_b = -qL^2/12$, genau die Tafelwerte. Auch das
Vorzeichen von $M_a$ passt: Ohne Einspannung würde sich das linke Ende im Uhrzeigersinn drehen, die
Einspannung hält gegen den Uhrzeigersinn dagegen.

\textbf{Die anderen Lastarten.}
\begin{itemize}
  \item Ein \textbf{Streckenmoment} $m(x)$ (gegen den Uhrzeigersinn positiv) leistet Arbeit auf der
        Verdrehung, $\int m\,\varphi\,dx = -\int m\,w'\,dx$. Mit Betti wird daraus
        $F_i = \int_0^L m\,S_i'\,dx$, also $Q_a = \int m\,N_1'\,dx$, $M_a = -\int m\,N_2'\,dx$ und
        entsprechend am Ende $b$. Bei konstantem $m$ ist $\int N_1'\,dx = N_1(L) - N_1(0) = -1$ und
        $\int N_2'\,dx = 0$: $Q_a = -m$, $Q_b = +m$ und keine Endmomente. Das ist das Kräftepaar aus
        Abschnitt~\ref{sec:wgv-vorzeichen} und der Grund, warum das Programm vom $m\,l/2$ der
        WGV-Tafel abweicht.
  \item Für eine \textbf{Längslast} $n(x)$ sind die Formfunktionen linear, $1 - \xi$ und $\xi$ (sie
        lösen $EA\,u'' = 0$ exakt), und es gilt $N_a = -\int n\,(1-\xi)\,dx$, $N_b = -\int n\,\xi\,dx$.
  \item \textbf{Temperatur} ist keine Last, sondern eine eingeprägte Verzerrung. Die
        Festeinspannkräfte stehen direkt da: $N_a = EA\,\varepsilon_T$, $N_b = -EA\,\varepsilon_T$,
        $M_a = EI\,\kappa_T$, $M_b = -EI\,\kappa_T$.
\end{itemize}

\textbf{Zwei Wege im Code.} Für konstante Lasten und Trapezlasten
$q = p_a(1-\xi) + p_b\,\xi$ sind die Integrale schon ausgerechnet:
\begin{align*}
  Q_a &= -\tfrac{L}{60}\,(21p_a + 9p_b),   & M_a &= \tfrac{L^2}{60}\,(3p_a + 2p_b), \\
  Q_b &= -\tfrac{L}{60}\,(9p_a + 21p_b),   & M_b &= -\tfrac{L^2}{60}\,(2p_a + 3p_b).
\end{align*}
Hängt die Last von $x$ ab (auch bei $m(x)$, $n(x)$ und den globalen Linienlasten), rechnet das CAS
genau die Integrale oben aus. Ein Beispiel ist $q(x) = q_0 \sin(\pi x/L)$. Für die Querkraft
braucht man $\int_0^L q\,N_1\,dx$, und das geht ohne Rechnung in drei Schritten:
\begin{enumerate}
  \item $N_1 + N_3 = 1$: Schiebt man beide Stabenden um $1$ nach unten, verschiebt sich der ganze
        Stab um $1$. Also ist $\int_0^L q\,N_1\,dx + \int_0^L q\,N_3\,dx = \int_0^L q\,dx$.
  \item $N_3$ ist $N_1$ gespiegelt, $N_3(\xi) = N_1(1-\xi)$. Ist die Last symmetrisch zur
        Stabmitte, so wie der Sinus, sind beide Integrale deshalb gleich groß.
  \item Also ist jedes die Hälfte: $\int_0^L q\,N_1\,dx = \frac12 \int_0^L q\,dx
        = \frac12 \cdot \frac{2q_0L}{\pi} = \frac{q_0L}{\pi}$, und $Q_a = -q_0L/\pi$.
\end{enumerate}
Anschaulich: Bei symmetrischer Last trägt jede Einspannung die halbe Last. Für das Moment
braucht man $\int_0^1 \sin(\pi\xi)\,(\xi - 2\xi^2 + \xi^3)\,d\xi = 2/\pi^3$, also
$M_a = 2q_0L^2/\pi^3 \approx 0{,}0645\,q_0L^2$.

\textbf{Warum das exakt ist.} Betti gilt für jede Last, und die $S_i$ sind exakte Lösungen der
homogenen Gleichung. Also sind die Festeinspannkräfte exakt, damit auch die Knotenverschiebungen aus
$K\,u = F$ und die Stabendkräfte $s = k\,v + p$. Nicht exakt wäre nur die Hermite-Interpolation
$\sum_j v_j S_j(x)$ im Feld, denn die kennt die Last nicht. Die benutzt das Programm aber nie als
Ergebnis (siehe unten).

\textbf{Zwei Zerlegungen derselben Biegelinie.} Die Biegelinie eines Stabes ist immer
Partikulärlösung plus homogene Lösung. Welche Partikulärlösung man nimmt, ist aber nicht
festgelegt: Zwei Partikulärlösungen unterscheiden sich immer um eine kubische Parabel, also um eine
homogene Lösung. Das Programm benutzt zwei davon, und beide beschreiben dieselbe Linie.

\begin{itemize}
  \item \emph{Die Anfangswert-Lösung} ist das Vierfachintegral der Querlast,
        \[
          EI\,w_p(x) = \int_0^x q(t)\,\frac{(x-t)^3}{6}\,dt ,
        \]
        mit $w_p = w_p' = w_p'' = w_p''' = 0$ bei $x = 0$, so wie ihr sie in \textbf{TM2} aufstellt.
        Dazu kommt die homogene Lösung, deren vier Konstanten die Randbedingungen erfüllen. Das
        Rand-LGS schreibt sie als $c_1 x^3/6 + c_2 x^2/2 + c_3 x + c_4$ und holt die Konstanten aus
        allen Rand- und Übergangsbedingungen (Abschnitt~\ref{sec:mathe-randlgs}). Der alte Export
        \code{v\_EI\_stab} schreibt sie mit den Hermite-Funktionen, und die Koeffizienten sind die
        Endverschiebungen aus dem WGV:
        \[
          EI\,w = EI\,w_p + EI w_a\,N_1 + EI w_a'\,N_2 + \bigl(EI w_b - EI w_p(L)\bigr)\,N_3
                  + \bigl(EI w_b' - EI w_p'(L)\bigr)\,N_4 .
        \]
        $w_p$ bringt am Ende $b$ schon eigene Werte mit, die ziehen die letzten beiden Koeffizienten
        wieder ab, damit die Summe dort genau die Endverschiebung hat. Die Koeffizienten sind
        Durchbiegung und Neigung, so wie ihr die Randbedingungen in \textbf{TM2} aufstellt. In der
        Zählung des WGV ist $w_a = v_1$, $w_a' = -\varphi_a = -v_2$, $w_b = v_3$ und $w_b' = -v_4$.
  \item \emph{Die Lösung des eingespannten Stabes} ist
        \[
          w_0(x) = w_p(x) - w_p(L)\,N_3(x) - w_p'(L)\,N_4(x) .
        \]
        Auch sie löst $EI\,w'''' = q$, hat aber an \emph{beiden} Enden $w_0 = w_0' = 0$: Das ist die
        Biegelinie des beidseitig eingespannten Stabes unter der Last. Mit ihr zerfällt die Linie in
        \[
          w(x) = w_0(x) + \sum_j v_j\,S_j(x) ,
        \]
        und die Koeffizienten des homogenen Teils sind genau die Endverschiebungen $v_j$, also die
        Unbekannten des WGV. Diese Zerlegung ist die Idee hinter $s = k\,v + p$: Der homogene Teil
        $\sum_j v_j S_j$ liefert die Endkräfte $k\,v$. $w_0$ verschiebt die Stabenden nicht
        ($w_0 = w_0' = 0$ an beiden Enden), stört also die Verträglichkeit an den Knoten nicht und
        liefert nur Endkräfte, und das sind genau die Festeinspannkräfte $p$. Mehr braucht das WGV
        von der Last nicht.
\end{itemize}

Bei Gleichlast ist $EI\,w_p = qx^4/24$ mit $EI\,w_p(L) = qL^4/24$ und $EI\,w_p'(L) = qL^3/6$, und
daraus wird
\[
  EI\,w_0 = \frac{qL^4}{24}\bigl(\xi^4 - (3\xi^2 - 2\xi^3) - 4(\xi^3 - \xi^2)\bigr)
          = \frac{q\,x^2 (L-x)^2}{24},
\]
die bekannte Biegelinie des eingespannten Balkens. Ihr Moment am linken Ende ist
$-EI\,w_0''(0) = -qL^2/12$ (Zug oben), und die Einspannung hält mit $M_a = +qL^2/12$ gegen den
Uhrzeigersinn dagegen. Der Satz von Betti erspart es, $w_0$ überhaupt auszurechnen: Die
Festeinspannkräfte kommen direkt aus den Integralen über $q\,S_i$.

Die fertige Biegelinie ist in beiden Zerlegungen dieselbe. Das sieht man, wenn man die Formel des
alten Exports umsortiert. Mit $w_a'\,N_2 = -v_2\,N_2 = v_2\,S_2$ und genauso am Ende $b$ wird aus ihr
\[
  EI\,w = \underbrace{EI\,\bigl(w_p - w_p(L)\,N_3 - w_p'(L)\,N_4\bigr)}_{EI\,w_0}
          + EI \sum_j v_j\,S_j ,
\]
also Buchstabe für Buchstabe die zweite Zerlegung, nur anders gruppiert. Ausgerechnet wird $w_0$
dabei nie. Für die Knotengleichungen braucht das WGV nur seine Endkräfte $p$, und die kommen nach
Betti direkt aus den Integralen über $q\,S_i$. Die Verläufe rechnet das Programm danach aus den
exakten Endkräften (Gleichgewicht) bzw.\ mit der Anfangswert-Lösung plus Hermite-Anteil, wie im
ersten Punkt. Die zweite Zerlegung dient hier also nur der Erklärung, warum $s = k\,v + p$ exakt ist.

\subsection{Gelenke und die Knotenverdrehung}\label{sec:mathe-gelenke}

Hier kommt die Warnung aus Abschnitt~\ref{sec:Eingabe_Knoten} her. Ein \textbf{M-Gelenk} wird nicht
herauskondensiert, sondern bekommt einen eigenen Freiheitsgrad: Das Stabende am Gelenk hat seine
eigene Verdrehung. Zu diesem Freiheitsgrad gehört eine eigene Gleichgewichtsgleichung, und an ihm
greift kein äußeres Moment an. Somit ist das Gelenkmoment null. Für die angezeigte Knotenverdrehung
heißt das:

\begin{itemize}
  \item \textbf{Angelenk} (Gelenk am Stabende, andere Stäbe biegesteif am Knoten): Die
        Knotenverdrehung ist die Verdrehung der \emph{biegesteif} angeschlossenen Stäbe. Der Stab
        mit dem Gelenk dreht sich anders, seine Endverdrehung steht im Stabmenü auf Seite~7
        ($\varphi_a$, $\varphi_b$), der Knick im Gelenk ist die Differenz.
  \item \textbf{Vollgelenk} (Knoten gelenkig): Der Knoten hat \emph{gar keine} eigene Verdrehung,
        jedes Stabende dreht sich für sich. Die $0$ im Knotenmenü ist nur ein Platzhalter, die
        echten Verdrehungen stehen wieder auf Seite~7 der Stäbe.
  \item \textbf{Nur Angelenke an einem Knoten}, aber kein Vollgelenk und kein Lager, das die
        Drehung hält: Dann hat der Knoten eine Verdrehung, an der nichts hängt. Das Programm
        meldet \emph{kinematisch}. Deshalb der Stolperstein im Anhang: In diesem Fall braucht ihr
        das Vollgelenk.
\end{itemize}

\textbf{N- und Q-Gelenke} werden dagegen herauskondensiert (siehe Dallmann Baustatik 3): Im Stab wird der entsprechende Anteil
der Steifigkeit gestrichen, die Festeinspannkräfte werden statisch gleichwertig
umverteilt, und nach der Rechnung wird die Klaffung im Gelenk wieder ausgerechnet. Die
Endverschiebungen auf Seite~7 enthalten sie. Was dabei gestrichen wird, betrifft immer nur die
freigegebene Richtung:

\begin{itemize}
  \item \textbf{N-Gelenk:} Die Dehnsteifigkeit fällt weg. Am Gelenk ist $N = 0$, im Stab bleibt
        nur der Anteil aus Längslasten. Die Biegung ist davon nicht betroffen.
  \item \textbf{Q-Gelenk:} Am Gelenk ist $Q = 0$, ohne Querlast also im ganzen Stab, und wegen
        $M' = Q$ bleibt von der Biegesteifigkeit nur ein konstantes Moment übrig. Die Normalkraft
        trägt der Stab ganz normal weiter.
\end{itemize}

%---------------------------------------------------------------------
\section{Bögen: Teilstäbe und Schnittgrößen über den Winkel}\label{sec:mathe-bogen}

\subsection{Geometrie und Teilstäbe}

Ein Bogen ist ein Kreisstück mit Radius $R$ über der Sehne $L$ vom Stabanfang zum Stabende. Sein
Öffnungswinkel ist $\alpha = 2\arcsin\frac{L}{2R}$, und jeder Punkt auf ihm hat einen Winkel
$\varphi \in [0, \alpha]$ ab dem Stabanfang. Für die Rechnung ersetzt das Programm den Bogen durch
$n$ gerade Teilstäbe (Obermenü Seite~3, Standard 20). Das erste und das letzte Teilstück liegen
\emph{tangential} am Bogen und sind so lang wie $R\alpha/n$, die übrigen Teilungspunkte liegen auf
dem Bogen. Die tangentialen Endstücke haben einen Zweck: Die Stabendkräfte des ersten Teilstücks
sind dann direkt Normal- und Querkraft \emph{am Bogen}, und von dort startet der Export.

Die Teilstäbe erben Steifigkeiten, Gelenke an den Bogenenden, Temperatur und die Lasten:

\begin{itemize}
  \item $q$ (radial) und $n$ (tangential) bekommt jeder Teilstab mit den Werten an seinen beiden
        Endpunkten, dazwischen linear. Eine Lastfunktion wertet das Programm an den
        Teilungspunkten aus, mit $x = \varphi$ im Winkelmodus des Rechners.
  \item Lasten je Bogenlänge ($p_X$, $p_Y$ ohne Projektion, $m$) werden je Teilstab mit
        $R\,\Delta\varphi / l$ skaliert ($l$ = Länge des Teilstabs). Damit trägt jeder Teilstab genau
        die Last seines Bogenstücks, und die Summe ist exakt, z.\,B.\ das Eigengewicht $gR\alpha$.
  \item Projizierte Lasten gehen je Grundriss- bzw.\ Höhenlänge des Teilstabs ein.
\end{itemize}

\textbf{Warum konstante radiale Lasten exakt sind.} Ein konstanter Druck $q$ senkrecht auf einen
beliebigen Weg von $A$ nach $B$ hat die Resultierende $q$ mal den um $90^\circ$ gedrehten Vektor
$B - A$. Sein Moment um einen Punkt $P$ ist $\frac q2\left(|B-P|^2 - |A-P|^2\right)$. Beides hängt nur
von den Endpunkten ab, nicht vom Weg dazwischen. Der Linienzug der Teilstäbe hat dieselben Endpunkte
wie der Bogen, also stimmen bei einem statisch bestimmten System die Lagerkräfte und die
Schnittgrößen an den Teilungspunkten auf dem Bogen exakt: Sie folgen dort allein aus dem
Gleichgewicht, und das sieht nur die Resultierenden. Wie biegsam die Stäbe sind, spielt dabei keine
Rolle. Bei einem statisch unbestimmten System hängt die Kraftverteilung dagegen auch von der
Steifigkeit ab, und der Linienzug ist etwas anders steif als der Bogen, dann gilt das nur bis auf
den Teilungsfehler (siehe \emph{Genauigkeit} unten).
Beim Eigengewicht stimmt dagegen nur die Summe exakt. Die Hebel weichen um Größenordnung
$(\alpha/n)^2$ ab, mit 20 Teilstäben bei einem Viertelkreis um etwa 0,04\,\%.

\subsection{Schnittgrößen über den Winkel}

Der Export schreibt die Schnittgrößen als Funktionen des Winkels $p$ ab Stabanfang. Dafür schneidet
er den Bogen bei $p$ und schreibt das Gleichgewicht des Stücks von $0$ bis $p$ auf. Am Anfang wirken
$N_A$, $Q_A$, $M_A$ aus der Rechnung (erstes, tangentiales Teilstück). Mit $\varepsilon = +1$, wenn
die gestrichelte Faser innen liegt ($R > 0$), und $\varepsilon = -1$, wenn sie außen liegt, gilt
ohne Last auf dem Bogen
\begin{align*}
  N(p) &= N_A\cos p + \varepsilon\,Q_A\sin p , \\
  Q(p) &= Q_A\cos p - \varepsilon\,N_A\sin p , \\
  M(p) &= M_A + Q_A\,R\sin p - \varepsilon\,N_A\,R\,(1-\cos p) .
\end{align*}
Die ersten beiden Zeilen drehen nur die Kraft vom Anfang in die Tangente bei $p$. In der dritten
stehen die Hebel: Der Schnittpunkt liegt um $R\sin p$ entlang der Anfangstangente und um
$R(1-\cos p)$ zum Mittelpunkt hin versetzt.
Hier sieht man auch, warum im Bogen Normalkraft und Biegung gekoppelt sind: $N_A$ hat am
gekrümmten Stab den Hebel $R(1-\cos p)$ und erzeugt ein Moment. Beim geraden Stab ($R \to \infty$
mit $R\,p = x$) geht dieser Hebel gegen $x^2/(2R) \to 0$, und $Q_A R\sin p$ wird zu $Q_A\,x$. Das
ist reines Gleichgewicht am unverformten Stab, also Theorie I.~Ordnung. Theorie II.~Ordnung hieße,
das Gleichgewicht am \emph{verformten} System aufzustellen (Knicken, Durchschlagen flacher Bögen),
das rechnet das Programm weder bei Bögen noch bei geraden Stäben.

Last auf dem Bogenstück kommt als Integral über den Lastwinkel $t$ von $0$ bis $p$ dazu (im Bogenmaß,
$ds = R\,dt$):
\begin{align*}
  \Delta N(p) &= -R\int_0^p \Bigl(\varepsilon\,q(t)\sin(p-t) + n(t)\cos(p-t)\Bigr)\,dt , \\
  \Delta Q(p) &= -R\int_0^p \Bigl(q(t)\cos(p-t) - \varepsilon\,n(t)\sin(p-t)\Bigr)\,dt , \\
  \Delta M(p) &= -R\int_0^p \Bigl(R\,q(t)\sin(p-t) - \varepsilon\,R\,n(t)\bigl(1-\cos(p-t)\bigr)
                + m(t)\Bigr)\,dt .
\end{align*}
Auch hier stehen die Hebel direkt da. Eine radiale Kraft beim Winkel $t$ geht durch den
Mittelpunkt, ihr Abstand zum Schnittpunkt bei $p$ ist $R\sin(p-t)$. Eine tangentiale Kraft liegt
auf der Tangente bei $t$, und der Schnittpunkt hat von dieser Tangente den Abstand $R(1-\cos(p-t))$.
Globale Linienlasten $p_X$, $p_Y$ gehen genauso ein, nur mit Hebeln aus den Koordinaten. Ihre
Richtung relativ zur Tangente ändert sich mit $p$.

\textbf{Winkelmodus.} Steht der Rechner auf DEG, laufen $p$ und $t$ in Grad. Dann ist
$ds = R\,\frac{\pi}{180}\,dt$, und genau diesen Faktor setzt das Programm vor die Integrale. Den
Modus erkennt es daran, ob $\cos(180) = -1$ ist. $\sin p$ und $\cos p$ rechnet der Rechner
ohnehin im eingestellten Modus, deshalb stimmen die Funktionen in jedem Modus.

\textbf{Genauigkeit.} Die Formeln sind das exakte Gleichgewicht am Kreisbogen. Ungenau sein können
nur die Startwerte $N_A$, $Q_A$, $M_A$, denn sie stammen aus der Rechnung mit Teilstäben. Bei einem
statisch bestimmten System mit Lasten, die die Teilstäbe exakt abbilden (Knotenlasten, konstante
radiale Last, projizierte Lasten), sind auch sie exakt. Sonst gehen die Abweichungen mit mehr
Teilstäben zurück: bei Eigengewicht, tangentialen Lasten und Lastfunktionen, und bei statisch
unbestimmten Systemen, in denen die Steifigkeit des Linienzugs etwas von der des Bogens abweicht.

%---------------------------------------------------------------------
\section{Starre Stäbe: Zwangsbedingungen und Nullraum}\label{sec:mathe-starr}

Warum nicht einfach $EA = 10^{15}$? Weil der Rechner mit etwa 16 gültigen Stellen rechnet. Steht
in derselben Matrix eine Feder mit $c = 1$ neben $10^{15}$, bleibt von der Feder nach der
Zerlegung praktisch nichts übrig. Das ist der Zahlenmüll aus Abschnitt~\ref{sec:starr}. Das
Programm macht es deshalb exakt.

\textbf{1. Zwangsbedingungen.} Ein starrer Anteil bekommt in der Steifigkeitsmatrix die Steifigkeit
\emph{null}. Stattdessen schreibt das Programm für jeden starren Stab auf, was \enquote{starr}
bedeutet, und zwar in den lokalen Endverschiebungen:
\begin{align*}
  \text{dehnsteif:}\quad & u_b - u_a = \varepsilon_T L, \\
  \text{biegesteif:}\quad & \varphi_a - \varphi_b = -\kappa_T L, \qquad
                            w_b - w_a + L\,\varphi_a = -\kappa_T \frac{L^2}{2} .
\end{align*}
Ohne Temperatur steht rechts null: gleiche Länge, gleiche Endverdrehung, Enden auf einer Geraden.
Alle Bedingungen zusammen, ins globale System gedreht, ergeben $C\,u = g$.

\textbf{2. Kern der Matrix (=Nullraum).} $C$ wird mit Gauß-Jordan auf reduzierte Stufenform gebracht. Jede Zeile hat
dann einen \emph{abhängigen} Freiheitsgrad (die Pivotspalte), der durch die übrigen, \emph{freien}
Freiheitsgrade ausgedrückt wird. Damit lässt sich jede erlaubte Verschiebung schreiben als
\[
  u = T\,u_r + u_0 , \qquad C\,T = 0 .
\]
Die Spalten von $T$ spannen den Nullraum bzw. den Kern von $C$ auf, also genau die Bewegungen, die die starren
Stäbe mitmachen. $u_r$ sind nur noch die freien Freiheitsgrade, $u_0$ ist der Temperaturanteil.

\textbf{3. Projizieren.} Eingesetzt in $K\,u = F$ und von links mit $T^T$ multipliziert:
\begin{equation} \label{eq:LGSred}
  \left(T^T K\,T\right) u_r = T^T\left(F - K\,u_0\right).
\end{equation}
Das ist nichts anderes als das Prinzip der virtuellen Verschiebungen: Als virtuelle
Verschiebungen sind nur noch solche zugelassen, die die Zwangsbedingungen erfüllen, also Spalten
von $T$. Das neue System ist kleiner als das alte und genauso gut konditioniert wie ein System
ohne starre Stäbe, weil nirgends eine riesige Zahl steht.

\textbf{4. Kräfte in den starren Stäben.} Nach Schritt~3 kennt man die Verschiebungen $u$, aber
noch keine Kraft in einem starren Stab, denn ihr Anteil in $K$ ist null. Setzt man $u$ in das
Gleichgewicht der Knoten ein, bleibt ein Rest
\[
  r = F - K\,u .
\]
Das sind die Knotenkräfte, die weder die elastischen Stäbe noch die Federn aufnehmen. Diesen Rest
müssen die starren Stäbe tragen.

Das widerspricht nicht $K\,u = F$. In $K$ fehlen die starren Stäbe, denn ihre Kraft lässt sich
nicht aus $u$ ausrechnen: Ein starrer Stab verformt sich bei jeder Kraft um null. Das vollständige
Gleichgewicht der Knoten lautet
\[
  K\,u + (\text{Kräfte der starren Stäbe}) = F ,
\]
und $r$ ist genau der fehlende Anteil. Das kennt man von den Lagern: Auch an einem gehaltenen
Freiheitsgrad bleibt ein Rest, und der ist die Auflagerkraft. Ein starrer Stab ist so etwas wie ein
inneres Lager, das statt einer absoluten eine relative Verschiebung festhält.

\textbf{Woher die $\lambda$ kommen.} Schritt~3 hat nicht $K\,u = F$ gelöst, sondern nur
\eqref{eq:LGSred}. Alles auf eine Seite gebracht und $u = u_0 + T\,u_r$ eingesetzt:
\begin{align*}
  0 &= T^T F - T^T K\,u_0 - T^T K\,T\,u_r \\
    &= T^T\bigl(F - K\,(u_0 + T\,u_r)\bigr) = T^T\,(F - K\,u) = T^T r .
\end{align*}
Die Zeilen von $T^T$ sind die Spalten von $T$, also eine Basis des Kerns von $C$. $T^T r = 0$ sagt
deshalb: $r$ steht auf jedem dieser Basisvektoren und damit auf dem ganzen Kern von $C$ senkrecht.

Daraus folgt, dass sich $r$ aus den Zeilen von $C$ zusammensetzen lässt. Das ist reine lineare
Algebra aus \textbf{HM1}. $C$ hat $m$ Zeilen (eine je Zwangsbedingung) und $n$ Spalten (eine je
Freiheitsgrad). Im $\mathbb R^n$ liegen damit zwei Unterräume:
\begin{itemize}
  \item $\operatorname{Kern} C$: alle $x$ mit $C\,x = 0$,
  \item $\operatorname{Bild} C^T$: alle Vektoren $C^T\lambda$, also alle Kombinationen der Zeilen
        von $C$.
\end{itemize}
Gezeigt wird $r \in \operatorname{Bild} C^T$, und zwar in vier Schritten.

\emph{(a) $r$ steht senkrecht auf $\operatorname{Kern} C$.} Das haben wir gerade gesehen.

\emph{(b) $\operatorname{Bild} C^T$ steht senkrecht auf $\operatorname{Kern} C$.} Für ein $x$ aus
$\operatorname{Kern} C$ und einen beliebigen Vektor $C^T\lambda$ gilt
\[
  (C^T\lambda)^T x = \lambda^T C\,x = \lambda^T 0 = 0 .
\]
Das erste Gleichheitszeichen ist nur die Rechenregel $(AB)^T = B^T A^T$. Insbesondere liegt außer
der Null kein Vektor in beiden Räumen, denn er müsste auf sich selbst senkrecht stehen.

\emph{(c) Zusammen füllen beide Räume den ganzen $\mathbb R^n$.} Nach der Dimensionsformel ist
$\dim\operatorname{Kern} C + \dim\operatorname{Bild} C = n$, und weil Zeilenrang gleich Spaltenrang
ist, gilt $\dim\operatorname{Bild} C^T = \dim\operatorname{Bild} C$. Also ist
$\dim\operatorname{Kern} C + \dim\operatorname{Bild} C^T = n$. Weil sich die beiden Räume nur in der
Null treffen, lässt sich jeder Vektor des $\mathbb R^n$ eindeutig als Summe aus einem Teil in
$\operatorname{Kern} C$ und einem Teil in $\operatorname{Bild} C^T$ schreiben.

\emph{(d) $r$ liegt in $\operatorname{Bild} C^T$.} Man zerlegt $r = r_K + r_B$ mit $r_K$ in
$\operatorname{Kern} C$ und $r_B$ in $\operatorname{Bild} C^T$. Nach (a) und (b) stehen $r$ und
$r_B$ senkrecht auf $\operatorname{Kern} C$, also auch $r_K = r - r_B$. $r_K$ liegt aber selbst in
$\operatorname{Kern} C$, steht also senkrecht auf sich selbst, und das kann nur der Nullvektor.
Damit ist $r = r_B$.

Ein Beispiel mit Zahlen: $30$ Freiheitsgrade und ein starrer Stab, der dehn- und biegesteif ist
($3$ unabhängige Zeilen). $\operatorname{Kern} C$ ist $27$-dimensional, das sind die $27$
Unbekannten $u_r$ in Schritt~3. $\operatorname{Bild} C^T$ ist $3$-dimensional, das sind die drei
$\lambda$.

Es gibt also Zahlen $\lambda_k$, eine je Zwangsbedingung, mit
\[
  r = C^T \lambda = \sum_k \lambda_k\,C_k^T ,
\]
wobei $C_k$ die $k$-te Zeile von $C$ ist. Weil $r$ nach dem vollständigen Gleichgewicht genau die
Kräfte der starren Stäbe sind, heißt das: Die starren Stäbe wirken mit $C^T\lambda$ auf die Knoten,
und jede Zwangsbedingung~$k$ trägt $\lambda_k\,C_k^T$ dazu bei. Das vollständige Gleichgewicht
lautet damit
\[
  K\,u + C^T\lambda = F .
\]
Mit $T^T$ multipliziert fällt $\lambda$ wegen $T^T C^T = (C\,T)^T = 0$ heraus, das ist Schritt~3.
Die übrigen Gleichungen erfüllt Schritt~4 mit der passenden Wahl von $\lambda$. Die $\lambda$ sind
übrigens die Lagrange-Multiplikatoren, die man bekommt, wenn man das Minimum der potentiellen
Energie $\frac12\,u^T K\,u - F^T u$ unter der Nebenbedingung $C\,u = g$ sucht.

\textbf{Mechanisch gelesen.} Das Skalarprodukt aus Kraft und Verschiebung ist die Arbeit,
\enquote{senkrecht} heißt also \enquote{keine Arbeit}. (a) sagt dann: Die Kräfte der starren Stäbe
leisten auf keiner Bewegung Arbeit, die die starren Stäbe erlauben. Genau das zeichnet eine
Zwangskraft aus. Eine reibungsfreie Führung leistet auf der Bewegung, die sie zulässt, auch keine
Arbeit, sie verhindert nur alle anderen. Physikalisch steckt diese Annahme schon in Schritt~3: Dort
sind als virtuelle Verschiebungen nur solche zugelassen, die die Zwangsbedingungen erfüllen. Alles
danach ist reine Mathematik.

Eindeutig sind die $\lambda$ nur, wenn die Zeilen von $C$ linear unabhängig sind, also
$\operatorname{Kern} C^T = \{0\}$ gilt. Dann gehört zu jeder Kraft aus $\operatorname{Bild} C^T$
genau ein $\lambda$. Ist $\operatorname{Kern} C^T$ größer, gibt es Kräfte in den starren Stäben,
die auf keinen Knoten wirken. Das sind die Eigenspannungszustände weiter unten.

\textbf{Warum die $\lambda$ Schnittgrößen sind.} Das folgt aus dem Gleichgewicht des einzelnen
starren Stabes, und zwar mit demselben Werkzeug wie eben: Kern, Bild und Dimensionsformel, diesmal
für einen Stab allein. Im lokalen System hat ein Stab die sechs Endverschiebungen
$v = (u_a, w_a, \varphi_a, u_b, w_b, \varphi_b)$ und sechs Endkräfte $s$ in denselben Richtungen.
Die drei Zeilen eines biege- und dehnsteifen Stabes bilden die $3 \times 6$-Matrix
\[
  C_\ell = \begin{pmatrix} -1 & 0 & 0 & 1 & 0 & 0 \\ 0 & 0 & 1 & 0 & 0 & -1 \\ 0 & -1 & L & 0 & 1 & 0 \end{pmatrix}
  \qquad
  \begin{matrix} \text{Längszeile: } u_b - u_a \\ \text{Drehzeile: } \varphi_a - \varphi_b \\
                 \text{Querzeile: } w_b - w_a + L\,\varphi_a \end{matrix}
\]

\emph{(a) Starr heißt: nur Starrkörperbewegungen.} Ein Stab kann sich auf drei Arten bewegen, ohne
sich zu verformen: längs verschieben, quer verschieben und um das Ende $a$ drehen. In
Endverschiebungen sind das die Spalten von
\[
  R = \begin{pmatrix} 1 & 0 & 0 & 1 & 0 & 0 \\ 0 & 1 & 0 & 0 & 1 & 0 \\ 0 & 0 & 1 & 0 & -L & 1 \end{pmatrix}^T
\]
(bei der Drehung um $\varphi = 1$ wandert das Ende $b$ wegen $\varphi = -w'$ um $w_b = -L$).
Einsetzen zeigt $C_\ell R = 0$: Starrkörperbewegungen verformen nichts, $\operatorname{Bild} R
\subseteq \operatorname{Kern} C_\ell$. Beide Matrizen haben Rang~$3$, nach der Dimensionsformel ist
also $\dim \operatorname{Kern} C_\ell = 6 - 3 = 3 = \dim \operatorname{Bild} R$, und damit
$\operatorname{Kern} C_\ell = \operatorname{Bild} R$. Die Zwangsbedingungen lassen genau die
Starrkörperbewegungen zu, nicht mehr und nicht weniger.

\emph{(b) Gleichgewicht heißt: keine Arbeit auf Starrkörperbewegungen.} Zuerst ein Stab ohne
Stablasten, die kommen in (d) dazu. Endkräfte $s$ sind genau dann im Gleichgewicht, wenn
\[
  R^T s = \begin{pmatrix} s_1 + s_4 \\ s_2 + s_5 \\ s_3 - L\,s_5 + s_6 \end{pmatrix} = 0 ,
\]
das sind Summe der Kräfte längs, Summe der Kräfte quer und Summe der Momente um $a$. Anders gelesen
ist das das PvV für einen starren Körper aus \textbf{TM1}: Auf jeder Starrkörperbewegung leisten die
Kräfte keine Arbeit. Die Endkräfte, die für sich im Gleichgewicht sind, bilden also den Raum
$\operatorname{Kern} R^T$.

\emph{(c) Dieser Raum wird genau von den Zeilen aufgespannt.} Wie beim Nullraum oben in zwei
Schritten: $R^T C_\ell^T = (C_\ell R)^T = 0$, also $\operatorname{Bild} C_\ell^T
\subseteq \operatorname{Kern} R^T$. Dimension: $\dim \operatorname{Bild} C_\ell^T = \operatorname{Rang}
C_\ell = 3$ und $\dim \operatorname{Kern} R^T = 6 - \operatorname{Rang} R = 3$. Also
\[
  \operatorname{Kern} R^T = \operatorname{Bild} C_\ell^T .
\]
Jede Gruppe von Endkräften im Gleichgewicht ist eine Kombination der drei Zeilen, und weil die
Zeilen linear unabhängig sind, auf genau eine Weise.

\emph{(d) Anwendung auf den starren Stab.} Auf den starren Stab wirken zwei Dinge: die Lasten auf
dem Stab ($q$, $n$, $m$) und die Endkräfte $s$, mit denen die Knoten auf seine Enden wirken. $s$
sind also nur die Endkräfte, die Lasten stehen getrennt. Zusammen sind beide im Gleichgewicht.
Die Festeinspannkräfte $p$ sind ebenfalls Endkräfte, die dieselben Lasten im Gleichgewicht
halten. Dass sie zu einem beidseitig eingespannten Stab gehören, spielt hier keine Rolle. Man
braucht nur \emph{irgendeinen} Satz Endkräfte, der die Lasten allein ausgleicht, und $p$ hat das
Programm aus dem WGV ohnehin. Zieht man die beiden Gleichgewichtsgruppen
\enquote{Lasten $+ s$} und \enquote{Lasten $+ p$} voneinander ab, fallen die Lasten heraus: $s - p$
ist eine Gruppe reiner Endkräfte im Gleichgewicht, also genau der Fall aus (b) ohne Stablasten.
Ausgeschrieben: Nach (b) gilt $R^T (s - p) = 0$, also liegt
$s - p$ in $\operatorname{Kern} R^T$, und das ist nach (c) $\operatorname{Bild} C_\ell^T$. Es gibt
also genau ein $\lambda_\ell \in \mathbb R^3$ mit $s - p = C_\ell^T \lambda_\ell$, umgestellt
\[
  s = p + C_\ell^T \lambda_\ell .
\]
Ins globale System gedreht und über alle starren Stäbe summiert sind das genau die Knotenkräfte
$C^T\lambda$ von oben, die $\lambda_\ell$ eines Stabes sind seine Einträge in $\lambda$. Ist ein Stab
nur dehnsteif, fehlen Dreh- und Querzeile, die Biegung liefert dann wie bei jedem Stab $k\,v$.
Ein N-Gelenk streicht die Längszeile, ein Q-Gelenk die Querzeile. Das hat zwei Gründe, die
zusammenpassen. Kinematisch kann sich das Stabende in der freigegebenen Richtung gegenüber dem
Knoten verschieben. Die Zeile würde aber den Knoten selbst festhalten, also mehr verlangen, als der
Stab tut. Und weil N- und Q-Gelenke herauskondensiert werden, hat das Stabende in dieser Richtung
keinen eigenen Freiheitsgrad, mit dem man die Zeile richtig schreiben könnte. Statisch ist die
zugehörige Schnittgröße am Gelenk null. Die Lasten trägt schon $p$, das für Stäbe mit N- oder
Q-Gelenk passend umverteilt ist, ein $\lambda$ für diese Zeile wäre also ohnehin null. Ein
M-Gelenk streicht dagegen nichts: Das Stabende hat dort einen eigenen Drehfreiheitsgrad
(Abschnitt~\ref{sec:mathe-gelenke}). Die Drehzeile verbindet die Verdrehung des Stabes mit diesem
Freiheitsgrad statt mit dem Knoten, und $M = 0$ am Gelenk folgt aus dessen Gleichung. $p$ darf man dabei unverändert nehmen, weil die
Festeinspannkräfte nach Abschnitt~\ref{sec:mathe-festeinspann} gar nicht von $EI$ abhängen.

\emph{(e) Was die drei Zahlen bedeuten.} Für jede Zeile liest man den Schnittgrößenverlauf des
unbelasteten Stabes direkt ab, mit $N(0) = -s_1$, $Q(0) = -s_2$, $M(0) = -s_3$ wie bei jedem Stab:
\begin{align*}
  \text{Längszeile:}\quad & s = \lambda_N\,(-1, 0, 0, 1, 0, 0)^T: && N(x) = \lambda_N, \\
  \text{Drehzeile:}\quad  & s = \lambda_M\,(0, 0, 1, 0, 0, -1)^T: && M(x) = -\lambda_M, \\
  \text{Querzeile:}\quad  & s = \lambda_Q\,(0, -1, L, 0, 1, 0)^T: && Q(x) = \lambda_Q, \quad
                                                                     M(x) = -\lambda_Q\,(L - x) .
\end{align*}
$\lambda_N$ ist also die Normalkraft, $\lambda_Q$ die Querkraft und $-\lambda_M$ das Moment am Ende
$b$. Zusammen $M(x) = -\lambda_M - \lambda_Q\,(L - x)$, dazu kommt der Anteil aus $p$. Mechanisch
passt das: Die Arbeit von $\lambda_k\,C_k^T$ auf den Endverschiebungen ist $\lambda_k\,(C_k\,v)$,
die Normalkraft arbeitet also auf der Längenänderung, das Moment auf der gegenseitigen Verdrehung
und die Querkraft darauf, wie weit das Ende $b$ von der Tangente am Ende $a$ abweicht. Das ist mit
dem Satz in Abschnitt~\ref{sec:starr} gemeint, die Schnittgrößen starrer Stäbe würden \enquote{aus
dem Gleichgewicht bestimmt}: Sie sind die $\lambda$, und die kommen allein aus dem Rest $r$, also aus
dem Gleichgewicht der Knoten, nicht aus einer Steifigkeit.

\textbf{Wie das Programm sie ausrechnet.} $r = C^T\lambda$ hat eine Gleichung je Freiheitsgrad, aber
nur eine Unbekannte je Zwangsbedingung, also meist mehr Gleichungen als Unbekannte. Weil $r$ nach
Schritt~3 eine Kombination der Zeilen von $C$ ist, geht es trotzdem exakt auf. Von links mit $C$
multipliziert wird daraus ein kleines quadratisches System
\[
  (C\,C^T)\,\lambda = C\,r ,
\]
das eindeutig lösbar ist, solange die Zeilen von $C$ linear unabhängig sind.

\textbf{Starre Teile, die sich gegenseitig halten.} Beispiele sind zwei dehnstarre Stäbe
nebeneinander zwischen denselben zwei Knoten oder ein biegestarrer Zweifeldträger (zwei starre
Stäbe über drei Lager). Dann sind einige Zeilen von $C$ linear abhängig. Das merkt schon die
Elimination in Schritt~2: Eine solche Zeile wird null, nachdem die früheren Pivotspalten abgezogen
sind. Das Programm protokolliert dabei, welche Kombination der ursprünglichen Zeilen null ergeben
hat:
\[
  \sum_k \mu_k\,C_k = 0 \qquad\Longleftrightarrow\qquad C^T\mu = 0 .
\]
Physikalisch ist $\mu$ ein Satz von Kräften in den starren Stäben, der auf keinen Knoten eine Kraft
ausübt, ein \emph{Eigenspannungszustand}. Bei den zwei Stäben nebeneinander ist das Zug~$1$ im
einen und Druck~$1$ im anderen.

Für die Kinematik ist eine abhängige Zeile harmlos, ohne Temperatur ist sie automatisch erfüllt. Für
das Gleichgewicht heißt sie: $r = C^T\lambda$ hat unendlich viele Lösungen
\[
  \lambda = \lambda_p + \sum_i t_i\,\mu_i .
\]
Das ist die Aufteilung, die man von jedem LGS kennt: eine partikuläre Lösung plus beliebige
Lösungen der homogenen Gleichung $C^T\mu = 0$. Und es ist genau die Aufteilung aus dem KGV in
\textbf{Statik}:
\begin{itemize}
  \item $\lambda_p$ rechnet das Programm mit $(C\,C^T)\,\lambda = C\,r$ nur über die unabhängigen
        Zeilen, die abhängigen bekommen $\lambda = 0$. Das ist, als würde man die überzähligen
        Zwangsbedingungen aufschneiden: ein statisch bestimmtes Hauptsystem mit seinem
        $0$-Zustand.
  \item Die $\mu_i$ sind die $1$-Zustände. Jedes $\mu_i$ hat an seiner abhängigen Zeile eine $1$,
        dazu kommen die Kräfte in den übrigen starren Stäben, die diese $1$ im Gleichgewicht halten.
  \item Die $t_i$ sind die statisch Überzähligen $X_i$.
\end{itemize}
Wie im KGV legt das Gleichgewicht die $X_i$ nicht fest. Dort kommen sie aus der Verträglichkeit,
und die braucht Nachgiebigkeiten. Genau die haben starre Stäbe nicht: Jede Wahl der $t_i$ ergibt
dieselben Verschiebungen, und die Starrkörpermechanik kann nicht sagen, welche stimmt. Das hängt
davon ab, wie starr die Stäbe \emph{im Verhältnis zueinander} wirklich sind.

\textbf{Der Grenzwert und die Ergänzungsenergie.} Das Programm trifft dieselbe Annahme wie der
Grenzwert im Rand-LGS: Alle starren Stäbe bekommen $EA = EI = 1/\varepsilon$ mit
$\varepsilon \to 0$. Dann verformen sie sich ein kleines bisschen. In den Größen, die die Zeilen
von $C$ messen, ist ihre elastische Verformung
\[
  C\,u - g = \varepsilon\,W\,\lambda ,
\]
wobei $W$ die Nachgiebigkeit der starren Stäbe mit $EA = EI = 1$ ist (blockweise je Stab, siehe
unten). Das ist das Elastizitätsgesetz der fast starren Stäbe.

Für das Folgende muss man zwei Dinge auseinanderhalten: $u$ sind die Knotenverschiebungen, $C\,u$
sind die Verformungen der starren Stäbe, gemessen in den Größen der Zwangsbedingungen
(Längenänderung, gegenseitige Verdrehung, Abweichung von der Tangente, siehe (e)). Die Stabkräfte
$\lambda$ arbeiten auf $C\,u$, ihre Knotenkräfte $C^T\lambda$ auf $u$, und beides ist dieselbe
Arbeit: $\lambda^T (C\,u) = (C^T\lambda)^T u$, das ist wieder nur $(AB)^T = B^T A^T$.

Jetzt kommt die Verträglichkeit, mit dem Arbeitssatz wie im KGV. Ein Eigenspannungszustand $\mu_i$
übt auf keinen Knoten eine Kraft aus, $C^T\mu_i = 0$. Seine Arbeit auf den wirklichen Verformungen
der starren Stäbe ist deshalb null:
\[
  \mu_i^T (C\,u) = (C^T\mu_i)^T u = 0^T u = 0 .
\]
Multipliziert man das Elastizitätsgesetz $C\,u - g = \varepsilon\,W\,\lambda$ von links mit
$\mu_i^T$, steht links genau diese Arbeit, und es bleibt
\[
  \underbrace{\mu_i^T C\,u}_{=\,0} - \mu_i^T g = \varepsilon\,\mu_i^T W\,\lambda
  \qquad\Longrightarrow\qquad
  \varepsilon\,\mu_i^T W\,\lambda + \mu_i^T g = 0 .
\]
Die Temperatur $g$ bleibt dabei stehen. Ohne Temperatur fällt der zweite Term weg, das $\varepsilon$ kürzt sich heraus, und mit
$\lambda = \lambda_p + \sum_j t_j\,\mu_j$ bleibt
\[
  \sum_j \underbrace{\mu_i^T W\,\mu_j}_{\delta_{ij}}\;t_j
  = -\underbrace{\mu_i^T W\,\lambda_p}_{\delta_{i0}} .
\]
Das sind die Elastizitätsgleichungen des KGV (hier ist $\delta_{ij}$ wieder eine $\delta$-Zahl und
kein Kronecker-Delta). Weil $\varepsilon$ herausfällt, hängt das Ergebnis nicht davon ab,
\emph{wie} starr die Stäbe sind, nur davon, dass sie alle gleich starr sind. Deshalb existiert der
Grenzwert. Passt die Temperatur dagegen nicht zu den abhängigen Zeilen ($\mu_i^T g \ne 0$), muss
$\varepsilon\,\mu_i^T W \lambda$ einen festen Wert annehmen, und für $\varepsilon \to 0$ wird
$\lambda$ unendlich groß. Das ist die unendliche Zwangskraft, zu der das Programm einen Hinweis
zeigt.

\textbf{Wo die Ergänzungsenergie hineinspielt.} Die Gleichungen oben sind genau die Bedingung dafür,
dass
\[
  \Pi^*(t) = \tfrac12\,\lambda^T W \lambda , \qquad \lambda = \lambda_p + \sum_i t_i\,\mu_i ,
\]
minimal wird, denn $\partial\Pi^*/\partial t_i = \mu_i^T W \lambda$. Das ist der Satz von Menabrea:
Unter allen Kraftzuständen im Gleichgewicht stellt sich der mit der kleinsten Ergänzungsenergie
ein. \enquote{Kleinste Ergänzungsenergie} ist dabei nur eine andere Schreibweise für
\enquote{verträglich}, das KGV in \textbf{Statik} rechnet mit dem Prinzip der virtuellen Kräfte
genau dasselbe. Die Ergänzungsenergie ist also kein zusätzliches Prinzip, sondern die
Verträglichkeit in Energieform. Genau genommen ist die Ergänzungsenergie der fast starren Stäbe
$\varepsilon\,\Pi^*$, denn ihre Nachgiebigkeit ist $\varepsilon\,W$ und nicht $W$. Ein fester
Faktor ändert aber nichts daran, wo das Minimum liegt, deshalb steht $\varepsilon$ oben gar nicht
erst in $\Pi^*$.

\textbf{$W$ ausrechnen} geht wie eine $\delta$-Zahl mit der Koppeltafel, mit $EA = EI = 1$ und für
jeden Stab für sich (verschiedene Stäbe koppeln nicht):
\begin{itemize}
  \item Längszeile: $N = 1$ konstant, $\int N^2\,dx = L$.
  \item Drehzeile: $M = 1$ konstant, $\int M^2\,dx = L$.
  \item Querzeile: $M$ linear von $L$ auf $0$, $\int M^2\,dx = L^3/3$.
  \item Dreh- mal Querzeile: $\int 1 \cdot (L - x)\,dx = L^2/2$. Längs- mal Biegezeile: $0$.
\end{itemize}
Die Querkraftverformung fehlt, wie überall im Programm (Euler-Bernoulli). Dass $EA$ und $EI$ beide
$1$ sind, ist eine Annahme und keine Physik. Sie spielt keine Rolle, solange in einem
Eigenspannungszustand nur dehnstarre oder nur biegestarre Zeilen vorkommen, denn dann kürzt sich der
gemeinsame Faktor. Mischen sich beide, hängt die Aufteilung vom angenommenen Verhältnis $EA : EI$
ab.

So teilen sich zwei dehnstarre Stäbe nebeneinander eine Last je zur Hälfte, und ein starrer
Zweifeldträger unter Gleichlast hat die bekannten Lagerkräfte $\frac38\,ql$, $\frac{10}8\,ql$,
$\frac38\,ql$, genau wie mit dem KGV und dem Stützmoment als $X_1$. Eine gelbe Box weist darauf
hin, dass diese Aufteilung angenommen ist, aber nur, wenn sie das Ergebnis auch ändert. Ist
$\lambda_p$ auf allen Zeilen null, die über $W$ mit einem $\mu_i$ gekoppelt sind, dann ist
$\delta_{i0} = \mu_i^T W \lambda_p = 0$ für \emph{jede} Wahl der Nachgiebigkeiten, also $t_i = 0$, egal wie
starr die Stäbe im Verhältnis zueinander sind. Ein Beispiel ist eine starre Platte, die oben und
unten auf Loslagern steht: Die Längskraft zwischen den beiden Lagern ist unbestimmt, aber ohne Last
in Plattenrichtung immer null.

\textbf{Das Gesamtbild.} Beide Aufteilungen in partikulär und homogen gehören zusammen. Die
Kinematik $C\,u = g$ und das Gleichgewicht $C^T\lambda = r$ benutzen dieselbe Matrix, einmal so und
einmal transponiert:

\begin{center}
\begin{tabular}{@{}p{2.9cm}p{5.9cm}p{5.9cm}@{}}
  \toprule
   & \textbf{Kinematik} $C\,u = g$ & \textbf{Gleichgewicht} $C^T\lambda = r$ \\
  \midrule
  partikulär       & $u_0$ (Temperatur) & $\lambda_p$ (Hauptsystem) \\
  homogen          & $T\,u_r$: Bewegungen, die die starren Stäbe erlauben
                   & $\sum_i t_i\,\mu_i$: Eigenspannungszustände \\
  festgelegt durch & Gleichgewicht (Schritt~3) & Verträglichkeit \\
  als Energie      & kleinste potentielle Energie & kleinste Ergänzungsenergie \\
  wie im           & WGV & KGV \\
  \bottomrule
\end{tabular}
\end{center}

Der Kern von $C$ sind die Bewegungen, die die starren Teile frei lassen. Der Kern von $C^T$ besteht
aus den Eigenspannungszuständen $\mu_i$, seine Dimension ist der Grad der statischen
Unbestimmtheit der starren Teile untereinander, also die Anzahl der $t_i$ (im KGV der $X_i$).

Die rote Box \emph{Starrmodus \& statisch unbestimmt} kommt übrigens schon dann, wenn das System
statisch unbestimmt ist und alle Stäbe $EA, EI \ge 10^7$ haben, also auch mit den
Standardwerten.

%---------------------------------------------------------------------
\section{Das Rand-LGS: die Biegelinie wie von Hand}\label{sec:mathe-randlgs}

Der Export \emph{Neue Biegelinie} rechnet unabhängig vom Weggrößenverfahren genau das, was ihr in
\textbf{TM2} von Hand macht, nur für jedes System und mit allen Bedingungen auf einmal.

\subsection{Ansatz je Stab}

Mit $\varphi = -w'$, $Q' = -q$, $M' = Q - m$ und $N' = -n$ setzt das Programm für jeden Stab an

\[
  EI\,w(x) = w_p(x) + c_1\frac{x^3}{6} + c_2\frac{x^2}{2} + c_3 x + c_4, \qquad
  EA\,u(x) = u_p(x) + C_1 x + C_2 .
\]

Das sind sechs Unbekannte je Stab. Die Partikulärlösungen baut es nicht aus einer Tabelle,
sondern als Integral, das der CAS-Kern exakt auswertet:

\begin{align*}
  w_p(x) &= \int_0^x q(t)\,\frac{(x-t)^3}{6}\,dt + \int_0^x m(t)\,\frac{(x-t)^2}{2}\,dt
            - EI\,\kappa_T\,\frac{x^2}{2}, \\
  u_p(x) &= -\int_0^x n(t)\,(x-t)\,dt + EA\,\varepsilon_T\,x .
\end{align*}

Die Idee dahinter: Man denkt sich die Streckenlast in viele kleine Einzellasten $q(t)\,dt$ an den
Stellen $t$ zerlegt. Bei einer Einzellast $1$ an der Stelle $t$ springt die Querkraft dort um $-1$
(wegen $Q' = -q$), und ab $t$ kommt der Anteil $EI\,w = (x-t)^3/6$ dazu. Das ist $c_1 x^3/6$ mit
$c_1 = 1$, nur um $t$ verschoben, davor ist der Anteil null. Weil alles linear ist, addieren sich
die Beiträge aller kleinen Einzellasten, und aus der Summe wird das Integral. Genauso gibt ein
Einzelmoment $1$ bei $t$ den Anteil $(x-t)^2/2$ in $EI\,w$ und eine Längskraft $1$ den Anteil
$-(x-t)$ in $EA\,u$. In der Mathematik heißt so ein
Integral \emph{Faltung}. Deshalb funktioniert es für jede Lastfunktion, für Trapezlasten und für
die globalen Linienlasten (die vorher auf Quer- und Längsrichtung des Stabes zerlegt werden). Und deshalb heißt die Integrationsvariable \code{t}, und
ihr könnt \code{t} nicht als Symbol verwenden.

Aus dem Ansatz folgen die Größen an beiden Stabenden als \emph{lineare} Ausdrücke in den
Konstanten, am Anfang z.\,B.\ $w = c_4/EI$, $\varphi = -c_3/EI$, $M = -c_2$, $Q = -c_1$,
$u = C_2/EA$, $N = C_1$.

\subsection{Bedingungen an den Knoten}

An jedem Knoten schreibt das Programm die Bedingungen, die ihr auch aufschreiben würdet, nur
systematisch:

\begin{itemize}
  \item \textbf{Kinematik:} Alle Stabenden ohne N- bzw.\ Q-Gelenk haben dieselbe
        Knotenverschiebung $\vec U$, also $u_e = \vec U\cdot\vec t_e$ und
        $w_e = \vec U\cdot\vec n_e$, alle Enden ohne M-Gelenk dieselbe Verdrehung. $\vec U$ selbst
        wird durch das erste passende Stabende ausgedrückt, damit keine zusätzlichen Unbekannten
        entstehen.
  \item \textbf{Gelenke:} Ein M-Gelenk gibt $M_e = 0$, ein N-Gelenk $N_e = 0$, ein Q-Gelenk
        $Q_e = 0$.
  \item \textbf{Lager:} In einer gehaltenen Richtung $\vec U\cdot\vec e = 0$, in einer freien
        Richtung das Gleichgewicht der Stabendkräfte, Federn und Lasten in dieser Richtung. Die
        Richtungen sind $\vec\xi = (\cos\alpha, -\sin\alpha)$ und $\vec\eta = (\sin\alpha, \cos\alpha)$,
        so gehen gedrehte Lager ohne Sonderfall.
  \item \textbf{Drehung:} Einspannung $\varphi = 0$, sonst $\sum M_e - c_m\,\varphi + M = 0$.
\end{itemize}

Am Ende muss die Anzahl der Gleichungen genau $6 \cdot (\text{Stäbe})$ sein (plus zwei je Knoten,
an dem alle Stäbe über N- oder Q-Gelenke hängen, siehe unten). Sonst kommt \emph{nicht quadratisch}.

\textbf{Warum solche Knoten eigene Verschiebungen brauchen.} Normalerweise wird $\vec U$ durch ein
Stabende ohne N- und Q-Gelenk ausgedrückt: Dessen $u_e$ und $w_e$ sind die beiden Komponenten von
$\vec U$ in Stab- und Querrichtung, damit ist $\vec U$ ganz bekannt. Ein Stabende mit N-Gelenk darf
sich aber längs gegen den Knoten verschieben, es liefert nur noch $w_e = \vec U\cdot\vec n_e$. Eines
mit Q-Gelenk liefert nur $u_e = \vec U\cdot\vec t_e$. Hat jedes Stabende am Knoten so ein Gelenk,
gibt es kein Ende, das beide Komponenten festlegt. Dann bekommt $\vec U$ die eigenen Unbekannten
$u_x(K)$, $u_y(K)$, und jedes Ende schreibt seine Bedingung gegen $\vec U$. Normalerweise
\enquote{verbraucht} das Bezugsende seine zwei Komponenten, um $\vec U$ festzulegen, und liefert
selbst keine Gleichung. Hier verbraucht keines etwas, daher die zwei Gleichungen mehr. Beispiel:
zwei Stäbe, die sich an einem Knoten mit je einem Q-Gelenk treffen. Jeder sagt nur, wie weit sich
der Knoten in seiner Längsrichtung bewegt, erst beide zusammen bestimmen $\vec U$.

\subsection{Lösen: erst wie von Hand, dann das CAS}

Ein LGS mit 30 symbolischen Unbekannten direkt in \code{simult} zu stecken, ist auf dem Rechner
langsam und liefert unübersichtliche Ausdrücke. Das Programm macht deshalb vorher, was ihr auch
machen würdet:

\begin{enumerate}
  \item Jede Gleichung mit nur \emph{einer} Unbekannten wird sofort aufgelöst und der Wert überall
        eingesetzt (z.\,B.\ $M = 0$ am Gelenk liefert direkt $c_2$). Das wiederholt sich, solange
        etwas passiert.
  \item Eine Unbekannte, die nur noch in einer einzigen Gleichung vorkommt, wird mit dieser
        Gleichung beiseitegelegt und am Schluss aus ihr berechnet.
  \item Gleichungen der Form $a\,x_j + b\,x_m = 0$ mit Zahlen $a$, $b$ werden zu $x_j = f\,x_m$
        und eingesetzt.
\end{enumerate}

Was übrig bleibt, zerfällt oft in unabhängige Blöcke (Unbekannte, die nie zusammen in einer
Gleichung stehen). Jeder Block wird einzeln mit \code{exact(simult(\dots))} gelöst, danach wird
rückwärts eingesetzt. Aus demselben Protokoll entsteht auch \code{randbed}: Eine Größe steht dort
nur dann als Wert, wenn ihre Konstanten in einem \emph{früheren} Schritt schon bekannt waren,
also genau in der Reihenfolge, in der ihr es auf dem Papier bekommen würdet.

\textbf{Die Aufteilung in \enquote{nur Biegelinie} und \enquote{nur Längslinie}} kommt aus
demselben Protokoll: Das Programm verfolgt für jede Konstante der Biegelinie zurück, aus welchen
Gleichungen sie bestimmt wurde. Tauchen darin keine Längslinien-Konstanten auf, die nicht schon
vorher feststanden, reichen die Bedingungen der Biegelinie allein.

Bei Rahmen mit dehnsteifen Stäben reicht das nicht. Dort bestimmt das LGS die Biegelinie zusammen
mit $C_1$ und $C_2$ der dehnsteifen Stäbe, weil z.\,B.\ die Seitenverschiebung $u$ des Riegels in
die Biegelinie der Stiele eingeht. Bei der direkten Methode ist $u = C_2$ und $N = C_1$ aber an
beiden Stabenden dieselbe Unbekannte (ohne Längslast und Temperatur). Das Programm nimmt deshalb
eine Bedingung, in der nur eine dieser Unbekannten steht, löst sie danach auf und setzt sie in die
anderen Bedingungen ein, genau wie von Hand. Mathematisch ist das ein Gauß-Schritt mit dieser
Unbekannten als Pivot, das LGS bleibt dabei regulär. Bleiben danach genau $4\cdot(\text{Stäbe})$
Bedingungen übrig, in denen nur noch Konstanten der Biegelinie stehen, genügen sie für $w$. Bei
dehnbaren Stäben ist $u(L) = (C_1 L + C_2)/EA$, dort hängt die Längsverschiebung am Stabende von
beiden Konstanten ab, und das Durchreichen geht nicht.

\subsection{Starre Stäbe im Rand-LGS: direkt oder Grenzwert}

Für einen starren Stab gibt es kein $EI\,w'' = -M$ mehr, trotzdem braucht das Programm sechs
Gleichungen. Zwei Wege:

\begin{itemize}
  \item \textbf{direkt} (Standard): Der biegesteife Stab bewegt sich als Starrkörper,
        $w(x) = c_4 + c_3 x - \kappa_T x^2/2$ und beim dehnsteifen $u(x) = C_2 + \varepsilon_T x$.
        $c_1$, $c_2$ und $C_1$ bleiben als reine Kraftgrößen stehen und werden aus dem
        Gleichgewicht bestimmt.

        Warum das geht: Auch beim elastischen Stab sind $c_1$, $c_2$ und $C_1$ eigentlich
        Kraftgrößen, nämlich die Schnittgrößen am Stabanfang, $Q(0) = -c_1$, $M(0) = -c_2$,
        $N = C_1$. In der Biegelinie stehen sie nur deshalb, weil die Krümmung aus dem Moment kommt:
        Sie gehen als $c_1 x^3/(6EI)$ und $c_2 x^2/(2EI)$ in $w$ ein, also immer geteilt durch
        $EI$. $c_3$ und $c_4$ dagegen sind echte Weggrößen ($c_3/EI = w'(0)$, $c_4/EI = w(0)$).
        Wird der Stab starr, $EI \to \infty$, verschwindet bei endlichen Kräften der Beitrag von
        $c_1$ und $c_2$ aus der Biegelinie, und $c_3$, $c_4$ schreibt man gleich ohne $EI$
        ($c_4 = w(0)$, $c_3 = w'(0)$). Die Schnittgrößen bleiben aber da: $Q(x)$, $M(x)$ und $N(x)$
        folgen aus $c_1$, $c_2$, $C_1$ und den Lasten allein über das Gleichgewicht am Stabstück
        ($Q' = -q$, $M' = Q$, $N' = -n$), ganz ohne Steifigkeit.

        Die Bedingungen an den Knoten sind dieselben wie beim elastischen Stab. Jede enthält aber
        entweder nur Kraftgrößen ($c_1$, $c_2$, $C_1$) oder nur Weggrößen ($c_3$, $c_4$, $C_2$) des
        starren Stabes: Kinematik und Lager mit gehaltener Richtung enthalten Wege, Gelenk-,
        Gleichgewichts- und Drehbedingungen Kräfte. Beispiel: starrer Kragarm, eingespannt am
        Anfang, Kraft $F$ (in $+w$) am freien Ende. Die Einspannung gibt $w(0) = c_4 = 0$ und
        $w'(0) = c_3 = 0$, also bewegt sich nichts. Das freie Ende gibt $Q(L) = -c_1 = F$ und
        $M(L) = -c_2 - c_1 L = 0$, also $c_1 = -F$, $c_2 = F L$ und $M(0) = -F L$. Das ist genau die
        Rechnung des statisch bestimmten Kragarms von Hand. In \code{randbed} heißen die
        Kraftgrößen eines starren Stabes deshalb $Q_i$, $M_i$, $N_i$ statt $-EI\,w_i'''$ usw.
  \item \textbf{Grenzwert}: Jeder starre Stab bekommt $EI = 1/\varepsilon$ bzw.\
        $EA = 1/\varepsilon$ (im CAS heißt $\varepsilon$ \code{rleps}), das System wird symbolisch in
        $\varepsilon$ gelöst und danach \code{limit(\dots, rleps, 0)} gebildet.
\end{itemize}

Warum überhaupt der Grenzwert? Das direkte Verfahren wird singulär, wenn sich starre Teile
\emph{gegenseitig} statisch unbestimmt halten. Beispiel: zwei dehnstarre Stäbe hintereinander
zwischen zwei Festlagern, Horizontallast in der Mitte. Beide Stäbe ändern ihre Länge nicht, also
ist $u = 0$ überall. Aber wie sich die Last auf die beiden Stäbe verteilt, sagt euch das
Gleichgewicht nicht, nur $N_1 - N_2 = F$. Mit $EA = 1/\varepsilon$ ist das System regulär, und
der Grenzwert $\varepsilon \to 0$ liefert die Aufteilung, die zwei \enquote{gleich starren} Stäben
entspricht. Genau in diesem Fall wechselt das Programm von selbst auf den Grenzwert. In
\code{randinfo} steht dann \emph{direkt singulär, starre Teile gegenseitig statisch unbestimmt}.

\subsection{Knotenverschiebungen aus den Randwerten}\label{sec:mathe-knotenverschiebung}

Für die Knotenverschiebung (Knotenmenü Seite~4) muss nichts Neues gelöst werden. Die
Kinematik-Bedingungen drücken $\vec U$ ja schon durch ein Referenz-Stabende aus, und zwar durch das erste
Stabende am Knoten ohne N- und Q-Gelenk. Dessen Längs- und Querverschiebung stehen nach dem Lösen
fest, $u(0)$ und $w(0)$ am Stabanfang bzw.\ $u(L)$ und $w(L)$ am Stabende. Zurück ins globale
System:
\[
  U_x = u\,t_x + w\,n_x, \qquad U_y = u\,t_y + w\,n_y, \qquad
  \vec t = \frac{(\Delta x, \Delta y)}{L}, \quad \vec n = \frac{(-\Delta y, \Delta x)}{L} .
\]
Das sind dieselben $\vec t$ und $\vec n$ wie in den Bedingungen, deshalb zählt $U_y$ wie die
$y$-Achse des Programms nach unten. Die Verdrehung kommt vom ersten Stabende ohne M-Gelenk,
\[
  \varphi = -w' = -\frac{w_p'(x) + c_1 x^2/2 + c_2 x + c_3}{EI},
\]
beim biegesteifen Stab $\varphi = \kappa_T x - c_3$. Hängen alle Stäbe über N- oder Q-Gelenke am
Knoten, hat das LGS für ihn ohnehin die eigenen Unbekannten $u_x(K)$, $u_y(K)$, die werden direkt
genommen.

Damit das beim zweiten Knoten nicht wieder ein ganzes LGS kostet, hebt das Programm die Linien des
letzten Rand-LGS mit einem Schlüssel auf: dem Fingerabdruck aller Eingaben (derselbe, an dem die
Hinweisbox erkennt, ob sich etwas geändert hat. Er enthält auch die Knoten- und Stabnummern) und
den Schaltern, die das Rand-LGS sonst beeinflussen (Methode für starre Stäbe, symbolischer Modus,
Zahlenformat). Passt der Schlüssel, werden nur die Randwerte eingesetzt, das sind drei kurze
CAS-Aufrufe statt eines LGS. Gleicher Fingerabdruck heißt wirklich gleiches System, auch wenn ein
Knoten zwischendurch gelöscht und genau so wieder eingebaut wurde. Werden dabei Nummern
vertauscht, ändert sich der Fingerabdruck, und das ist auch richtig so: Dann gehören die Namen
\code{wv1}, \code{wv2}, \dots\ zu anderen Stäben.

%---------------------------------------------------------------------
\section{Symbolischer Modus: rechnen mit Platzhaltern}\label{sec:mathe-symbolisch}

Die Finite-Elemente-Rechnung für die Diagramme ist eine reine Zahlenrechnung. Symbolisch
invertiert der Rechner keine Steifigkeitsmatrix in vernünftiger Zeit. Der Trick:

\begin{enumerate}
  \item Jedes Symbol bekommt einen \textbf{Platzhalterwert}, standardmäßig~1. Die Länge $l$
        bekommt, wenn die Geometrie kein eigenes Symbol hat, den Rasterabstand.
  \item \textbf{Superposition:} Alle Lasten werden auf null gesetzt und dann \emph{einzeln} wieder
        eingeschaltet. Für jede Last rechnet das Programm die Stabendkräfte und teilt sie durch den
        Platzhalter der Last und die passende Potenz von $l$. Übrig bleibt ein reiner
        Zahlenfaktor~$c$:
        \[
          N = c\,F, \quad M = c\,F\,l \quad\text{(Einzellast)}, \qquad
          N = c\,q\,l, \quad M = c\,q\,l^2 \quad\text{(Streckenlast)} .
        \]
        Die Steifigkeitsmatrix wird dafür nur einmal zerlegt, jede Last kostet dann nur noch ein
        Vorwärts- und Rückwärtseinsetzen.
  \item Der Zahlenfaktor wird \textbf{zurück in einen exakten Ausdruck} verwandelt (siehe unten)
        und alle Anteile werden addiert.
\end{enumerate}

Das funktioniert, weil die Rechnung linear ist und alle Längen Vielfache von $l$ sind. Genau hier
kommen die Einschränkungen her: Eine Temperatur passt nicht in das Schema \enquote{Last mal
Länge hoch $k$} (deshalb fehlt sie in den Diagrammwerten; die Exporte \emph{Neue Biegelinie} und
\emph{Neue Längslinie}, also \code{wneu}, \code{uneu}, \code{wv}, \code{uv}, haben sie, weil das
Rand-LGS gar nicht mit Superposition rechnet), und ein Zahlen-$EI$
neben einem symbolischen $EI$ macht die Steifigkeitsverhältnisse vom Platzhalterwert abhängig
(deshalb der gelbe Hinweis). Die Stelle des Maximums in einem Feld sucht das Programm numerisch
und rastet sie auf typische Werte wie $1/\sqrt 3$ oder einfache Brüche ein.

\textbf{Aus einer Zahl wird wieder ein Bruch.} Aus $0{,}708333\ldots$ soll wieder
$\frac{17}{24}$ werden. Dafür entwickelt das Programm die Zahl in einen \emph{Kettenbruch}
$a_0 + \cfrac{1}{a_1 + \cfrac{1}{a_2 + \dots}}$ und nimmt den ersten Näherungsbruch, der bis
auf $10^{-12}$ stimmt (Nenner höchstens 10\,000). Klappt das nicht und steht das Zahlenformat auf
\emph{Wurzel/Bruch}, versucht es dasselbe mit dem \emph{Quadrat} der Zahl und zieht aus Zähler und
Nenner die Quadratfaktoren heraus. So wird aus $0{,}7071\ldots$ wieder $\frac{\sqrt 2}{2}$ und
aus $86{,}60\ldots$ ein $50\sqrt 3$. Scheitert beides, sucht es zuletzt noch einen Bruch mit
kleinem Nenner (höchstens 100), der nur bis auf die eingestellte Toleranz $10^{-x}$ stimmen muss
(Obermenü Seite~3, \emph{CAS Bruch-Tol.}). Das ist die einzige Stelle, an der diese Einstellung
wirkt, und gerechnet wird dabei in Lua, nicht im CAS. Erst wenn auch das nicht klappt, bleibt eine
Dezimalzahl stehen. Im Format \emph{Dezimal} fallen alle drei Versuche weg.

%---------------------------------------------------------------------
\section{Kraftgrößenverfahren: Hauptsystem und \texorpdfstring{$\delta$}{delta}-Zahlen}\label{sec:mathe-kgv}

\textbf{Das Hauptsystem} sucht das Programm \enquote{gierig}: Es löst der Reihe nach
Einspannmomente, Lagerkräfte, Momente an biegesteifen Stabanschlüssen und zuletzt unbelastete
Pendelstäbe. Es behält jede Lösung nur, wenn das System danach noch tragfähig ist (die
Steifigkeitsmatrix lässt sich noch zerlegen). Sobald nichts mehr geht, ist das System statisch
bestimmt, und die Anzahl der gelösten Bindungen ist der Grad $n$.

\textbf{Die $\delta$-Zahlen} integriert das Programm nicht über $\int M_i M_j/(EI)\,dx$ (so rechnet
ihr sie in \textbf{Statik} mit dem Prinzip der virtuellen Kräfte), sondern es
rechnet für jeden Zustand ($X_j = 1$ bzw.\ Lastzustand) die Verformung des Hauptsystems mit dem
Weggrößenverfahren und liest die zu $X_i$ gehörende Klaffung ab: Verschiebung am gelösten Lager,
gegenseitige Verdrehung am eingefügten Gelenk, Längenänderung am gelösten Pendelstab. Nach
Maxwell--Betti ist das dasselbe wie das Arbeitsintegral, inklusive der Anteile aus Normalkraft,
Federn und endlichem $EA$, die man von Hand gern vergisst. Ein gelöster Pendelstab ist im
Hauptsystem ganz entfernt. Seine eigene Längenänderung $L/(EA)$ fehlt deshalb in den
Knotenverschiebungen und wird zu $\delta_{jj}$ addiert. Bei gedrehten Lagern liegen die
Knotenfreiheitsgrade schon in Lagerrichtung ($u_\xi$, $u_\eta$), die Klaffung eines gelösten Lagers
ist genau dieser Freiheitsgrad, und die Einheitskräfte eines gelösten Pendelstabes werden in
diese Richtungen zerlegt.

%---------------------------------------------------------------------
\section{Polplan, Verschiebungszustand und Arbeitssatz}\label{sec:mathe-polplan}

\subsection{Scheiben finden}

Zuerst muss das Programm wissen, welche Stäbe zusammen eine starre Scheibe bilden. Das passiert
in zwei Durchgängen:

\begin{enumerate}
  \item \textbf{Ohne Lager:} An jedem Knoten ohne Vollgelenk fasst das Programm alle Stäbe
        zusammen, deren Enden dort kein Gelenk haben (\emph{Union-Find}: jede Gruppe hat einen
        Vertreter, zwei Gruppen werden vereinigt, indem man einen Vertreter unter den anderen
        hängt). Mit diesen Scheiben und nur den inneren Gelenken läuft der Polplan einmal durch.
        Findet er zwei Scheiben, die gegeneinander unverschieblich sind, z.\,B.\ die drei Stäbe
        eines Dreigelenk-Dreiecks, werden sie zu einer Scheibe zusammengefasst.
  \item \textbf{Mit Lagern:} Stäbe, die biegesteif an einer Einspannung hängen (oder an einem
        Festlager mit Drehfeder), gehören zur Scheibe~0, dem Boden. Mit diesen zusammengefassten
        Scheiben und den Lagern entsteht der Polplan, den ihr seht.
\end{enumerate}

\subsection{Pole und Polstrahlen}

An jedem Knoten schaut das Programm, welche Scheiben sich dort treffen und welche Relativbewegung
zwischen ihnen möglich ist:

\begin{longtable}{p{0.34\linewidth}p{0.60\linewidth}}
  \toprule
  \textbf{Verbindung} & \textbf{Folge im Polplan} \\
  \midrule \endhead
  M-Gelenk, Festlager & Pol im Knoten (Drehung möglich, keine Verschiebung) \\
  Rollenlager & Polstrahl durch den Knoten senkrecht zur Rollrichtung, also entlang der
    Lagerkraft \\
  N- oder Q-Gelenk & Pol im Unendlichen (reine Verschiebung längs bzw.\ quer zum Stab) \\
  Einspannung, biegesteifer Anschluss & kein Pol: die Scheiben sind starr verbunden \\
  \bottomrule
\end{longtable}

Danach wiederholt das Programm (höchstens zehnmal) zwei Schritte, bis nichts Neues mehr dazukommt:

\begin{enumerate}
  \item \textbf{Satz von den drei Polen:} Für je drei Scheiben $i$, $j$, $k$ liegen die Pole
        $(i,j)$, $(j,k)$ und $(i,k)$ auf einer Geraden. Sind zwei davon bekannt, entsteht ein
        Polstrahl für den dritten. Liegt einer der beiden im Unendlichen, geht der Strahl durch
        den anderen in dessen Richtung.
  \item \textbf{Strahlen schneiden:} Zwei Polstrahlen derselben Scheibenpaarung schneiden sich im
        Pol. Sind sie parallel, liegt der Pol im Unendlichen in dieser Richtung.
\end{enumerate}

Findet es für eine Paarung zwei verschiedene Pole, oder liegt ein Pol nicht auf einem Strahl, auf
dem er liegen müsste, sind die beiden Scheiben gegeneinander unverschieblich (im Polplan rot). Im
Durchgang ohne Lager werden solche Paare zu einer Scheibe zusammengefasst (siehe oben).

\subsection{Der Verschiebungszustand}

Eine Scheibe, die sich um ihren Hauptpol $(c_x, c_y)$ um den kleinen Winkel $\varphi$ dreht,
verschiebt jeden Punkt um $\vec r_\perp \varphi$, also senkrecht zum Polstrahl vom Pol zu diesem
Punkt. Das Programm gibt einer Startscheibe $\varphi = 1$ (bevorzugt eine, deren Hauptpol auf einem
Knoten liegt) und reicht die Bewegung über die Nebenpole weiter: Der Nebenpol $(A,B)$ ist der Punkt,
der auf beiden Scheiben dieselbe Verschiebung hat. Kennt man die Verschiebung dieses Punktes aus
Scheibe~$A$ und den Hauptpol von Scheibe~$B$, folgt $\varphi_B$. Hat eine Scheibe ihren Hauptpol im
Unendlichen, verschiebt sie sich nur.

\textbf{Verschobene Scheiben unterwegs.} Die Kette reißt an einer Scheibe, die sich nur verschiebt,
nicht ab. Bei ihr bewegt sich jeder Punkt um denselben Vektor $\vec v$, also auch der Nebenpol zur
nächsten Scheibe~$B$. Damit ist die Verschiebung des Nebenpols bekannt, und es geht weiter wie bei
einer drehenden Scheibe: Liegt der Hauptpol von $B$ im Endlichen, folgt $\varphi_B$ aus $\vec v$,
liegt er im Unendlichen, verschiebt sich auch $B$ nur. Liegt der Nebenpol $(A,B)$ selbst im
Unendlichen, drehen beide Scheiben gleich, $\varphi_B = \varphi_A = 0$: $B$ verschiebt sich dann auch
nur, oder sie steht still, wenn ihr Hauptpol im Endlichen liegt. Eine reine Verschiebung ist nur der
Grenzfall einer Drehung mit dem Pol im Unendlichen.

\textbf{Nebenpol gleich Hauptpol.} Der Hauptpol einer Scheibe ist der Punkt, der sich auf ihr nicht
bewegt. Fällt der Nebenpol $(A,B)$ mit dem Hauptpol von $B$ zusammen, bewegt sich dieser Punkt also
auf $B$ nicht, und weil er auf $A$ dieselbe Verschiebung haben muss, auch auf $A$ nicht. Kommt aus
der Startscheibe $A$ trotzdem eine Verschiebung ungleich null an, passt das nicht zusammen: Der Punkt
kann nicht zugleich stehen und sich bewegen. Dann kann sich $A$ gar nicht bewegen, die Startscheibe
war falsch gewählt, und das Programm nimmt die nächste.

Ein Beispiel ist euer Dreigelenktragwerk mit dem Gelenk $G$ senkrecht über dem linken Lager $A$,
bei dem ihr das linke Festlager zum Loslager macht (verschieblich in $x$). Das rechte Festlager $B$
ist der Hauptpol $(\mathrm{II})$ der rechten Scheibe, das Gelenk der Nebenpol $(\mathrm{I,II})$. Der
Hauptpol $(\mathrm{I})$ der linken Scheibe liegt auf der Senkrechten durch $A$ (Loslager) und nach
dem Satz von den drei Polen auf der Geraden durch $(\mathrm{II})$ und $(\mathrm{I,II})$, also auf
$BG$. Beide Geraden schneiden sich in $G$, damit ist $G$ zugleich $(\mathrm{I})$ und
$(\mathrm{I,II})$. Startet das Programm mit der rechten Scheibe, dreht sie um $B$ und schiebt $G$
seitlich weg, aber auf der linken Scheibe steht $G$ als Hauptpol still: Widerspruch, die rechte
Scheibe kann sich nicht bewegen. Mit der linken Scheibe als Start geht es: Sie dreht um $G$, ihr
Lager $A$ wandert waagerecht (senkrecht zum Polstrahl $GA$, und das lässt das Loslager zu). $G$
selbst bewegt sich dabei nicht. Die rechte Scheibe hat damit zwei Punkte, die stehen bleiben ($B$
und $G$), und bewegt sich deshalb gar nicht, $\varphi_\mathrm{II} = 0$. Das ist kein Fehler,
sondern die Kinematik dieses Systems: Nur die linke Scheibe kippt um das Gelenk.

\subsection{Der Arbeitssatz}\label{sec:mathe-arbeitssatz}

Nach dem Lösen einer Bindung ist das System eine kinematische Kette mit genau einem Freiheitsgrad.
Für jede virtuelle Bewegung, die die übrigen Bindungen zulassen, gilt
\[
  \delta W = \sum_k \vec F_k \cdot \delta\vec u_k + \sum_m M_m\,\delta\varphi_m = 0 ,
\]
und weil es nur einen Freiheitsgrad gibt, sind alle Wege Vielfache eines einzigen Parameters. Der
Export (Abschnitt~\ref{sec:arbeitssatz}) macht daraus drei Schritte.

\textbf{Weg einer Kraft (\code{pvv\_kin}, Spalte~2).} Eine Kraft mit Richtung
$\vec e = (e_x, e_y)$ ($y$ nach unten) greift im Punkt $P$ auf Scheibe~$i$ an. Ihr Weg ist die
Verschiebung von $P$ in Richtung $\vec e$:
\begin{itemize}
  \item Scheibe~$i$ dreht um ihren Hauptpol $(c_x, c_y)$ um $\varphi_i$ (gegen den Uhrzeigersinn).
        Jeder Punkt bewegt sich senkrecht zum Polstrahl vom Pol zu ihm, projiziert auf $\vec e$ ist das
        \[
          \delta_e = \varphi_i\,\bigl[e_x\,(P_y - c_y) - e_y\,(P_x - c_x)\bigr] .
        \]
        Die eckige Klammer ist der Hebel der Kraft um den Pol.
  \item Scheibe~$i$ verschiebt sich nur (Hauptpol im Unendlichen) um $\delta_i$ in Richtung
        $\vec t_i$: $\delta_e = \delta_i\,(\vec e\cdot\vec t_i)$.
  \item Ein Moment auf Scheibe~$i$ arbeitet am Drehwinkel: $\delta\varphi = \varphi_i$.
\end{itemize}
Verschieben \emph{und} drehen zugleich gibt es nicht als eigenen Fall: Eine kleine ebene
Starrkörperbewegung ist immer eine Drehung um einen Punkt, eben den Pol. Eine Verschiebung plus
Drehung ist nur eine Drehung um einen anderen Pol, die reine Verschiebung der Grenzfall mit dem Pol
im Unendlichen.
Streckenlasten dürfen exakt durch ihre Resultierende $R$ ersetzt werden. Auf einem geraden Stab
der Länge $L$, der zu einer starren Scheibe gehört, ändert sich die virtuelle Verschiebung linear
mit der Stabkoordinate $x$, denn ein Starrkörper verschiebt jeden Punkt linear mit seinem Ort:
$\delta\vec u(x) = \delta\vec u_A + \frac{x}{L}\,(\delta\vec u_B - \delta\vec u_A)$. Damit ist die
Arbeit der Streckenlast
\[
  \int_0^L q(x)\,\vec e\cdot\delta\vec u(x)\,dx
  = \vec e\cdot\delta\vec u_A \underbrace{\int_0^L q\,dx}_{R}
    + \vec e\cdot\frac{\delta\vec u_B - \delta\vec u_A}{L} \underbrace{\int_0^L q\,x\,dx}_{S = R\,x_R}
  = R\;\vec e\cdot\delta\vec u(x_R) ,
\]
also die Arbeit der Resultierenden $R$ an ihrem Angriffspunkt $x_R = S/R$. Auf einem Bogen ist die
Verschiebung nicht linear in der Bogenlänge, das Ergebnis gilt aber genauso: Dort wird die Last
durch ihre Resultierende (Kraft und Richtung) auf ihrer Wirkungslinie ersetzt. Die freigesetzte
Größe ist eine äußere Last in der gezeichneten Richtung, bei einem gelösten Gelenk arbeitet sie an
der gegenseitigen Verdrehung bzw.\ Verschiebung der beiden Seiten.

\textbf{Verhältnis der Scheiben (\code{pvv\_pol}).} Der Nebenpol $P_{AB}$ ist der Punkt, der auf
Scheibe~$A$ und Scheibe~$B$ dieselbe Verschiebung hat. Schreibt man diese Verschiebung mit der
Formel oben einmal über $A$ und einmal über $B$ hin, entsteht eine lineare Beziehung
$\varphi_B = \kappa_{AB}\,\varphi_A$. Die Verschiebung ist ein Vektor mit $x$- und $y$-Komponente,
das sind eigentlich zwei Gleichungen. Beide liefern aber dasselbe $\kappa_{AB}$: Über $A$ steht die
Verschiebung senkrecht auf dem Polstrahl von $(A)$ zum Nebenpol, über $B$ senkrecht auf dem von
$(B)$, und weil $(A)$, $(A,B)$ und $(B)$ auf einer Geraden liegen, sind beide Richtungen gleich.
Deshalb reicht eine Komponente. Das Programm nimmt die betragsmäßig größere, damit es nicht durch
eine Zahl nahe null teilt. Liegt der Nebenpol zweier drehender Scheiben im Unendlichen, drehen beide
gleich, $\varphi_A = \varphi_B$.

Von der Bezugsscheibe aus (der Startscheibe des Verschiebungszustands, ihr Winkel heißt
$\varphi_0$) laufen diese Beziehungen über die Nebenpole durch alle Scheiben und liefern
$\varphi_i = \kappa_i\,\varphi_0$. Das sind die unteren Zeilen von \code{pvv\_kin}. Gibt es mehr
Nebenpole als nötig, sind die übrigen Proben. Als Bezug nimmt das Programm eine Scheibe, die sich
dreht, wenn es eine gibt. Eine Scheibe, die stehen bleibt, bekommt einfach $\kappa_i = 0$, ihre
Lasten leisten keine Arbeit. Verschieben sich alle bewegten Scheiben nur (alle ihre Hauptpole im
Unendlichen), gibt es keinen Drehwinkel als Bezug. Dann ist die Bezugsgröße die Verschiebung $\delta_0$ der Startscheibe, und
alle anderen Bewegungen werden als Vielfache von $\delta_0$ geschrieben (siehe unten).

\textbf{Einsetzen (\code{w\_eq2}).} Mit $\varphi_i = \kappa_i\,\varphi_0$ wird jeder Weg zu
$h_k\,\kappa_{i(k)}\,\varphi_0$ (Spalte~3 von \code{pvv\_kin}). Dabei ist $i(k)$ die Scheibe, an der
die Kraft $F_k$ angreift, und $h_k$ ihr Hebel um den Pol dieser Scheibe, also die eckige Klammer
von oben, $h_k = e_x\,(P_y - c_y) - e_y\,(P_x - c_x)$. Das ist der senkrechte Abstand der
Wirkungslinie vom Pol mit Vorzeichen, nicht der Abstand des Angriffspunktes. Verschiebt sich die
Scheibe nur, tritt an die Stelle von $h_k\,\varphi_i$ der Weg $\delta_i\,(\vec e\cdot\vec t_i)$ aus
der Liste oben, $h_k$ ist dann der Richtungskosinus $\vec e\cdot\vec t_i$. Der Arbeitssatz zerfällt in
\[
  \delta W = \varphi_0 \Bigl(\sum_k F_k\,h_k\,\kappa_{i(k)} + \sum_m M_m\,\kappa_{i(m)}\Bigr) = 0 .
\]
Da $\varphi_0$ beliebig ist, muss die Klammer verschwinden. Sie ist linear in der freigesetzten
Größe, also eine Gleichung mit einer Unbekannten. Verschiebt sich die Bezugsscheibe nur, steht
$\delta_0$ statt $\varphi_0$ da, und die Faktoren der Drehwinkel haben die Einheit $1/$Länge.
Intern zählt das Programm Drehwinkel im Uhrzeigersinn. Im Export ist das Vorzeichen umgedreht,
damit alles zur Konvention der Vorlesung (gegen den Uhrzeigersinn positiv) passt.

%---------------------------------------------------------------------
\section{Gleichgewicht der Teilsysteme (LGS-Matrix)}\label{sec:mathe-gleichgewicht}

Für den Export \emph{LGS-Matrix} schneidet das Programm das System an allen Gelenken auf, genau
so, wie Seite~2 der Explosionszeichnung es zeigt. Beide benutzen dieselben drei Regeln im Code:
\code{gelenkKnoten} (welche Stäbe hängen mit Gelenk, welche biegesteif am Knoten, wie viele
Teilsysteme treffen sich), \code{gelenkRichtung} (positive Richtung der Gelenkgrößen) und
\code{lastStabKnoten} (an welchem Teilsystem Lasten und Lager eines Gelenkknotens hängen). Jedes
Teilsystem liefert drei Gleichungen: Kräfte nach rechts, nach oben und Momente gegen den
Uhrzeigersinn. Unbekannt sind die Lagerreaktionen und die Größen, die ein Gelenk überträgt: beim
M-Gelenk die beiden Kraftkomponenten $Gx$, $Gy$, beim N-Gelenk $Gq$ und $Gm$, beim Q-Gelenk $Gn$
und $Gm$. Treffen sich mehr als zwei Teilsysteme in einem Gelenk, bekommt jeder Gelenkstab eigene
Gelenkkräfte, gibt es dort keinen biegesteifen Rest, sagen zwei zusätzliche Gleichungen, dass der
Gelenkbolzen selbst im Gleichgewicht ist. Bei nur zwei Teilsystemen braucht der Bolzen keine
eigenen Gleichungen: Die Gelenkkraft steht in beiden mit entgegengesetztem Vorzeichen (actio =
reactio), damit ist sein Gleichgewicht schon erfüllt. Hängt am Gelenk ein biegesteifer Rest, gehört
der Bolzen zu ihm. Ein Beispiel ist eine Stütze, die durch den Knoten durchläuft, und zwei Riegel,
die gelenkig an ihr hängen. Die Gelenkkräfte der Riegel greifen dann direkt an der Stütze an und
stehen in deren drei Gleichungen. Der Bolzen ist kein eigener Körper, also gibt es für ihn keine
eigenen Gleichungen, sein Gleichgewicht steckt im Gleichgewicht der Stütze. Erst wenn alle Stäbe
gelenkig am Knoten hängen, gehört der Bolzen zu keinem Teilsystem und braucht seine zwei
Gleichungen selbst. Im symbolischen Modus stehen Lasten mit ihren Namen und
Hebelarme als Vielfache von $l$ in der Matrix.

Den \textbf{Momentenbezugspunkt} wählt das Programm so, wie ihr es in \textbf{TM1} lernt: Es
sucht den Punkt, durch den möglichst viele unbekannte Kräfte gehen, also die senkrechte Linie mit den
meisten unbekannten Vertikalkräften und die waagerechte mit den meisten unbekannten
Horizontalkräften, am liebsten ein Knoten, der auf beiden liegt. Die Momentengleichung wird dadurch
so kurz wie möglich.

Die Werte auf Seite~1 der Explosionszeichnung kommen aus dem Weggrößenverfahren. Löst ihr das LGS
(mit den oben beschriebenen Befehlen oder von Hand), müssen dieselben Zahlen herauskommen. Das ist
eure Kontrolle.

\textbf{Lösungsweg (\code{lgs\_weg}):} Gesucht ist eine Reihenfolge, in der jede Gleichung möglichst
nur eine neue Unbekannte enthält. Dafür braucht das Programm mehr Gleichungen als die in \code{lgs}.
Alle sind Linearkombinationen der Zeilen aus \code{lgs}, also nichts Neues, nur anders
zusammengefasst.

\textbf{Anderer Momentenpunkt.} Eine Kraft mit den Komponenten $F_\rightarrow$ (nach rechts) und
$F_\uparrow$ (nach oben) im Punkt $(x, y)$ hat um $P$ das Moment
$M_P = (x - x_P)\,F_\uparrow - (y - y_P)\,F_\rightarrow$ (gegen den Uhrzeigersinn positiv, hier $y$
nach oben). Die Pfeile sind dieselben wie in den Namen der Kraftzeilen von \code{lgs}. Um einen
anderen Punkt $Q$ ist es
\[
  M_Q = (x - x_Q)\,F_\uparrow - (y - y_Q)\,F_\rightarrow
      = M_P + (x_P - x_Q)\,F_\uparrow - (y_P - y_Q)\,F_\rightarrow .
\]
Über alle Kräfte eines Teilsystems summiert (reine Momente hängen nicht vom Punkt ab) heißt das
\[
  \textstyle\sum M_Q = \sum M_P + (x_P - x_Q) \sum F_\uparrow - (y_P - y_Q) \sum F_\rightarrow .
\]
Die Momentengleichung um $Q$ ist also die Momentenzeile plus Vielfache der beiden Kraftzeilen
desselben Teilsystems, die Faktoren sind Hebelarme. Im Programm zeigt $y$ nach unten, dort steht
deshalb $+(y_P - y_Q)$.

\textbf{Teilsysteme zusammenfassen.} Für eine Gruppe werden die Kraftzeilen addiert und die
Momentenzeilen, jede vorher auf denselben Punkt $Q$ umgerechnet. Eine Gelenkkraft zwischen zwei
Teilsystemen der Gruppe steht in beiden mit entgegengesetztem Vorzeichen und greift am selben Punkt
an. Sie fällt deshalb aus den Kräften \emph{und} aus den Momenten heraus. Alle Teilsysteme zusammen
sind das Gesamtsystem.

\textbf{Kräfte quer zu einer Richtung.} Zeigt eine Unbekannte in Richtung $(d_x, d_y)$ (wieder mit $y$
nach unten), dann enthält $d_y \sum F_\rightarrow + d_x \sum F_\uparrow$ (normiert) diese Unbekannte
nicht mehr.

\textbf{Warum höchstens drei Gleichungen je Scheibe?} Jede Gleichgewichtsbedingung einer Scheibe ist
entweder eine Kraftsumme in eine Richtung $\vec e$, also $e_x \sum F_\rightarrow + e_y \sum
F_\uparrow$, oder eine Momentensumme um einen Punkt, nach der Formel oben $\sum M_P$ plus Vielfache
der beiden Kraftsummen. Alle liegen also in dem Raum, den $\sum F_\rightarrow$, $\sum F_\uparrow$ und
$\sum M_P$ aufspannen, und der hat höchstens die Dimension~3. Mehr als drei linear unabhängige
Gleichungen bekommt ihr deshalb nie, egal um wie viele Punkte ihr die Momente bildet. Drei
Momentengleichungen um $A$, $B$ und $C$ sind genau dann unabhängig, wenn die drei Punkte nicht auf
einer Geraden liegen. Liegen sie auf einer Geraden, ist die dritte eine Kombination der beiden
anderen. Ein Durchlaufträger auf drei Lagern bleibt deshalb einfach statisch unbestimmt, auch wenn
man um jedes Lager die Momente bildet. Anschaulich (PvV): Jede Gleichung gehört zu einer virtuellen
Starrkörperbewegung, einer Verschiebung in Richtung $\vec e$ bzw.\ einer Drehung um $P$, und eine
Scheibe hat in der Ebene nur drei Freiheitsgrade.

\textbf{Welche Gruppen und Punkte?} Das Programm sucht die Gleichungen nicht gezielt, sondern stellt
alle aus einem festen Vorrat auf und wählt danach aus:

\begin{itemize}
  \item \textbf{Gruppen:} jedes Teilsystem einzeln, das Gesamtsystem und bei Rahmen alle
        zusammenhängenden Gruppen. Zwei Teilsysteme hängen zusammen, wenn sie eine
        Gelenkunbekannte teilen. Ohne gemeinsame Gelenkkraft hebt sich beim Zusammenfassen nichts
        auf. Die Gruppen entstehen, indem an jede Gruppe ein Nachbar angehängt wird, bis die
        Zahl unter \emph{Max. Teilgruppen} erreicht ist (ab Werk 30). Im Fachwerk gibt es nur die Knoten und das
        Gesamtsystem, denn eine Gruppe von Knoten wäre ein Ritterschnitt.
  \item \textbf{Momentenpunkte $Q$:} Knoten und Schnittpunkte zweier Wirkungslinien von
        Unbekannten bzw.\ Pendelachsen (innerhalb der doppelten Systemgröße), aber je Gruppe nur
        die, an denen eine \emph{äußere} Kraft der Gruppe herausfällt. Äußere Kräfte sind die, die
        nach dem Zusammenfassen noch in der Kraftsumme stehen. Gelenkkräfte zwischen zwei
        Teilsystemen der Gruppe heben sich auf und zählen nicht. Ein Knoten zählt, wenn die
        Wirkungslinie einer äußeren Kraft durch ihn geht, ein Schnittpunkt, wenn beide Linien zu
        äußeren Kräften gehören, und er wird nach diesem Paar benannt. Für \code{TS 1+2} ist das
        Gelenk zwischen TS~1 und TS~2 deshalb nur dann ein Kandidat, wenn noch eine andere äußere
        Kraft durch es hindurchgeht. $Q$ wird also nicht vorher ausgesucht: Jeder dieser Punkte
        liefert eine Gleichung, und die Auswahl unten nimmt die, in der nur noch eine Unbekannte
        übrig ist.
  \item \textbf{Kräfte quer} zu jeder schrägen äußeren Unbekannten und jeder schrägen Pendelachse.
\end{itemize}

Gleichungen, die sich nur um einen Faktor unterscheiden, bleiben einmal übrig, und zwar in der
einfachsten Form. Weil dabei nur mit Zahlen gerechnet wird, dauert das auch bei einigen hundert
Kandidaten kaum Zeit.

\textbf{Benennung.} Der Name nennt die Art, den Punkt und die Gruppe: \code{$\Sigma$M P5 TS 1+2}
sind die Momente um Knoten~5 an den Teilsystemen 1 und 2 zusammen, \code{$\Sigma$M S(A\_h,B\_v)
Ges} die Momente des Gesamtsystems (\code{Ges}) um den Schnittpunkt der Wirkungslinien von $A_h$ und
$B_v$, \code{$\Sigma$F quer zu B\_h TS 2} die Kräfte an Teilsystem~2 quer zu $B_h$. Neue Gleichungen
bekommen die nächste freie Nummer nach denen aus \code{lgs}.

\textbf{Pendelstäbe.} Ein Teilsystem eines Rahmens ohne Lasten, an dem nur an zwei Punkten $U$ und
$V$ Anschlusskräfte angreifen (Lager- bzw.\ Gelenkkräfte, keine Momente), nennen wir hier
Pendelstab. Ist es ein gerader Stab, ist das der Pendelstab aus \textbf{TM1} mit reiner
Normalkraft. Ist es geknickt (z.\,B.\ ein unbelasteter Winkel zwischen zwei Gelenken), hat es zwar
Momente und Querkräfte, für das Gleichgewicht verhält es sich aber genauso: Die Anschlusskraft liegt
auf der Verbindungslinie $UV$. Nur darum geht es hier. Sein
Momentengleichgewicht um $U$ enthält nur die beiden Kraftkomponenten bei $V$ und sagt: Die Kraft bei
$V$ liegt auf der Achse $UV$. Liegt der Punkt $Q$ einer anderen Momentengleichung auf dieser Achse,
stehen die beiden Komponenten bei $V$ dort im selben Verhältnis wie im Pendelgleichgewicht, nur mit
einem Faktor $\lambda$. Zieht man $\lambda$-mal das Pendelgleichgewicht ab, fallen beide heraus. Das
Programm prüft das für jede Gleichung und jeden Pendelstab mit den Zahlen und nimmt die Achse als
weitere Wirkungslinie für die Schnittpunkte dazu. So findet sich auch beim Dreigelenkrahmen mit
verschieden hohen Lagern und einer unbelasteten Hälfte ein Punkt, der direkt eine Lagerkraft
liefert. Ist die Achse waagerecht oder senkrecht, braucht es den Umweg nicht, dann hat das
Pendelgleichgewicht selbst schon nur eine Unbekannte. Dieselbe Gleichung entsteht oft auf zwei
Wegen, etwa am Gesamtsystem (Lagerkräfte über das Pendelgleichgewicht um das Gelenkende
herausgerechnet) und an der anderen Scheibe allein (Gelenkkräfte über das Pendelgleichgewicht um
das Lager herausgerechnet). Das Programm behält dann den Namen mit dem Gesamtsystem, weil das der
Vorstellung entspricht, dass die Lagerkraft auf der Pendelachse wirkt.

\textbf{Auswahl.} Dann geht es schrittweise: Gibt es eine Gleichung mit genau einer noch unbekannten
Größe, ist das der nächste Schritt. Unter mehreren gewinnt die einfachste: erst die aus \code{lgs},
dann andere Momentenpunkte am selben Teilsystem, dann das Gesamtsystem, dann Zusammenfassungen und
Schnittpunkte, zuletzt Gleichungen, in denen ein Pendelgleichgewicht steckt. Bei Gleichstand gewinnt
die mit weniger Unbekannten, danach die, die zuerst aufgestellt wurde (Knoten vor Schnittpunkten,
nach Nummern). Lagergrößen (Lager und Federn) haben dabei Vorrang, aber nur innerhalb derselben
Stufe: Ein Einzelschritt für eine Lagergröße geht vor jedem anderen Einzelschritt. Gibt es keinen,
nimmt das Programm einen Einzelschritt, nach dem eine Lagergröße als Einzelschritt folgt (einen
Schritt vorausgeschaut), sonst den einfachsten. Ein Block kommt erst, wenn es gar keinen
Einzelschritt mehr gibt. Gibt es keinen, sucht das Programm zwei Gleichungen mit denselben zwei neuen
Unbekannten, die voneinander unabhängig sind (Rang~2), danach drei mit drei. Unter gleich großen
Blöcken gewinnt der mit mehr Lagergrößen, dann der aus den einfacheren Gleichungen. Dafür hat jede
Gleichung eine Stufe (\emph{Klasse}) nach der Reihenfolge oben: $0$ für eine Zeile aus \code{lgs},
höher für andere Momentenpunkte, das Gesamtsystem, Zusammenfassungen, Schnittpunkte und zuletzt
Gleichungen mit Pendelgleichgewicht. Verglichen wird die Summe der Stufen der Gleichungen im Block,
die kleinere gewinnt. Erst wenn auch das nicht geht, löst es alle übrigen Unbekannten als einen
Block aus Gleichungen von \code{lgs}. Ein solcher Block wird durch Einsetzen gelöst: Die Gleichung
mit den wenigsten Unbekannten wird nach einer davon aufgelöst und in die übrigen eingesetzt, das
wiederholt sich, bis die letzte Gleichung nur noch eine Unbekannte hat. Dann wird von hinten
zurückeingesetzt. Bei gleich vielen Unbekannten nimmt das Programm am liebsten eine Unbekannte, die
in ihrer Gleichung den Faktor $+1$ oder $-1$ hat, denn dann entsteht beim Auflösen kein Bruch.

Die Auswahl läuft mit Zahlen, denn die Faktoren vor den Unbekannten hängen nur von der Geometrie ab,
auch im symbolischen Modus. Die Texte und Ergebnisse in \code{lgs\_weg} rechnet danach das CAS
exakt aus, sie sind also dieselben wie die Lösung des ganzen LGS.

\textbf{Option LGS wie Lösungsweg.} Der Lösungsweg benutzt genau $n$ Gleichungen für $n$
Unbekannte. Ordnet man sie nach den Schritten, hat ihre Matrix Blockdreiecksform: Jeder Schritt
enthält außer schon bekannten Größen nur seine neuen Unbekannten, und deren Block ist regulär
(Einzelschritt: Faktor ungleich null, Block: voller Rang). Damit ist die ganze Matrix regulär. Das
Programm lässt die Gleichungen aus \code{lgs} stehen, die der Weg benutzt, und setzt jede
abgeleitete Gleichung auf den Platz einer Gleichung, die er nicht benutzt, am liebsten auf eine,
aus der sie mit entsteht. Danach prüft es den Rang des neuen Systems noch einmal mit Zahlen. Allgemein
gilt: Ersetzt man Zeile $r$ durch $\sum_i f_i\, r_i$, bleibt die Lösung eindeutig genau dann, wenn
$f_r \neq 0$ ist (die Determinante ändert sich um den Faktor $f_r$).

\textbf{Gegenprüfung mit der FEM.} Für jede Unbekannte bestimmt das Programm aus den FEM-Ergebnissen
den Wert mit demselben Namen und derselben Richtung wie auf Seite~2: Lagerreaktionen direkt,
Gelenkgrößen aus den Stabendkräften des angeschlossenen Stabes (ohne Knotenlast und Lagerkraft, wenn
die zu diesem Stab gehören), Stabkräfte im Fachwerk aus der Normalkraft. Das vergleicht es mit dem
Wert aus \code{lgs\_weg}, im symbolischen Modus mit denselben Ersatzwerten für die Symbole, mit denen
die FEM gerechnet hat. Weicht etwas um mehr als $10^{-4}$ des größten Wertes ab, kommt die Warnung.

\textbf{Resultierende von Streckenlasten:} $R = \int_0^L q\,dx$ und das statische Moment
$S = \int_0^L q\,x\,dx$ um den Stabanfang, die Wirkungslinie liegt bei $x_R = S/R$. Für Trapezlasten
sind das die bekannten Formeln, für Lastfunktionen integriert das Programm mit der Simpsonregel
über elf Stützstellen. Die ist für $q\,x$ bis zu einem quadratischen $q$ exakt.

%---------------------------------------------------------------------
\section{Querschnittswerte}\label{sec:mathe-qswerte}

\textbf{Vielecke} (Rechteck, Dreieck, Trapez, Polygon, dicke Linie) rechnet das Programm alle mit
derselben Formel, der Gaußschen Trapezformel. Mit den Eckpunkten $(y_i, z_i)$ gegen den
Uhrzeigersinn und $c_i = y_i z_{i+1} - y_{i+1} z_i$ ist
\[
  A = \tfrac12 \sum c_i, \qquad
  y_S = \frac{\sum (y_i + y_{i+1})\,c_i}{6A}, \qquad
  \int z^2\,dA = \tfrac1{12} \sum \left(z_i^2 + z_i z_{i+1} + z_{i+1}^2\right) c_i ,
\]
entsprechend für $z_S$, $\int y^2\,dA$ und $\int yz\,dA$. Das ist exakt, es ist der Gaußsche
Integralsatz, angewandt auf die Randkanten. Deshalb ist es auch egal, wie kompliziert das Vieleck
ist, solange es sich nicht selbst schneidet.

\textbf{Kreisausschnitt und -abschnitt} mit halbem Öffnungswinkel $\alpha$: Der Schwerpunkt liegt
auf der Winkelhalbierenden im Abstand $\frac{2R\sin\alpha}{3\alpha}$ (Ausschnitt) bzw.\
$\frac{4R\sin^3\alpha}{3(2\alpha - \sin 2\alpha)}$ (Abschnitt). Die Trägheitsmomente werden erst um
die Winkelhalbierende bestimmt (beim Abschnitt wird das Dreieck zum Mittelpunkt abgezogen) und dann
ins KOS gedreht.

\textbf{Dünnwandige Linien} sind Mittellinie mal Dicke: $A = L\,t$,
$I_y = t L^3 \sin^2\beta/12$ usw. Den Eigenanteil $L t^3/12$ lässt das Programm weg, genau wie
die Vorlesung, weil er von höherer Ordnung in $t$ ist. \textbf{Dünnwandige Kreise und Bögen}
ebenso: $A = R\,\Delta\varphi\,t$, $I_y = R^3 t \int \sin^2\varphi\,d\varphi$ um den Mittelpunkt
(Vollkreis $\pi R^3 t$), Schwerpunkt des Bogens im Abstand $R\,\frac{\sin(\Delta\varphi/2)}{\Delta\varphi/2}$
auf der Winkelhalbierenden.

\textbf{Löcher} gehen mit negativer Fläche in die Summen ein, danach setzt der Satz von Steiner alles
zusammen. \textbf{Numerisch addierte} Elemente legt das Programm auf ein Raster von $40 \times 40$
Zellen über ihr umschreibendes Rechteck und zählt nur die Zellen, die nicht schon in einem anderen
Element liegen. Deshalb die Warnung, das ist eine Näherung.

\textbf{Hauptachsen:} Mit $m = \frac{I_y + I_z}2$ und
$d = \sqrt{\left(\frac{I_y - I_z}2\right)^2 + I_{yz}^2}$ ist $I_{1,2} = m \pm d$ und
\[
  \alpha = \tfrac12 \operatorname{atan2}(2 I_{yz},\ I_y - I_z) .
\]
Das \code{atan2} (der Arkustangens, der das Vorzeichen von Zähler und Nenner kennt) sorgt dafür,
dass $\alpha$ immer zur Achse~1 mit $I_1$ gehört. Mit dem einfachen $\tan 2\alpha$ aus der
Formelsammlung müsstet ihr das von Hand prüfen, dort kann auch die Achse mit $I_2$ herauskommen.

\textbf{Widerstandsmomente} $W = I/e$ nehmen als Randabstand $e$ den größten Abstand der Elemente
vom Schwerpunkt. Bei dünnwandigen Elementen ist das die Mittellinie, $t/2$ kommt nicht dazu.

%---------------------------------------------------------------------
\section{Normalspannung und neutrale Faser}\label{sec:mathe-sigma}

Mit $\Delta = I_y I_z - I_{yz}^2$ und $y$, $z$ vom Schwerpunkt aus rechnet das Programm immer mit
\[
  \sigma_x = \frac NA + \frac{M_y I_z - M_z I_{yz}}{\Delta}\,z
                      - \frac{M_z I_y - M_y I_{yz}}{\Delta}\,y .
\]
Für $I_{yz} = 0$ bleibt das bekannte $\frac NA + \frac{M_y}{I_y} z - \frac{M_z}{I_z} y$ übrig. Die
Formel kommt direkt aus $M_y = \int \sigma z\,dA$ und $M_z = -\int \sigma y\,dA$ mit dem
linearen Ansatz $\sigma = a + b\,y + c\,z$, also zwei Gleichungen für $b$ und $c$.

Weil $\sigma_x$ linear in $y$ und $z$ ist, speichert das Programm nur den Gradienten $(c_y, c_z)$ und
den Wert $c_0$ im Ursprung des gezeigten KOS. Daraus entstehen beide Gleichungen in der Tabelle: die
KOS-Form direkt, die HAS-Form, indem der Gradient um den Hauptachsenwinkel gedreht wird. Die
neutrale Faser ist die Gerade $c_0 + c_y y + c_z z = 0$.

Die \textbf{Extremwerte} sucht das Programm nur an den Eckpunkten des Querschnitts (Bögen alle
$10^\circ$ abgetastet). Eine lineare Funktion hat ihr Maximum immer in einer Ecke der konvexen
Hülle. Bei dünnwandigen Profilen sind das die Endpunkte der Mittellinien, also wieder die
Näherung aus der Vorlesung.

%---------------------------------------------------------------------
\section{Schubfluss in dünnwandigen Profilen}\label{sec:mathe-schub}

\subsection{Offene Profile}

Aus den dünnwandigen Linien baut das Programm einen \emph{Profilgraph}: Die Linien werden an
Kreuzungen, Berührpunkten und an den Schwerachsen geteilt, Punkte, die zusammenfallen, werden ein
Knoten. \textbf{Kreise und Bögen} kommen als Kette kurzer Sehnen dazu, geteilt an ihren Enden, an
Endpunkten anderer Elemente auf dem Bogen, an Kreuzungen mit geraden Linien und an den Schwerachsen,
dazwischen gleichmäßig mit höchstens $15^\circ$ je Sehne. Die Sehnen bestimmen nur die Knoten und
den Zusammenhang (Zellen, Symmetrie). Jede Sehne kennt ihren Bogen (Mittelpunkt, Radius, Winkel
$\varphi_1$ bis $\varphi_2$), und gerechnet wird auf dem Bogen: Länge $R\,|\varphi_2 - \varphi_1|$
statt der Sehne, Stützstellen auf dem Bogen. Dann läuft es los:

\begin{itemize}
  \item An jedem \textbf{freien Ende} ist der Schubfluss null, dort beginnt ein Pfad mit
        $S_y = S_z = 0$.
  \item Entlang eines Segments werden die statischen Momente aufintegriert,
        $S_z \mathrel{+}= (y - y_S)\,t\,\Delta s$ und $S_y \mathrel{+}= (z - z_S)\,t\,\Delta s$. Weil
        der Integrand auf einer geraden Linie linear ist, ist das mit der Mittelpunktsregel exakt.
        $S$ selbst ist quadratisch in $s$. Auf einem Bogen mit Mittelpunkt $(y_C, z_C)$ ist
        $ds = R\,d\varphi$, und das Integral ist geschlossen lösbar:
        \[
          \Delta S_z = t R \left[(y_C - y_S)\,\Delta\varphi + R\,(\sin\varphi_b - \sin\varphi_a)\right],
          \quad
          \Delta S_y = t R \left[(z_C - z_S)\,\Delta\varphi - R\,(\cos\varphi_b - \cos\varphi_a)\right]
        \]
        (für den Lauf gegen den Uhrzeigersinn, sonst mit umgekehrtem Vorzeichen).
  \item An einer \textbf{Verzweigung} gilt das Knotengleichgewicht: Kommen alle Äste bis auf einen an,
        startet der letzte mit der Summe der ankommenden statischen Momente.
\end{itemize}

Aus den statischen Momenten folgt der Schubfluss mit der gekoppelten Formel
\[
  \tau\,t = -\frac{Q_y\,(I_y S_z + I_{yz} S_y) + Q_z\,(I_z S_y + I_{yz} S_z)}{I_y I_z - I_{yz}^2} ,
\]
die für $I_{yz} = 0$ in das bekannte $\tau = -\frac{Q_z S_y}{I_y t}$ bzw.\ $-\frac{Q_y S_z}{I_z t}$
zerfällt. $\tau$ ist positiv in Laufrichtung $s$.

\subsection{Eine geschlossene Zelle: woher kommt \texorpdfstring{$q_0$}{q0}?}\label{sec:mathe-zelle}

Eine Zelle hat kein freies Ende, an dem der Schubfluss null sein müsste. Das Programm schneidet sie
deshalb an einer Stelle auf, rechnet wie beim offenen Profil und bekommt einen Verlauf, der sich
vom richtigen nur um einen \textbf{konstanten Umlaufanteil} $q_0$ unterscheidet. Diesen bestimmt
es auf zwei Wegen:

\textbf{(a) Symmetrie.} Liegt eine Symmetrieachse parallel zur Querkraft, ist der richtige Verlauf
antisymmetrisch und auf der Achse null. Das Programm schneidet die Zelle deshalb möglichst auf der
Achse auf. Klappt das nicht genau, verschiebt es den Umlaufanteil so, dass ein Punkt und sein
Spiegelpunkt entgegengesetzt gleiche Werte haben. Die Symmetrie selbst prüft es, indem es jede
Linie an 13 Stellen abtastet (jede Bogensehne an 3) und schaut, ob der gespiegelte Punkt wieder auf
einer Linie gleicher Dicke liegt. Weil Bögen auch an den Schwerachsen geteilt und die Stücke
gleichmäßig in Sehnen zerlegt werden, sind die Sehnen spiegelbildlicher Bögen selbst
spiegelbildlich.

\textbf{(b) Verträglichkeit.} Ohne passende Symmetrieachse kommt die Bedingung, dass sich die Zelle
unter einer Querkraft im Schubmittelpunkt nicht verdrillt:
\[
  \oint \frac{q}{G\,t}\,ds = 0 .
\]
Der Kniff ist, diese Bedingung nicht für die gerade eingegebene Querkraft zu lösen, sondern für
die statischen Momente selbst. Dazu sortiert man die Formel für den Schubfluss nach $S_z$ und $S_y$:
\[
  q = -\frac{a\,S_z + b\,S_y}{\Delta} , \qquad
  \begin{pmatrix} a \\ b \end{pmatrix}
  = \underbrace{\begin{pmatrix} I_y & I_{yz} \\ I_{yz} & I_z \end{pmatrix}}_{J}
    \begin{pmatrix} Q_y \\ Q_z \end{pmatrix} ,
\]
mit $\Delta = I_y I_z - I_{yz}^2 = \det J$. $a$, $b$, $G$ und $\Delta$ sind entlang der Zelle
konstant, also wird die Verträglichkeit zu
\[
  a \oint \frac{S_z}{t}\,ds + b \oint \frac{S_y}{t}\,ds = 0 .
\]
\emph{Warum $J$ regulär sein muss.} Die statischen Momente hängen nur vom Querschnitt ab, nicht von
der Last. Sollen sie für \emph{jede} Querkraft den richtigen Schubfluss liefern, muss die Gleichung
oben für jedes $(Q_y, Q_z)$ gelten. Weil $J$ regulär ist, gibt es zu jedem gewünschten Paar
$(a, b)$ eine Querkraft, die es erzeugt, nämlich $(Q_y, Q_z)^T = J^{-1}(a, b)^T$. Man darf also
$(a, b) = (1, 0)$ und $(a, b) = (0, 1)$ einsetzen und bekommt die beiden Bedingungen einzeln:
\[
  \oint \frac{S_z}{t}\,ds = 0 \quad\text{und}\quad \oint \frac{S_y}{t}\,ds = 0 .
\]
Wäre $J$ singulär, lägen alle Paare $(a, b)$, die überhaupt vorkommen können, auf einer Geraden durch
den Ursprung. Dann wüsste man nur, dass eine bestimmte Kombination der beiden Umlaufintegrale
verschwindet, und könnte die beiden statischen Momente nicht einzeln festlegen.

Regulär ist $J$ bei jedem echten Querschnitt: Nach der Cauchy-Schwarz-Ungleichung aus \textbf{HM1}
ist $I_{yz}^2 = \bigl(\int y z\,dA\bigr)^2 \le \int y^2\,dA \int z^2\,dA = I_z I_y$, also
$\Delta \ge 0$. Gleichheit gibt es nur, wenn die ganze Fläche auf einer Geraden durch den Schwerpunkt
liegt, und dann versagt schon die Biegeformel, in der $\Delta$ im Nenner steht.

Streng genommen braucht man die Regularität nur für diese Richtung des Schlusses. Dass zwei
einzeln verschobene statische Momente für jede Querkraft den richtigen Schubfluss liefern, folgt
schon daraus, dass $q_0 = -\oint (q_{\text{offen}}/t)\,ds \big/ \oint ds/t$ linear in $a$ und $b$
ist, also in einen Teil zu $S_z$ und einen zu $S_y$ zerfällt. Die Regularität sagt zusätzlich, dass
es keine andere Wahl der beiden Konstanten gibt. Deshalb sind auch die angezeigten $S$-Verläufe
eindeutig.

Das Programm verschiebt also jedes statische Moment in der Zelle um seine eigene Konstante $s_0$.
Aus $\oint (S + s_0)/t\,ds = 0$ folgt
\[
  s_0 = -\,\frac{\oint S/t\;ds}{\oint ds/t}
\]
(beide Integrale im Umlaufsinn gegen den Uhrzeigersinn). Damit stimmen nicht nur $\tau$, sondern
auch die angezeigten $S$-Verläufe. Integriert wird mit der Simpsonregel, die für das quadratische
$S/t$ exakt ist. Auf einem Bogen ist $S$ eine glatte Winkelfunktion, mit 16 Teilen je Sehne
von höchstens $15^\circ$ ist Simpson dort auf etwa $10^{-10}$ genau. Beim symmetrischen Profil kommt
so übrigens genau die antisymmetrische Lösung aus (a) heraus.

\subsection{Schubmittelpunkt}

Der Schubmittelpunkt ist der Punkt, durch den die Resultierende des Schubflusses geht. Das
Programm rechnet den Schubfluss für die Einheitslasten $Q_y = 1$ und $Q_z = 1$ (bei einer Zelle
inklusive $q_0$) und dessen Moment um den Schwerpunkt,
\[
  M_q = \int \left[(y - y_S)\,t_z - (z - z_S)\,t_y\right] q\,ds ,
\]
mit der Tangentenrichtung $(t_y, t_z)$ der Linie. Weil dasselbe Moment aus der Querkraft im
Schubmittelpunkt kommen muss, $M_q = (y_M - y_S)\,Q_z - (z_M - z_S)\,Q_y$, ist
\[
  y_M = y_S + M_q(Q_z = 1), \qquad z_M = z_S - M_q(Q_y = 1) .
\]
Der Integrand ist kubisch in $s$, die Simpsonregel also wieder exakt (auf Bögen, mit der
Bogentangente, wieder auf etwa $10^{-10}$). Die Momententabelle zeigt genau diese Summe, aufgeteilt
nach Elementen. Gegen die Handrechnung: offener Halbkreis $e = 4R/\pi$, D-Profil aus Halbkreis und
Platte und Kreisrohr stimmen auf etwa $10^{-9}$.

\subsection{Massive Querschnitte}

Für $\tau = \frac{Q\,S}{I\,b}$ teilt das Programm den Querschnitt in etwa 160 waagerechte (bzw.\
senkrechte) Streifen. Die Breite $b(z)$ kommt exakt aus allen Elementen: Löcher abgezogen, Polygone
über die Schnittpunkte der Kanten, Kreise über $2\sqrt{R^2 - \Delta^2}$, Kreisausschnitt und
-abschnitt als Kreis geschnitten mit den Halbebenen der Randstrahlen bzw.\ der Sehne. Die
Streifengrenzen liegen auf allen Knickstellen von $b$ (Ecken, Rechteckkanten, Kreisscheitel,
Lochränder) und auf der Schwerachse, dazwischen gleich weit. In jedem Streifen ist $b$ damit glatt,
und $S(z)$ wird mit drei Gauß-Punkten je Streifen von oben aufsummiert (an Knickstellen fein
unterteilt, wegen der Wurzel am Kreisscheitel).

$\tau$ wertet das Programm auf den Streifengrenzen aus, also auch genau am oberen und unteren Rand.
Dort ist $S = 0$, damit $\tau = 0$. Der kleine Rest, den die Quadratur am zweiten Rand übrig lässt,
wird anteilig nach der Fläche abgezogen. An Spitzen und am Kreisscheitel geht $b \to 0$, aber $S$
schneller, der Grenzwert $\tau \to 0$ wird direkt eingesetzt. An einem Breitensprung (Flansch und
Steg, Lochrand) gibt es zwei Stützstellen, eine mit der Breite darunter und eine mit der darüber,
der Verlauf zeigt also den Sprung. Früher wurde $\tau$ in der Streifenmitte ausgewertet, aber mit
$S$ bis zur Streifenunterkante: Der oberste Wert war deshalb nie null (beim Rechteck etwa
$2{,}5\,\%$ von $\tau_{max}$), und der Rand selbst kam gar nicht vor.

Wirken $Q_y$ und $Q_z$ gleichzeitig, ist das angezeigte $\tau_{max}$
$\sqrt{\tau_{y,max}^2 + \tau_{z,max}^2}$, also eine Abschätzung nach oben, denn die beiden Maxima
liegen meist nicht an derselben Stelle.

%---------------------------------------------------------------------
\section{Torsion und Verwölbung}\label{sec:mathe-torsion}

\textbf{Torsion:} Beim offenen Profil ist $I_T = \sum \frac{\xi}{3} L\,t^3$ und
$\tau = M_T\,t/I_T$. Bei einer Zelle kommt die Fläche $A_m$ der Mittellinie aus der Gaußschen
Trapezformel und $\oint ds/t = \sum L_i/t_i$, damit $I_T = 4A_m^2/\oint ds/t$ und
$\tau = M_T/(2A_m t)$. Läuft die Zelle über einen Bogen, nimmt die Trapezformel die Sehnen, und je
Sehne mit dem Winkel $\Delta\varphi$ kommt das Kreissegment
\[
  \Delta A = \frac{R^2}{2}\left(\Delta\varphi - \sin\Delta\varphi\right)
\]
zwischen Sehne und Bogen dazu (mit dem Vorzeichen des Umlaufs). $L_i$ ist dort die Bogenlänge
$R\,\Delta\varphi$. Beides ist exakt, die Sehnen sind also nur eine Hilfe für die Topologie.

\textbf{Zelle mit offenen Ästen} (z.\,B.\ ein Kasten mit überstehenden Gurten): Die Zelle trägt
nach Bredt das ganze $M_T$, und $I_T$ ist ihr Bredt-Wert. Daraus folgt die Verdrillung
$\vartheta = M_T/(G\,I_T)$, und die offenen Äste müssen sich genauso verdrillen wie die Zelle. Ein
Ast der Dicke $t$ hat dann am Rand $\tau = G\,\vartheta\,t = M_T\,t/I_T$. Streng genommen nehmen die
Äste damit zusätzlich das Moment $G\,\vartheta\,I_{T,\text{offen}}$ auf, und alle Schubspannungen
zusammen ergeben $M_T\,(1 + I_{T,\text{offen}}/I_{T,\text{Zelle}})$ statt $M_T$. Das ist eine
Näherung: $I_{T,\text{offen}}$ wächst mit $t^3$, $I_{T,\text{Zelle}}$ mit $A_m^2$, und der
Unterschied liegt bei üblichen Profilen im Promillebereich. Bei einem Kasten $200 \times 100$ mit
$t = 10$ und zwei $50$ langen Gurten derselben Dicke sind es etwa $0{,}1\,\%$. Genauso lässt Bredt
den kleinen offenen Anteil $\frac13 L\,t^3$ der Zellwände selbst weg. Exakt im Gleichgewicht wäre
$I_T = I_{T,\text{Zelle}} + I_{T,\text{offen}}$ mit einer Aufteilung von $M_T$ im Verhältnis der
beiden Steifigkeiten.

Ein
dünnwandiges Kreisrohr ist die Zelle mit $A_m = \pi R^2$ und $\oint ds/t = 2\pi R/t$, also
$I_T = 2\pi R^3 t$. Das Programm erkennt das Rohr aus einem einzelnen Kreis gesondert. Über die
Sehnenkette käme dasselbe heraus, ein Rohr aus zwei Halbkreisen rechnet so. Das Torsionsmoment
aus platzierten Kräften ist das Moment um den \emph{Schubmittelpunkt}, denn um den dreht sich der
Querschnitt.

\textbf{Verwölbung:} Auf einem geraden Segment ist der Fahrstrahlabstand $r_\perp$ linear, damit
ist die Verwölbung nach Formelsammlung geschlossen integrierbar:
\[
  u_x(s) = u_{x,0} + \frac{q_T}{G\,t}\,s - \vartheta\left[r_1 s + (r_2 - r_1)\frac{s^2}{2L}\right],
  \qquad q_T = \frac{M_T}{2A_m}, \quad \vartheta = \frac{M_T}{G\,I_T} ,
\]
wobei der Bredt-Anteil $q_T/(Gt)$ nur in der Zelle auftritt (mit Vorzeichen je nach Umlaufsinn).
Auf einem Bogen ist $r_\perp$ nicht linear, das Integral aber ebenfalls geschlossen lösbar. Mit
dem Pol $(y_M, z_M)$ und dem Winkel $\varphi$ ab dem Sehnenanfang $\varphi_1$ ist
\[
  \int r_\perp\,ds = R\left[(y_C - y_M)(\sin\varphi - \sin\varphi_1)
    - (z_C - z_M)(\cos\varphi - \cos\varphi_1) + R\,(\varphi - \varphi_1)\right] .
\]
Ein Kreisrohr verwölbt sich damit nicht: Mit dem Pol im Mittelpunkt ist $r_\perp = R$, und der
Bredt-Anteil hebt $R\,\vartheta$ genau auf.
Das Programm läuft von einem Knoten auf einer Symmetrieachse aus (dort ist $u_x$ bei symmetrischen
Profilen null) über den Profilgraph und setzt die Werte Segment für Segment fort. Beim Schließen der
Zelle muss es wieder beim Startwert ankommen, weil $\oint r_\perp\,ds = 2A_m$ ist. Das ist die
eingebaute Probe. Zum Schluss wird die Konstante so gewählt, dass $\int u_x\,t\,ds = 0$ ist, also
keine mittlere Längsverschiebung übrig bleibt.

%---------------------------------------------------------------------
\section{Kernfläche}\label{sec:mathe-kern}

Die Kernfläche ist die Menge aller Lastpunkte, deren neutrale Faser den Querschnitt nicht
schneidet. Gerechnet wird rückwärts:

\begin{enumerate}
  \item Das Programm sammelt die Eckpunkte des Querschnitts (Bögen abgetastet) und bildet ihre
        \textbf{konvexe Hülle}. Einbuchtungen spielen keine Rolle, denn eine Gerade, die die Hülle
        berührt, berührt auch den Querschnitt nicht innen.
  \item Jede Kante der Hülle ist eine mögliche \textbf{neutrale Faser}. Mit dem nach außen
        zeigenden Einheitsvektor $\vec n$ der Kante und ihrem Abstand $d$ vom Schwerpunkt gehört
        dazu der Lastpunkt
        \[
          \vec e = -\frac{1}{A\,d}
          \begin{pmatrix} I_z & -I_{yz} \\ -I_{yz} & I_y \end{pmatrix} \vec n .
        \]
        Für ein Rechteck und die obere Kante ist $\vec n = (0,1)$, $d = h/2$, und es kommt
        $e_z = -h/6$ heraus, das bekannte Kerndrittel.
  \item Die Kernpunkte aller Kanten, der Reihe nach verbunden, sind die Kernfläche. An Bögen legt
        das Programm alle $5^\circ$ eine Tangente an, beim Vollkreis 36 Stück.
\end{enumerate}

\textbf{Woher die Matrix kommt.} Die Formel für $\vec e$ ist die Formel für $\sigma_x$ aus
Abschnitt~\ref{sec:mathe-sigma}, rückwärts gelesen. Eine Normalkraft $N$ im Punkt
$\vec e = (e_y, e_z)$, gemessen vom Schwerpunkt, erzeugt $M_y = N e_z$ und $M_z = -N e_y$. In
$\sigma_x$ eingesetzt und nach $y$ und $z$ sortiert:
\[
  \sigma_x = \frac NA + \frac N\Delta\Bigl[\bigl(I_y e_y + I_{yz} e_z\bigr)\,y
                                       + \bigl(I_{yz} e_y + I_z e_z\bigr)\,z\Bigr]
           = \frac NA\Bigl(1 + \frac A\Delta\,\vec x^T J\,\vec e\Bigr) ,
\]
mit $\vec x = (y, z)$ und der Trägheitsmatrix $J = \begin{pmatrix} I_y & I_{yz} \\ I_{yz} & I_z
\end{pmatrix}$, $\Delta = \det J$. Die neutrale Faser ist $\sigma_x = 0$, also die Gerade
\[
  \vec x^T \Bigl(\frac A\Delta\,J\,\vec e\Bigr) = -1 .
\]
Die Kante mit dem nach außen zeigenden Einheitsvektor $\vec n$ und dem Abstand $d$ vom Schwerpunkt
ist die Gerade $\vec x^T \vec n = d$, oder genauso $\vec x^T(-\vec n/d) = -1$. Eine Gerade, die
nicht durch den Schwerpunkt geht, lässt sich nur auf eine Weise in der Form $\vec x^T \vec g = -1$
schreiben. Die neutrale Faser ist also genau dann die Kante, wenn die beiden Vektoren gleich sind:
\[
  \frac A\Delta\,J\,\vec e = -\frac{\vec n}{d}
  \qquad\Longrightarrow\qquad
  \vec e = -\frac{1}{A\,d}\,\bigl(\Delta\,J^{-1}\bigr)\,\vec n ,
  \qquad
  \Delta\,J^{-1} = \begin{pmatrix} I_z & -I_{yz} \\ -I_{yz} & I_y \end{pmatrix} .
\]
Die Matrix in der Formel ist also die Inverse der Trägheitsmatrix ohne den Vorfaktor $1/\det J$
(Diagonale vertauscht, Nebendiagonale mit anderem Vorzeichen). Auch hier braucht man, dass $J$
regulär ist, siehe Abschnitt~\ref{sec:mathe-zelle}.

\textbf{Warum man die Kernpunkte geradlinig verbinden darf.} Dreht man die neutrale Faser um eine
Ecke $\vec P$ der Hülle von einer Kante zur nächsten, gehen alle diese Geraden durch $\vec P$, also
gilt $\vec P^T(\vec n/d) = 1$. Der Vektor $\vec n/d$ wandert damit auf einer Geraden, und weil
$\vec e$ linear von $\vec n/d$ abhängt, wandert auch der Lastpunkt auf einer Geraden. Jeder Ecke der
Hülle entspricht so eine gerade Kante der Kernfläche, jeder Kante der Hülle eine Ecke.

Im Hauptachsensystem fällt die Kopplung weg ($I_{\eta\zeta} = 0$). Schreibt man die Kante mit
ihren Achsenabschnitten als $\eta/a_\eta + \zeta/a_\zeta = 1$, ist $\vec n/d = (1/a_\eta,
1/a_\zeta)$, und aus der Formel wird die bekannte Form $\eta_K = -i_\zeta^2/a_\eta$,
$\zeta_K = -i_\eta^2/a_\zeta$ mit $i^2 = I/A$. Deshalb beschriftet das
Programm die umhüllenden Geraden mit ihren Abschnitten auf den Hauptachsen, gemessen vom
Schwerpunkt, und nicht mit Koordinaten im gezeichneten KOS.

%---------------------------------------------------------------------
\section{Spannungskreis: vom Schnitt zum Zustand}\label{sec:mathe-kreis}

\textbf{Lineare Gleichungen.} Mit $\sigma_m$, $R_c$ und $R_s$ als Unbekannten ist jede Angabe mit
bekanntem Winkel eine Zeile
\[
  \sigma = \sigma_m + R_c\cos 2\alpha + R_s\sin 2\alpha, \qquad
  \tau = -R_c\sin 2\alpha + R_s\cos 2\alpha .
\]
\textbf{Punkte ohne Winkel} erfüllen $(\sigma - \sigma_m)^2 + \tau^2 = R_c^2 + R_s^2$. Zieht man
diese Gleichung für zwei solche Punkte voneinander ab, fallen alle quadratischen Terme weg und es
bleibt eine lineare Gleichung für $\sigma_m$ allein:
$2(\sigma_2 - \sigma_1)\,\sigma_m = \sigma_2^2 + \tau_2^2 - \sigma_1^2 - \tau_1^2$. Bei mehreren
Punkten ohne Winkel zieht das Programm die Gleichung des ersten Punktes von jeder anderen ab. Jede
Differenz ist linear und kommt zu den linearen Zeilen. Quadratisch bleibt damit nur noch eine
Gleichung, die des ersten Punktes selbst: Die Information der anderen Punkte steckt vollständig in
den Differenzen.

\textbf{Rang.} Die linearen Zeilen löst das Programm mit Gauß-Jordan und Spaltenpivotsuche und
zählt dabei, wie viele Pivotelemente es findet (Toleranz relativ zum größten Eintrag). Bleibt in
einer Zeile ohne Pivot rechts etwas übrig, widersprechen sich die Angaben.

\begin{itemize}
  \item \textbf{Rang 3:} Der Zustand ist eindeutig. Ein Punkt ohne Winkel muss dann nur noch auf
        dem Kreis liegen (Probe).
  \item \textbf{Rang 2 und ein Punkt ohne Winkel:} Die Lösungen der linearen Zeilen bilden eine
        Gerade $\vec x = \vec x_0 + t\,\vec d$ im $(\sigma_m, R_c, R_s)$-Raum. In die quadratische
        Gleichung eingesetzt, wird daraus $a t^2 + b t + c = 0$ mit keiner, einer oder zwei Lösungen.
        Genau daher kommt \emph{zwei Zustände möglich}.
  \item \textbf{Weniger:} Es fehlen Angaben.
\end{itemize}

Schnitte ohne Winkel mit nur $\sigma$ oder nur $\tau$ geben keine Gleichung, aber eine
Ungleichung: $|\sigma - \sigma_m| \le R$ bzw.\ $|\tau| \le R$. Damit sortiert das Programm von zwei
möglichen Zuständen den aus, auf dessen Kreis der Wert nicht passt.

\textbf{Die Spannungsscheibe} ist das Vieleck, dessen Kanten alle einen Kreis mit Radius $r$
berühren. Jeder Schnitt $i$ mit dem Normaleneinheitsvektor $\vec n_i$ gibt eine Kante: die Gerade
$\vec n_i \cdot \vec x = r$, senkrecht zu $\vec n_i$ im Abstand $r$ vom Mittelpunkt. Zur Scheibe
gehört alles auf der Seite des Mittelpunktes, also die Halbebene $\vec n_i \cdot \vec x \le r$. Die
Scheibe ist die Menge der Punkte, die in allen diesen Halbebenen zugleich liegen. Liegen zwei
benachbarte Normalen $160^\circ$ oder mehr auseinander, wäre die Scheibe
offen (oder eine sehr spitze Nadel). Dann ergänzt das Programm den $x$- oder $y$-Schnitt, dessen
Normale am nächsten an der Mitte der größten Lücke liegt, und das so oft, bis die Lücke kleiner ist.
 eingebunden.

\part{Die Mathematik hinter den Features}\label{sec:Mathe}

Vorab will ich erwähnen, dass Teile dieses Abschnitts nicht meine eigene Arbeit sind. Gerade wenn
es darum geht, das Programm effizient und robust zu schreiben, oder wenn sich das Skript Methoden der
weiterführenden Mathematik bzw.\ Numerik bedient, stammen Ideen von KI, z.\,B.\ die
$LDL^T$-Zerlegung (ich bin selbst erst im 2.~Semester, deswegen halten sich meine mathematischen
Kenntnisse noch in Grenzen). Das gilt übrigens auch für große Teile der Umsetzung. Aus diesem Grund
halte ich die Erklärungen zu diesen Aspekten hier etwas kürzer.

Nichtsdestotrotz ist dieser Teil für alle, die wissen wollen, \emph{warum} das Programm so rechnet, wie es rechnet.
Zum Benutzen braucht ihr ihn nicht. Grundlegende Sachen wie Hebelarme, Steiner und Gleichgewicht seien hier nicht weiter erläutert. Hier
werden die Stellen behandelt, an denen das Programm mehr macht als die Vorlesung oder es anders macht,
weil ein Rechner eben nicht wie ein Mensch rechnet. Ich habe versucht, zu jedem Punkt dazuzuschreiben,
was er für euch in der Klausur bedeutet.

Ein paar Konventionen, die überall gelten:

\begin{itemize}
  \item Im \textbf{Tragwerksskript} zeigt $x$ nach \emph{rechts} und $y$ nach \emph{unten}. Drehwinkel,
        Momente und Lastwinkel zählen trotzdem so, wie ihr sie seht: gegen den Uhrzeigersinn
        positiv.
  \item Jeder Stab hat seine eigene Koordinate $x \in [0, L]$ vom Anfangs- zum Endknoten, die
        Tangente $\vec t = (\Delta x, \Delta y)/L$ und die Normale $\vec n = (-\Delta y, \Delta x)/L$.
        Die Schnittgrößen folgen der TM1-Konvention mit gestrichelter Faser auf der Seite von
        $\vec n$.
  \item Im \textbf{Querschnittsskript} gilt die Konvention aus dem Gross (TM2):
        $I_y = \int z^2\,dA$, $I_z = \int y^2\,dA$ und $I_{yz} = -\int yz\,dA$, also mit Minus.
\end{itemize}

%---------------------------------------------------------------------
\section{Das Weggrößenverfahren im Programm}\label{sec:mathe-wgv}

Alles, was ihr im Tragwerksskript seht (Verläufe, Auflagerkräfte, Verformungen,
Explosionszeichnung), kommt aus einer einzigen Rechnung: dem Weggrößenverfahren, also genau dem
allgemeinen WGV aus \textbf{Statik}. Unbekannt sind die Verschiebungen der Knoten, nicht die
Kräfte.

\textbf{Unbekannte:} Jeder Knoten hat drei Freiheitsgrade, $u_\xi$, $u_\eta$ und $\varphi$. Bei
einem ungedrehten Lager ist $\xi = x$ und $\eta = y$, bei einem gedrehten Lager zeigen sie in die
Lagerrichtungen. Ein Lager \emph{streicht} einfach den gehaltenen Freiheitsgrad.

\textbf{Stabsteifigkeit:} Für jeden Stab steht die bekannte $6 \times 6$-Matrix aus $EA/L$,
$12EI/L^3$, $6EI/L^2$, $4EI/L$ und $2EI/L$. Sie verknüpft die sechs Endverschiebungen
$(u_a, w_a, \varphi_a, u_b, w_b, \varphi_b)$ im lokalen System des Stabes mit den sechs
Endkräften. Zwei Drehungen bringen sie ins globale System: die Stabrichtung und (bei einem
gedrehten Lager) die Lagerrichtung,
\[
  K_e = T_K\,T^T\,k\,T\,T_K^T .
\]
So ist ein gedrehtes Lager kein Sonderfall, sondern nur eine andere Basis am Knoten. Matrix,
Transformation und Drehrichtungen sind dieselben wie in der WGV-Formelsammlung aus Statik, die
FEM-Formelsammlung hat dieselben Elemente, zählt die Querrichtung aber nach oben. Die
Unterschiede stehen in Abschnitt~\ref{sec:wgv-vorzeichen}.

\textbf{Lasten:} Knotenlasten gehen direkt in den Lastvektor (vorher um den Lastwinkel gedreht).
Streckenlasten gehen als Festeinspannkräfte ein. Das sind die Endkräfte, die ein
beidseitig eingespannter Stab unter der Last hätte, z.\,B.\ $\mp qL/2$ und $\pm qL^2/12$ bei
Gleichlast (die Volleinspannmomente aus Tafel~A9 der Formelsammlung, mit denselben Vorzeichen).
Im Gesamtsystem stehen sie mit umgekehrtem Vorzeichen, also als Knotenersatzlasten. Für Trapezlasten stehen die entsprechenden Formeln im Code, Lastfunktionen wie
\code{x\^{}2} integriert das Programm mit den Hermite-Formfunktionen
$N_1 = 1 - 3\xi^2 + 2\xi^3$, $N_2 = L(\xi - 2\xi^2 + \xi^3)$ usw.\ (arbeitsäquivalente
Knotenlasten, Abschnitt~\ref{sec:mathe-festeinspann}). Temperatur wirkt wie eine Last $EA\,\varepsilon_T$ bzw.\ $EI\,\kappa_T$ an den
Stabenden. Globale Linienlasten $p_X(x)$, $p_Y(x)$ werden vorher in Quer- und Längsanteil zerlegt,
bei einer projizierten Last ist $x$ die globale Koordinate ab dem Stabanfang, $x_{global} = f\,x$
mit dem Projektionsfaktor $f = |\Delta y|/L$ bzw.\ $|\Delta x|/L$, und je Stablänge wirkt
$f\,p(f\,x)$.

\textbf{Lösen:} Das Gesamtsystem $K\,u = F$ wird mit einer $LDL^T$-Zerlegung gelöst. Wird dabei
ein Diagonalelement $D_j$ (fast) null, fehlt dem System in dieser Richtung jede Steifigkeit. Es
ist ein Mechanismus, und ihr bekommt \emph{System ist kinematisch!}.

\textbf{Zurückrechnen:} Aus den Knotenverschiebungen folgen die Endkräfte jedes Stabes,
$s = k\,v + p$ (Steifigkeit mal Endverschiebung plus Festeinspannkraft), und aus dem Anfangswert
durch Integrieren der Lasten die Verläufe, z.\,B.\
$M(x) = M_a + Q_a x - m\,x - \int_0^x q(t)\,(x-t)\,dt$. Das ist exakt, nicht genähert. Deshalb
braucht jeder Bereich zwischen Lasten, Lagern und Gelenken nur einen einzigen Stab.

\textbf{Auflagerkräfte:} Die Summe der Stabendkräfte an einem Knoten minus der Knotenlast ist die
Kraft, die Lager und Federn auf den Knoten ausüben. Davon zeigt das Programm
$R_\xi = \vec R\cdot\vec\xi$ und $R_\eta = -\vec R\cdot\vec\eta$, deshalb zeigt $R_\eta$ nach oben,
während die Knotenverschiebung \code{ü\_y} auf derselben Menüseite nach unten positiv ist.

\subsection{Festeinspannkräfte beliebiger Lasten}\label{sec:mathe-festeinspann}

Für das Gesamtsystem braucht das Programm bzw. das WGV von einer Streckenlast nur, was sie an den Stabenden
bewirkt, die Festeinspannkräfte. Für Gleich- und Trapezlasten stehen diese In Tafeln oder in der FOS, für eine
Lastfunktion wie $q(x) = x^2$ oder $q_0 \sin(\pi x/L)$ nicht. Das Programm rechnet sie deshalb für
jede Lastfunktion mit den Hermite-Funktionen aus, und zwar exakt, nicht genähert.

\textbf{Die Hermite-Funktionen.} Im lokalen System läuft $x$ von $0$ bis $L$, $\xi = x/L$, $w$ zeigt
zur gestrichelten Faser und $\varphi = -w'$ gegen den Uhrzeigersinn. Die vier Funktionen sind
\begin{align*}
  N_1 &= 1 - 3\xi^2 + 2\xi^3       && w(0) = 1, \\
  N_2 &= L\,(\xi - 2\xi^2 + \xi^3) && w'(0) = 1, \text{ also } \varphi_a = -1, \\
  N_3 &= 3\xi^2 - 2\xi^3           && w(L) = 1, \\
  N_4 &= L\,(-\xi^2 + \xi^3)       && w'(L) = 1, \text{ also } \varphi_b = -1,
\end{align*}
und die jeweils drei anderen Endwerte sind null. Jede ist eine kubische Parabel, löst also
$EI\,w'''' = 0$ exakt und stellt somit eine homogene Lösung der DGL dar.
Gleichzeitig erfüllt sie die vier von ihr geforderten Randbedingungen. Damit ist sie genau die Biegelinie des
beidseitig eingespannten Stabes, dem man diese eine Einheitsverformung aufzwingt.

Das WGV zählt die vier Endverschiebungen eines Stabes aber als $v = (w_a, \varphi_a, w_b, \varphi_b)$,
mit $\varphi$ gegen den Uhrzeigersinn positiv. Die Einheitsverformung zum Freiheitsgrad $j$ in
dieser Zählung heißt im Folgenden $S_j$ (S für Formfunktion, engl.\ \emph{shape}):
\[
  S_1 = N_1, \qquad S_2 = -N_2, \qquad S_3 = N_3, \qquad S_4 = -N_4 .
\]
Bei den Drehungen kehrt sich das Vorzeichen um, weil $N_2$ und $N_4$ zu $\varphi = -1$ gehören.

\textbf{Beweisidee mit dem Satz von Betti.} Die folgende Herleitung soll die Idee zeigen und
ist kein formal vollständiger Beweis. So nennen wir alle $F_i$ Kräfte und alle Verschiebungen Wege,
obwohl bei den Drehungen Momente und Winkel gemeint sind, und die Einheitsverschiebung \enquote{$1$}
hat eigentlich eine Einheit (m bzw.\ rad). Das formal sauber mitzuschleppen ändert am Ergebnis
nichts, macht die Rechnung aber deutlich schwerer zu lesen.

Die Festeinspannkräfte $F_i$ sind die Kräfte und
Momente, die die Einspannungen auf den belasteten Stab ausüben, jeweils in positiver Richtung des
Freiheitsgrades (das $p$ in $s = k\,v + p$). Man schneidet den Stab frei. Die Einspannungen sind
dann nur noch Kräfte an den Stabenden, und man hat \emph{ein} System mit zwei Lastgruppen:

\begin{itemize}
  \item Lastgruppe~1: die Last $q(x)$ und die Endkräfte $F_j$, die die Enden festhalten (gesucht).
        Die Biegelinie ist $w_1$ mit $w_1 = w_1' = 0$ an beiden Enden, alle Endverschiebungen sind
        also $v_j^{(1)} = 0$.
  \item Lastgruppe~2: nur die Endkräfte $k_{ji}$, die den Freiheitsgrad $i$ um $1$ verschieben und
        alle anderen festhalten. Die Biegelinie ist $S_i(x)$, die Endverschiebungen sind
        $v_j^{(2)} = \delta_{ij}$ (Kronecker-Delta: $1$ für $j = i$, sonst $0$).
\end{itemize}

Beide Gruppen sind für sich im Gleichgewicht, deshalb stört es nicht, dass der freigeschnittene
Stab nicht gelagert ist. Nach Betti leistet Lastgruppe~1 auf den Wegen von Lastgruppe~2 dieselbe
Arbeit wie Lastgruppe~2 auf den Wegen von Lastgruppe~1:
\[
  \int_0^L q(x)\,S_i(x)\,dx + F_j\,v_j^{(2)} = k_{ji}\,v_j^{(1)} .
\]
Über das doppelte $j$ wird auf beiden Seiten summiert (Einsteinsche Summenkonvention). Setzt man
die Verschiebungen ein, bleibt links $F_j\,\delta_{ij} = F_i$ und rechts null. Die innere Arbeit
fehlt dabei nicht, sie ist der Grund für die Gleichheit:
\[
  W_a^{1 \to 2} = \int_0^L M_1\,\kappa_2\,dx = \int_0^L \frac{M_1 M_2}{EI}\,dx
  = \int_0^L M_2\,\kappa_1\,dx = W_a^{2 \to 1} .
\]
Das erste Gleichheitszeichen ist der Arbeitssatz (äußere gleich innere Arbeit) mit zwei
verschiedenen Zuständen: die Kräfte und Momente $M_1$ aus Lastgruppe~1, die im Gleichgewicht sind,
und die Verformung $\kappa_2$ aus Lastgruppe~2, die verträglich ist. Einer der beiden Zustände
spielt dabei die Rolle des virtuellen. Im PvV aus \textbf{TM1} ist es die Verschiebung (hier die
von Lastgruppe~2), beim Prinzip der virtuellen Kräfte aus \textbf{Statik} wäre es die Kraftgruppe.
Der Satz selbst gilt in beiden Lesarten, solange der eine Zustand im Gleichgewicht und der andere
verträglich ist. Das letzte Gleichheitszeichen ist derselbe Satz mit vertauschten Rollen: Kräfte
aus Lastgruppe~2, Verformung aus Lastgruppe~1. In der Mitte steht nur $\kappa = M/EI$. Betti
schreibt die beiden Enden dieser Kette hin. Die Kette enthält nur die Biegearbeit. Das Programm rechnet nach Euler-Bernoulli, der Stab ist
also schubstarr ($\gamma = 0$), und die Querkraft leistet keine innere Arbeit. Nach Timoschenko käme
in der Mitte noch $\int_0^L Q_1 Q_2/(GA_s)\,dx$ dazu. Betti gilt dann weiterhin, nur wären die $S_i$
andere Kurven, die auch von $GA_s$ abhängen.

Dass die gemischte innere Arbeit hier null ist, obwohl die Schnittgrößen aus Lastgruppe~2 auf
der Biegung aus Lastgruppe~1 lokal sehr wohl Arbeit leisten, sieht man auch direkt: $M_2$ ist ohne
Streckenlast linear, also $M_2'' = 0$, und zweimal partiell integriert bleibt
\[
  \int_0^L M_2\,\kappa_1\,dx = -\int_0^L M_2\,w_1''\,dx
  = \Bigl[-M_2\,w_1' + M_2'\,w_1\Bigr]_0^L - \int_0^L M_2''\,w_1\,dx = 0 ,
\]
weil $w_1$ und $w_1'$ an beiden Enden verschwinden. Also gilt für jede Lastfunktion
\[
  F_i = -\int_0^L q(x)\,S_i(x)\,dx .
\]
Mit $S_1 = N_1$, $S_2 = -N_2$, $S_3 = N_3$ und $S_4 = -N_4$ ergibt das genau die Formeln im
Programm:
\begin{align*}
  Q_a &= -\int_0^L q\,N_1\,dx, & M_a &= \int_0^L q\,N_2\,dx, \\
  Q_b &= -\int_0^L q\,N_3\,dx, & M_b &= \int_0^L q\,N_4\,dx.
\end{align*}

\textbf{Anschaulich} ist $S_i(x)$ die Einflusslinie der Festeinspannkraft $i$ (Satz von
Müller-Breslau): Eine Einzelkraft $P$ an der Stelle $a$ erzeugt $F_i = -P\,S_i(a)$, und das Integral
summiert nur die vielen kleinen Einzellasten $q\,dx$ auf.

\textbf{Kontrolle mit Gleichlast:} $\int_0^L N_1\,dx = L/2$, $\int_0^L N_2\,dx = L^2/12$ und
$\int_0^L N_4\,dx = -L^2/12$. Damit ist $Q_a = Q_b = -qL/2$ (die Last zeigt nach unten, die
Einspannung drückt dagegen) und $M_a = qL^2/12$, $M_b = -qL^2/12$, genau die Tafelwerte. Auch das
Vorzeichen von $M_a$ passt: Ohne Einspannung würde sich das linke Ende im Uhrzeigersinn drehen, die
Einspannung hält gegen den Uhrzeigersinn dagegen.

\textbf{Die anderen Lastarten.}
\begin{itemize}
  \item Ein \textbf{Streckenmoment} $m(x)$ (gegen den Uhrzeigersinn positiv) leistet Arbeit auf der
        Verdrehung, $\int m\,\varphi\,dx = -\int m\,w'\,dx$. Mit Betti wird daraus
        $F_i = \int_0^L m\,S_i'\,dx$, also $Q_a = \int m\,N_1'\,dx$, $M_a = -\int m\,N_2'\,dx$ und
        entsprechend am Ende $b$. Bei konstantem $m$ ist $\int N_1'\,dx = N_1(L) - N_1(0) = -1$ und
        $\int N_2'\,dx = 0$: $Q_a = -m$, $Q_b = +m$ und keine Endmomente. Das ist das Kräftepaar aus
        Abschnitt~\ref{sec:wgv-vorzeichen} und der Grund, warum das Programm vom $m\,l/2$ der
        WGV-Tafel abweicht.
  \item Für eine \textbf{Längslast} $n(x)$ sind die Formfunktionen linear, $1 - \xi$ und $\xi$ (sie
        lösen $EA\,u'' = 0$ exakt), und es gilt $N_a = -\int n\,(1-\xi)\,dx$, $N_b = -\int n\,\xi\,dx$.
  \item \textbf{Temperatur} ist keine Last, sondern eine eingeprägte Verzerrung. Die
        Festeinspannkräfte stehen direkt da: $N_a = EA\,\varepsilon_T$, $N_b = -EA\,\varepsilon_T$,
        $M_a = EI\,\kappa_T$, $M_b = -EI\,\kappa_T$.
\end{itemize}

\textbf{Zwei Wege im Code.} Für konstante Lasten und Trapezlasten
$q = p_a(1-\xi) + p_b\,\xi$ sind die Integrale schon ausgerechnet:
\begin{align*}
  Q_a &= -\tfrac{L}{60}\,(21p_a + 9p_b),   & M_a &= \tfrac{L^2}{60}\,(3p_a + 2p_b), \\
  Q_b &= -\tfrac{L}{60}\,(9p_a + 21p_b),   & M_b &= -\tfrac{L^2}{60}\,(2p_a + 3p_b).
\end{align*}
Hängt die Last von $x$ ab (auch bei $m(x)$, $n(x)$ und den globalen Linienlasten), rechnet das CAS
genau die Integrale oben aus. Ein Beispiel ist $q(x) = q_0 \sin(\pi x/L)$. Für die Querkraft
braucht man $\int_0^L q\,N_1\,dx$, und das geht ohne Rechnung in drei Schritten:
\begin{enumerate}
  \item $N_1 + N_3 = 1$: Schiebt man beide Stabenden um $1$ nach unten, verschiebt sich der ganze
        Stab um $1$. Also ist $\int_0^L q\,N_1\,dx + \int_0^L q\,N_3\,dx = \int_0^L q\,dx$.
  \item $N_3$ ist $N_1$ gespiegelt, $N_3(\xi) = N_1(1-\xi)$. Ist die Last symmetrisch zur
        Stabmitte, so wie der Sinus, sind beide Integrale deshalb gleich groß.
  \item Also ist jedes die Hälfte: $\int_0^L q\,N_1\,dx = \frac12 \int_0^L q\,dx
        = \frac12 \cdot \frac{2q_0L}{\pi} = \frac{q_0L}{\pi}$, und $Q_a = -q_0L/\pi$.
\end{enumerate}
Anschaulich: Bei symmetrischer Last trägt jede Einspannung die halbe Last. Für das Moment
braucht man $\int_0^1 \sin(\pi\xi)\,(\xi - 2\xi^2 + \xi^3)\,d\xi = 2/\pi^3$, also
$M_a = 2q_0L^2/\pi^3 \approx 0{,}0645\,q_0L^2$.

\textbf{Warum das exakt ist.} Betti gilt für jede Last, und die $S_i$ sind exakte Lösungen der
homogenen Gleichung. Also sind die Festeinspannkräfte exakt, damit auch die Knotenverschiebungen aus
$K\,u = F$ und die Stabendkräfte $s = k\,v + p$. Nicht exakt wäre nur die Hermite-Interpolation
$\sum_j v_j S_j(x)$ im Feld, denn die kennt die Last nicht. Die benutzt das Programm aber nie als
Ergebnis (siehe unten).

\textbf{Zwei Zerlegungen derselben Biegelinie.} Die Biegelinie eines Stabes ist immer
Partikulärlösung plus homogene Lösung. Welche Partikulärlösung man nimmt, ist aber nicht
festgelegt: Zwei Partikulärlösungen unterscheiden sich immer um eine kubische Parabel, also um eine
homogene Lösung. Das Programm benutzt zwei davon, und beide beschreiben dieselbe Linie.

\begin{itemize}
  \item \emph{Die Anfangswert-Lösung} ist das Vierfachintegral der Querlast,
        \[
          EI\,w_p(x) = \int_0^x q(t)\,\frac{(x-t)^3}{6}\,dt ,
        \]
        mit $w_p = w_p' = w_p'' = w_p''' = 0$ bei $x = 0$, so wie ihr sie in \textbf{TM2} aufstellt.
        Dazu kommt die homogene Lösung, deren vier Konstanten die Randbedingungen erfüllen. Das
        Rand-LGS schreibt sie als $c_1 x^3/6 + c_2 x^2/2 + c_3 x + c_4$ und holt die Konstanten aus
        allen Rand- und Übergangsbedingungen (Abschnitt~\ref{sec:mathe-randlgs}). Der alte Export
        \code{v\_EI\_stab} schreibt sie mit den Hermite-Funktionen, und die Koeffizienten sind die
        Endverschiebungen aus dem WGV:
        \[
          EI\,w = EI\,w_p + EI w_a\,N_1 + EI w_a'\,N_2 + \bigl(EI w_b - EI w_p(L)\bigr)\,N_3
                  + \bigl(EI w_b' - EI w_p'(L)\bigr)\,N_4 .
        \]
        $w_p$ bringt am Ende $b$ schon eigene Werte mit, die ziehen die letzten beiden Koeffizienten
        wieder ab, damit die Summe dort genau die Endverschiebung hat. Die Koeffizienten sind
        Durchbiegung und Neigung, so wie ihr die Randbedingungen in \textbf{TM2} aufstellt. In der
        Zählung des WGV ist $w_a = v_1$, $w_a' = -\varphi_a = -v_2$, $w_b = v_3$ und $w_b' = -v_4$.
  \item \emph{Die Lösung des eingespannten Stabes} ist
        \[
          w_0(x) = w_p(x) - w_p(L)\,N_3(x) - w_p'(L)\,N_4(x) .
        \]
        Auch sie löst $EI\,w'''' = q$, hat aber an \emph{beiden} Enden $w_0 = w_0' = 0$: Das ist die
        Biegelinie des beidseitig eingespannten Stabes unter der Last. Mit ihr zerfällt die Linie in
        \[
          w(x) = w_0(x) + \sum_j v_j\,S_j(x) ,
        \]
        und die Koeffizienten des homogenen Teils sind genau die Endverschiebungen $v_j$, also die
        Unbekannten des WGV. Diese Zerlegung ist die Idee hinter $s = k\,v + p$: Der homogene Teil
        $\sum_j v_j S_j$ liefert die Endkräfte $k\,v$. $w_0$ verschiebt die Stabenden nicht
        ($w_0 = w_0' = 0$ an beiden Enden), stört also die Verträglichkeit an den Knoten nicht und
        liefert nur Endkräfte, und das sind genau die Festeinspannkräfte $p$. Mehr braucht das WGV
        von der Last nicht.
\end{itemize}

Bei Gleichlast ist $EI\,w_p = qx^4/24$ mit $EI\,w_p(L) = qL^4/24$ und $EI\,w_p'(L) = qL^3/6$, und
daraus wird
\[
  EI\,w_0 = \frac{qL^4}{24}\bigl(\xi^4 - (3\xi^2 - 2\xi^3) - 4(\xi^3 - \xi^2)\bigr)
          = \frac{q\,x^2 (L-x)^2}{24},
\]
die bekannte Biegelinie des eingespannten Balkens. Ihr Moment am linken Ende ist
$-EI\,w_0''(0) = -qL^2/12$ (Zug oben), und die Einspannung hält mit $M_a = +qL^2/12$ gegen den
Uhrzeigersinn dagegen. Der Satz von Betti erspart es, $w_0$ überhaupt auszurechnen: Die
Festeinspannkräfte kommen direkt aus den Integralen über $q\,S_i$.

Die fertige Biegelinie ist in beiden Zerlegungen dieselbe. Das sieht man, wenn man die Formel des
alten Exports umsortiert. Mit $w_a'\,N_2 = -v_2\,N_2 = v_2\,S_2$ und genauso am Ende $b$ wird aus ihr
\[
  EI\,w = \underbrace{EI\,\bigl(w_p - w_p(L)\,N_3 - w_p'(L)\,N_4\bigr)}_{EI\,w_0}
          + EI \sum_j v_j\,S_j ,
\]
also Buchstabe für Buchstabe die zweite Zerlegung, nur anders gruppiert. Ausgerechnet wird $w_0$
dabei nie. Für die Knotengleichungen braucht das WGV nur seine Endkräfte $p$, und die kommen nach
Betti direkt aus den Integralen über $q\,S_i$. Die Verläufe rechnet das Programm danach aus den
exakten Endkräften (Gleichgewicht) bzw.\ mit der Anfangswert-Lösung plus Hermite-Anteil, wie im
ersten Punkt. Die zweite Zerlegung dient hier also nur der Erklärung, warum $s = k\,v + p$ exakt ist.

\subsection{Gelenke und die Knotenverdrehung}\label{sec:mathe-gelenke}

Hier kommt die Warnung aus Abschnitt~\ref{sec:Eingabe_Knoten} her. Ein \textbf{M-Gelenk} wird nicht
herauskondensiert, sondern bekommt einen eigenen Freiheitsgrad: Das Stabende am Gelenk hat seine
eigene Verdrehung. Zu diesem Freiheitsgrad gehört eine eigene Gleichgewichtsgleichung, und an ihm
greift kein äußeres Moment an. Somit ist das Gelenkmoment null. Für die angezeigte Knotenverdrehung
heißt das:

\begin{itemize}
  \item \textbf{Angelenk} (Gelenk am Stabende, andere Stäbe biegesteif am Knoten): Die
        Knotenverdrehung ist die Verdrehung der \emph{biegesteif} angeschlossenen Stäbe. Der Stab
        mit dem Gelenk dreht sich anders, seine Endverdrehung steht im Stabmenü auf Seite~7
        ($\varphi_a$, $\varphi_b$), der Knick im Gelenk ist die Differenz.
  \item \textbf{Vollgelenk} (Knoten gelenkig): Der Knoten hat \emph{gar keine} eigene Verdrehung,
        jedes Stabende dreht sich für sich. Die $0$ im Knotenmenü ist nur ein Platzhalter, die
        echten Verdrehungen stehen wieder auf Seite~7 der Stäbe.
  \item \textbf{Nur Angelenke an einem Knoten}, aber kein Vollgelenk und kein Lager, das die
        Drehung hält: Dann hat der Knoten eine Verdrehung, an der nichts hängt. Das Programm
        meldet \emph{kinematisch}. Deshalb der Stolperstein im Anhang: In diesem Fall braucht ihr
        das Vollgelenk.
\end{itemize}

\textbf{N- und Q-Gelenke} werden dagegen herauskondensiert (siehe Dallmann Baustatik 3): Im Stab wird der entsprechende Anteil
der Steifigkeit gestrichen, die Festeinspannkräfte werden statisch gleichwertig
umverteilt, und nach der Rechnung wird die Klaffung im Gelenk wieder ausgerechnet. Die
Endverschiebungen auf Seite~7 enthalten sie. Was dabei gestrichen wird, betrifft immer nur die
freigegebene Richtung:

\begin{itemize}
  \item \textbf{N-Gelenk:} Die Dehnsteifigkeit fällt weg. Am Gelenk ist $N = 0$, im Stab bleibt
        nur der Anteil aus Längslasten. Die Biegung ist davon nicht betroffen.
  \item \textbf{Q-Gelenk:} Am Gelenk ist $Q = 0$, ohne Querlast also im ganzen Stab, und wegen
        $M' = Q$ bleibt von der Biegesteifigkeit nur ein konstantes Moment übrig. Die Normalkraft
        trägt der Stab ganz normal weiter.
\end{itemize}

%---------------------------------------------------------------------
\section{Bögen: Teilstäbe und Schnittgrößen über den Winkel}\label{sec:mathe-bogen}

\subsection{Geometrie und Teilstäbe}

Ein Bogen ist ein Kreisstück mit Radius $R$ über der Sehne $L$ vom Stabanfang zum Stabende. Sein
Öffnungswinkel ist $\alpha = 2\arcsin\frac{L}{2R}$, und jeder Punkt auf ihm hat einen Winkel
$\varphi \in [0, \alpha]$ ab dem Stabanfang. Für die Rechnung ersetzt das Programm den Bogen durch
$n$ gerade Teilstäbe (Obermenü Seite~3, Standard 20). Das erste und das letzte Teilstück liegen
\emph{tangential} am Bogen und sind so lang wie $R\alpha/n$, die übrigen Teilungspunkte liegen auf
dem Bogen. Die tangentialen Endstücke haben einen Zweck: Die Stabendkräfte des ersten Teilstücks
sind dann direkt Normal- und Querkraft \emph{am Bogen}, und von dort startet der Export.

Die Teilstäbe erben Steifigkeiten, Gelenke an den Bogenenden, Temperatur und die Lasten:

\begin{itemize}
  \item $q$ (radial) und $n$ (tangential) bekommt jeder Teilstab mit den Werten an seinen beiden
        Endpunkten, dazwischen linear. Eine Lastfunktion wertet das Programm an den
        Teilungspunkten aus, mit $x = \varphi$ im Winkelmodus des Rechners.
  \item Lasten je Bogenlänge ($p_X$, $p_Y$ ohne Projektion, $m$) werden je Teilstab mit
        $R\,\Delta\varphi / l$ skaliert ($l$ = Länge des Teilstabs). Damit trägt jeder Teilstab genau
        die Last seines Bogenstücks, und die Summe ist exakt, z.\,B.\ das Eigengewicht $gR\alpha$.
  \item Projizierte Lasten gehen je Grundriss- bzw.\ Höhenlänge des Teilstabs ein.
\end{itemize}

\textbf{Warum konstante radiale Lasten exakt sind.} Ein konstanter Druck $q$ senkrecht auf einen
beliebigen Weg von $A$ nach $B$ hat die Resultierende $q$ mal den um $90^\circ$ gedrehten Vektor
$B - A$. Sein Moment um einen Punkt $P$ ist $\frac q2\left(|B-P|^2 - |A-P|^2\right)$. Beides hängt nur
von den Endpunkten ab, nicht vom Weg dazwischen. Der Linienzug der Teilstäbe hat dieselben Endpunkte
wie der Bogen, also stimmen bei einem statisch bestimmten System die Lagerkräfte und die
Schnittgrößen an den Teilungspunkten auf dem Bogen exakt: Sie folgen dort allein aus dem
Gleichgewicht, und das sieht nur die Resultierenden. Wie biegsam die Stäbe sind, spielt dabei keine
Rolle. Bei einem statisch unbestimmten System hängt die Kraftverteilung dagegen auch von der
Steifigkeit ab, und der Linienzug ist etwas anders steif als der Bogen, dann gilt das nur bis auf
den Teilungsfehler (siehe \emph{Genauigkeit} unten).
Beim Eigengewicht stimmt dagegen nur die Summe exakt. Die Hebel weichen um Größenordnung
$(\alpha/n)^2$ ab, mit 20 Teilstäben bei einem Viertelkreis um etwa 0,04\,\%.

\subsection{Schnittgrößen über den Winkel}

Der Export schreibt die Schnittgrößen als Funktionen des Winkels $p$ ab Stabanfang. Dafür schneidet
er den Bogen bei $p$ und schreibt das Gleichgewicht des Stücks von $0$ bis $p$ auf. Am Anfang wirken
$N_A$, $Q_A$, $M_A$ aus der Rechnung (erstes, tangentiales Teilstück). Mit $\varepsilon = +1$, wenn
die gestrichelte Faser innen liegt ($R > 0$), und $\varepsilon = -1$, wenn sie außen liegt, gilt
ohne Last auf dem Bogen
\begin{align*}
  N(p) &= N_A\cos p + \varepsilon\,Q_A\sin p , \\
  Q(p) &= Q_A\cos p - \varepsilon\,N_A\sin p , \\
  M(p) &= M_A + Q_A\,R\sin p - \varepsilon\,N_A\,R\,(1-\cos p) .
\end{align*}
Die ersten beiden Zeilen drehen nur die Kraft vom Anfang in die Tangente bei $p$. In der dritten
stehen die Hebel: Der Schnittpunkt liegt um $R\sin p$ entlang der Anfangstangente und um
$R(1-\cos p)$ zum Mittelpunkt hin versetzt.
Hier sieht man auch, warum im Bogen Normalkraft und Biegung gekoppelt sind: $N_A$ hat am
gekrümmten Stab den Hebel $R(1-\cos p)$ und erzeugt ein Moment. Beim geraden Stab ($R \to \infty$
mit $R\,p = x$) geht dieser Hebel gegen $x^2/(2R) \to 0$, und $Q_A R\sin p$ wird zu $Q_A\,x$. Das
ist reines Gleichgewicht am unverformten Stab, also Theorie I.~Ordnung. Theorie II.~Ordnung hieße,
das Gleichgewicht am \emph{verformten} System aufzustellen (Knicken, Durchschlagen flacher Bögen),
das rechnet das Programm weder bei Bögen noch bei geraden Stäben.

Last auf dem Bogenstück kommt als Integral über den Lastwinkel $t$ von $0$ bis $p$ dazu (im Bogenmaß,
$ds = R\,dt$):
\begin{align*}
  \Delta N(p) &= -R\int_0^p \Bigl(\varepsilon\,q(t)\sin(p-t) + n(t)\cos(p-t)\Bigr)\,dt , \\
  \Delta Q(p) &= -R\int_0^p \Bigl(q(t)\cos(p-t) - \varepsilon\,n(t)\sin(p-t)\Bigr)\,dt , \\
  \Delta M(p) &= -R\int_0^p \Bigl(R\,q(t)\sin(p-t) - \varepsilon\,R\,n(t)\bigl(1-\cos(p-t)\bigr)
                + m(t)\Bigr)\,dt .
\end{align*}
Auch hier stehen die Hebel direkt da. Eine radiale Kraft beim Winkel $t$ geht durch den
Mittelpunkt, ihr Abstand zum Schnittpunkt bei $p$ ist $R\sin(p-t)$. Eine tangentiale Kraft liegt
auf der Tangente bei $t$, und der Schnittpunkt hat von dieser Tangente den Abstand $R(1-\cos(p-t))$.
Globale Linienlasten $p_X$, $p_Y$ gehen genauso ein, nur mit Hebeln aus den Koordinaten. Ihre
Richtung relativ zur Tangente ändert sich mit $p$.

\textbf{Winkelmodus.} Steht der Rechner auf DEG, laufen $p$ und $t$ in Grad. Dann ist
$ds = R\,\frac{\pi}{180}\,dt$, und genau diesen Faktor setzt das Programm vor die Integrale. Den
Modus erkennt es daran, ob $\cos(180) = -1$ ist. $\sin p$ und $\cos p$ rechnet der Rechner
ohnehin im eingestellten Modus, deshalb stimmen die Funktionen in jedem Modus.

\textbf{Genauigkeit.} Die Formeln sind das exakte Gleichgewicht am Kreisbogen. Ungenau sein können
nur die Startwerte $N_A$, $Q_A$, $M_A$, denn sie stammen aus der Rechnung mit Teilstäben. Bei einem
statisch bestimmten System mit Lasten, die die Teilstäbe exakt abbilden (Knotenlasten, konstante
radiale Last, projizierte Lasten), sind auch sie exakt. Sonst gehen die Abweichungen mit mehr
Teilstäben zurück: bei Eigengewicht, tangentialen Lasten und Lastfunktionen, und bei statisch
unbestimmten Systemen, in denen die Steifigkeit des Linienzugs etwas von der des Bogens abweicht.

%---------------------------------------------------------------------
\section{Starre Stäbe: Zwangsbedingungen und Nullraum}\label{sec:mathe-starr}

Warum nicht einfach $EA = 10^{15}$? Weil der Rechner mit etwa 16 gültigen Stellen rechnet. Steht
in derselben Matrix eine Feder mit $c = 1$ neben $10^{15}$, bleibt von der Feder nach der
Zerlegung praktisch nichts übrig. Das ist der Zahlenmüll aus Abschnitt~\ref{sec:starr}. Das
Programm macht es deshalb exakt.

\textbf{1. Zwangsbedingungen.} Ein starrer Anteil bekommt in der Steifigkeitsmatrix die Steifigkeit
\emph{null}. Stattdessen schreibt das Programm für jeden starren Stab auf, was \enquote{starr}
bedeutet, und zwar in den lokalen Endverschiebungen:
\begin{align*}
  \text{dehnsteif:}\quad & u_b - u_a = \varepsilon_T L, \\
  \text{biegesteif:}\quad & \varphi_a - \varphi_b = -\kappa_T L, \qquad
                            w_b - w_a + L\,\varphi_a = -\kappa_T \frac{L^2}{2} .
\end{align*}
Ohne Temperatur steht rechts null: gleiche Länge, gleiche Endverdrehung, Enden auf einer Geraden.
Alle Bedingungen zusammen, ins globale System gedreht, ergeben $C\,u = g$.

\textbf{2. Kern der Matrix (=Nullraum).} $C$ wird mit Gauß-Jordan auf reduzierte Stufenform gebracht. Jede Zeile hat
dann einen \emph{abhängigen} Freiheitsgrad (die Pivotspalte), der durch die übrigen, \emph{freien}
Freiheitsgrade ausgedrückt wird. Damit lässt sich jede erlaubte Verschiebung schreiben als
\[
  u = T\,u_r + u_0 , \qquad C\,T = 0 .
\]
Die Spalten von $T$ spannen den Nullraum bzw. den Kern von $C$ auf, also genau die Bewegungen, die die starren
Stäbe mitmachen. $u_r$ sind nur noch die freien Freiheitsgrade, $u_0$ ist der Temperaturanteil.

\textbf{3. Projizieren.} Eingesetzt in $K\,u = F$ und von links mit $T^T$ multipliziert:
\begin{equation} \label{eq:LGSred}
  \left(T^T K\,T\right) u_r = T^T\left(F - K\,u_0\right).
\end{equation}
Das ist nichts anderes als das Prinzip der virtuellen Verschiebungen: Als virtuelle
Verschiebungen sind nur noch solche zugelassen, die die Zwangsbedingungen erfüllen, also Spalten
von $T$. Das neue System ist kleiner als das alte und genauso gut konditioniert wie ein System
ohne starre Stäbe, weil nirgends eine riesige Zahl steht.

\textbf{4. Kräfte in den starren Stäben.} Nach Schritt~3 kennt man die Verschiebungen $u$, aber
noch keine Kraft in einem starren Stab, denn ihr Anteil in $K$ ist null. Setzt man $u$ in das
Gleichgewicht der Knoten ein, bleibt ein Rest
\[
  r = F - K\,u .
\]
Das sind die Knotenkräfte, die weder die elastischen Stäbe noch die Federn aufnehmen. Diesen Rest
müssen die starren Stäbe tragen.

Das widerspricht nicht $K\,u = F$. In $K$ fehlen die starren Stäbe, denn ihre Kraft lässt sich
nicht aus $u$ ausrechnen: Ein starrer Stab verformt sich bei jeder Kraft um null. Das vollständige
Gleichgewicht der Knoten lautet
\[
  K\,u + (\text{Kräfte der starren Stäbe}) = F ,
\]
und $r$ ist genau der fehlende Anteil. Das kennt man von den Lagern: Auch an einem gehaltenen
Freiheitsgrad bleibt ein Rest, und der ist die Auflagerkraft. Ein starrer Stab ist so etwas wie ein
inneres Lager, das statt einer absoluten eine relative Verschiebung festhält.

\textbf{Woher die $\lambda$ kommen.} Schritt~3 hat nicht $K\,u = F$ gelöst, sondern nur
\eqref{eq:LGSred}. Alles auf eine Seite gebracht und $u = u_0 + T\,u_r$ eingesetzt:
\begin{align*}
  0 &= T^T F - T^T K\,u_0 - T^T K\,T\,u_r \\
    &= T^T\bigl(F - K\,(u_0 + T\,u_r)\bigr) = T^T\,(F - K\,u) = T^T r .
\end{align*}
Die Zeilen von $T^T$ sind die Spalten von $T$, also eine Basis des Kerns von $C$. $T^T r = 0$ sagt
deshalb: $r$ steht auf jedem dieser Basisvektoren und damit auf dem ganzen Kern von $C$ senkrecht.

Daraus folgt, dass sich $r$ aus den Zeilen von $C$ zusammensetzen lässt. Das ist reine lineare
Algebra aus \textbf{HM1}. $C$ hat $m$ Zeilen (eine je Zwangsbedingung) und $n$ Spalten (eine je
Freiheitsgrad). Im $\mathbb R^n$ liegen damit zwei Unterräume:
\begin{itemize}
  \item $\operatorname{Kern} C$: alle $x$ mit $C\,x = 0$,
  \item $\operatorname{Bild} C^T$: alle Vektoren $C^T\lambda$, also alle Kombinationen der Zeilen
        von $C$.
\end{itemize}
Gezeigt wird $r \in \operatorname{Bild} C^T$, und zwar in vier Schritten.

\emph{(a) $r$ steht senkrecht auf $\operatorname{Kern} C$.} Das haben wir gerade gesehen.

\emph{(b) $\operatorname{Bild} C^T$ steht senkrecht auf $\operatorname{Kern} C$.} Für ein $x$ aus
$\operatorname{Kern} C$ und einen beliebigen Vektor $C^T\lambda$ gilt
\[
  (C^T\lambda)^T x = \lambda^T C\,x = \lambda^T 0 = 0 .
\]
Das erste Gleichheitszeichen ist nur die Rechenregel $(AB)^T = B^T A^T$. Insbesondere liegt außer
der Null kein Vektor in beiden Räumen, denn er müsste auf sich selbst senkrecht stehen.

\emph{(c) Zusammen füllen beide Räume den ganzen $\mathbb R^n$.} Nach der Dimensionsformel ist
$\dim\operatorname{Kern} C + \dim\operatorname{Bild} C = n$, und weil Zeilenrang gleich Spaltenrang
ist, gilt $\dim\operatorname{Bild} C^T = \dim\operatorname{Bild} C$. Also ist
$\dim\operatorname{Kern} C + \dim\operatorname{Bild} C^T = n$. Weil sich die beiden Räume nur in der
Null treffen, lässt sich jeder Vektor des $\mathbb R^n$ eindeutig als Summe aus einem Teil in
$\operatorname{Kern} C$ und einem Teil in $\operatorname{Bild} C^T$ schreiben.

\emph{(d) $r$ liegt in $\operatorname{Bild} C^T$.} Man zerlegt $r = r_K + r_B$ mit $r_K$ in
$\operatorname{Kern} C$ und $r_B$ in $\operatorname{Bild} C^T$. Nach (a) und (b) stehen $r$ und
$r_B$ senkrecht auf $\operatorname{Kern} C$, also auch $r_K = r - r_B$. $r_K$ liegt aber selbst in
$\operatorname{Kern} C$, steht also senkrecht auf sich selbst, und das kann nur der Nullvektor.
Damit ist $r = r_B$.

Ein Beispiel mit Zahlen: $30$ Freiheitsgrade und ein starrer Stab, der dehn- und biegesteif ist
($3$ unabhängige Zeilen). $\operatorname{Kern} C$ ist $27$-dimensional, das sind die $27$
Unbekannten $u_r$ in Schritt~3. $\operatorname{Bild} C^T$ ist $3$-dimensional, das sind die drei
$\lambda$.

Es gibt also Zahlen $\lambda_k$, eine je Zwangsbedingung, mit
\[
  r = C^T \lambda = \sum_k \lambda_k\,C_k^T ,
\]
wobei $C_k$ die $k$-te Zeile von $C$ ist. Weil $r$ nach dem vollständigen Gleichgewicht genau die
Kräfte der starren Stäbe sind, heißt das: Die starren Stäbe wirken mit $C^T\lambda$ auf die Knoten,
und jede Zwangsbedingung~$k$ trägt $\lambda_k\,C_k^T$ dazu bei. Das vollständige Gleichgewicht
lautet damit
\[
  K\,u + C^T\lambda = F .
\]
Mit $T^T$ multipliziert fällt $\lambda$ wegen $T^T C^T = (C\,T)^T = 0$ heraus, das ist Schritt~3.
Die übrigen Gleichungen erfüllt Schritt~4 mit der passenden Wahl von $\lambda$. Die $\lambda$ sind
übrigens die Lagrange-Multiplikatoren, die man bekommt, wenn man das Minimum der potentiellen
Energie $\frac12\,u^T K\,u - F^T u$ unter der Nebenbedingung $C\,u = g$ sucht.

\textbf{Mechanisch gelesen.} Das Skalarprodukt aus Kraft und Verschiebung ist die Arbeit,
\enquote{senkrecht} heißt also \enquote{keine Arbeit}. (a) sagt dann: Die Kräfte der starren Stäbe
leisten auf keiner Bewegung Arbeit, die die starren Stäbe erlauben. Genau das zeichnet eine
Zwangskraft aus. Eine reibungsfreie Führung leistet auf der Bewegung, die sie zulässt, auch keine
Arbeit, sie verhindert nur alle anderen. Physikalisch steckt diese Annahme schon in Schritt~3: Dort
sind als virtuelle Verschiebungen nur solche zugelassen, die die Zwangsbedingungen erfüllen. Alles
danach ist reine Mathematik.

Eindeutig sind die $\lambda$ nur, wenn die Zeilen von $C$ linear unabhängig sind, also
$\operatorname{Kern} C^T = \{0\}$ gilt. Dann gehört zu jeder Kraft aus $\operatorname{Bild} C^T$
genau ein $\lambda$. Ist $\operatorname{Kern} C^T$ größer, gibt es Kräfte in den starren Stäben,
die auf keinen Knoten wirken. Das sind die Eigenspannungszustände weiter unten.

\textbf{Warum die $\lambda$ Schnittgrößen sind.} Das folgt aus dem Gleichgewicht des einzelnen
starren Stabes, und zwar mit demselben Werkzeug wie eben: Kern, Bild und Dimensionsformel, diesmal
für einen Stab allein. Im lokalen System hat ein Stab die sechs Endverschiebungen
$v = (u_a, w_a, \varphi_a, u_b, w_b, \varphi_b)$ und sechs Endkräfte $s$ in denselben Richtungen.
Die drei Zeilen eines biege- und dehnsteifen Stabes bilden die $3 \times 6$-Matrix
\[
  C_\ell = \begin{pmatrix} -1 & 0 & 0 & 1 & 0 & 0 \\ 0 & 0 & 1 & 0 & 0 & -1 \\ 0 & -1 & L & 0 & 1 & 0 \end{pmatrix}
  \qquad
  \begin{matrix} \text{Längszeile: } u_b - u_a \\ \text{Drehzeile: } \varphi_a - \varphi_b \\
                 \text{Querzeile: } w_b - w_a + L\,\varphi_a \end{matrix}
\]

\emph{(a) Starr heißt: nur Starrkörperbewegungen.} Ein Stab kann sich auf drei Arten bewegen, ohne
sich zu verformen: längs verschieben, quer verschieben und um das Ende $a$ drehen. In
Endverschiebungen sind das die Spalten von
\[
  R = \begin{pmatrix} 1 & 0 & 0 & 1 & 0 & 0 \\ 0 & 1 & 0 & 0 & 1 & 0 \\ 0 & 0 & 1 & 0 & -L & 1 \end{pmatrix}^T
\]
(bei der Drehung um $\varphi = 1$ wandert das Ende $b$ wegen $\varphi = -w'$ um $w_b = -L$).
Einsetzen zeigt $C_\ell R = 0$: Starrkörperbewegungen verformen nichts, $\operatorname{Bild} R
\subseteq \operatorname{Kern} C_\ell$. Beide Matrizen haben Rang~$3$, nach der Dimensionsformel ist
also $\dim \operatorname{Kern} C_\ell = 6 - 3 = 3 = \dim \operatorname{Bild} R$, und damit
$\operatorname{Kern} C_\ell = \operatorname{Bild} R$. Die Zwangsbedingungen lassen genau die
Starrkörperbewegungen zu, nicht mehr und nicht weniger.

\emph{(b) Gleichgewicht heißt: keine Arbeit auf Starrkörperbewegungen.} Zuerst ein Stab ohne
Stablasten, die kommen in (d) dazu. Endkräfte $s$ sind genau dann im Gleichgewicht, wenn
\[
  R^T s = \begin{pmatrix} s_1 + s_4 \\ s_2 + s_5 \\ s_3 - L\,s_5 + s_6 \end{pmatrix} = 0 ,
\]
das sind Summe der Kräfte längs, Summe der Kräfte quer und Summe der Momente um $a$. Anders gelesen
ist das das PvV für einen starren Körper aus \textbf{TM1}: Auf jeder Starrkörperbewegung leisten die
Kräfte keine Arbeit. Die Endkräfte, die für sich im Gleichgewicht sind, bilden also den Raum
$\operatorname{Kern} R^T$.

\emph{(c) Dieser Raum wird genau von den Zeilen aufgespannt.} Wie beim Nullraum oben in zwei
Schritten: $R^T C_\ell^T = (C_\ell R)^T = 0$, also $\operatorname{Bild} C_\ell^T
\subseteq \operatorname{Kern} R^T$. Dimension: $\dim \operatorname{Bild} C_\ell^T = \operatorname{Rang}
C_\ell = 3$ und $\dim \operatorname{Kern} R^T = 6 - \operatorname{Rang} R = 3$. Also
\[
  \operatorname{Kern} R^T = \operatorname{Bild} C_\ell^T .
\]
Jede Gruppe von Endkräften im Gleichgewicht ist eine Kombination der drei Zeilen, und weil die
Zeilen linear unabhängig sind, auf genau eine Weise.

\emph{(d) Anwendung auf den starren Stab.} Auf den starren Stab wirken zwei Dinge: die Lasten auf
dem Stab ($q$, $n$, $m$) und die Endkräfte $s$, mit denen die Knoten auf seine Enden wirken. $s$
sind also nur die Endkräfte, die Lasten stehen getrennt. Zusammen sind beide im Gleichgewicht.
Die Festeinspannkräfte $p$ sind ebenfalls Endkräfte, die dieselben Lasten im Gleichgewicht
halten. Dass sie zu einem beidseitig eingespannten Stab gehören, spielt hier keine Rolle. Man
braucht nur \emph{irgendeinen} Satz Endkräfte, der die Lasten allein ausgleicht, und $p$ hat das
Programm aus dem WGV ohnehin. Zieht man die beiden Gleichgewichtsgruppen
\enquote{Lasten $+ s$} und \enquote{Lasten $+ p$} voneinander ab, fallen die Lasten heraus: $s - p$
ist eine Gruppe reiner Endkräfte im Gleichgewicht, also genau der Fall aus (b) ohne Stablasten.
Ausgeschrieben: Nach (b) gilt $R^T (s - p) = 0$, also liegt
$s - p$ in $\operatorname{Kern} R^T$, und das ist nach (c) $\operatorname{Bild} C_\ell^T$. Es gibt
also genau ein $\lambda_\ell \in \mathbb R^3$ mit $s - p = C_\ell^T \lambda_\ell$, umgestellt
\[
  s = p + C_\ell^T \lambda_\ell .
\]
Ins globale System gedreht und über alle starren Stäbe summiert sind das genau die Knotenkräfte
$C^T\lambda$ von oben, die $\lambda_\ell$ eines Stabes sind seine Einträge in $\lambda$. Ist ein Stab
nur dehnsteif, fehlen Dreh- und Querzeile, die Biegung liefert dann wie bei jedem Stab $k\,v$.
Ein N-Gelenk streicht die Längszeile, ein Q-Gelenk die Querzeile. Das hat zwei Gründe, die
zusammenpassen. Kinematisch kann sich das Stabende in der freigegebenen Richtung gegenüber dem
Knoten verschieben. Die Zeile würde aber den Knoten selbst festhalten, also mehr verlangen, als der
Stab tut. Und weil N- und Q-Gelenke herauskondensiert werden, hat das Stabende in dieser Richtung
keinen eigenen Freiheitsgrad, mit dem man die Zeile richtig schreiben könnte. Statisch ist die
zugehörige Schnittgröße am Gelenk null. Die Lasten trägt schon $p$, das für Stäbe mit N- oder
Q-Gelenk passend umverteilt ist, ein $\lambda$ für diese Zeile wäre also ohnehin null. Ein
M-Gelenk streicht dagegen nichts: Das Stabende hat dort einen eigenen Drehfreiheitsgrad
(Abschnitt~\ref{sec:mathe-gelenke}). Die Drehzeile verbindet die Verdrehung des Stabes mit diesem
Freiheitsgrad statt mit dem Knoten, und $M = 0$ am Gelenk folgt aus dessen Gleichung. $p$ darf man dabei unverändert nehmen, weil die
Festeinspannkräfte nach Abschnitt~\ref{sec:mathe-festeinspann} gar nicht von $EI$ abhängen.

\emph{(e) Was die drei Zahlen bedeuten.} Für jede Zeile liest man den Schnittgrößenverlauf des
unbelasteten Stabes direkt ab, mit $N(0) = -s_1$, $Q(0) = -s_2$, $M(0) = -s_3$ wie bei jedem Stab:
\begin{align*}
  \text{Längszeile:}\quad & s = \lambda_N\,(-1, 0, 0, 1, 0, 0)^T: && N(x) = \lambda_N, \\
  \text{Drehzeile:}\quad  & s = \lambda_M\,(0, 0, 1, 0, 0, -1)^T: && M(x) = -\lambda_M, \\
  \text{Querzeile:}\quad  & s = \lambda_Q\,(0, -1, L, 0, 1, 0)^T: && Q(x) = \lambda_Q, \quad
                                                                     M(x) = -\lambda_Q\,(L - x) .
\end{align*}
$\lambda_N$ ist also die Normalkraft, $\lambda_Q$ die Querkraft und $-\lambda_M$ das Moment am Ende
$b$. Zusammen $M(x) = -\lambda_M - \lambda_Q\,(L - x)$, dazu kommt der Anteil aus $p$. Mechanisch
passt das: Die Arbeit von $\lambda_k\,C_k^T$ auf den Endverschiebungen ist $\lambda_k\,(C_k\,v)$,
die Normalkraft arbeitet also auf der Längenänderung, das Moment auf der gegenseitigen Verdrehung
und die Querkraft darauf, wie weit das Ende $b$ von der Tangente am Ende $a$ abweicht. Das ist mit
dem Satz in Abschnitt~\ref{sec:starr} gemeint, die Schnittgrößen starrer Stäbe würden \enquote{aus
dem Gleichgewicht bestimmt}: Sie sind die $\lambda$, und die kommen allein aus dem Rest $r$, also aus
dem Gleichgewicht der Knoten, nicht aus einer Steifigkeit.

\textbf{Wie das Programm sie ausrechnet.} $r = C^T\lambda$ hat eine Gleichung je Freiheitsgrad, aber
nur eine Unbekannte je Zwangsbedingung, also meist mehr Gleichungen als Unbekannte. Weil $r$ nach
Schritt~3 eine Kombination der Zeilen von $C$ ist, geht es trotzdem exakt auf. Von links mit $C$
multipliziert wird daraus ein kleines quadratisches System
\[
  (C\,C^T)\,\lambda = C\,r ,
\]
das eindeutig lösbar ist, solange die Zeilen von $C$ linear unabhängig sind.

\textbf{Starre Teile, die sich gegenseitig halten.} Beispiele sind zwei dehnstarre Stäbe
nebeneinander zwischen denselben zwei Knoten oder ein biegestarrer Zweifeldträger (zwei starre
Stäbe über drei Lager). Dann sind einige Zeilen von $C$ linear abhängig. Das merkt schon die
Elimination in Schritt~2: Eine solche Zeile wird null, nachdem die früheren Pivotspalten abgezogen
sind. Das Programm protokolliert dabei, welche Kombination der ursprünglichen Zeilen null ergeben
hat:
\[
  \sum_k \mu_k\,C_k = 0 \qquad\Longleftrightarrow\qquad C^T\mu = 0 .
\]
Physikalisch ist $\mu$ ein Satz von Kräften in den starren Stäben, der auf keinen Knoten eine Kraft
ausübt, ein \emph{Eigenspannungszustand}. Bei den zwei Stäben nebeneinander ist das Zug~$1$ im
einen und Druck~$1$ im anderen.

Für die Kinematik ist eine abhängige Zeile harmlos, ohne Temperatur ist sie automatisch erfüllt. Für
das Gleichgewicht heißt sie: $r = C^T\lambda$ hat unendlich viele Lösungen
\[
  \lambda = \lambda_p + \sum_i t_i\,\mu_i .
\]
Das ist die Aufteilung, die man von jedem LGS kennt: eine partikuläre Lösung plus beliebige
Lösungen der homogenen Gleichung $C^T\mu = 0$. Und es ist genau die Aufteilung aus dem KGV in
\textbf{Statik}:
\begin{itemize}
  \item $\lambda_p$ rechnet das Programm mit $(C\,C^T)\,\lambda = C\,r$ nur über die unabhängigen
        Zeilen, die abhängigen bekommen $\lambda = 0$. Das ist, als würde man die überzähligen
        Zwangsbedingungen aufschneiden: ein statisch bestimmtes Hauptsystem mit seinem
        $0$-Zustand.
  \item Die $\mu_i$ sind die $1$-Zustände. Jedes $\mu_i$ hat an seiner abhängigen Zeile eine $1$,
        dazu kommen die Kräfte in den übrigen starren Stäben, die diese $1$ im Gleichgewicht halten.
  \item Die $t_i$ sind die statisch Überzähligen $X_i$.
\end{itemize}
Wie im KGV legt das Gleichgewicht die $X_i$ nicht fest. Dort kommen sie aus der Verträglichkeit,
und die braucht Nachgiebigkeiten. Genau die haben starre Stäbe nicht: Jede Wahl der $t_i$ ergibt
dieselben Verschiebungen, und die Starrkörpermechanik kann nicht sagen, welche stimmt. Das hängt
davon ab, wie starr die Stäbe \emph{im Verhältnis zueinander} wirklich sind.

\textbf{Der Grenzwert und die Ergänzungsenergie.} Das Programm trifft dieselbe Annahme wie der
Grenzwert im Rand-LGS: Alle starren Stäbe bekommen $EA = EI = 1/\varepsilon$ mit
$\varepsilon \to 0$. Dann verformen sie sich ein kleines bisschen. In den Größen, die die Zeilen
von $C$ messen, ist ihre elastische Verformung
\[
  C\,u - g = \varepsilon\,W\,\lambda ,
\]
wobei $W$ die Nachgiebigkeit der starren Stäbe mit $EA = EI = 1$ ist (blockweise je Stab, siehe
unten). Das ist das Elastizitätsgesetz der fast starren Stäbe.

Für das Folgende muss man zwei Dinge auseinanderhalten: $u$ sind die Knotenverschiebungen, $C\,u$
sind die Verformungen der starren Stäbe, gemessen in den Größen der Zwangsbedingungen
(Längenänderung, gegenseitige Verdrehung, Abweichung von der Tangente, siehe (e)). Die Stabkräfte
$\lambda$ arbeiten auf $C\,u$, ihre Knotenkräfte $C^T\lambda$ auf $u$, und beides ist dieselbe
Arbeit: $\lambda^T (C\,u) = (C^T\lambda)^T u$, das ist wieder nur $(AB)^T = B^T A^T$.

Jetzt kommt die Verträglichkeit, mit dem Arbeitssatz wie im KGV. Ein Eigenspannungszustand $\mu_i$
übt auf keinen Knoten eine Kraft aus, $C^T\mu_i = 0$. Seine Arbeit auf den wirklichen Verformungen
der starren Stäbe ist deshalb null:
\[
  \mu_i^T (C\,u) = (C^T\mu_i)^T u = 0^T u = 0 .
\]
Multipliziert man das Elastizitätsgesetz $C\,u - g = \varepsilon\,W\,\lambda$ von links mit
$\mu_i^T$, steht links genau diese Arbeit, und es bleibt
\[
  \underbrace{\mu_i^T C\,u}_{=\,0} - \mu_i^T g = \varepsilon\,\mu_i^T W\,\lambda
  \qquad\Longrightarrow\qquad
  \varepsilon\,\mu_i^T W\,\lambda + \mu_i^T g = 0 .
\]
Die Temperatur $g$ bleibt dabei stehen. Ohne Temperatur fällt der zweite Term weg, das $\varepsilon$ kürzt sich heraus, und mit
$\lambda = \lambda_p + \sum_j t_j\,\mu_j$ bleibt
\[
  \sum_j \underbrace{\mu_i^T W\,\mu_j}_{\delta_{ij}}\;t_j
  = -\underbrace{\mu_i^T W\,\lambda_p}_{\delta_{i0}} .
\]
Das sind die Elastizitätsgleichungen des KGV (hier ist $\delta_{ij}$ wieder eine $\delta$-Zahl und
kein Kronecker-Delta). Weil $\varepsilon$ herausfällt, hängt das Ergebnis nicht davon ab,
\emph{wie} starr die Stäbe sind, nur davon, dass sie alle gleich starr sind. Deshalb existiert der
Grenzwert. Passt die Temperatur dagegen nicht zu den abhängigen Zeilen ($\mu_i^T g \ne 0$), muss
$\varepsilon\,\mu_i^T W \lambda$ einen festen Wert annehmen, und für $\varepsilon \to 0$ wird
$\lambda$ unendlich groß. Das ist die unendliche Zwangskraft, zu der das Programm einen Hinweis
zeigt.

\textbf{Wo die Ergänzungsenergie hineinspielt.} Die Gleichungen oben sind genau die Bedingung dafür,
dass
\[
  \Pi^*(t) = \tfrac12\,\lambda^T W \lambda , \qquad \lambda = \lambda_p + \sum_i t_i\,\mu_i ,
\]
minimal wird, denn $\partial\Pi^*/\partial t_i = \mu_i^T W \lambda$. Das ist der Satz von Menabrea:
Unter allen Kraftzuständen im Gleichgewicht stellt sich der mit der kleinsten Ergänzungsenergie
ein. \enquote{Kleinste Ergänzungsenergie} ist dabei nur eine andere Schreibweise für
\enquote{verträglich}, das KGV in \textbf{Statik} rechnet mit dem Prinzip der virtuellen Kräfte
genau dasselbe. Die Ergänzungsenergie ist also kein zusätzliches Prinzip, sondern die
Verträglichkeit in Energieform. Genau genommen ist die Ergänzungsenergie der fast starren Stäbe
$\varepsilon\,\Pi^*$, denn ihre Nachgiebigkeit ist $\varepsilon\,W$ und nicht $W$. Ein fester
Faktor ändert aber nichts daran, wo das Minimum liegt, deshalb steht $\varepsilon$ oben gar nicht
erst in $\Pi^*$.

\textbf{$W$ ausrechnen} geht wie eine $\delta$-Zahl mit der Koppeltafel, mit $EA = EI = 1$ und für
jeden Stab für sich (verschiedene Stäbe koppeln nicht):
\begin{itemize}
  \item Längszeile: $N = 1$ konstant, $\int N^2\,dx = L$.
  \item Drehzeile: $M = 1$ konstant, $\int M^2\,dx = L$.
  \item Querzeile: $M$ linear von $L$ auf $0$, $\int M^2\,dx = L^3/3$.
  \item Dreh- mal Querzeile: $\int 1 \cdot (L - x)\,dx = L^2/2$. Längs- mal Biegezeile: $0$.
\end{itemize}
Die Querkraftverformung fehlt, wie überall im Programm (Euler-Bernoulli). Dass $EA$ und $EI$ beide
$1$ sind, ist eine Annahme und keine Physik. Sie spielt keine Rolle, solange in einem
Eigenspannungszustand nur dehnstarre oder nur biegestarre Zeilen vorkommen, denn dann kürzt sich der
gemeinsame Faktor. Mischen sich beide, hängt die Aufteilung vom angenommenen Verhältnis $EA : EI$
ab.

So teilen sich zwei dehnstarre Stäbe nebeneinander eine Last je zur Hälfte, und ein starrer
Zweifeldträger unter Gleichlast hat die bekannten Lagerkräfte $\frac38\,ql$, $\frac{10}8\,ql$,
$\frac38\,ql$, genau wie mit dem KGV und dem Stützmoment als $X_1$. Eine gelbe Box weist darauf
hin, dass diese Aufteilung angenommen ist, aber nur, wenn sie das Ergebnis auch ändert. Ist
$\lambda_p$ auf allen Zeilen null, die über $W$ mit einem $\mu_i$ gekoppelt sind, dann ist
$\delta_{i0} = \mu_i^T W \lambda_p = 0$ für \emph{jede} Wahl der Nachgiebigkeiten, also $t_i = 0$, egal wie
starr die Stäbe im Verhältnis zueinander sind. Ein Beispiel ist eine starre Platte, die oben und
unten auf Loslagern steht: Die Längskraft zwischen den beiden Lagern ist unbestimmt, aber ohne Last
in Plattenrichtung immer null.

\textbf{Das Gesamtbild.} Beide Aufteilungen in partikulär und homogen gehören zusammen. Die
Kinematik $C\,u = g$ und das Gleichgewicht $C^T\lambda = r$ benutzen dieselbe Matrix, einmal so und
einmal transponiert:

\begin{center}
\begin{tabular}{@{}p{2.9cm}p{5.9cm}p{5.9cm}@{}}
  \toprule
   & \textbf{Kinematik} $C\,u = g$ & \textbf{Gleichgewicht} $C^T\lambda = r$ \\
  \midrule
  partikulär       & $u_0$ (Temperatur) & $\lambda_p$ (Hauptsystem) \\
  homogen          & $T\,u_r$: Bewegungen, die die starren Stäbe erlauben
                   & $\sum_i t_i\,\mu_i$: Eigenspannungszustände \\
  festgelegt durch & Gleichgewicht (Schritt~3) & Verträglichkeit \\
  als Energie      & kleinste potentielle Energie & kleinste Ergänzungsenergie \\
  wie im           & WGV & KGV \\
  \bottomrule
\end{tabular}
\end{center}

Der Kern von $C$ sind die Bewegungen, die die starren Teile frei lassen. Der Kern von $C^T$ besteht
aus den Eigenspannungszuständen $\mu_i$, seine Dimension ist der Grad der statischen
Unbestimmtheit der starren Teile untereinander, also die Anzahl der $t_i$ (im KGV der $X_i$).

Die rote Box \emph{Starrmodus \& statisch unbestimmt} kommt übrigens schon dann, wenn das System
statisch unbestimmt ist und alle Stäbe $EA, EI \ge 10^7$ haben, also auch mit den
Standardwerten.

%---------------------------------------------------------------------
\section{Das Rand-LGS: die Biegelinie wie von Hand}\label{sec:mathe-randlgs}

Der Export \emph{Neue Biegelinie} rechnet unabhängig vom Weggrößenverfahren genau das, was ihr in
\textbf{TM2} von Hand macht, nur für jedes System und mit allen Bedingungen auf einmal.

\subsection{Ansatz je Stab}

Mit $\varphi = -w'$, $Q' = -q$, $M' = Q - m$ und $N' = -n$ setzt das Programm für jeden Stab an

\[
  EI\,w(x) = w_p(x) + c_1\frac{x^3}{6} + c_2\frac{x^2}{2} + c_3 x + c_4, \qquad
  EA\,u(x) = u_p(x) + C_1 x + C_2 .
\]

Das sind sechs Unbekannte je Stab. Die Partikulärlösungen baut es nicht aus einer Tabelle,
sondern als Integral, das der CAS-Kern exakt auswertet:

\begin{align*}
  w_p(x) &= \int_0^x q(t)\,\frac{(x-t)^3}{6}\,dt + \int_0^x m(t)\,\frac{(x-t)^2}{2}\,dt
            - EI\,\kappa_T\,\frac{x^2}{2}, \\
  u_p(x) &= -\int_0^x n(t)\,(x-t)\,dt + EA\,\varepsilon_T\,x .
\end{align*}

Die Idee dahinter: Man denkt sich die Streckenlast in viele kleine Einzellasten $q(t)\,dt$ an den
Stellen $t$ zerlegt. Bei einer Einzellast $1$ an der Stelle $t$ springt die Querkraft dort um $-1$
(wegen $Q' = -q$), und ab $t$ kommt der Anteil $EI\,w = (x-t)^3/6$ dazu. Das ist $c_1 x^3/6$ mit
$c_1 = 1$, nur um $t$ verschoben, davor ist der Anteil null. Weil alles linear ist, addieren sich
die Beiträge aller kleinen Einzellasten, und aus der Summe wird das Integral. Genauso gibt ein
Einzelmoment $1$ bei $t$ den Anteil $(x-t)^2/2$ in $EI\,w$ und eine Längskraft $1$ den Anteil
$-(x-t)$ in $EA\,u$. In der Mathematik heißt so ein
Integral \emph{Faltung}. Deshalb funktioniert es für jede Lastfunktion, für Trapezlasten und für
die globalen Linienlasten (die vorher auf Quer- und Längsrichtung des Stabes zerlegt werden). Und deshalb heißt die Integrationsvariable \code{t}, und
ihr könnt \code{t} nicht als Symbol verwenden.

Aus dem Ansatz folgen die Größen an beiden Stabenden als \emph{lineare} Ausdrücke in den
Konstanten, am Anfang z.\,B.\ $w = c_4/EI$, $\varphi = -c_3/EI$, $M = -c_2$, $Q = -c_1$,
$u = C_2/EA$, $N = C_1$.

\subsection{Bedingungen an den Knoten}

An jedem Knoten schreibt das Programm die Bedingungen, die ihr auch aufschreiben würdet, nur
systematisch:

\begin{itemize}
  \item \textbf{Kinematik:} Alle Stabenden ohne N- bzw.\ Q-Gelenk haben dieselbe
        Knotenverschiebung $\vec U$, also $u_e = \vec U\cdot\vec t_e$ und
        $w_e = \vec U\cdot\vec n_e$, alle Enden ohne M-Gelenk dieselbe Verdrehung. $\vec U$ selbst
        wird durch das erste passende Stabende ausgedrückt, damit keine zusätzlichen Unbekannten
        entstehen.
  \item \textbf{Gelenke:} Ein M-Gelenk gibt $M_e = 0$, ein N-Gelenk $N_e = 0$, ein Q-Gelenk
        $Q_e = 0$.
  \item \textbf{Lager:} In einer gehaltenen Richtung $\vec U\cdot\vec e = 0$, in einer freien
        Richtung das Gleichgewicht der Stabendkräfte, Federn und Lasten in dieser Richtung. Die
        Richtungen sind $\vec\xi = (\cos\alpha, -\sin\alpha)$ und $\vec\eta = (\sin\alpha, \cos\alpha)$,
        so gehen gedrehte Lager ohne Sonderfall.
  \item \textbf{Drehung:} Einspannung $\varphi = 0$, sonst $\sum M_e - c_m\,\varphi + M = 0$.
\end{itemize}

Am Ende muss die Anzahl der Gleichungen genau $6 \cdot (\text{Stäbe})$ sein (plus zwei je Knoten,
an dem alle Stäbe über N- oder Q-Gelenke hängen, siehe unten). Sonst kommt \emph{nicht quadratisch}.

\textbf{Warum solche Knoten eigene Verschiebungen brauchen.} Normalerweise wird $\vec U$ durch ein
Stabende ohne N- und Q-Gelenk ausgedrückt: Dessen $u_e$ und $w_e$ sind die beiden Komponenten von
$\vec U$ in Stab- und Querrichtung, damit ist $\vec U$ ganz bekannt. Ein Stabende mit N-Gelenk darf
sich aber längs gegen den Knoten verschieben, es liefert nur noch $w_e = \vec U\cdot\vec n_e$. Eines
mit Q-Gelenk liefert nur $u_e = \vec U\cdot\vec t_e$. Hat jedes Stabende am Knoten so ein Gelenk,
gibt es kein Ende, das beide Komponenten festlegt. Dann bekommt $\vec U$ die eigenen Unbekannten
$u_x(K)$, $u_y(K)$, und jedes Ende schreibt seine Bedingung gegen $\vec U$. Normalerweise
\enquote{verbraucht} das Bezugsende seine zwei Komponenten, um $\vec U$ festzulegen, und liefert
selbst keine Gleichung. Hier verbraucht keines etwas, daher die zwei Gleichungen mehr. Beispiel:
zwei Stäbe, die sich an einem Knoten mit je einem Q-Gelenk treffen. Jeder sagt nur, wie weit sich
der Knoten in seiner Längsrichtung bewegt, erst beide zusammen bestimmen $\vec U$.

\subsection{Lösen: erst wie von Hand, dann das CAS}

Ein LGS mit 30 symbolischen Unbekannten direkt in \code{simult} zu stecken, ist auf dem Rechner
langsam und liefert unübersichtliche Ausdrücke. Das Programm macht deshalb vorher, was ihr auch
machen würdet:

\begin{enumerate}
  \item Jede Gleichung mit nur \emph{einer} Unbekannten wird sofort aufgelöst und der Wert überall
        eingesetzt (z.\,B.\ $M = 0$ am Gelenk liefert direkt $c_2$). Das wiederholt sich, solange
        etwas passiert.
  \item Eine Unbekannte, die nur noch in einer einzigen Gleichung vorkommt, wird mit dieser
        Gleichung beiseitegelegt und am Schluss aus ihr berechnet.
  \item Gleichungen der Form $a\,x_j + b\,x_m = 0$ mit Zahlen $a$, $b$ werden zu $x_j = f\,x_m$
        und eingesetzt.
\end{enumerate}

Was übrig bleibt, zerfällt oft in unabhängige Blöcke (Unbekannte, die nie zusammen in einer
Gleichung stehen). Jeder Block wird einzeln mit \code{exact(simult(\dots))} gelöst, danach wird
rückwärts eingesetzt. Aus demselben Protokoll entsteht auch \code{randbed}: Eine Größe steht dort
nur dann als Wert, wenn ihre Konstanten in einem \emph{früheren} Schritt schon bekannt waren,
also genau in der Reihenfolge, in der ihr es auf dem Papier bekommen würdet.

\textbf{Die Aufteilung in \enquote{nur Biegelinie} und \enquote{nur Längslinie}} kommt aus
demselben Protokoll: Das Programm verfolgt für jede Konstante der Biegelinie zurück, aus welchen
Gleichungen sie bestimmt wurde. Tauchen darin keine Längslinien-Konstanten auf, die nicht schon
vorher feststanden, reichen die Bedingungen der Biegelinie allein.

Bei Rahmen mit dehnsteifen Stäben reicht das nicht. Dort bestimmt das LGS die Biegelinie zusammen
mit $C_1$ und $C_2$ der dehnsteifen Stäbe, weil z.\,B.\ die Seitenverschiebung $u$ des Riegels in
die Biegelinie der Stiele eingeht. Bei der direkten Methode ist $u = C_2$ und $N = C_1$ aber an
beiden Stabenden dieselbe Unbekannte (ohne Längslast und Temperatur). Das Programm nimmt deshalb
eine Bedingung, in der nur eine dieser Unbekannten steht, löst sie danach auf und setzt sie in die
anderen Bedingungen ein, genau wie von Hand. Mathematisch ist das ein Gauß-Schritt mit dieser
Unbekannten als Pivot, das LGS bleibt dabei regulär. Bleiben danach genau $4\cdot(\text{Stäbe})$
Bedingungen übrig, in denen nur noch Konstanten der Biegelinie stehen, genügen sie für $w$. Bei
dehnbaren Stäben ist $u(L) = (C_1 L + C_2)/EA$, dort hängt die Längsverschiebung am Stabende von
beiden Konstanten ab, und das Durchreichen geht nicht.

\subsection{Starre Stäbe im Rand-LGS: direkt oder Grenzwert}

Für einen starren Stab gibt es kein $EI\,w'' = -M$ mehr, trotzdem braucht das Programm sechs
Gleichungen. Zwei Wege:

\begin{itemize}
  \item \textbf{direkt} (Standard): Der biegesteife Stab bewegt sich als Starrkörper,
        $w(x) = c_4 + c_3 x - \kappa_T x^2/2$ und beim dehnsteifen $u(x) = C_2 + \varepsilon_T x$.
        $c_1$, $c_2$ und $C_1$ bleiben als reine Kraftgrößen stehen und werden aus dem
        Gleichgewicht bestimmt.

        Warum das geht: Auch beim elastischen Stab sind $c_1$, $c_2$ und $C_1$ eigentlich
        Kraftgrößen, nämlich die Schnittgrößen am Stabanfang, $Q(0) = -c_1$, $M(0) = -c_2$,
        $N = C_1$. In der Biegelinie stehen sie nur deshalb, weil die Krümmung aus dem Moment kommt:
        Sie gehen als $c_1 x^3/(6EI)$ und $c_2 x^2/(2EI)$ in $w$ ein, also immer geteilt durch
        $EI$. $c_3$ und $c_4$ dagegen sind echte Weggrößen ($c_3/EI = w'(0)$, $c_4/EI = w(0)$).
        Wird der Stab starr, $EI \to \infty$, verschwindet bei endlichen Kräften der Beitrag von
        $c_1$ und $c_2$ aus der Biegelinie, und $c_3$, $c_4$ schreibt man gleich ohne $EI$
        ($c_4 = w(0)$, $c_3 = w'(0)$). Die Schnittgrößen bleiben aber da: $Q(x)$, $M(x)$ und $N(x)$
        folgen aus $c_1$, $c_2$, $C_1$ und den Lasten allein über das Gleichgewicht am Stabstück
        ($Q' = -q$, $M' = Q$, $N' = -n$), ganz ohne Steifigkeit.

        Die Bedingungen an den Knoten sind dieselben wie beim elastischen Stab. Jede enthält aber
        entweder nur Kraftgrößen ($c_1$, $c_2$, $C_1$) oder nur Weggrößen ($c_3$, $c_4$, $C_2$) des
        starren Stabes: Kinematik und Lager mit gehaltener Richtung enthalten Wege, Gelenk-,
        Gleichgewichts- und Drehbedingungen Kräfte. Beispiel: starrer Kragarm, eingespannt am
        Anfang, Kraft $F$ (in $+w$) am freien Ende. Die Einspannung gibt $w(0) = c_4 = 0$ und
        $w'(0) = c_3 = 0$, also bewegt sich nichts. Das freie Ende gibt $Q(L) = -c_1 = F$ und
        $M(L) = -c_2 - c_1 L = 0$, also $c_1 = -F$, $c_2 = F L$ und $M(0) = -F L$. Das ist genau die
        Rechnung des statisch bestimmten Kragarms von Hand. In \code{randbed} heißen die
        Kraftgrößen eines starren Stabes deshalb $Q_i$, $M_i$, $N_i$ statt $-EI\,w_i'''$ usw.
  \item \textbf{Grenzwert}: Jeder starre Stab bekommt $EI = 1/\varepsilon$ bzw.\
        $EA = 1/\varepsilon$ (im CAS heißt $\varepsilon$ \code{rleps}), das System wird symbolisch in
        $\varepsilon$ gelöst und danach \code{limit(\dots, rleps, 0)} gebildet.
\end{itemize}

Warum überhaupt der Grenzwert? Das direkte Verfahren wird singulär, wenn sich starre Teile
\emph{gegenseitig} statisch unbestimmt halten. Beispiel: zwei dehnstarre Stäbe hintereinander
zwischen zwei Festlagern, Horizontallast in der Mitte. Beide Stäbe ändern ihre Länge nicht, also
ist $u = 0$ überall. Aber wie sich die Last auf die beiden Stäbe verteilt, sagt euch das
Gleichgewicht nicht, nur $N_1 - N_2 = F$. Mit $EA = 1/\varepsilon$ ist das System regulär, und
der Grenzwert $\varepsilon \to 0$ liefert die Aufteilung, die zwei \enquote{gleich starren} Stäben
entspricht. Genau in diesem Fall wechselt das Programm von selbst auf den Grenzwert. In
\code{randinfo} steht dann \emph{direkt singulär, starre Teile gegenseitig statisch unbestimmt}.

\subsection{Knotenverschiebungen aus den Randwerten}\label{sec:mathe-knotenverschiebung}

Für die Knotenverschiebung (Knotenmenü Seite~4) muss nichts Neues gelöst werden. Die
Kinematik-Bedingungen drücken $\vec U$ ja schon durch ein Referenz-Stabende aus, und zwar durch das erste
Stabende am Knoten ohne N- und Q-Gelenk. Dessen Längs- und Querverschiebung stehen nach dem Lösen
fest, $u(0)$ und $w(0)$ am Stabanfang bzw.\ $u(L)$ und $w(L)$ am Stabende. Zurück ins globale
System:
\[
  U_x = u\,t_x + w\,n_x, \qquad U_y = u\,t_y + w\,n_y, \qquad
  \vec t = \frac{(\Delta x, \Delta y)}{L}, \quad \vec n = \frac{(-\Delta y, \Delta x)}{L} .
\]
Das sind dieselben $\vec t$ und $\vec n$ wie in den Bedingungen, deshalb zählt $U_y$ wie die
$y$-Achse des Programms nach unten. Die Verdrehung kommt vom ersten Stabende ohne M-Gelenk,
\[
  \varphi = -w' = -\frac{w_p'(x) + c_1 x^2/2 + c_2 x + c_3}{EI},
\]
beim biegesteifen Stab $\varphi = \kappa_T x - c_3$. Hängen alle Stäbe über N- oder Q-Gelenke am
Knoten, hat das LGS für ihn ohnehin die eigenen Unbekannten $u_x(K)$, $u_y(K)$, die werden direkt
genommen.

Damit das beim zweiten Knoten nicht wieder ein ganzes LGS kostet, hebt das Programm die Linien des
letzten Rand-LGS mit einem Schlüssel auf: dem Fingerabdruck aller Eingaben (derselbe, an dem die
Hinweisbox erkennt, ob sich etwas geändert hat. Er enthält auch die Knoten- und Stabnummern) und
den Schaltern, die das Rand-LGS sonst beeinflussen (Methode für starre Stäbe, symbolischer Modus,
Zahlenformat). Passt der Schlüssel, werden nur die Randwerte eingesetzt, das sind drei kurze
CAS-Aufrufe statt eines LGS. Gleicher Fingerabdruck heißt wirklich gleiches System, auch wenn ein
Knoten zwischendurch gelöscht und genau so wieder eingebaut wurde. Werden dabei Nummern
vertauscht, ändert sich der Fingerabdruck, und das ist auch richtig so: Dann gehören die Namen
\code{wv1}, \code{wv2}, \dots\ zu anderen Stäben.

%---------------------------------------------------------------------
\section{Symbolischer Modus: rechnen mit Platzhaltern}\label{sec:mathe-symbolisch}

Die Finite-Elemente-Rechnung für die Diagramme ist eine reine Zahlenrechnung. Symbolisch
invertiert der Rechner keine Steifigkeitsmatrix in vernünftiger Zeit. Der Trick:

\begin{enumerate}
  \item Jedes Symbol bekommt einen \textbf{Platzhalterwert}, standardmäßig~1. Die Länge $l$
        bekommt, wenn die Geometrie kein eigenes Symbol hat, den Rasterabstand.
  \item \textbf{Superposition:} Alle Lasten werden auf null gesetzt und dann \emph{einzeln} wieder
        eingeschaltet. Für jede Last rechnet das Programm die Stabendkräfte und teilt sie durch den
        Platzhalter der Last und die passende Potenz von $l$. Übrig bleibt ein reiner
        Zahlenfaktor~$c$:
        \[
          N = c\,F, \quad M = c\,F\,l \quad\text{(Einzellast)}, \qquad
          N = c\,q\,l, \quad M = c\,q\,l^2 \quad\text{(Streckenlast)} .
        \]
        Die Steifigkeitsmatrix wird dafür nur einmal zerlegt, jede Last kostet dann nur noch ein
        Vorwärts- und Rückwärtseinsetzen.
  \item Der Zahlenfaktor wird \textbf{zurück in einen exakten Ausdruck} verwandelt (siehe unten)
        und alle Anteile werden addiert.
\end{enumerate}

Das funktioniert, weil die Rechnung linear ist und alle Längen Vielfache von $l$ sind. Genau hier
kommen die Einschränkungen her: Eine Temperatur passt nicht in das Schema \enquote{Last mal
Länge hoch $k$} (deshalb fehlt sie in den Diagrammwerten; die Exporte \emph{Neue Biegelinie} und
\emph{Neue Längslinie}, also \code{wneu}, \code{uneu}, \code{wv}, \code{uv}, haben sie, weil das
Rand-LGS gar nicht mit Superposition rechnet), und ein Zahlen-$EI$
neben einem symbolischen $EI$ macht die Steifigkeitsverhältnisse vom Platzhalterwert abhängig
(deshalb der gelbe Hinweis). Die Stelle des Maximums in einem Feld sucht das Programm numerisch
und rastet sie auf typische Werte wie $1/\sqrt 3$ oder einfache Brüche ein.

\textbf{Aus einer Zahl wird wieder ein Bruch.} Aus $0{,}708333\ldots$ soll wieder
$\frac{17}{24}$ werden. Dafür entwickelt das Programm die Zahl in einen \emph{Kettenbruch}
$a_0 + \cfrac{1}{a_1 + \cfrac{1}{a_2 + \dots}}$ und nimmt den ersten Näherungsbruch, der bis
auf $10^{-12}$ stimmt (Nenner höchstens 10\,000). Klappt das nicht und steht das Zahlenformat auf
\emph{Wurzel/Bruch}, versucht es dasselbe mit dem \emph{Quadrat} der Zahl und zieht aus Zähler und
Nenner die Quadratfaktoren heraus. So wird aus $0{,}7071\ldots$ wieder $\frac{\sqrt 2}{2}$ und
aus $86{,}60\ldots$ ein $50\sqrt 3$. Scheitert beides, sucht es zuletzt noch einen Bruch mit
kleinem Nenner (höchstens 100), der nur bis auf die eingestellte Toleranz $10^{-x}$ stimmen muss
(Obermenü Seite~3, \emph{CAS Bruch-Tol.}). Das ist die einzige Stelle, an der diese Einstellung
wirkt, und gerechnet wird dabei in Lua, nicht im CAS. Erst wenn auch das nicht klappt, bleibt eine
Dezimalzahl stehen. Im Format \emph{Dezimal} fallen alle drei Versuche weg.

%---------------------------------------------------------------------
\section{Kraftgrößenverfahren: Hauptsystem und \texorpdfstring{$\delta$}{delta}-Zahlen}\label{sec:mathe-kgv}

\textbf{Das Hauptsystem} sucht das Programm \enquote{gierig}: Es löst der Reihe nach
Einspannmomente, Lagerkräfte, Momente an biegesteifen Stabanschlüssen und zuletzt unbelastete
Pendelstäbe. Es behält jede Lösung nur, wenn das System danach noch tragfähig ist (die
Steifigkeitsmatrix lässt sich noch zerlegen). Sobald nichts mehr geht, ist das System statisch
bestimmt, und die Anzahl der gelösten Bindungen ist der Grad $n$.

\textbf{Die $\delta$-Zahlen} integriert das Programm nicht über $\int M_i M_j/(EI)\,dx$ (so rechnet
ihr sie in \textbf{Statik} mit dem Prinzip der virtuellen Kräfte), sondern es
rechnet für jeden Zustand ($X_j = 1$ bzw.\ Lastzustand) die Verformung des Hauptsystems mit dem
Weggrößenverfahren und liest die zu $X_i$ gehörende Klaffung ab: Verschiebung am gelösten Lager,
gegenseitige Verdrehung am eingefügten Gelenk, Längenänderung am gelösten Pendelstab. Nach
Maxwell--Betti ist das dasselbe wie das Arbeitsintegral, inklusive der Anteile aus Normalkraft,
Federn und endlichem $EA$, die man von Hand gern vergisst. Ein gelöster Pendelstab ist im
Hauptsystem ganz entfernt. Seine eigene Längenänderung $L/(EA)$ fehlt deshalb in den
Knotenverschiebungen und wird zu $\delta_{jj}$ addiert. Bei gedrehten Lagern liegen die
Knotenfreiheitsgrade schon in Lagerrichtung ($u_\xi$, $u_\eta$), die Klaffung eines gelösten Lagers
ist genau dieser Freiheitsgrad, und die Einheitskräfte eines gelösten Pendelstabes werden in
diese Richtungen zerlegt.

%---------------------------------------------------------------------
\section{Polplan, Verschiebungszustand und Arbeitssatz}\label{sec:mathe-polplan}

\subsection{Scheiben finden}

Zuerst muss das Programm wissen, welche Stäbe zusammen eine starre Scheibe bilden. Das passiert
in zwei Durchgängen:

\begin{enumerate}
  \item \textbf{Ohne Lager:} An jedem Knoten ohne Vollgelenk fasst das Programm alle Stäbe
        zusammen, deren Enden dort kein Gelenk haben (\emph{Union-Find}: jede Gruppe hat einen
        Vertreter, zwei Gruppen werden vereinigt, indem man einen Vertreter unter den anderen
        hängt). Mit diesen Scheiben und nur den inneren Gelenken läuft der Polplan einmal durch.
        Findet er zwei Scheiben, die gegeneinander unverschieblich sind, z.\,B.\ die drei Stäbe
        eines Dreigelenk-Dreiecks, werden sie zu einer Scheibe zusammengefasst.
  \item \textbf{Mit Lagern:} Stäbe, die biegesteif an einer Einspannung hängen (oder an einem
        Festlager mit Drehfeder), gehören zur Scheibe~0, dem Boden. Mit diesen zusammengefassten
        Scheiben und den Lagern entsteht der Polplan, den ihr seht.
\end{enumerate}

\subsection{Pole und Polstrahlen}

An jedem Knoten schaut das Programm, welche Scheiben sich dort treffen und welche Relativbewegung
zwischen ihnen möglich ist:

\begin{longtable}{p{0.34\linewidth}p{0.60\linewidth}}
  \toprule
  \textbf{Verbindung} & \textbf{Folge im Polplan} \\
  \midrule \endhead
  M-Gelenk, Festlager & Pol im Knoten (Drehung möglich, keine Verschiebung) \\
  Rollenlager & Polstrahl durch den Knoten senkrecht zur Rollrichtung, also entlang der
    Lagerkraft \\
  N- oder Q-Gelenk & Pol im Unendlichen (reine Verschiebung längs bzw.\ quer zum Stab) \\
  Einspannung, biegesteifer Anschluss & kein Pol: die Scheiben sind starr verbunden \\
  \bottomrule
\end{longtable}

Danach wiederholt das Programm (höchstens zehnmal) zwei Schritte, bis nichts Neues mehr dazukommt:

\begin{enumerate}
  \item \textbf{Satz von den drei Polen:} Für je drei Scheiben $i$, $j$, $k$ liegen die Pole
        $(i,j)$, $(j,k)$ und $(i,k)$ auf einer Geraden. Sind zwei davon bekannt, entsteht ein
        Polstrahl für den dritten. Liegt einer der beiden im Unendlichen, geht der Strahl durch
        den anderen in dessen Richtung.
  \item \textbf{Strahlen schneiden:} Zwei Polstrahlen derselben Scheibenpaarung schneiden sich im
        Pol. Sind sie parallel, liegt der Pol im Unendlichen in dieser Richtung.
\end{enumerate}

Findet es für eine Paarung zwei verschiedene Pole, oder liegt ein Pol nicht auf einem Strahl, auf
dem er liegen müsste, sind die beiden Scheiben gegeneinander unverschieblich (im Polplan rot). Im
Durchgang ohne Lager werden solche Paare zu einer Scheibe zusammengefasst (siehe oben).

\subsection{Der Verschiebungszustand}

Eine Scheibe, die sich um ihren Hauptpol $(c_x, c_y)$ um den kleinen Winkel $\varphi$ dreht,
verschiebt jeden Punkt um $\vec r_\perp \varphi$, also senkrecht zum Polstrahl vom Pol zu diesem
Punkt. Das Programm gibt einer Startscheibe $\varphi = 1$ (bevorzugt eine, deren Hauptpol auf einem
Knoten liegt) und reicht die Bewegung über die Nebenpole weiter: Der Nebenpol $(A,B)$ ist der Punkt,
der auf beiden Scheiben dieselbe Verschiebung hat. Kennt man die Verschiebung dieses Punktes aus
Scheibe~$A$ und den Hauptpol von Scheibe~$B$, folgt $\varphi_B$. Hat eine Scheibe ihren Hauptpol im
Unendlichen, verschiebt sie sich nur.

\textbf{Verschobene Scheiben unterwegs.} Die Kette reißt an einer Scheibe, die sich nur verschiebt,
nicht ab. Bei ihr bewegt sich jeder Punkt um denselben Vektor $\vec v$, also auch der Nebenpol zur
nächsten Scheibe~$B$. Damit ist die Verschiebung des Nebenpols bekannt, und es geht weiter wie bei
einer drehenden Scheibe: Liegt der Hauptpol von $B$ im Endlichen, folgt $\varphi_B$ aus $\vec v$,
liegt er im Unendlichen, verschiebt sich auch $B$ nur. Liegt der Nebenpol $(A,B)$ selbst im
Unendlichen, drehen beide Scheiben gleich, $\varphi_B = \varphi_A = 0$: $B$ verschiebt sich dann auch
nur, oder sie steht still, wenn ihr Hauptpol im Endlichen liegt. Eine reine Verschiebung ist nur der
Grenzfall einer Drehung mit dem Pol im Unendlichen.

\textbf{Nebenpol gleich Hauptpol.} Der Hauptpol einer Scheibe ist der Punkt, der sich auf ihr nicht
bewegt. Fällt der Nebenpol $(A,B)$ mit dem Hauptpol von $B$ zusammen, bewegt sich dieser Punkt also
auf $B$ nicht, und weil er auf $A$ dieselbe Verschiebung haben muss, auch auf $A$ nicht. Kommt aus
der Startscheibe $A$ trotzdem eine Verschiebung ungleich null an, passt das nicht zusammen: Der Punkt
kann nicht zugleich stehen und sich bewegen. Dann kann sich $A$ gar nicht bewegen, die Startscheibe
war falsch gewählt, und das Programm nimmt die nächste.

Ein Beispiel ist euer Dreigelenktragwerk mit dem Gelenk $G$ senkrecht über dem linken Lager $A$,
bei dem ihr das linke Festlager zum Loslager macht (verschieblich in $x$). Das rechte Festlager $B$
ist der Hauptpol $(\mathrm{II})$ der rechten Scheibe, das Gelenk der Nebenpol $(\mathrm{I,II})$. Der
Hauptpol $(\mathrm{I})$ der linken Scheibe liegt auf der Senkrechten durch $A$ (Loslager) und nach
dem Satz von den drei Polen auf der Geraden durch $(\mathrm{II})$ und $(\mathrm{I,II})$, also auf
$BG$. Beide Geraden schneiden sich in $G$, damit ist $G$ zugleich $(\mathrm{I})$ und
$(\mathrm{I,II})$. Startet das Programm mit der rechten Scheibe, dreht sie um $B$ und schiebt $G$
seitlich weg, aber auf der linken Scheibe steht $G$ als Hauptpol still: Widerspruch, die rechte
Scheibe kann sich nicht bewegen. Mit der linken Scheibe als Start geht es: Sie dreht um $G$, ihr
Lager $A$ wandert waagerecht (senkrecht zum Polstrahl $GA$, und das lässt das Loslager zu). $G$
selbst bewegt sich dabei nicht. Die rechte Scheibe hat damit zwei Punkte, die stehen bleiben ($B$
und $G$), und bewegt sich deshalb gar nicht, $\varphi_\mathrm{II} = 0$. Das ist kein Fehler,
sondern die Kinematik dieses Systems: Nur die linke Scheibe kippt um das Gelenk.

\subsection{Der Arbeitssatz}\label{sec:mathe-arbeitssatz}

Nach dem Lösen einer Bindung ist das System eine kinematische Kette mit genau einem Freiheitsgrad.
Für jede virtuelle Bewegung, die die übrigen Bindungen zulassen, gilt
\[
  \delta W = \sum_k \vec F_k \cdot \delta\vec u_k + \sum_m M_m\,\delta\varphi_m = 0 ,
\]
und weil es nur einen Freiheitsgrad gibt, sind alle Wege Vielfache eines einzigen Parameters. Der
Export (Abschnitt~\ref{sec:arbeitssatz}) macht daraus drei Schritte.

\textbf{Weg einer Kraft (\code{pvv\_kin}, Spalte~2).} Eine Kraft mit Richtung
$\vec e = (e_x, e_y)$ ($y$ nach unten) greift im Punkt $P$ auf Scheibe~$i$ an. Ihr Weg ist die
Verschiebung von $P$ in Richtung $\vec e$:
\begin{itemize}
  \item Scheibe~$i$ dreht um ihren Hauptpol $(c_x, c_y)$ um $\varphi_i$ (gegen den Uhrzeigersinn).
        Jeder Punkt bewegt sich senkrecht zum Polstrahl vom Pol zu ihm, projiziert auf $\vec e$ ist das
        \[
          \delta_e = \varphi_i\,\bigl[e_x\,(P_y - c_y) - e_y\,(P_x - c_x)\bigr] .
        \]
        Die eckige Klammer ist der Hebel der Kraft um den Pol.
  \item Scheibe~$i$ verschiebt sich nur (Hauptpol im Unendlichen) um $\delta_i$ in Richtung
        $\vec t_i$: $\delta_e = \delta_i\,(\vec e\cdot\vec t_i)$.
  \item Ein Moment auf Scheibe~$i$ arbeitet am Drehwinkel: $\delta\varphi = \varphi_i$.
\end{itemize}
Verschieben \emph{und} drehen zugleich gibt es nicht als eigenen Fall: Eine kleine ebene
Starrkörperbewegung ist immer eine Drehung um einen Punkt, eben den Pol. Eine Verschiebung plus
Drehung ist nur eine Drehung um einen anderen Pol, die reine Verschiebung der Grenzfall mit dem Pol
im Unendlichen.
Streckenlasten dürfen exakt durch ihre Resultierende $R$ ersetzt werden. Auf einem geraden Stab
der Länge $L$, der zu einer starren Scheibe gehört, ändert sich die virtuelle Verschiebung linear
mit der Stabkoordinate $x$, denn ein Starrkörper verschiebt jeden Punkt linear mit seinem Ort:
$\delta\vec u(x) = \delta\vec u_A + \frac{x}{L}\,(\delta\vec u_B - \delta\vec u_A)$. Damit ist die
Arbeit der Streckenlast
\[
  \int_0^L q(x)\,\vec e\cdot\delta\vec u(x)\,dx
  = \vec e\cdot\delta\vec u_A \underbrace{\int_0^L q\,dx}_{R}
    + \vec e\cdot\frac{\delta\vec u_B - \delta\vec u_A}{L} \underbrace{\int_0^L q\,x\,dx}_{S = R\,x_R}
  = R\;\vec e\cdot\delta\vec u(x_R) ,
\]
also die Arbeit der Resultierenden $R$ an ihrem Angriffspunkt $x_R = S/R$. Auf einem Bogen ist die
Verschiebung nicht linear in der Bogenlänge, das Ergebnis gilt aber genauso: Dort wird die Last
durch ihre Resultierende (Kraft und Richtung) auf ihrer Wirkungslinie ersetzt. Die freigesetzte
Größe ist eine äußere Last in der gezeichneten Richtung, bei einem gelösten Gelenk arbeitet sie an
der gegenseitigen Verdrehung bzw.\ Verschiebung der beiden Seiten.

\textbf{Verhältnis der Scheiben (\code{pvv\_pol}).} Der Nebenpol $P_{AB}$ ist der Punkt, der auf
Scheibe~$A$ und Scheibe~$B$ dieselbe Verschiebung hat. Schreibt man diese Verschiebung mit der
Formel oben einmal über $A$ und einmal über $B$ hin, entsteht eine lineare Beziehung
$\varphi_B = \kappa_{AB}\,\varphi_A$. Die Verschiebung ist ein Vektor mit $x$- und $y$-Komponente,
das sind eigentlich zwei Gleichungen. Beide liefern aber dasselbe $\kappa_{AB}$: Über $A$ steht die
Verschiebung senkrecht auf dem Polstrahl von $(A)$ zum Nebenpol, über $B$ senkrecht auf dem von
$(B)$, und weil $(A)$, $(A,B)$ und $(B)$ auf einer Geraden liegen, sind beide Richtungen gleich.
Deshalb reicht eine Komponente. Das Programm nimmt die betragsmäßig größere, damit es nicht durch
eine Zahl nahe null teilt. Liegt der Nebenpol zweier drehender Scheiben im Unendlichen, drehen beide
gleich, $\varphi_A = \varphi_B$.

Von der Bezugsscheibe aus (der Startscheibe des Verschiebungszustands, ihr Winkel heißt
$\varphi_0$) laufen diese Beziehungen über die Nebenpole durch alle Scheiben und liefern
$\varphi_i = \kappa_i\,\varphi_0$. Das sind die unteren Zeilen von \code{pvv\_kin}. Gibt es mehr
Nebenpole als nötig, sind die übrigen Proben. Als Bezug nimmt das Programm eine Scheibe, die sich
dreht, wenn es eine gibt. Eine Scheibe, die stehen bleibt, bekommt einfach $\kappa_i = 0$, ihre
Lasten leisten keine Arbeit. Verschieben sich alle bewegten Scheiben nur (alle ihre Hauptpole im
Unendlichen), gibt es keinen Drehwinkel als Bezug. Dann ist die Bezugsgröße die Verschiebung $\delta_0$ der Startscheibe, und
alle anderen Bewegungen werden als Vielfache von $\delta_0$ geschrieben (siehe unten).

\textbf{Einsetzen (\code{w\_eq2}).} Mit $\varphi_i = \kappa_i\,\varphi_0$ wird jeder Weg zu
$h_k\,\kappa_{i(k)}\,\varphi_0$ (Spalte~3 von \code{pvv\_kin}). Dabei ist $i(k)$ die Scheibe, an der
die Kraft $F_k$ angreift, und $h_k$ ihr Hebel um den Pol dieser Scheibe, also die eckige Klammer
von oben, $h_k = e_x\,(P_y - c_y) - e_y\,(P_x - c_x)$. Das ist der senkrechte Abstand der
Wirkungslinie vom Pol mit Vorzeichen, nicht der Abstand des Angriffspunktes. Verschiebt sich die
Scheibe nur, tritt an die Stelle von $h_k\,\varphi_i$ der Weg $\delta_i\,(\vec e\cdot\vec t_i)$ aus
der Liste oben, $h_k$ ist dann der Richtungskosinus $\vec e\cdot\vec t_i$. Der Arbeitssatz zerfällt in
\[
  \delta W = \varphi_0 \Bigl(\sum_k F_k\,h_k\,\kappa_{i(k)} + \sum_m M_m\,\kappa_{i(m)}\Bigr) = 0 .
\]
Da $\varphi_0$ beliebig ist, muss die Klammer verschwinden. Sie ist linear in der freigesetzten
Größe, also eine Gleichung mit einer Unbekannten. Verschiebt sich die Bezugsscheibe nur, steht
$\delta_0$ statt $\varphi_0$ da, und die Faktoren der Drehwinkel haben die Einheit $1/$Länge.
Intern zählt das Programm Drehwinkel im Uhrzeigersinn. Im Export ist das Vorzeichen umgedreht,
damit alles zur Konvention der Vorlesung (gegen den Uhrzeigersinn positiv) passt.

%---------------------------------------------------------------------
\section{Gleichgewicht der Teilsysteme (LGS-Matrix)}\label{sec:mathe-gleichgewicht}

Für den Export \emph{LGS-Matrix} schneidet das Programm das System an allen Gelenken auf, genau
so, wie Seite~2 der Explosionszeichnung es zeigt. Beide benutzen dieselben drei Regeln im Code:
\code{gelenkKnoten} (welche Stäbe hängen mit Gelenk, welche biegesteif am Knoten, wie viele
Teilsysteme treffen sich), \code{gelenkRichtung} (positive Richtung der Gelenkgrößen) und
\code{lastStabKnoten} (an welchem Teilsystem Lasten und Lager eines Gelenkknotens hängen). Jedes
Teilsystem liefert drei Gleichungen: Kräfte nach rechts, nach oben und Momente gegen den
Uhrzeigersinn. Unbekannt sind die Lagerreaktionen und die Größen, die ein Gelenk überträgt: beim
M-Gelenk die beiden Kraftkomponenten $Gx$, $Gy$, beim N-Gelenk $Gq$ und $Gm$, beim Q-Gelenk $Gn$
und $Gm$. Treffen sich mehr als zwei Teilsysteme in einem Gelenk, bekommt jeder Gelenkstab eigene
Gelenkkräfte, gibt es dort keinen biegesteifen Rest, sagen zwei zusätzliche Gleichungen, dass der
Gelenkbolzen selbst im Gleichgewicht ist. Bei nur zwei Teilsystemen braucht der Bolzen keine
eigenen Gleichungen: Die Gelenkkraft steht in beiden mit entgegengesetztem Vorzeichen (actio =
reactio), damit ist sein Gleichgewicht schon erfüllt. Hängt am Gelenk ein biegesteifer Rest, gehört
der Bolzen zu ihm. Ein Beispiel ist eine Stütze, die durch den Knoten durchläuft, und zwei Riegel,
die gelenkig an ihr hängen. Die Gelenkkräfte der Riegel greifen dann direkt an der Stütze an und
stehen in deren drei Gleichungen. Der Bolzen ist kein eigener Körper, also gibt es für ihn keine
eigenen Gleichungen, sein Gleichgewicht steckt im Gleichgewicht der Stütze. Erst wenn alle Stäbe
gelenkig am Knoten hängen, gehört der Bolzen zu keinem Teilsystem und braucht seine zwei
Gleichungen selbst. Im symbolischen Modus stehen Lasten mit ihren Namen und
Hebelarme als Vielfache von $l$ in der Matrix.

Den \textbf{Momentenbezugspunkt} wählt das Programm so, wie ihr es in \textbf{TM1} lernt: Es
sucht den Punkt, durch den möglichst viele unbekannte Kräfte gehen, also die senkrechte Linie mit den
meisten unbekannten Vertikalkräften und die waagerechte mit den meisten unbekannten
Horizontalkräften, am liebsten ein Knoten, der auf beiden liegt. Die Momentengleichung wird dadurch
so kurz wie möglich.

Die Werte auf Seite~1 der Explosionszeichnung kommen aus dem Weggrößenverfahren. Löst ihr das LGS
(mit den oben beschriebenen Befehlen oder von Hand), müssen dieselben Zahlen herauskommen. Das ist
eure Kontrolle.

\textbf{Lösungsweg (\code{lgs\_weg}):} Gesucht ist eine Reihenfolge, in der jede Gleichung möglichst
nur eine neue Unbekannte enthält. Dafür braucht das Programm mehr Gleichungen als die in \code{lgs}.
Alle sind Linearkombinationen der Zeilen aus \code{lgs}, also nichts Neues, nur anders
zusammengefasst.

\textbf{Anderer Momentenpunkt.} Eine Kraft mit den Komponenten $F_\rightarrow$ (nach rechts) und
$F_\uparrow$ (nach oben) im Punkt $(x, y)$ hat um $P$ das Moment
$M_P = (x - x_P)\,F_\uparrow - (y - y_P)\,F_\rightarrow$ (gegen den Uhrzeigersinn positiv, hier $y$
nach oben). Die Pfeile sind dieselben wie in den Namen der Kraftzeilen von \code{lgs}. Um einen
anderen Punkt $Q$ ist es
\[
  M_Q = (x - x_Q)\,F_\uparrow - (y - y_Q)\,F_\rightarrow
      = M_P + (x_P - x_Q)\,F_\uparrow - (y_P - y_Q)\,F_\rightarrow .
\]
Über alle Kräfte eines Teilsystems summiert (reine Momente hängen nicht vom Punkt ab) heißt das
\[
  \textstyle\sum M_Q = \sum M_P + (x_P - x_Q) \sum F_\uparrow - (y_P - y_Q) \sum F_\rightarrow .
\]
Die Momentengleichung um $Q$ ist also die Momentenzeile plus Vielfache der beiden Kraftzeilen
desselben Teilsystems, die Faktoren sind Hebelarme. Im Programm zeigt $y$ nach unten, dort steht
deshalb $+(y_P - y_Q)$.

\textbf{Teilsysteme zusammenfassen.} Für eine Gruppe werden die Kraftzeilen addiert und die
Momentenzeilen, jede vorher auf denselben Punkt $Q$ umgerechnet. Eine Gelenkkraft zwischen zwei
Teilsystemen der Gruppe steht in beiden mit entgegengesetztem Vorzeichen und greift am selben Punkt
an. Sie fällt deshalb aus den Kräften \emph{und} aus den Momenten heraus. Alle Teilsysteme zusammen
sind das Gesamtsystem.

\textbf{Kräfte quer zu einer Richtung.} Zeigt eine Unbekannte in Richtung $(d_x, d_y)$ (wieder mit $y$
nach unten), dann enthält $d_y \sum F_\rightarrow + d_x \sum F_\uparrow$ (normiert) diese Unbekannte
nicht mehr.

\textbf{Warum höchstens drei Gleichungen je Scheibe?} Jede Gleichgewichtsbedingung einer Scheibe ist
entweder eine Kraftsumme in eine Richtung $\vec e$, also $e_x \sum F_\rightarrow + e_y \sum
F_\uparrow$, oder eine Momentensumme um einen Punkt, nach der Formel oben $\sum M_P$ plus Vielfache
der beiden Kraftsummen. Alle liegen also in dem Raum, den $\sum F_\rightarrow$, $\sum F_\uparrow$ und
$\sum M_P$ aufspannen, und der hat höchstens die Dimension~3. Mehr als drei linear unabhängige
Gleichungen bekommt ihr deshalb nie, egal um wie viele Punkte ihr die Momente bildet. Drei
Momentengleichungen um $A$, $B$ und $C$ sind genau dann unabhängig, wenn die drei Punkte nicht auf
einer Geraden liegen. Liegen sie auf einer Geraden, ist die dritte eine Kombination der beiden
anderen. Ein Durchlaufträger auf drei Lagern bleibt deshalb einfach statisch unbestimmt, auch wenn
man um jedes Lager die Momente bildet. Anschaulich (PvV): Jede Gleichung gehört zu einer virtuellen
Starrkörperbewegung, einer Verschiebung in Richtung $\vec e$ bzw.\ einer Drehung um $P$, und eine
Scheibe hat in der Ebene nur drei Freiheitsgrade.

\textbf{Welche Gruppen und Punkte?} Das Programm sucht die Gleichungen nicht gezielt, sondern stellt
alle aus einem festen Vorrat auf und wählt danach aus:

\begin{itemize}
  \item \textbf{Gruppen:} jedes Teilsystem einzeln, das Gesamtsystem und bei Rahmen alle
        zusammenhängenden Gruppen. Zwei Teilsysteme hängen zusammen, wenn sie eine
        Gelenkunbekannte teilen. Ohne gemeinsame Gelenkkraft hebt sich beim Zusammenfassen nichts
        auf. Die Gruppen entstehen, indem an jede Gruppe ein Nachbar angehängt wird, bis die
        Zahl unter \emph{Max. Teilgruppen} erreicht ist (ab Werk 30). Im Fachwerk gibt es nur die Knoten und das
        Gesamtsystem, denn eine Gruppe von Knoten wäre ein Ritterschnitt.
  \item \textbf{Momentenpunkte $Q$:} Knoten und Schnittpunkte zweier Wirkungslinien von
        Unbekannten bzw.\ Pendelachsen (innerhalb der doppelten Systemgröße), aber je Gruppe nur
        die, an denen eine \emph{äußere} Kraft der Gruppe herausfällt. Äußere Kräfte sind die, die
        nach dem Zusammenfassen noch in der Kraftsumme stehen. Gelenkkräfte zwischen zwei
        Teilsystemen der Gruppe heben sich auf und zählen nicht. Ein Knoten zählt, wenn die
        Wirkungslinie einer äußeren Kraft durch ihn geht, ein Schnittpunkt, wenn beide Linien zu
        äußeren Kräften gehören, und er wird nach diesem Paar benannt. Für \code{TS 1+2} ist das
        Gelenk zwischen TS~1 und TS~2 deshalb nur dann ein Kandidat, wenn noch eine andere äußere
        Kraft durch es hindurchgeht. $Q$ wird also nicht vorher ausgesucht: Jeder dieser Punkte
        liefert eine Gleichung, und die Auswahl unten nimmt die, in der nur noch eine Unbekannte
        übrig ist.
  \item \textbf{Kräfte quer} zu jeder schrägen äußeren Unbekannten und jeder schrägen Pendelachse.
\end{itemize}

Gleichungen, die sich nur um einen Faktor unterscheiden, bleiben einmal übrig, und zwar in der
einfachsten Form. Weil dabei nur mit Zahlen gerechnet wird, dauert das auch bei einigen hundert
Kandidaten kaum Zeit.

\textbf{Benennung.} Der Name nennt die Art, den Punkt und die Gruppe: \code{$\Sigma$M P5 TS 1+2}
sind die Momente um Knoten~5 an den Teilsystemen 1 und 2 zusammen, \code{$\Sigma$M S(A\_h,B\_v)
Ges} die Momente des Gesamtsystems (\code{Ges}) um den Schnittpunkt der Wirkungslinien von $A_h$ und
$B_v$, \code{$\Sigma$F quer zu B\_h TS 2} die Kräfte an Teilsystem~2 quer zu $B_h$. Neue Gleichungen
bekommen die nächste freie Nummer nach denen aus \code{lgs}.

\textbf{Pendelstäbe.} Ein Teilsystem eines Rahmens ohne Lasten, an dem nur an zwei Punkten $U$ und
$V$ Anschlusskräfte angreifen (Lager- bzw.\ Gelenkkräfte, keine Momente), nennen wir hier
Pendelstab. Ist es ein gerader Stab, ist das der Pendelstab aus \textbf{TM1} mit reiner
Normalkraft. Ist es geknickt (z.\,B.\ ein unbelasteter Winkel zwischen zwei Gelenken), hat es zwar
Momente und Querkräfte, für das Gleichgewicht verhält es sich aber genauso: Die Anschlusskraft liegt
auf der Verbindungslinie $UV$. Nur darum geht es hier. Sein
Momentengleichgewicht um $U$ enthält nur die beiden Kraftkomponenten bei $V$ und sagt: Die Kraft bei
$V$ liegt auf der Achse $UV$. Liegt der Punkt $Q$ einer anderen Momentengleichung auf dieser Achse,
stehen die beiden Komponenten bei $V$ dort im selben Verhältnis wie im Pendelgleichgewicht, nur mit
einem Faktor $\lambda$. Zieht man $\lambda$-mal das Pendelgleichgewicht ab, fallen beide heraus. Das
Programm prüft das für jede Gleichung und jeden Pendelstab mit den Zahlen und nimmt die Achse als
weitere Wirkungslinie für die Schnittpunkte dazu. So findet sich auch beim Dreigelenkrahmen mit
verschieden hohen Lagern und einer unbelasteten Hälfte ein Punkt, der direkt eine Lagerkraft
liefert. Ist die Achse waagerecht oder senkrecht, braucht es den Umweg nicht, dann hat das
Pendelgleichgewicht selbst schon nur eine Unbekannte. Dieselbe Gleichung entsteht oft auf zwei
Wegen, etwa am Gesamtsystem (Lagerkräfte über das Pendelgleichgewicht um das Gelenkende
herausgerechnet) und an der anderen Scheibe allein (Gelenkkräfte über das Pendelgleichgewicht um
das Lager herausgerechnet). Das Programm behält dann den Namen mit dem Gesamtsystem, weil das der
Vorstellung entspricht, dass die Lagerkraft auf der Pendelachse wirkt.

\textbf{Auswahl.} Dann geht es schrittweise: Gibt es eine Gleichung mit genau einer noch unbekannten
Größe, ist das der nächste Schritt. Unter mehreren gewinnt die einfachste: erst die aus \code{lgs},
dann andere Momentenpunkte am selben Teilsystem, dann das Gesamtsystem, dann Zusammenfassungen und
Schnittpunkte, zuletzt Gleichungen, in denen ein Pendelgleichgewicht steckt. Bei Gleichstand gewinnt
die mit weniger Unbekannten, danach die, die zuerst aufgestellt wurde (Knoten vor Schnittpunkten,
nach Nummern). Lagergrößen (Lager und Federn) haben dabei Vorrang, aber nur innerhalb derselben
Stufe: Ein Einzelschritt für eine Lagergröße geht vor jedem anderen Einzelschritt. Gibt es keinen,
nimmt das Programm einen Einzelschritt, nach dem eine Lagergröße als Einzelschritt folgt (einen
Schritt vorausgeschaut), sonst den einfachsten. Ein Block kommt erst, wenn es gar keinen
Einzelschritt mehr gibt. Gibt es keinen, sucht das Programm zwei Gleichungen mit denselben zwei neuen
Unbekannten, die voneinander unabhängig sind (Rang~2), danach drei mit drei. Unter gleich großen
Blöcken gewinnt der mit mehr Lagergrößen, dann der aus den einfacheren Gleichungen. Dafür hat jede
Gleichung eine Stufe (\emph{Klasse}) nach der Reihenfolge oben: $0$ für eine Zeile aus \code{lgs},
höher für andere Momentenpunkte, das Gesamtsystem, Zusammenfassungen, Schnittpunkte und zuletzt
Gleichungen mit Pendelgleichgewicht. Verglichen wird die Summe der Stufen der Gleichungen im Block,
die kleinere gewinnt. Erst wenn auch das nicht geht, löst es alle übrigen Unbekannten als einen
Block aus Gleichungen von \code{lgs}. Ein solcher Block wird durch Einsetzen gelöst: Die Gleichung
mit den wenigsten Unbekannten wird nach einer davon aufgelöst und in die übrigen eingesetzt, das
wiederholt sich, bis die letzte Gleichung nur noch eine Unbekannte hat. Dann wird von hinten
zurückeingesetzt. Bei gleich vielen Unbekannten nimmt das Programm am liebsten eine Unbekannte, die
in ihrer Gleichung den Faktor $+1$ oder $-1$ hat, denn dann entsteht beim Auflösen kein Bruch.

Die Auswahl läuft mit Zahlen, denn die Faktoren vor den Unbekannten hängen nur von der Geometrie ab,
auch im symbolischen Modus. Die Texte und Ergebnisse in \code{lgs\_weg} rechnet danach das CAS
exakt aus, sie sind also dieselben wie die Lösung des ganzen LGS.

\textbf{Option LGS wie Lösungsweg.} Der Lösungsweg benutzt genau $n$ Gleichungen für $n$
Unbekannte. Ordnet man sie nach den Schritten, hat ihre Matrix Blockdreiecksform: Jeder Schritt
enthält außer schon bekannten Größen nur seine neuen Unbekannten, und deren Block ist regulär
(Einzelschritt: Faktor ungleich null, Block: voller Rang). Damit ist die ganze Matrix regulär. Das
Programm lässt die Gleichungen aus \code{lgs} stehen, die der Weg benutzt, und setzt jede
abgeleitete Gleichung auf den Platz einer Gleichung, die er nicht benutzt, am liebsten auf eine,
aus der sie mit entsteht. Danach prüft es den Rang des neuen Systems noch einmal mit Zahlen. Allgemein
gilt: Ersetzt man Zeile $r$ durch $\sum_i f_i\, r_i$, bleibt die Lösung eindeutig genau dann, wenn
$f_r \neq 0$ ist (die Determinante ändert sich um den Faktor $f_r$).

\textbf{Gegenprüfung mit der FEM.} Für jede Unbekannte bestimmt das Programm aus den FEM-Ergebnissen
den Wert mit demselben Namen und derselben Richtung wie auf Seite~2: Lagerreaktionen direkt,
Gelenkgrößen aus den Stabendkräften des angeschlossenen Stabes (ohne Knotenlast und Lagerkraft, wenn
die zu diesem Stab gehören), Stabkräfte im Fachwerk aus der Normalkraft. Das vergleicht es mit dem
Wert aus \code{lgs\_weg}, im symbolischen Modus mit denselben Ersatzwerten für die Symbole, mit denen
die FEM gerechnet hat. Weicht etwas um mehr als $10^{-4}$ des größten Wertes ab, kommt die Warnung.

\textbf{Resultierende von Streckenlasten:} $R = \int_0^L q\,dx$ und das statische Moment
$S = \int_0^L q\,x\,dx$ um den Stabanfang, die Wirkungslinie liegt bei $x_R = S/R$. Für Trapezlasten
sind das die bekannten Formeln, für Lastfunktionen integriert das Programm mit der Simpsonregel
über elf Stützstellen. Die ist für $q\,x$ bis zu einem quadratischen $q$ exakt.

%---------------------------------------------------------------------
\section{Querschnittswerte}\label{sec:mathe-qswerte}

\textbf{Vielecke} (Rechteck, Dreieck, Trapez, Polygon, dicke Linie) rechnet das Programm alle mit
derselben Formel, der Gaußschen Trapezformel. Mit den Eckpunkten $(y_i, z_i)$ gegen den
Uhrzeigersinn und $c_i = y_i z_{i+1} - y_{i+1} z_i$ ist
\[
  A = \tfrac12 \sum c_i, \qquad
  y_S = \frac{\sum (y_i + y_{i+1})\,c_i}{6A}, \qquad
  \int z^2\,dA = \tfrac1{12} \sum \left(z_i^2 + z_i z_{i+1} + z_{i+1}^2\right) c_i ,
\]
entsprechend für $z_S$, $\int y^2\,dA$ und $\int yz\,dA$. Das ist exakt, es ist der Gaußsche
Integralsatz, angewandt auf die Randkanten. Deshalb ist es auch egal, wie kompliziert das Vieleck
ist, solange es sich nicht selbst schneidet.

\textbf{Kreisausschnitt und -abschnitt} mit halbem Öffnungswinkel $\alpha$: Der Schwerpunkt liegt
auf der Winkelhalbierenden im Abstand $\frac{2R\sin\alpha}{3\alpha}$ (Ausschnitt) bzw.\
$\frac{4R\sin^3\alpha}{3(2\alpha - \sin 2\alpha)}$ (Abschnitt). Die Trägheitsmomente werden erst um
die Winkelhalbierende bestimmt (beim Abschnitt wird das Dreieck zum Mittelpunkt abgezogen) und dann
ins KOS gedreht.

\textbf{Dünnwandige Linien} sind Mittellinie mal Dicke: $A = L\,t$,
$I_y = t L^3 \sin^2\beta/12$ usw. Den Eigenanteil $L t^3/12$ lässt das Programm weg, genau wie
die Vorlesung, weil er von höherer Ordnung in $t$ ist. \textbf{Dünnwandige Kreise und Bögen}
ebenso: $A = R\,\Delta\varphi\,t$, $I_y = R^3 t \int \sin^2\varphi\,d\varphi$ um den Mittelpunkt
(Vollkreis $\pi R^3 t$), Schwerpunkt des Bogens im Abstand $R\,\frac{\sin(\Delta\varphi/2)}{\Delta\varphi/2}$
auf der Winkelhalbierenden.

\textbf{Löcher} gehen mit negativer Fläche in die Summen ein, danach setzt der Satz von Steiner alles
zusammen. \textbf{Numerisch addierte} Elemente legt das Programm auf ein Raster von $40 \times 40$
Zellen über ihr umschreibendes Rechteck und zählt nur die Zellen, die nicht schon in einem anderen
Element liegen. Deshalb die Warnung, das ist eine Näherung.

\textbf{Hauptachsen:} Mit $m = \frac{I_y + I_z}2$ und
$d = \sqrt{\left(\frac{I_y - I_z}2\right)^2 + I_{yz}^2}$ ist $I_{1,2} = m \pm d$ und
\[
  \alpha = \tfrac12 \operatorname{atan2}(2 I_{yz},\ I_y - I_z) .
\]
Das \code{atan2} (der Arkustangens, der das Vorzeichen von Zähler und Nenner kennt) sorgt dafür,
dass $\alpha$ immer zur Achse~1 mit $I_1$ gehört. Mit dem einfachen $\tan 2\alpha$ aus der
Formelsammlung müsstet ihr das von Hand prüfen, dort kann auch die Achse mit $I_2$ herauskommen.

\textbf{Widerstandsmomente} $W = I/e$ nehmen als Randabstand $e$ den größten Abstand der Elemente
vom Schwerpunkt. Bei dünnwandigen Elementen ist das die Mittellinie, $t/2$ kommt nicht dazu.

%---------------------------------------------------------------------
\section{Normalspannung und neutrale Faser}\label{sec:mathe-sigma}

Mit $\Delta = I_y I_z - I_{yz}^2$ und $y$, $z$ vom Schwerpunkt aus rechnet das Programm immer mit
\[
  \sigma_x = \frac NA + \frac{M_y I_z - M_z I_{yz}}{\Delta}\,z
                      - \frac{M_z I_y - M_y I_{yz}}{\Delta}\,y .
\]
Für $I_{yz} = 0$ bleibt das bekannte $\frac NA + \frac{M_y}{I_y} z - \frac{M_z}{I_z} y$ übrig. Die
Formel kommt direkt aus $M_y = \int \sigma z\,dA$ und $M_z = -\int \sigma y\,dA$ mit dem
linearen Ansatz $\sigma = a + b\,y + c\,z$, also zwei Gleichungen für $b$ und $c$.

Weil $\sigma_x$ linear in $y$ und $z$ ist, speichert das Programm nur den Gradienten $(c_y, c_z)$ und
den Wert $c_0$ im Ursprung des gezeigten KOS. Daraus entstehen beide Gleichungen in der Tabelle: die
KOS-Form direkt, die HAS-Form, indem der Gradient um den Hauptachsenwinkel gedreht wird. Die
neutrale Faser ist die Gerade $c_0 + c_y y + c_z z = 0$.

Die \textbf{Extremwerte} sucht das Programm nur an den Eckpunkten des Querschnitts (Bögen alle
$10^\circ$ abgetastet). Eine lineare Funktion hat ihr Maximum immer in einer Ecke der konvexen
Hülle. Bei dünnwandigen Profilen sind das die Endpunkte der Mittellinien, also wieder die
Näherung aus der Vorlesung.

%---------------------------------------------------------------------
\section{Schubfluss in dünnwandigen Profilen}\label{sec:mathe-schub}

\subsection{Offene Profile}

Aus den dünnwandigen Linien baut das Programm einen \emph{Profilgraph}: Die Linien werden an
Kreuzungen, Berührpunkten und an den Schwerachsen geteilt, Punkte, die zusammenfallen, werden ein
Knoten. \textbf{Kreise und Bögen} kommen als Kette kurzer Sehnen dazu, geteilt an ihren Enden, an
Endpunkten anderer Elemente auf dem Bogen, an Kreuzungen mit geraden Linien und an den Schwerachsen,
dazwischen gleichmäßig mit höchstens $15^\circ$ je Sehne. Die Sehnen bestimmen nur die Knoten und
den Zusammenhang (Zellen, Symmetrie). Jede Sehne kennt ihren Bogen (Mittelpunkt, Radius, Winkel
$\varphi_1$ bis $\varphi_2$), und gerechnet wird auf dem Bogen: Länge $R\,|\varphi_2 - \varphi_1|$
statt der Sehne, Stützstellen auf dem Bogen. Dann läuft es los:

\begin{itemize}
  \item An jedem \textbf{freien Ende} ist der Schubfluss null, dort beginnt ein Pfad mit
        $S_y = S_z = 0$.
  \item Entlang eines Segments werden die statischen Momente aufintegriert,
        $S_z \mathrel{+}= (y - y_S)\,t\,\Delta s$ und $S_y \mathrel{+}= (z - z_S)\,t\,\Delta s$. Weil
        der Integrand auf einer geraden Linie linear ist, ist das mit der Mittelpunktsregel exakt.
        $S$ selbst ist quadratisch in $s$. Auf einem Bogen mit Mittelpunkt $(y_C, z_C)$ ist
        $ds = R\,d\varphi$, und das Integral ist geschlossen lösbar:
        \[
          \Delta S_z = t R \left[(y_C - y_S)\,\Delta\varphi + R\,(\sin\varphi_b - \sin\varphi_a)\right],
          \quad
          \Delta S_y = t R \left[(z_C - z_S)\,\Delta\varphi - R\,(\cos\varphi_b - \cos\varphi_a)\right]
        \]
        (für den Lauf gegen den Uhrzeigersinn, sonst mit umgekehrtem Vorzeichen).
  \item An einer \textbf{Verzweigung} gilt das Knotengleichgewicht: Kommen alle Äste bis auf einen an,
        startet der letzte mit der Summe der ankommenden statischen Momente.
\end{itemize}

Aus den statischen Momenten folgt der Schubfluss mit der gekoppelten Formel
\[
  \tau\,t = -\frac{Q_y\,(I_y S_z + I_{yz} S_y) + Q_z\,(I_z S_y + I_{yz} S_z)}{I_y I_z - I_{yz}^2} ,
\]
die für $I_{yz} = 0$ in das bekannte $\tau = -\frac{Q_z S_y}{I_y t}$ bzw.\ $-\frac{Q_y S_z}{I_z t}$
zerfällt. $\tau$ ist positiv in Laufrichtung $s$.

\subsection{Eine geschlossene Zelle: woher kommt \texorpdfstring{$q_0$}{q0}?}\label{sec:mathe-zelle}

Eine Zelle hat kein freies Ende, an dem der Schubfluss null sein müsste. Das Programm schneidet sie
deshalb an einer Stelle auf, rechnet wie beim offenen Profil und bekommt einen Verlauf, der sich
vom richtigen nur um einen \textbf{konstanten Umlaufanteil} $q_0$ unterscheidet. Diesen bestimmt
es auf zwei Wegen:

\textbf{(a) Symmetrie.} Liegt eine Symmetrieachse parallel zur Querkraft, ist der richtige Verlauf
antisymmetrisch und auf der Achse null. Das Programm schneidet die Zelle deshalb möglichst auf der
Achse auf. Klappt das nicht genau, verschiebt es den Umlaufanteil so, dass ein Punkt und sein
Spiegelpunkt entgegengesetzt gleiche Werte haben. Die Symmetrie selbst prüft es, indem es jede
Linie an 13 Stellen abtastet (jede Bogensehne an 3) und schaut, ob der gespiegelte Punkt wieder auf
einer Linie gleicher Dicke liegt. Weil Bögen auch an den Schwerachsen geteilt und die Stücke
gleichmäßig in Sehnen zerlegt werden, sind die Sehnen spiegelbildlicher Bögen selbst
spiegelbildlich.

\textbf{(b) Verträglichkeit.} Ohne passende Symmetrieachse kommt die Bedingung, dass sich die Zelle
unter einer Querkraft im Schubmittelpunkt nicht verdrillt:
\[
  \oint \frac{q}{G\,t}\,ds = 0 .
\]
Der Kniff ist, diese Bedingung nicht für die gerade eingegebene Querkraft zu lösen, sondern für
die statischen Momente selbst. Dazu sortiert man die Formel für den Schubfluss nach $S_z$ und $S_y$:
\[
  q = -\frac{a\,S_z + b\,S_y}{\Delta} , \qquad
  \begin{pmatrix} a \\ b \end{pmatrix}
  = \underbrace{\begin{pmatrix} I_y & I_{yz} \\ I_{yz} & I_z \end{pmatrix}}_{J}
    \begin{pmatrix} Q_y \\ Q_z \end{pmatrix} ,
\]
mit $\Delta = I_y I_z - I_{yz}^2 = \det J$. $a$, $b$, $G$ und $\Delta$ sind entlang der Zelle
konstant, also wird die Verträglichkeit zu
\[
  a \oint \frac{S_z}{t}\,ds + b \oint \frac{S_y}{t}\,ds = 0 .
\]
\emph{Warum $J$ regulär sein muss.} Die statischen Momente hängen nur vom Querschnitt ab, nicht von
der Last. Sollen sie für \emph{jede} Querkraft den richtigen Schubfluss liefern, muss die Gleichung
oben für jedes $(Q_y, Q_z)$ gelten. Weil $J$ regulär ist, gibt es zu jedem gewünschten Paar
$(a, b)$ eine Querkraft, die es erzeugt, nämlich $(Q_y, Q_z)^T = J^{-1}(a, b)^T$. Man darf also
$(a, b) = (1, 0)$ und $(a, b) = (0, 1)$ einsetzen und bekommt die beiden Bedingungen einzeln:
\[
  \oint \frac{S_z}{t}\,ds = 0 \quad\text{und}\quad \oint \frac{S_y}{t}\,ds = 0 .
\]
Wäre $J$ singulär, lägen alle Paare $(a, b)$, die überhaupt vorkommen können, auf einer Geraden durch
den Ursprung. Dann wüsste man nur, dass eine bestimmte Kombination der beiden Umlaufintegrale
verschwindet, und könnte die beiden statischen Momente nicht einzeln festlegen.

Regulär ist $J$ bei jedem echten Querschnitt: Nach der Cauchy-Schwarz-Ungleichung aus \textbf{HM1}
ist $I_{yz}^2 = \bigl(\int y z\,dA\bigr)^2 \le \int y^2\,dA \int z^2\,dA = I_z I_y$, also
$\Delta \ge 0$. Gleichheit gibt es nur, wenn die ganze Fläche auf einer Geraden durch den Schwerpunkt
liegt, und dann versagt schon die Biegeformel, in der $\Delta$ im Nenner steht.

Streng genommen braucht man die Regularität nur für diese Richtung des Schlusses. Dass zwei
einzeln verschobene statische Momente für jede Querkraft den richtigen Schubfluss liefern, folgt
schon daraus, dass $q_0 = -\oint (q_{\text{offen}}/t)\,ds \big/ \oint ds/t$ linear in $a$ und $b$
ist, also in einen Teil zu $S_z$ und einen zu $S_y$ zerfällt. Die Regularität sagt zusätzlich, dass
es keine andere Wahl der beiden Konstanten gibt. Deshalb sind auch die angezeigten $S$-Verläufe
eindeutig.

Das Programm verschiebt also jedes statische Moment in der Zelle um seine eigene Konstante $s_0$.
Aus $\oint (S + s_0)/t\,ds = 0$ folgt
\[
  s_0 = -\,\frac{\oint S/t\;ds}{\oint ds/t}
\]
(beide Integrale im Umlaufsinn gegen den Uhrzeigersinn). Damit stimmen nicht nur $\tau$, sondern
auch die angezeigten $S$-Verläufe. Integriert wird mit der Simpsonregel, die für das quadratische
$S/t$ exakt ist. Auf einem Bogen ist $S$ eine glatte Winkelfunktion, mit 16 Teilen je Sehne
von höchstens $15^\circ$ ist Simpson dort auf etwa $10^{-10}$ genau. Beim symmetrischen Profil kommt
so übrigens genau die antisymmetrische Lösung aus (a) heraus.

\subsection{Schubmittelpunkt}

Der Schubmittelpunkt ist der Punkt, durch den die Resultierende des Schubflusses geht. Das
Programm rechnet den Schubfluss für die Einheitslasten $Q_y = 1$ und $Q_z = 1$ (bei einer Zelle
inklusive $q_0$) und dessen Moment um den Schwerpunkt,
\[
  M_q = \int \left[(y - y_S)\,t_z - (z - z_S)\,t_y\right] q\,ds ,
\]
mit der Tangentenrichtung $(t_y, t_z)$ der Linie. Weil dasselbe Moment aus der Querkraft im
Schubmittelpunkt kommen muss, $M_q = (y_M - y_S)\,Q_z - (z_M - z_S)\,Q_y$, ist
\[
  y_M = y_S + M_q(Q_z = 1), \qquad z_M = z_S - M_q(Q_y = 1) .
\]
Der Integrand ist kubisch in $s$, die Simpsonregel also wieder exakt (auf Bögen, mit der
Bogentangente, wieder auf etwa $10^{-10}$). Die Momententabelle zeigt genau diese Summe, aufgeteilt
nach Elementen. Gegen die Handrechnung: offener Halbkreis $e = 4R/\pi$, D-Profil aus Halbkreis und
Platte und Kreisrohr stimmen auf etwa $10^{-9}$.

\subsection{Massive Querschnitte}

Für $\tau = \frac{Q\,S}{I\,b}$ teilt das Programm den Querschnitt in etwa 160 waagerechte (bzw.\
senkrechte) Streifen. Die Breite $b(z)$ kommt exakt aus allen Elementen: Löcher abgezogen, Polygone
über die Schnittpunkte der Kanten, Kreise über $2\sqrt{R^2 - \Delta^2}$, Kreisausschnitt und
-abschnitt als Kreis geschnitten mit den Halbebenen der Randstrahlen bzw.\ der Sehne. Die
Streifengrenzen liegen auf allen Knickstellen von $b$ (Ecken, Rechteckkanten, Kreisscheitel,
Lochränder) und auf der Schwerachse, dazwischen gleich weit. In jedem Streifen ist $b$ damit glatt,
und $S(z)$ wird mit drei Gauß-Punkten je Streifen von oben aufsummiert (an Knickstellen fein
unterteilt, wegen der Wurzel am Kreisscheitel).

$\tau$ wertet das Programm auf den Streifengrenzen aus, also auch genau am oberen und unteren Rand.
Dort ist $S = 0$, damit $\tau = 0$. Der kleine Rest, den die Quadratur am zweiten Rand übrig lässt,
wird anteilig nach der Fläche abgezogen. An Spitzen und am Kreisscheitel geht $b \to 0$, aber $S$
schneller, der Grenzwert $\tau \to 0$ wird direkt eingesetzt. An einem Breitensprung (Flansch und
Steg, Lochrand) gibt es zwei Stützstellen, eine mit der Breite darunter und eine mit der darüber,
der Verlauf zeigt also den Sprung. Früher wurde $\tau$ in der Streifenmitte ausgewertet, aber mit
$S$ bis zur Streifenunterkante: Der oberste Wert war deshalb nie null (beim Rechteck etwa
$2{,}5\,\%$ von $\tau_{max}$), und der Rand selbst kam gar nicht vor.

Wirken $Q_y$ und $Q_z$ gleichzeitig, ist das angezeigte $\tau_{max}$
$\sqrt{\tau_{y,max}^2 + \tau_{z,max}^2}$, also eine Abschätzung nach oben, denn die beiden Maxima
liegen meist nicht an derselben Stelle.

%---------------------------------------------------------------------
\section{Torsion und Verwölbung}\label{sec:mathe-torsion}

\textbf{Torsion:} Beim offenen Profil ist $I_T = \sum \frac{\xi}{3} L\,t^3$ und
$\tau = M_T\,t/I_T$. Bei einer Zelle kommt die Fläche $A_m$ der Mittellinie aus der Gaußschen
Trapezformel und $\oint ds/t = \sum L_i/t_i$, damit $I_T = 4A_m^2/\oint ds/t$ und
$\tau = M_T/(2A_m t)$. Läuft die Zelle über einen Bogen, nimmt die Trapezformel die Sehnen, und je
Sehne mit dem Winkel $\Delta\varphi$ kommt das Kreissegment
\[
  \Delta A = \frac{R^2}{2}\left(\Delta\varphi - \sin\Delta\varphi\right)
\]
zwischen Sehne und Bogen dazu (mit dem Vorzeichen des Umlaufs). $L_i$ ist dort die Bogenlänge
$R\,\Delta\varphi$. Beides ist exakt, die Sehnen sind also nur eine Hilfe für die Topologie.

\textbf{Zelle mit offenen Ästen} (z.\,B.\ ein Kasten mit überstehenden Gurten): Die Zelle trägt
nach Bredt das ganze $M_T$, und $I_T$ ist ihr Bredt-Wert. Daraus folgt die Verdrillung
$\vartheta = M_T/(G\,I_T)$, und die offenen Äste müssen sich genauso verdrillen wie die Zelle. Ein
Ast der Dicke $t$ hat dann am Rand $\tau = G\,\vartheta\,t = M_T\,t/I_T$. Streng genommen nehmen die
Äste damit zusätzlich das Moment $G\,\vartheta\,I_{T,\text{offen}}$ auf, und alle Schubspannungen
zusammen ergeben $M_T\,(1 + I_{T,\text{offen}}/I_{T,\text{Zelle}})$ statt $M_T$. Das ist eine
Näherung: $I_{T,\text{offen}}$ wächst mit $t^3$, $I_{T,\text{Zelle}}$ mit $A_m^2$, und der
Unterschied liegt bei üblichen Profilen im Promillebereich. Bei einem Kasten $200 \times 100$ mit
$t = 10$ und zwei $50$ langen Gurten derselben Dicke sind es etwa $0{,}1\,\%$. Genauso lässt Bredt
den kleinen offenen Anteil $\frac13 L\,t^3$ der Zellwände selbst weg. Exakt im Gleichgewicht wäre
$I_T = I_{T,\text{Zelle}} + I_{T,\text{offen}}$ mit einer Aufteilung von $M_T$ im Verhältnis der
beiden Steifigkeiten.

Ein
dünnwandiges Kreisrohr ist die Zelle mit $A_m = \pi R^2$ und $\oint ds/t = 2\pi R/t$, also
$I_T = 2\pi R^3 t$. Das Programm erkennt das Rohr aus einem einzelnen Kreis gesondert. Über die
Sehnenkette käme dasselbe heraus, ein Rohr aus zwei Halbkreisen rechnet so. Das Torsionsmoment
aus platzierten Kräften ist das Moment um den \emph{Schubmittelpunkt}, denn um den dreht sich der
Querschnitt.

\textbf{Verwölbung:} Auf einem geraden Segment ist der Fahrstrahlabstand $r_\perp$ linear, damit
ist die Verwölbung nach Formelsammlung geschlossen integrierbar:
\[
  u_x(s) = u_{x,0} + \frac{q_T}{G\,t}\,s - \vartheta\left[r_1 s + (r_2 - r_1)\frac{s^2}{2L}\right],
  \qquad q_T = \frac{M_T}{2A_m}, \quad \vartheta = \frac{M_T}{G\,I_T} ,
\]
wobei der Bredt-Anteil $q_T/(Gt)$ nur in der Zelle auftritt (mit Vorzeichen je nach Umlaufsinn).
Auf einem Bogen ist $r_\perp$ nicht linear, das Integral aber ebenfalls geschlossen lösbar. Mit
dem Pol $(y_M, z_M)$ und dem Winkel $\varphi$ ab dem Sehnenanfang $\varphi_1$ ist
\[
  \int r_\perp\,ds = R\left[(y_C - y_M)(\sin\varphi - \sin\varphi_1)
    - (z_C - z_M)(\cos\varphi - \cos\varphi_1) + R\,(\varphi - \varphi_1)\right] .
\]
Ein Kreisrohr verwölbt sich damit nicht: Mit dem Pol im Mittelpunkt ist $r_\perp = R$, und der
Bredt-Anteil hebt $R\,\vartheta$ genau auf.
Das Programm läuft von einem Knoten auf einer Symmetrieachse aus (dort ist $u_x$ bei symmetrischen
Profilen null) über den Profilgraph und setzt die Werte Segment für Segment fort. Beim Schließen der
Zelle muss es wieder beim Startwert ankommen, weil $\oint r_\perp\,ds = 2A_m$ ist. Das ist die
eingebaute Probe. Zum Schluss wird die Konstante so gewählt, dass $\int u_x\,t\,ds = 0$ ist, also
keine mittlere Längsverschiebung übrig bleibt.

%---------------------------------------------------------------------
\section{Kernfläche}\label{sec:mathe-kern}

Die Kernfläche ist die Menge aller Lastpunkte, deren neutrale Faser den Querschnitt nicht
schneidet. Gerechnet wird rückwärts:

\begin{enumerate}
  \item Das Programm sammelt die Eckpunkte des Querschnitts (Bögen abgetastet) und bildet ihre
        \textbf{konvexe Hülle}. Einbuchtungen spielen keine Rolle, denn eine Gerade, die die Hülle
        berührt, berührt auch den Querschnitt nicht innen.
  \item Jede Kante der Hülle ist eine mögliche \textbf{neutrale Faser}. Mit dem nach außen
        zeigenden Einheitsvektor $\vec n$ der Kante und ihrem Abstand $d$ vom Schwerpunkt gehört
        dazu der Lastpunkt
        \[
          \vec e = -\frac{1}{A\,d}
          \begin{pmatrix} I_z & -I_{yz} \\ -I_{yz} & I_y \end{pmatrix} \vec n .
        \]
        Für ein Rechteck und die obere Kante ist $\vec n = (0,1)$, $d = h/2$, und es kommt
        $e_z = -h/6$ heraus, das bekannte Kerndrittel.
  \item Die Kernpunkte aller Kanten, der Reihe nach verbunden, sind die Kernfläche. An Bögen legt
        das Programm alle $5^\circ$ eine Tangente an, beim Vollkreis 36 Stück.
\end{enumerate}

\textbf{Woher die Matrix kommt.} Die Formel für $\vec e$ ist die Formel für $\sigma_x$ aus
Abschnitt~\ref{sec:mathe-sigma}, rückwärts gelesen. Eine Normalkraft $N$ im Punkt
$\vec e = (e_y, e_z)$, gemessen vom Schwerpunkt, erzeugt $M_y = N e_z$ und $M_z = -N e_y$. In
$\sigma_x$ eingesetzt und nach $y$ und $z$ sortiert:
\[
  \sigma_x = \frac NA + \frac N\Delta\Bigl[\bigl(I_y e_y + I_{yz} e_z\bigr)\,y
                                       + \bigl(I_{yz} e_y + I_z e_z\bigr)\,z\Bigr]
           = \frac NA\Bigl(1 + \frac A\Delta\,\vec x^T J\,\vec e\Bigr) ,
\]
mit $\vec x = (y, z)$ und der Trägheitsmatrix $J = \begin{pmatrix} I_y & I_{yz} \\ I_{yz} & I_z
\end{pmatrix}$, $\Delta = \det J$. Die neutrale Faser ist $\sigma_x = 0$, also die Gerade
\[
  \vec x^T \Bigl(\frac A\Delta\,J\,\vec e\Bigr) = -1 .
\]
Die Kante mit dem nach außen zeigenden Einheitsvektor $\vec n$ und dem Abstand $d$ vom Schwerpunkt
ist die Gerade $\vec x^T \vec n = d$, oder genauso $\vec x^T(-\vec n/d) = -1$. Eine Gerade, die
nicht durch den Schwerpunkt geht, lässt sich nur auf eine Weise in der Form $\vec x^T \vec g = -1$
schreiben. Die neutrale Faser ist also genau dann die Kante, wenn die beiden Vektoren gleich sind:
\[
  \frac A\Delta\,J\,\vec e = -\frac{\vec n}{d}
  \qquad\Longrightarrow\qquad
  \vec e = -\frac{1}{A\,d}\,\bigl(\Delta\,J^{-1}\bigr)\,\vec n ,
  \qquad
  \Delta\,J^{-1} = \begin{pmatrix} I_z & -I_{yz} \\ -I_{yz} & I_y \end{pmatrix} .
\]
Die Matrix in der Formel ist also die Inverse der Trägheitsmatrix ohne den Vorfaktor $1/\det J$
(Diagonale vertauscht, Nebendiagonale mit anderem Vorzeichen). Auch hier braucht man, dass $J$
regulär ist, siehe Abschnitt~\ref{sec:mathe-zelle}.

\textbf{Warum man die Kernpunkte geradlinig verbinden darf.} Dreht man die neutrale Faser um eine
Ecke $\vec P$ der Hülle von einer Kante zur nächsten, gehen alle diese Geraden durch $\vec P$, also
gilt $\vec P^T(\vec n/d) = 1$. Der Vektor $\vec n/d$ wandert damit auf einer Geraden, und weil
$\vec e$ linear von $\vec n/d$ abhängt, wandert auch der Lastpunkt auf einer Geraden. Jeder Ecke der
Hülle entspricht so eine gerade Kante der Kernfläche, jeder Kante der Hülle eine Ecke.

Im Hauptachsensystem fällt die Kopplung weg ($I_{\eta\zeta} = 0$). Schreibt man die Kante mit
ihren Achsenabschnitten als $\eta/a_\eta + \zeta/a_\zeta = 1$, ist $\vec n/d = (1/a_\eta,
1/a_\zeta)$, und aus der Formel wird die bekannte Form $\eta_K = -i_\zeta^2/a_\eta$,
$\zeta_K = -i_\eta^2/a_\zeta$ mit $i^2 = I/A$. Deshalb beschriftet das
Programm die umhüllenden Geraden mit ihren Abschnitten auf den Hauptachsen, gemessen vom
Schwerpunkt, und nicht mit Koordinaten im gezeichneten KOS.

%---------------------------------------------------------------------
\section{Spannungskreis: vom Schnitt zum Zustand}\label{sec:mathe-kreis}

\textbf{Lineare Gleichungen.} Mit $\sigma_m$, $R_c$ und $R_s$ als Unbekannten ist jede Angabe mit
bekanntem Winkel eine Zeile
\[
  \sigma = \sigma_m + R_c\cos 2\alpha + R_s\sin 2\alpha, \qquad
  \tau = -R_c\sin 2\alpha + R_s\cos 2\alpha .
\]
\textbf{Punkte ohne Winkel} erfüllen $(\sigma - \sigma_m)^2 + \tau^2 = R_c^2 + R_s^2$. Zieht man
diese Gleichung für zwei solche Punkte voneinander ab, fallen alle quadratischen Terme weg und es
bleibt eine lineare Gleichung für $\sigma_m$ allein:
$2(\sigma_2 - \sigma_1)\,\sigma_m = \sigma_2^2 + \tau_2^2 - \sigma_1^2 - \tau_1^2$. Bei mehreren
Punkten ohne Winkel zieht das Programm die Gleichung des ersten Punktes von jeder anderen ab. Jede
Differenz ist linear und kommt zu den linearen Zeilen. Quadratisch bleibt damit nur noch eine
Gleichung, die des ersten Punktes selbst: Die Information der anderen Punkte steckt vollständig in
den Differenzen.

\textbf{Rang.} Die linearen Zeilen löst das Programm mit Gauß-Jordan und Spaltenpivotsuche und
zählt dabei, wie viele Pivotelemente es findet (Toleranz relativ zum größten Eintrag). Bleibt in
einer Zeile ohne Pivot rechts etwas übrig, widersprechen sich die Angaben.

\begin{itemize}
  \item \textbf{Rang 3:} Der Zustand ist eindeutig. Ein Punkt ohne Winkel muss dann nur noch auf
        dem Kreis liegen (Probe).
  \item \textbf{Rang 2 und ein Punkt ohne Winkel:} Die Lösungen der linearen Zeilen bilden eine
        Gerade $\vec x = \vec x_0 + t\,\vec d$ im $(\sigma_m, R_c, R_s)$-Raum. In die quadratische
        Gleichung eingesetzt, wird daraus $a t^2 + b t + c = 0$ mit keiner, einer oder zwei Lösungen.
        Genau daher kommt \emph{zwei Zustände möglich}.
  \item \textbf{Weniger:} Es fehlen Angaben.
\end{itemize}

Schnitte ohne Winkel mit nur $\sigma$ oder nur $\tau$ geben keine Gleichung, aber eine
Ungleichung: $|\sigma - \sigma_m| \le R$ bzw.\ $|\tau| \le R$. Damit sortiert das Programm von zwei
möglichen Zuständen den aus, auf dessen Kreis der Wert nicht passt.

\textbf{Die Spannungsscheibe} ist das Vieleck, dessen Kanten alle einen Kreis mit Radius $r$
berühren. Jeder Schnitt $i$ mit dem Normaleneinheitsvektor $\vec n_i$ gibt eine Kante: die Gerade
$\vec n_i \cdot \vec x = r$, senkrecht zu $\vec n_i$ im Abstand $r$ vom Mittelpunkt. Zur Scheibe
gehört alles auf der Seite des Mittelpunktes, also die Halbebene $\vec n_i \cdot \vec x \le r$. Die
Scheibe ist die Menge der Punkte, die in allen diesen Halbebenen zugleich liegen. Liegen zwei
benachbarte Normalen $160^\circ$ oder mehr auseinander, wäre die Scheibe
offen (oder eine sehr spitze Nadel). Dann ergänzt das Programm den $x$- oder $y$-Schnitt, dessen
Normale am nächsten an der Mitte der größten Lücke liegt, und das so oft, bis die Lücke kleiner ist.
 eingebunden.

\part{Die Mathematik hinter den Features}\label{sec:Mathe}

Vorab will ich erwähnen, dass Teile dieses Abschnitts nicht meine eigene Arbeit sind. Gerade wenn
es darum geht, das Programm effizient und robust zu schreiben, oder wenn sich das Skript Methoden der
weiterführenden Mathematik bzw.\ Numerik bedient, stammen Ideen von KI, z.\,B.\ die
$LDL^T$-Zerlegung (ich bin selbst erst im 2.~Semester, deswegen halten sich meine mathematischen
Kenntnisse noch in Grenzen). Das gilt übrigens auch für große Teile der Umsetzung. Aus diesem Grund
halte ich die Erklärungen zu diesen Aspekten hier etwas kürzer.

Nichtsdestotrotz ist dieser Teil für alle, die wissen wollen, \emph{warum} das Programm so rechnet, wie es rechnet.
Zum Benutzen braucht ihr ihn nicht. Grundlegende Sachen wie Hebelarme, Steiner und Gleichgewicht seien hier nicht weiter erläutert. Hier
werden die Stellen behandelt, an denen das Programm mehr macht als die Vorlesung oder es anders macht,
weil ein Rechner eben nicht wie ein Mensch rechnet. Ich habe versucht, zu jedem Punkt dazuzuschreiben,
was er für euch in der Klausur bedeutet.

Ein paar Konventionen, die überall gelten:

\begin{itemize}
  \item Im \textbf{Tragwerksskript} zeigt $x$ nach \emph{rechts} und $y$ nach \emph{unten}. Drehwinkel,
        Momente und Lastwinkel zählen trotzdem so, wie ihr sie seht: gegen den Uhrzeigersinn
        positiv.
  \item Jeder Stab hat seine eigene Koordinate $x \in [0, L]$ vom Anfangs- zum Endknoten, die
        Tangente $\vec t = (\Delta x, \Delta y)/L$ und die Normale $\vec n = (-\Delta y, \Delta x)/L$.
        Die Schnittgrößen folgen der TM1-Konvention mit gestrichelter Faser auf der Seite von
        $\vec n$.
  \item Im \textbf{Querschnittsskript} gilt die Konvention aus dem Gross (TM2):
        $I_y = \int z^2\,dA$, $I_z = \int y^2\,dA$ und $I_{yz} = -\int yz\,dA$, also mit Minus.
\end{itemize}

%---------------------------------------------------------------------
\section{Das Weggrößenverfahren im Programm}\label{sec:mathe-wgv}

Alles, was ihr im Tragwerksskript seht (Verläufe, Auflagerkräfte, Verformungen,
Explosionszeichnung), kommt aus einer einzigen Rechnung: dem Weggrößenverfahren, also genau dem
allgemeinen WGV aus \textbf{Statik}. Unbekannt sind die Verschiebungen der Knoten, nicht die
Kräfte.

\textbf{Unbekannte:} Jeder Knoten hat drei Freiheitsgrade, $u_\xi$, $u_\eta$ und $\varphi$. Bei
einem ungedrehten Lager ist $\xi = x$ und $\eta = y$, bei einem gedrehten Lager zeigen sie in die
Lagerrichtungen. Ein Lager \emph{streicht} einfach den gehaltenen Freiheitsgrad.

\textbf{Stabsteifigkeit:} Für jeden Stab steht die bekannte $6 \times 6$-Matrix aus $EA/L$,
$12EI/L^3$, $6EI/L^2$, $4EI/L$ und $2EI/L$. Sie verknüpft die sechs Endverschiebungen
$(u_a, w_a, \varphi_a, u_b, w_b, \varphi_b)$ im lokalen System des Stabes mit den sechs
Endkräften. Zwei Drehungen bringen sie ins globale System: die Stabrichtung und (bei einem
gedrehten Lager) die Lagerrichtung,
\[
  K_e = T_K\,T^T\,k\,T\,T_K^T .
\]
So ist ein gedrehtes Lager kein Sonderfall, sondern nur eine andere Basis am Knoten. Matrix,
Transformation und Drehrichtungen sind dieselben wie in der WGV-Formelsammlung aus Statik, die
FEM-Formelsammlung hat dieselben Elemente, zählt die Querrichtung aber nach oben. Die
Unterschiede stehen in Abschnitt~\ref{sec:wgv-vorzeichen}.

\textbf{Lasten:} Knotenlasten gehen direkt in den Lastvektor (vorher um den Lastwinkel gedreht).
Streckenlasten gehen als Festeinspannkräfte ein. Das sind die Endkräfte, die ein
beidseitig eingespannter Stab unter der Last hätte, z.\,B.\ $\mp qL/2$ und $\pm qL^2/12$ bei
Gleichlast (die Volleinspannmomente aus Tafel~A9 der Formelsammlung, mit denselben Vorzeichen).
Im Gesamtsystem stehen sie mit umgekehrtem Vorzeichen, also als Knotenersatzlasten. Für Trapezlasten stehen die entsprechenden Formeln im Code, Lastfunktionen wie
\code{x\^{}2} integriert das Programm mit den Hermite-Formfunktionen
$N_1 = 1 - 3\xi^2 + 2\xi^3$, $N_2 = L(\xi - 2\xi^2 + \xi^3)$ usw.\ (arbeitsäquivalente
Knotenlasten, Abschnitt~\ref{sec:mathe-festeinspann}). Temperatur wirkt wie eine Last $EA\,\varepsilon_T$ bzw.\ $EI\,\kappa_T$ an den
Stabenden. Globale Linienlasten $p_X(x)$, $p_Y(x)$ werden vorher in Quer- und Längsanteil zerlegt,
bei einer projizierten Last ist $x$ die globale Koordinate ab dem Stabanfang, $x_{global} = f\,x$
mit dem Projektionsfaktor $f = |\Delta y|/L$ bzw.\ $|\Delta x|/L$, und je Stablänge wirkt
$f\,p(f\,x)$.

\textbf{Lösen:} Das Gesamtsystem $K\,u = F$ wird mit einer $LDL^T$-Zerlegung gelöst. Wird dabei
ein Diagonalelement $D_j$ (fast) null, fehlt dem System in dieser Richtung jede Steifigkeit. Es
ist ein Mechanismus, und ihr bekommt \emph{System ist kinematisch!}.

\textbf{Zurückrechnen:} Aus den Knotenverschiebungen folgen die Endkräfte jedes Stabes,
$s = k\,v + p$ (Steifigkeit mal Endverschiebung plus Festeinspannkraft), und aus dem Anfangswert
durch Integrieren der Lasten die Verläufe, z.\,B.\
$M(x) = M_a + Q_a x - m\,x - \int_0^x q(t)\,(x-t)\,dt$. Das ist exakt, nicht genähert. Deshalb
braucht jeder Bereich zwischen Lasten, Lagern und Gelenken nur einen einzigen Stab.

\textbf{Auflagerkräfte:} Die Summe der Stabendkräfte an einem Knoten minus der Knotenlast ist die
Kraft, die Lager und Federn auf den Knoten ausüben. Davon zeigt das Programm
$R_\xi = \vec R\cdot\vec\xi$ und $R_\eta = -\vec R\cdot\vec\eta$, deshalb zeigt $R_\eta$ nach oben,
während die Knotenverschiebung \code{ü\_y} auf derselben Menüseite nach unten positiv ist.

\subsection{Festeinspannkräfte beliebiger Lasten}\label{sec:mathe-festeinspann}

Für das Gesamtsystem braucht das Programm bzw. das WGV von einer Streckenlast nur, was sie an den Stabenden
bewirkt, die Festeinspannkräfte. Für Gleich- und Trapezlasten stehen diese In Tafeln oder in der FOS, für eine
Lastfunktion wie $q(x) = x^2$ oder $q_0 \sin(\pi x/L)$ nicht. Das Programm rechnet sie deshalb für
jede Lastfunktion mit den Hermite-Funktionen aus, und zwar exakt, nicht genähert.

\textbf{Die Hermite-Funktionen.} Im lokalen System läuft $x$ von $0$ bis $L$, $\xi = x/L$, $w$ zeigt
zur gestrichelten Faser und $\varphi = -w'$ gegen den Uhrzeigersinn. Die vier Funktionen sind
\begin{align*}
  N_1 &= 1 - 3\xi^2 + 2\xi^3       && w(0) = 1, \\
  N_2 &= L\,(\xi - 2\xi^2 + \xi^3) && w'(0) = 1, \text{ also } \varphi_a = -1, \\
  N_3 &= 3\xi^2 - 2\xi^3           && w(L) = 1, \\
  N_4 &= L\,(-\xi^2 + \xi^3)       && w'(L) = 1, \text{ also } \varphi_b = -1,
\end{align*}
und die jeweils drei anderen Endwerte sind null. Jede ist eine kubische Parabel, löst also
$EI\,w'''' = 0$ exakt und stellt somit eine homogene Lösung der DGL dar.
Gleichzeitig erfüllt sie die vier von ihr geforderten Randbedingungen. Damit ist sie genau die Biegelinie des
beidseitig eingespannten Stabes, dem man diese eine Einheitsverformung aufzwingt.

Das WGV zählt die vier Endverschiebungen eines Stabes aber als $v = (w_a, \varphi_a, w_b, \varphi_b)$,
mit $\varphi$ gegen den Uhrzeigersinn positiv. Die Einheitsverformung zum Freiheitsgrad $j$ in
dieser Zählung heißt im Folgenden $S_j$ (S für Formfunktion, engl.\ \emph{shape}):
\[
  S_1 = N_1, \qquad S_2 = -N_2, \qquad S_3 = N_3, \qquad S_4 = -N_4 .
\]
Bei den Drehungen kehrt sich das Vorzeichen um, weil $N_2$ und $N_4$ zu $\varphi = -1$ gehören.

\textbf{Beweisidee mit dem Satz von Betti.} Die folgende Herleitung soll die Idee zeigen und
ist kein formal vollständiger Beweis. So nennen wir alle $F_i$ Kräfte und alle Verschiebungen Wege,
obwohl bei den Drehungen Momente und Winkel gemeint sind, und die Einheitsverschiebung \enquote{$1$}
hat eigentlich eine Einheit (m bzw.\ rad). Das formal sauber mitzuschleppen ändert am Ergebnis
nichts, macht die Rechnung aber deutlich schwerer zu lesen.

Die Festeinspannkräfte $F_i$ sind die Kräfte und
Momente, die die Einspannungen auf den belasteten Stab ausüben, jeweils in positiver Richtung des
Freiheitsgrades (das $p$ in $s = k\,v + p$). Man schneidet den Stab frei. Die Einspannungen sind
dann nur noch Kräfte an den Stabenden, und man hat \emph{ein} System mit zwei Lastgruppen:

\begin{itemize}
  \item Lastgruppe~1: die Last $q(x)$ und die Endkräfte $F_j$, die die Enden festhalten (gesucht).
        Die Biegelinie ist $w_1$ mit $w_1 = w_1' = 0$ an beiden Enden, alle Endverschiebungen sind
        also $v_j^{(1)} = 0$.
  \item Lastgruppe~2: nur die Endkräfte $k_{ji}$, die den Freiheitsgrad $i$ um $1$ verschieben und
        alle anderen festhalten. Die Biegelinie ist $S_i(x)$, die Endverschiebungen sind
        $v_j^{(2)} = \delta_{ij}$ (Kronecker-Delta: $1$ für $j = i$, sonst $0$).
\end{itemize}

Beide Gruppen sind für sich im Gleichgewicht, deshalb stört es nicht, dass der freigeschnittene
Stab nicht gelagert ist. Nach Betti leistet Lastgruppe~1 auf den Wegen von Lastgruppe~2 dieselbe
Arbeit wie Lastgruppe~2 auf den Wegen von Lastgruppe~1:
\[
  \int_0^L q(x)\,S_i(x)\,dx + F_j\,v_j^{(2)} = k_{ji}\,v_j^{(1)} .
\]
Über das doppelte $j$ wird auf beiden Seiten summiert (Einsteinsche Summenkonvention). Setzt man
die Verschiebungen ein, bleibt links $F_j\,\delta_{ij} = F_i$ und rechts null. Die innere Arbeit
fehlt dabei nicht, sie ist der Grund für die Gleichheit:
\[
  W_a^{1 \to 2} = \int_0^L M_1\,\kappa_2\,dx = \int_0^L \frac{M_1 M_2}{EI}\,dx
  = \int_0^L M_2\,\kappa_1\,dx = W_a^{2 \to 1} .
\]
Das erste Gleichheitszeichen ist der Arbeitssatz (äußere gleich innere Arbeit) mit zwei
verschiedenen Zuständen: die Kräfte und Momente $M_1$ aus Lastgruppe~1, die im Gleichgewicht sind,
und die Verformung $\kappa_2$ aus Lastgruppe~2, die verträglich ist. Einer der beiden Zustände
spielt dabei die Rolle des virtuellen. Im PvV aus \textbf{TM1} ist es die Verschiebung (hier die
von Lastgruppe~2), beim Prinzip der virtuellen Kräfte aus \textbf{Statik} wäre es die Kraftgruppe.
Der Satz selbst gilt in beiden Lesarten, solange der eine Zustand im Gleichgewicht und der andere
verträglich ist. Das letzte Gleichheitszeichen ist derselbe Satz mit vertauschten Rollen: Kräfte
aus Lastgruppe~2, Verformung aus Lastgruppe~1. In der Mitte steht nur $\kappa = M/EI$. Betti
schreibt die beiden Enden dieser Kette hin. Die Kette enthält nur die Biegearbeit. Das Programm rechnet nach Euler-Bernoulli, der Stab ist
also schubstarr ($\gamma = 0$), und die Querkraft leistet keine innere Arbeit. Nach Timoschenko käme
in der Mitte noch $\int_0^L Q_1 Q_2/(GA_s)\,dx$ dazu. Betti gilt dann weiterhin, nur wären die $S_i$
andere Kurven, die auch von $GA_s$ abhängen.

Dass die gemischte innere Arbeit hier null ist, obwohl die Schnittgrößen aus Lastgruppe~2 auf
der Biegung aus Lastgruppe~1 lokal sehr wohl Arbeit leisten, sieht man auch direkt: $M_2$ ist ohne
Streckenlast linear, also $M_2'' = 0$, und zweimal partiell integriert bleibt
\[
  \int_0^L M_2\,\kappa_1\,dx = -\int_0^L M_2\,w_1''\,dx
  = \Bigl[-M_2\,w_1' + M_2'\,w_1\Bigr]_0^L - \int_0^L M_2''\,w_1\,dx = 0 ,
\]
weil $w_1$ und $w_1'$ an beiden Enden verschwinden. Also gilt für jede Lastfunktion
\[
  F_i = -\int_0^L q(x)\,S_i(x)\,dx .
\]
Mit $S_1 = N_1$, $S_2 = -N_2$, $S_3 = N_3$ und $S_4 = -N_4$ ergibt das genau die Formeln im
Programm:
\begin{align*}
  Q_a &= -\int_0^L q\,N_1\,dx, & M_a &= \int_0^L q\,N_2\,dx, \\
  Q_b &= -\int_0^L q\,N_3\,dx, & M_b &= \int_0^L q\,N_4\,dx.
\end{align*}

\textbf{Anschaulich} ist $S_i(x)$ die Einflusslinie der Festeinspannkraft $i$ (Satz von
Müller-Breslau): Eine Einzelkraft $P$ an der Stelle $a$ erzeugt $F_i = -P\,S_i(a)$, und das Integral
summiert nur die vielen kleinen Einzellasten $q\,dx$ auf.

\textbf{Kontrolle mit Gleichlast:} $\int_0^L N_1\,dx = L/2$, $\int_0^L N_2\,dx = L^2/12$ und
$\int_0^L N_4\,dx = -L^2/12$. Damit ist $Q_a = Q_b = -qL/2$ (die Last zeigt nach unten, die
Einspannung drückt dagegen) und $M_a = qL^2/12$, $M_b = -qL^2/12$, genau die Tafelwerte. Auch das
Vorzeichen von $M_a$ passt: Ohne Einspannung würde sich das linke Ende im Uhrzeigersinn drehen, die
Einspannung hält gegen den Uhrzeigersinn dagegen.

\textbf{Die anderen Lastarten.}
\begin{itemize}
  \item Ein \textbf{Streckenmoment} $m(x)$ (gegen den Uhrzeigersinn positiv) leistet Arbeit auf der
        Verdrehung, $\int m\,\varphi\,dx = -\int m\,w'\,dx$. Mit Betti wird daraus
        $F_i = \int_0^L m\,S_i'\,dx$, also $Q_a = \int m\,N_1'\,dx$, $M_a = -\int m\,N_2'\,dx$ und
        entsprechend am Ende $b$. Bei konstantem $m$ ist $\int N_1'\,dx = N_1(L) - N_1(0) = -1$ und
        $\int N_2'\,dx = 0$: $Q_a = -m$, $Q_b = +m$ und keine Endmomente. Das ist das Kräftepaar aus
        Abschnitt~\ref{sec:wgv-vorzeichen} und der Grund, warum das Programm vom $m\,l/2$ der
        WGV-Tafel abweicht.
  \item Für eine \textbf{Längslast} $n(x)$ sind die Formfunktionen linear, $1 - \xi$ und $\xi$ (sie
        lösen $EA\,u'' = 0$ exakt), und es gilt $N_a = -\int n\,(1-\xi)\,dx$, $N_b = -\int n\,\xi\,dx$.
  \item \textbf{Temperatur} ist keine Last, sondern eine eingeprägte Verzerrung. Die
        Festeinspannkräfte stehen direkt da: $N_a = EA\,\varepsilon_T$, $N_b = -EA\,\varepsilon_T$,
        $M_a = EI\,\kappa_T$, $M_b = -EI\,\kappa_T$.
\end{itemize}

\textbf{Zwei Wege im Code.} Für konstante Lasten und Trapezlasten
$q = p_a(1-\xi) + p_b\,\xi$ sind die Integrale schon ausgerechnet:
\begin{align*}
  Q_a &= -\tfrac{L}{60}\,(21p_a + 9p_b),   & M_a &= \tfrac{L^2}{60}\,(3p_a + 2p_b), \\
  Q_b &= -\tfrac{L}{60}\,(9p_a + 21p_b),   & M_b &= -\tfrac{L^2}{60}\,(2p_a + 3p_b).
\end{align*}
Hängt die Last von $x$ ab (auch bei $m(x)$, $n(x)$ und den globalen Linienlasten), rechnet das CAS
genau die Integrale oben aus. Ein Beispiel ist $q(x) = q_0 \sin(\pi x/L)$. Für die Querkraft
braucht man $\int_0^L q\,N_1\,dx$, und das geht ohne Rechnung in drei Schritten:
\begin{enumerate}
  \item $N_1 + N_3 = 1$: Schiebt man beide Stabenden um $1$ nach unten, verschiebt sich der ganze
        Stab um $1$. Also ist $\int_0^L q\,N_1\,dx + \int_0^L q\,N_3\,dx = \int_0^L q\,dx$.
  \item $N_3$ ist $N_1$ gespiegelt, $N_3(\xi) = N_1(1-\xi)$. Ist die Last symmetrisch zur
        Stabmitte, so wie der Sinus, sind beide Integrale deshalb gleich groß.
  \item Also ist jedes die Hälfte: $\int_0^L q\,N_1\,dx = \frac12 \int_0^L q\,dx
        = \frac12 \cdot \frac{2q_0L}{\pi} = \frac{q_0L}{\pi}$, und $Q_a = -q_0L/\pi$.
\end{enumerate}
Anschaulich: Bei symmetrischer Last trägt jede Einspannung die halbe Last. Für das Moment
braucht man $\int_0^1 \sin(\pi\xi)\,(\xi - 2\xi^2 + \xi^3)\,d\xi = 2/\pi^3$, also
$M_a = 2q_0L^2/\pi^3 \approx 0{,}0645\,q_0L^2$.

\textbf{Warum das exakt ist.} Betti gilt für jede Last, und die $S_i$ sind exakte Lösungen der
homogenen Gleichung. Also sind die Festeinspannkräfte exakt, damit auch die Knotenverschiebungen aus
$K\,u = F$ und die Stabendkräfte $s = k\,v + p$. Nicht exakt wäre nur die Hermite-Interpolation
$\sum_j v_j S_j(x)$ im Feld, denn die kennt die Last nicht. Die benutzt das Programm aber nie als
Ergebnis (siehe unten).

\textbf{Zwei Zerlegungen derselben Biegelinie.} Die Biegelinie eines Stabes ist immer
Partikulärlösung plus homogene Lösung. Welche Partikulärlösung man nimmt, ist aber nicht
festgelegt: Zwei Partikulärlösungen unterscheiden sich immer um eine kubische Parabel, also um eine
homogene Lösung. Das Programm benutzt zwei davon, und beide beschreiben dieselbe Linie.

\begin{itemize}
  \item \emph{Die Anfangswert-Lösung} ist das Vierfachintegral der Querlast,
        \[
          EI\,w_p(x) = \int_0^x q(t)\,\frac{(x-t)^3}{6}\,dt ,
        \]
        mit $w_p = w_p' = w_p'' = w_p''' = 0$ bei $x = 0$, so wie ihr sie in \textbf{TM2} aufstellt.
        Dazu kommt die homogene Lösung, deren vier Konstanten die Randbedingungen erfüllen. Das
        Rand-LGS schreibt sie als $c_1 x^3/6 + c_2 x^2/2 + c_3 x + c_4$ und holt die Konstanten aus
        allen Rand- und Übergangsbedingungen (Abschnitt~\ref{sec:mathe-randlgs}). Der alte Export
        \code{v\_EI\_stab} schreibt sie mit den Hermite-Funktionen, und die Koeffizienten sind die
        Endverschiebungen aus dem WGV:
        \[
          EI\,w = EI\,w_p + EI w_a\,N_1 + EI w_a'\,N_2 + \bigl(EI w_b - EI w_p(L)\bigr)\,N_3
                  + \bigl(EI w_b' - EI w_p'(L)\bigr)\,N_4 .
        \]
        $w_p$ bringt am Ende $b$ schon eigene Werte mit, die ziehen die letzten beiden Koeffizienten
        wieder ab, damit die Summe dort genau die Endverschiebung hat. Die Koeffizienten sind
        Durchbiegung und Neigung, so wie ihr die Randbedingungen in \textbf{TM2} aufstellt. In der
        Zählung des WGV ist $w_a = v_1$, $w_a' = -\varphi_a = -v_2$, $w_b = v_3$ und $w_b' = -v_4$.
  \item \emph{Die Lösung des eingespannten Stabes} ist
        \[
          w_0(x) = w_p(x) - w_p(L)\,N_3(x) - w_p'(L)\,N_4(x) .
        \]
        Auch sie löst $EI\,w'''' = q$, hat aber an \emph{beiden} Enden $w_0 = w_0' = 0$: Das ist die
        Biegelinie des beidseitig eingespannten Stabes unter der Last. Mit ihr zerfällt die Linie in
        \[
          w(x) = w_0(x) + \sum_j v_j\,S_j(x) ,
        \]
        und die Koeffizienten des homogenen Teils sind genau die Endverschiebungen $v_j$, also die
        Unbekannten des WGV. Diese Zerlegung ist die Idee hinter $s = k\,v + p$: Der homogene Teil
        $\sum_j v_j S_j$ liefert die Endkräfte $k\,v$. $w_0$ verschiebt die Stabenden nicht
        ($w_0 = w_0' = 0$ an beiden Enden), stört also die Verträglichkeit an den Knoten nicht und
        liefert nur Endkräfte, und das sind genau die Festeinspannkräfte $p$. Mehr braucht das WGV
        von der Last nicht.
\end{itemize}

Bei Gleichlast ist $EI\,w_p = qx^4/24$ mit $EI\,w_p(L) = qL^4/24$ und $EI\,w_p'(L) = qL^3/6$, und
daraus wird
\[
  EI\,w_0 = \frac{qL^4}{24}\bigl(\xi^4 - (3\xi^2 - 2\xi^3) - 4(\xi^3 - \xi^2)\bigr)
          = \frac{q\,x^2 (L-x)^2}{24},
\]
die bekannte Biegelinie des eingespannten Balkens. Ihr Moment am linken Ende ist
$-EI\,w_0''(0) = -qL^2/12$ (Zug oben), und die Einspannung hält mit $M_a = +qL^2/12$ gegen den
Uhrzeigersinn dagegen. Der Satz von Betti erspart es, $w_0$ überhaupt auszurechnen: Die
Festeinspannkräfte kommen direkt aus den Integralen über $q\,S_i$.

Die fertige Biegelinie ist in beiden Zerlegungen dieselbe. Das sieht man, wenn man die Formel des
alten Exports umsortiert. Mit $w_a'\,N_2 = -v_2\,N_2 = v_2\,S_2$ und genauso am Ende $b$ wird aus ihr
\[
  EI\,w = \underbrace{EI\,\bigl(w_p - w_p(L)\,N_3 - w_p'(L)\,N_4\bigr)}_{EI\,w_0}
          + EI \sum_j v_j\,S_j ,
\]
also Buchstabe für Buchstabe die zweite Zerlegung, nur anders gruppiert. Ausgerechnet wird $w_0$
dabei nie. Für die Knotengleichungen braucht das WGV nur seine Endkräfte $p$, und die kommen nach
Betti direkt aus den Integralen über $q\,S_i$. Die Verläufe rechnet das Programm danach aus den
exakten Endkräften (Gleichgewicht) bzw.\ mit der Anfangswert-Lösung plus Hermite-Anteil, wie im
ersten Punkt. Die zweite Zerlegung dient hier also nur der Erklärung, warum $s = k\,v + p$ exakt ist.

\subsection{Gelenke und die Knotenverdrehung}\label{sec:mathe-gelenke}

Hier kommt die Warnung aus Abschnitt~\ref{sec:Eingabe_Knoten} her. Ein \textbf{M-Gelenk} wird nicht
herauskondensiert, sondern bekommt einen eigenen Freiheitsgrad: Das Stabende am Gelenk hat seine
eigene Verdrehung. Zu diesem Freiheitsgrad gehört eine eigene Gleichgewichtsgleichung, und an ihm
greift kein äußeres Moment an. Somit ist das Gelenkmoment null. Für die angezeigte Knotenverdrehung
heißt das:

\begin{itemize}
  \item \textbf{Angelenk} (Gelenk am Stabende, andere Stäbe biegesteif am Knoten): Die
        Knotenverdrehung ist die Verdrehung der \emph{biegesteif} angeschlossenen Stäbe. Der Stab
        mit dem Gelenk dreht sich anders, seine Endverdrehung steht im Stabmenü auf Seite~7
        ($\varphi_a$, $\varphi_b$), der Knick im Gelenk ist die Differenz.
  \item \textbf{Vollgelenk} (Knoten gelenkig): Der Knoten hat \emph{gar keine} eigene Verdrehung,
        jedes Stabende dreht sich für sich. Die $0$ im Knotenmenü ist nur ein Platzhalter, die
        echten Verdrehungen stehen wieder auf Seite~7 der Stäbe.
  \item \textbf{Nur Angelenke an einem Knoten}, aber kein Vollgelenk und kein Lager, das die
        Drehung hält: Dann hat der Knoten eine Verdrehung, an der nichts hängt. Das Programm
        meldet \emph{kinematisch}. Deshalb der Stolperstein im Anhang: In diesem Fall braucht ihr
        das Vollgelenk.
\end{itemize}

\textbf{N- und Q-Gelenke} werden dagegen herauskondensiert (siehe Dallmann Baustatik 3): Im Stab wird der entsprechende Anteil
der Steifigkeit gestrichen, die Festeinspannkräfte werden statisch gleichwertig
umverteilt, und nach der Rechnung wird die Klaffung im Gelenk wieder ausgerechnet. Die
Endverschiebungen auf Seite~7 enthalten sie. Was dabei gestrichen wird, betrifft immer nur die
freigegebene Richtung:

\begin{itemize}
  \item \textbf{N-Gelenk:} Die Dehnsteifigkeit fällt weg. Am Gelenk ist $N = 0$, im Stab bleibt
        nur der Anteil aus Längslasten. Die Biegung ist davon nicht betroffen.
  \item \textbf{Q-Gelenk:} Am Gelenk ist $Q = 0$, ohne Querlast also im ganzen Stab, und wegen
        $M' = Q$ bleibt von der Biegesteifigkeit nur ein konstantes Moment übrig. Die Normalkraft
        trägt der Stab ganz normal weiter.
\end{itemize}

%---------------------------------------------------------------------
\section{Bögen: Teilstäbe und Schnittgrößen über den Winkel}\label{sec:mathe-bogen}

\subsection{Geometrie und Teilstäbe}

Ein Bogen ist ein Kreisstück mit Radius $R$ über der Sehne $L$ vom Stabanfang zum Stabende. Sein
Öffnungswinkel ist $\alpha = 2\arcsin\frac{L}{2R}$, und jeder Punkt auf ihm hat einen Winkel
$\varphi \in [0, \alpha]$ ab dem Stabanfang. Für die Rechnung ersetzt das Programm den Bogen durch
$n$ gerade Teilstäbe (Obermenü Seite~3, Standard 20). Das erste und das letzte Teilstück liegen
\emph{tangential} am Bogen und sind so lang wie $R\alpha/n$, die übrigen Teilungspunkte liegen auf
dem Bogen. Die tangentialen Endstücke haben einen Zweck: Die Stabendkräfte des ersten Teilstücks
sind dann direkt Normal- und Querkraft \emph{am Bogen}, und von dort startet der Export.

Die Teilstäbe erben Steifigkeiten, Gelenke an den Bogenenden, Temperatur und die Lasten:

\begin{itemize}
  \item $q$ (radial) und $n$ (tangential) bekommt jeder Teilstab mit den Werten an seinen beiden
        Endpunkten, dazwischen linear. Eine Lastfunktion wertet das Programm an den
        Teilungspunkten aus, mit $x = \varphi$ im Winkelmodus des Rechners.
  \item Lasten je Bogenlänge ($p_X$, $p_Y$ ohne Projektion, $m$) werden je Teilstab mit
        $R\,\Delta\varphi / l$ skaliert ($l$ = Länge des Teilstabs). Damit trägt jeder Teilstab genau
        die Last seines Bogenstücks, und die Summe ist exakt, z.\,B.\ das Eigengewicht $gR\alpha$.
  \item Projizierte Lasten gehen je Grundriss- bzw.\ Höhenlänge des Teilstabs ein.
\end{itemize}

\textbf{Warum konstante radiale Lasten exakt sind.} Ein konstanter Druck $q$ senkrecht auf einen
beliebigen Weg von $A$ nach $B$ hat die Resultierende $q$ mal den um $90^\circ$ gedrehten Vektor
$B - A$. Sein Moment um einen Punkt $P$ ist $\frac q2\left(|B-P|^2 - |A-P|^2\right)$. Beides hängt nur
von den Endpunkten ab, nicht vom Weg dazwischen. Der Linienzug der Teilstäbe hat dieselben Endpunkte
wie der Bogen, also stimmen bei einem statisch bestimmten System die Lagerkräfte und die
Schnittgrößen an den Teilungspunkten auf dem Bogen exakt: Sie folgen dort allein aus dem
Gleichgewicht, und das sieht nur die Resultierenden. Wie biegsam die Stäbe sind, spielt dabei keine
Rolle. Bei einem statisch unbestimmten System hängt die Kraftverteilung dagegen auch von der
Steifigkeit ab, und der Linienzug ist etwas anders steif als der Bogen, dann gilt das nur bis auf
den Teilungsfehler (siehe \emph{Genauigkeit} unten).
Beim Eigengewicht stimmt dagegen nur die Summe exakt. Die Hebel weichen um Größenordnung
$(\alpha/n)^2$ ab, mit 20 Teilstäben bei einem Viertelkreis um etwa 0,04\,\%.

\subsection{Schnittgrößen über den Winkel}

Der Export schreibt die Schnittgrößen als Funktionen des Winkels $p$ ab Stabanfang. Dafür schneidet
er den Bogen bei $p$ und schreibt das Gleichgewicht des Stücks von $0$ bis $p$ auf. Am Anfang wirken
$N_A$, $Q_A$, $M_A$ aus der Rechnung (erstes, tangentiales Teilstück). Mit $\varepsilon = +1$, wenn
die gestrichelte Faser innen liegt ($R > 0$), und $\varepsilon = -1$, wenn sie außen liegt, gilt
ohne Last auf dem Bogen
\begin{align*}
  N(p) &= N_A\cos p + \varepsilon\,Q_A\sin p , \\
  Q(p) &= Q_A\cos p - \varepsilon\,N_A\sin p , \\
  M(p) &= M_A + Q_A\,R\sin p - \varepsilon\,N_A\,R\,(1-\cos p) .
\end{align*}
Die ersten beiden Zeilen drehen nur die Kraft vom Anfang in die Tangente bei $p$. In der dritten
stehen die Hebel: Der Schnittpunkt liegt um $R\sin p$ entlang der Anfangstangente und um
$R(1-\cos p)$ zum Mittelpunkt hin versetzt.
Hier sieht man auch, warum im Bogen Normalkraft und Biegung gekoppelt sind: $N_A$ hat am
gekrümmten Stab den Hebel $R(1-\cos p)$ und erzeugt ein Moment. Beim geraden Stab ($R \to \infty$
mit $R\,p = x$) geht dieser Hebel gegen $x^2/(2R) \to 0$, und $Q_A R\sin p$ wird zu $Q_A\,x$. Das
ist reines Gleichgewicht am unverformten Stab, also Theorie I.~Ordnung. Theorie II.~Ordnung hieße,
das Gleichgewicht am \emph{verformten} System aufzustellen (Knicken, Durchschlagen flacher Bögen),
das rechnet das Programm weder bei Bögen noch bei geraden Stäben.

Last auf dem Bogenstück kommt als Integral über den Lastwinkel $t$ von $0$ bis $p$ dazu (im Bogenmaß,
$ds = R\,dt$):
\begin{align*}
  \Delta N(p) &= -R\int_0^p \Bigl(\varepsilon\,q(t)\sin(p-t) + n(t)\cos(p-t)\Bigr)\,dt , \\
  \Delta Q(p) &= -R\int_0^p \Bigl(q(t)\cos(p-t) - \varepsilon\,n(t)\sin(p-t)\Bigr)\,dt , \\
  \Delta M(p) &= -R\int_0^p \Bigl(R\,q(t)\sin(p-t) - \varepsilon\,R\,n(t)\bigl(1-\cos(p-t)\bigr)
                + m(t)\Bigr)\,dt .
\end{align*}
Auch hier stehen die Hebel direkt da. Eine radiale Kraft beim Winkel $t$ geht durch den
Mittelpunkt, ihr Abstand zum Schnittpunkt bei $p$ ist $R\sin(p-t)$. Eine tangentiale Kraft liegt
auf der Tangente bei $t$, und der Schnittpunkt hat von dieser Tangente den Abstand $R(1-\cos(p-t))$.
Globale Linienlasten $p_X$, $p_Y$ gehen genauso ein, nur mit Hebeln aus den Koordinaten. Ihre
Richtung relativ zur Tangente ändert sich mit $p$.

\textbf{Winkelmodus.} Steht der Rechner auf DEG, laufen $p$ und $t$ in Grad. Dann ist
$ds = R\,\frac{\pi}{180}\,dt$, und genau diesen Faktor setzt das Programm vor die Integrale. Den
Modus erkennt es daran, ob $\cos(180) = -1$ ist. $\sin p$ und $\cos p$ rechnet der Rechner
ohnehin im eingestellten Modus, deshalb stimmen die Funktionen in jedem Modus.

\textbf{Genauigkeit.} Die Formeln sind das exakte Gleichgewicht am Kreisbogen. Ungenau sein können
nur die Startwerte $N_A$, $Q_A$, $M_A$, denn sie stammen aus der Rechnung mit Teilstäben. Bei einem
statisch bestimmten System mit Lasten, die die Teilstäbe exakt abbilden (Knotenlasten, konstante
radiale Last, projizierte Lasten), sind auch sie exakt. Sonst gehen die Abweichungen mit mehr
Teilstäben zurück: bei Eigengewicht, tangentialen Lasten und Lastfunktionen, und bei statisch
unbestimmten Systemen, in denen die Steifigkeit des Linienzugs etwas von der des Bogens abweicht.

%---------------------------------------------------------------------
\section{Starre Stäbe: Zwangsbedingungen und Nullraum}\label{sec:mathe-starr}

Warum nicht einfach $EA = 10^{15}$? Weil der Rechner mit etwa 16 gültigen Stellen rechnet. Steht
in derselben Matrix eine Feder mit $c = 1$ neben $10^{15}$, bleibt von der Feder nach der
Zerlegung praktisch nichts übrig. Das ist der Zahlenmüll aus Abschnitt~\ref{sec:starr}. Das
Programm macht es deshalb exakt.

\textbf{1. Zwangsbedingungen.} Ein starrer Anteil bekommt in der Steifigkeitsmatrix die Steifigkeit
\emph{null}. Stattdessen schreibt das Programm für jeden starren Stab auf, was \enquote{starr}
bedeutet, und zwar in den lokalen Endverschiebungen:
\begin{align*}
  \text{dehnsteif:}\quad & u_b - u_a = \varepsilon_T L, \\
  \text{biegesteif:}\quad & \varphi_a - \varphi_b = -\kappa_T L, \qquad
                            w_b - w_a + L\,\varphi_a = -\kappa_T \frac{L^2}{2} .
\end{align*}
Ohne Temperatur steht rechts null: gleiche Länge, gleiche Endverdrehung, Enden auf einer Geraden.
Alle Bedingungen zusammen, ins globale System gedreht, ergeben $C\,u = g$.

\textbf{2. Kern der Matrix (=Nullraum).} $C$ wird mit Gauß-Jordan auf reduzierte Stufenform gebracht. Jede Zeile hat
dann einen \emph{abhängigen} Freiheitsgrad (die Pivotspalte), der durch die übrigen, \emph{freien}
Freiheitsgrade ausgedrückt wird. Damit lässt sich jede erlaubte Verschiebung schreiben als
\[
  u = T\,u_r + u_0 , \qquad C\,T = 0 .
\]
Die Spalten von $T$ spannen den Nullraum bzw. den Kern von $C$ auf, also genau die Bewegungen, die die starren
Stäbe mitmachen. $u_r$ sind nur noch die freien Freiheitsgrade, $u_0$ ist der Temperaturanteil.

\textbf{3. Projizieren.} Eingesetzt in $K\,u = F$ und von links mit $T^T$ multipliziert:
\begin{equation} \label{eq:LGSred}
  \left(T^T K\,T\right) u_r = T^T\left(F - K\,u_0\right).
\end{equation}
Das ist nichts anderes als das Prinzip der virtuellen Verschiebungen: Als virtuelle
Verschiebungen sind nur noch solche zugelassen, die die Zwangsbedingungen erfüllen, also Spalten
von $T$. Das neue System ist kleiner als das alte und genauso gut konditioniert wie ein System
ohne starre Stäbe, weil nirgends eine riesige Zahl steht.

\textbf{4. Kräfte in den starren Stäben.} Nach Schritt~3 kennt man die Verschiebungen $u$, aber
noch keine Kraft in einem starren Stab, denn ihr Anteil in $K$ ist null. Setzt man $u$ in das
Gleichgewicht der Knoten ein, bleibt ein Rest
\[
  r = F - K\,u .
\]
Das sind die Knotenkräfte, die weder die elastischen Stäbe noch die Federn aufnehmen. Diesen Rest
müssen die starren Stäbe tragen.

Das widerspricht nicht $K\,u = F$. In $K$ fehlen die starren Stäbe, denn ihre Kraft lässt sich
nicht aus $u$ ausrechnen: Ein starrer Stab verformt sich bei jeder Kraft um null. Das vollständige
Gleichgewicht der Knoten lautet
\[
  K\,u + (\text{Kräfte der starren Stäbe}) = F ,
\]
und $r$ ist genau der fehlende Anteil. Das kennt man von den Lagern: Auch an einem gehaltenen
Freiheitsgrad bleibt ein Rest, und der ist die Auflagerkraft. Ein starrer Stab ist so etwas wie ein
inneres Lager, das statt einer absoluten eine relative Verschiebung festhält.

\textbf{Woher die $\lambda$ kommen.} Schritt~3 hat nicht $K\,u = F$ gelöst, sondern nur
\eqref{eq:LGSred}. Alles auf eine Seite gebracht und $u = u_0 + T\,u_r$ eingesetzt:
\begin{align*}
  0 &= T^T F - T^T K\,u_0 - T^T K\,T\,u_r \\
    &= T^T\bigl(F - K\,(u_0 + T\,u_r)\bigr) = T^T\,(F - K\,u) = T^T r .
\end{align*}
Die Zeilen von $T^T$ sind die Spalten von $T$, also eine Basis des Kerns von $C$. $T^T r = 0$ sagt
deshalb: $r$ steht auf jedem dieser Basisvektoren und damit auf dem ganzen Kern von $C$ senkrecht.

Daraus folgt, dass sich $r$ aus den Zeilen von $C$ zusammensetzen lässt. Das ist reine lineare
Algebra aus \textbf{HM1}. $C$ hat $m$ Zeilen (eine je Zwangsbedingung) und $n$ Spalten (eine je
Freiheitsgrad). Im $\mathbb R^n$ liegen damit zwei Unterräume:
\begin{itemize}
  \item $\operatorname{Kern} C$: alle $x$ mit $C\,x = 0$,
  \item $\operatorname{Bild} C^T$: alle Vektoren $C^T\lambda$, also alle Kombinationen der Zeilen
        von $C$.
\end{itemize}
Gezeigt wird $r \in \operatorname{Bild} C^T$, und zwar in vier Schritten.

\emph{(a) $r$ steht senkrecht auf $\operatorname{Kern} C$.} Das haben wir gerade gesehen.

\emph{(b) $\operatorname{Bild} C^T$ steht senkrecht auf $\operatorname{Kern} C$.} Für ein $x$ aus
$\operatorname{Kern} C$ und einen beliebigen Vektor $C^T\lambda$ gilt
\[
  (C^T\lambda)^T x = \lambda^T C\,x = \lambda^T 0 = 0 .
\]
Das erste Gleichheitszeichen ist nur die Rechenregel $(AB)^T = B^T A^T$. Insbesondere liegt außer
der Null kein Vektor in beiden Räumen, denn er müsste auf sich selbst senkrecht stehen.

\emph{(c) Zusammen füllen beide Räume den ganzen $\mathbb R^n$.} Nach der Dimensionsformel ist
$\dim\operatorname{Kern} C + \dim\operatorname{Bild} C = n$, und weil Zeilenrang gleich Spaltenrang
ist, gilt $\dim\operatorname{Bild} C^T = \dim\operatorname{Bild} C$. Also ist
$\dim\operatorname{Kern} C + \dim\operatorname{Bild} C^T = n$. Weil sich die beiden Räume nur in der
Null treffen, lässt sich jeder Vektor des $\mathbb R^n$ eindeutig als Summe aus einem Teil in
$\operatorname{Kern} C$ und einem Teil in $\operatorname{Bild} C^T$ schreiben.

\emph{(d) $r$ liegt in $\operatorname{Bild} C^T$.} Man zerlegt $r = r_K + r_B$ mit $r_K$ in
$\operatorname{Kern} C$ und $r_B$ in $\operatorname{Bild} C^T$. Nach (a) und (b) stehen $r$ und
$r_B$ senkrecht auf $\operatorname{Kern} C$, also auch $r_K = r - r_B$. $r_K$ liegt aber selbst in
$\operatorname{Kern} C$, steht also senkrecht auf sich selbst, und das kann nur der Nullvektor.
Damit ist $r = r_B$.

Ein Beispiel mit Zahlen: $30$ Freiheitsgrade und ein starrer Stab, der dehn- und biegesteif ist
($3$ unabhängige Zeilen). $\operatorname{Kern} C$ ist $27$-dimensional, das sind die $27$
Unbekannten $u_r$ in Schritt~3. $\operatorname{Bild} C^T$ ist $3$-dimensional, das sind die drei
$\lambda$.

Es gibt also Zahlen $\lambda_k$, eine je Zwangsbedingung, mit
\[
  r = C^T \lambda = \sum_k \lambda_k\,C_k^T ,
\]
wobei $C_k$ die $k$-te Zeile von $C$ ist. Weil $r$ nach dem vollständigen Gleichgewicht genau die
Kräfte der starren Stäbe sind, heißt das: Die starren Stäbe wirken mit $C^T\lambda$ auf die Knoten,
und jede Zwangsbedingung~$k$ trägt $\lambda_k\,C_k^T$ dazu bei. Das vollständige Gleichgewicht
lautet damit
\[
  K\,u + C^T\lambda = F .
\]
Mit $T^T$ multipliziert fällt $\lambda$ wegen $T^T C^T = (C\,T)^T = 0$ heraus, das ist Schritt~3.
Die übrigen Gleichungen erfüllt Schritt~4 mit der passenden Wahl von $\lambda$. Die $\lambda$ sind
übrigens die Lagrange-Multiplikatoren, die man bekommt, wenn man das Minimum der potentiellen
Energie $\frac12\,u^T K\,u - F^T u$ unter der Nebenbedingung $C\,u = g$ sucht.

\textbf{Mechanisch gelesen.} Das Skalarprodukt aus Kraft und Verschiebung ist die Arbeit,
\enquote{senkrecht} heißt also \enquote{keine Arbeit}. (a) sagt dann: Die Kräfte der starren Stäbe
leisten auf keiner Bewegung Arbeit, die die starren Stäbe erlauben. Genau das zeichnet eine
Zwangskraft aus. Eine reibungsfreie Führung leistet auf der Bewegung, die sie zulässt, auch keine
Arbeit, sie verhindert nur alle anderen. Physikalisch steckt diese Annahme schon in Schritt~3: Dort
sind als virtuelle Verschiebungen nur solche zugelassen, die die Zwangsbedingungen erfüllen. Alles
danach ist reine Mathematik.

Eindeutig sind die $\lambda$ nur, wenn die Zeilen von $C$ linear unabhängig sind, also
$\operatorname{Kern} C^T = \{0\}$ gilt. Dann gehört zu jeder Kraft aus $\operatorname{Bild} C^T$
genau ein $\lambda$. Ist $\operatorname{Kern} C^T$ größer, gibt es Kräfte in den starren Stäben,
die auf keinen Knoten wirken. Das sind die Eigenspannungszustände weiter unten.

\textbf{Warum die $\lambda$ Schnittgrößen sind.} Das folgt aus dem Gleichgewicht des einzelnen
starren Stabes, und zwar mit demselben Werkzeug wie eben: Kern, Bild und Dimensionsformel, diesmal
für einen Stab allein. Im lokalen System hat ein Stab die sechs Endverschiebungen
$v = (u_a, w_a, \varphi_a, u_b, w_b, \varphi_b)$ und sechs Endkräfte $s$ in denselben Richtungen.
Die drei Zeilen eines biege- und dehnsteifen Stabes bilden die $3 \times 6$-Matrix
\[
  C_\ell = \begin{pmatrix} -1 & 0 & 0 & 1 & 0 & 0 \\ 0 & 0 & 1 & 0 & 0 & -1 \\ 0 & -1 & L & 0 & 1 & 0 \end{pmatrix}
  \qquad
  \begin{matrix} \text{Längszeile: } u_b - u_a \\ \text{Drehzeile: } \varphi_a - \varphi_b \\
                 \text{Querzeile: } w_b - w_a + L\,\varphi_a \end{matrix}
\]

\emph{(a) Starr heißt: nur Starrkörperbewegungen.} Ein Stab kann sich auf drei Arten bewegen, ohne
sich zu verformen: längs verschieben, quer verschieben und um das Ende $a$ drehen. In
Endverschiebungen sind das die Spalten von
\[
  R = \begin{pmatrix} 1 & 0 & 0 & 1 & 0 & 0 \\ 0 & 1 & 0 & 0 & 1 & 0 \\ 0 & 0 & 1 & 0 & -L & 1 \end{pmatrix}^T
\]
(bei der Drehung um $\varphi = 1$ wandert das Ende $b$ wegen $\varphi = -w'$ um $w_b = -L$).
Einsetzen zeigt $C_\ell R = 0$: Starrkörperbewegungen verformen nichts, $\operatorname{Bild} R
\subseteq \operatorname{Kern} C_\ell$. Beide Matrizen haben Rang~$3$, nach der Dimensionsformel ist
also $\dim \operatorname{Kern} C_\ell = 6 - 3 = 3 = \dim \operatorname{Bild} R$, und damit
$\operatorname{Kern} C_\ell = \operatorname{Bild} R$. Die Zwangsbedingungen lassen genau die
Starrkörperbewegungen zu, nicht mehr und nicht weniger.

\emph{(b) Gleichgewicht heißt: keine Arbeit auf Starrkörperbewegungen.} Zuerst ein Stab ohne
Stablasten, die kommen in (d) dazu. Endkräfte $s$ sind genau dann im Gleichgewicht, wenn
\[
  R^T s = \begin{pmatrix} s_1 + s_4 \\ s_2 + s_5 \\ s_3 - L\,s_5 + s_6 \end{pmatrix} = 0 ,
\]
das sind Summe der Kräfte längs, Summe der Kräfte quer und Summe der Momente um $a$. Anders gelesen
ist das das PvV für einen starren Körper aus \textbf{TM1}: Auf jeder Starrkörperbewegung leisten die
Kräfte keine Arbeit. Die Endkräfte, die für sich im Gleichgewicht sind, bilden also den Raum
$\operatorname{Kern} R^T$.

\emph{(c) Dieser Raum wird genau von den Zeilen aufgespannt.} Wie beim Nullraum oben in zwei
Schritten: $R^T C_\ell^T = (C_\ell R)^T = 0$, also $\operatorname{Bild} C_\ell^T
\subseteq \operatorname{Kern} R^T$. Dimension: $\dim \operatorname{Bild} C_\ell^T = \operatorname{Rang}
C_\ell = 3$ und $\dim \operatorname{Kern} R^T = 6 - \operatorname{Rang} R = 3$. Also
\[
  \operatorname{Kern} R^T = \operatorname{Bild} C_\ell^T .
\]
Jede Gruppe von Endkräften im Gleichgewicht ist eine Kombination der drei Zeilen, und weil die
Zeilen linear unabhängig sind, auf genau eine Weise.

\emph{(d) Anwendung auf den starren Stab.} Auf den starren Stab wirken zwei Dinge: die Lasten auf
dem Stab ($q$, $n$, $m$) und die Endkräfte $s$, mit denen die Knoten auf seine Enden wirken. $s$
sind also nur die Endkräfte, die Lasten stehen getrennt. Zusammen sind beide im Gleichgewicht.
Die Festeinspannkräfte $p$ sind ebenfalls Endkräfte, die dieselben Lasten im Gleichgewicht
halten. Dass sie zu einem beidseitig eingespannten Stab gehören, spielt hier keine Rolle. Man
braucht nur \emph{irgendeinen} Satz Endkräfte, der die Lasten allein ausgleicht, und $p$ hat das
Programm aus dem WGV ohnehin. Zieht man die beiden Gleichgewichtsgruppen
\enquote{Lasten $+ s$} und \enquote{Lasten $+ p$} voneinander ab, fallen die Lasten heraus: $s - p$
ist eine Gruppe reiner Endkräfte im Gleichgewicht, also genau der Fall aus (b) ohne Stablasten.
Ausgeschrieben: Nach (b) gilt $R^T (s - p) = 0$, also liegt
$s - p$ in $\operatorname{Kern} R^T$, und das ist nach (c) $\operatorname{Bild} C_\ell^T$. Es gibt
also genau ein $\lambda_\ell \in \mathbb R^3$ mit $s - p = C_\ell^T \lambda_\ell$, umgestellt
\[
  s = p + C_\ell^T \lambda_\ell .
\]
Ins globale System gedreht und über alle starren Stäbe summiert sind das genau die Knotenkräfte
$C^T\lambda$ von oben, die $\lambda_\ell$ eines Stabes sind seine Einträge in $\lambda$. Ist ein Stab
nur dehnsteif, fehlen Dreh- und Querzeile, die Biegung liefert dann wie bei jedem Stab $k\,v$.
Ein N-Gelenk streicht die Längszeile, ein Q-Gelenk die Querzeile. Das hat zwei Gründe, die
zusammenpassen. Kinematisch kann sich das Stabende in der freigegebenen Richtung gegenüber dem
Knoten verschieben. Die Zeile würde aber den Knoten selbst festhalten, also mehr verlangen, als der
Stab tut. Und weil N- und Q-Gelenke herauskondensiert werden, hat das Stabende in dieser Richtung
keinen eigenen Freiheitsgrad, mit dem man die Zeile richtig schreiben könnte. Statisch ist die
zugehörige Schnittgröße am Gelenk null. Die Lasten trägt schon $p$, das für Stäbe mit N- oder
Q-Gelenk passend umverteilt ist, ein $\lambda$ für diese Zeile wäre also ohnehin null. Ein
M-Gelenk streicht dagegen nichts: Das Stabende hat dort einen eigenen Drehfreiheitsgrad
(Abschnitt~\ref{sec:mathe-gelenke}). Die Drehzeile verbindet die Verdrehung des Stabes mit diesem
Freiheitsgrad statt mit dem Knoten, und $M = 0$ am Gelenk folgt aus dessen Gleichung. $p$ darf man dabei unverändert nehmen, weil die
Festeinspannkräfte nach Abschnitt~\ref{sec:mathe-festeinspann} gar nicht von $EI$ abhängen.

\emph{(e) Was die drei Zahlen bedeuten.} Für jede Zeile liest man den Schnittgrößenverlauf des
unbelasteten Stabes direkt ab, mit $N(0) = -s_1$, $Q(0) = -s_2$, $M(0) = -s_3$ wie bei jedem Stab:
\begin{align*}
  \text{Längszeile:}\quad & s = \lambda_N\,(-1, 0, 0, 1, 0, 0)^T: && N(x) = \lambda_N, \\
  \text{Drehzeile:}\quad  & s = \lambda_M\,(0, 0, 1, 0, 0, -1)^T: && M(x) = -\lambda_M, \\
  \text{Querzeile:}\quad  & s = \lambda_Q\,(0, -1, L, 0, 1, 0)^T: && Q(x) = \lambda_Q, \quad
                                                                     M(x) = -\lambda_Q\,(L - x) .
\end{align*}
$\lambda_N$ ist also die Normalkraft, $\lambda_Q$ die Querkraft und $-\lambda_M$ das Moment am Ende
$b$. Zusammen $M(x) = -\lambda_M - \lambda_Q\,(L - x)$, dazu kommt der Anteil aus $p$. Mechanisch
passt das: Die Arbeit von $\lambda_k\,C_k^T$ auf den Endverschiebungen ist $\lambda_k\,(C_k\,v)$,
die Normalkraft arbeitet also auf der Längenänderung, das Moment auf der gegenseitigen Verdrehung
und die Querkraft darauf, wie weit das Ende $b$ von der Tangente am Ende $a$ abweicht. Das ist mit
dem Satz in Abschnitt~\ref{sec:starr} gemeint, die Schnittgrößen starrer Stäbe würden \enquote{aus
dem Gleichgewicht bestimmt}: Sie sind die $\lambda$, und die kommen allein aus dem Rest $r$, also aus
dem Gleichgewicht der Knoten, nicht aus einer Steifigkeit.

\textbf{Wie das Programm sie ausrechnet.} $r = C^T\lambda$ hat eine Gleichung je Freiheitsgrad, aber
nur eine Unbekannte je Zwangsbedingung, also meist mehr Gleichungen als Unbekannte. Weil $r$ nach
Schritt~3 eine Kombination der Zeilen von $C$ ist, geht es trotzdem exakt auf. Von links mit $C$
multipliziert wird daraus ein kleines quadratisches System
\[
  (C\,C^T)\,\lambda = C\,r ,
\]
das eindeutig lösbar ist, solange die Zeilen von $C$ linear unabhängig sind.

\textbf{Starre Teile, die sich gegenseitig halten.} Beispiele sind zwei dehnstarre Stäbe
nebeneinander zwischen denselben zwei Knoten oder ein biegestarrer Zweifeldträger (zwei starre
Stäbe über drei Lager). Dann sind einige Zeilen von $C$ linear abhängig. Das merkt schon die
Elimination in Schritt~2: Eine solche Zeile wird null, nachdem die früheren Pivotspalten abgezogen
sind. Das Programm protokolliert dabei, welche Kombination der ursprünglichen Zeilen null ergeben
hat:
\[
  \sum_k \mu_k\,C_k = 0 \qquad\Longleftrightarrow\qquad C^T\mu = 0 .
\]
Physikalisch ist $\mu$ ein Satz von Kräften in den starren Stäben, der auf keinen Knoten eine Kraft
ausübt, ein \emph{Eigenspannungszustand}. Bei den zwei Stäben nebeneinander ist das Zug~$1$ im
einen und Druck~$1$ im anderen.

Für die Kinematik ist eine abhängige Zeile harmlos, ohne Temperatur ist sie automatisch erfüllt. Für
das Gleichgewicht heißt sie: $r = C^T\lambda$ hat unendlich viele Lösungen
\[
  \lambda = \lambda_p + \sum_i t_i\,\mu_i .
\]
Das ist die Aufteilung, die man von jedem LGS kennt: eine partikuläre Lösung plus beliebige
Lösungen der homogenen Gleichung $C^T\mu = 0$. Und es ist genau die Aufteilung aus dem KGV in
\textbf{Statik}:
\begin{itemize}
  \item $\lambda_p$ rechnet das Programm mit $(C\,C^T)\,\lambda = C\,r$ nur über die unabhängigen
        Zeilen, die abhängigen bekommen $\lambda = 0$. Das ist, als würde man die überzähligen
        Zwangsbedingungen aufschneiden: ein statisch bestimmtes Hauptsystem mit seinem
        $0$-Zustand.
  \item Die $\mu_i$ sind die $1$-Zustände. Jedes $\mu_i$ hat an seiner abhängigen Zeile eine $1$,
        dazu kommen die Kräfte in den übrigen starren Stäben, die diese $1$ im Gleichgewicht halten.
  \item Die $t_i$ sind die statisch Überzähligen $X_i$.
\end{itemize}
Wie im KGV legt das Gleichgewicht die $X_i$ nicht fest. Dort kommen sie aus der Verträglichkeit,
und die braucht Nachgiebigkeiten. Genau die haben starre Stäbe nicht: Jede Wahl der $t_i$ ergibt
dieselben Verschiebungen, und die Starrkörpermechanik kann nicht sagen, welche stimmt. Das hängt
davon ab, wie starr die Stäbe \emph{im Verhältnis zueinander} wirklich sind.

\textbf{Der Grenzwert und die Ergänzungsenergie.} Das Programm trifft dieselbe Annahme wie der
Grenzwert im Rand-LGS: Alle starren Stäbe bekommen $EA = EI = 1/\varepsilon$ mit
$\varepsilon \to 0$. Dann verformen sie sich ein kleines bisschen. In den Größen, die die Zeilen
von $C$ messen, ist ihre elastische Verformung
\[
  C\,u - g = \varepsilon\,W\,\lambda ,
\]
wobei $W$ die Nachgiebigkeit der starren Stäbe mit $EA = EI = 1$ ist (blockweise je Stab, siehe
unten). Das ist das Elastizitätsgesetz der fast starren Stäbe.

Für das Folgende muss man zwei Dinge auseinanderhalten: $u$ sind die Knotenverschiebungen, $C\,u$
sind die Verformungen der starren Stäbe, gemessen in den Größen der Zwangsbedingungen
(Längenänderung, gegenseitige Verdrehung, Abweichung von der Tangente, siehe (e)). Die Stabkräfte
$\lambda$ arbeiten auf $C\,u$, ihre Knotenkräfte $C^T\lambda$ auf $u$, und beides ist dieselbe
Arbeit: $\lambda^T (C\,u) = (C^T\lambda)^T u$, das ist wieder nur $(AB)^T = B^T A^T$.

Jetzt kommt die Verträglichkeit, mit dem Arbeitssatz wie im KGV. Ein Eigenspannungszustand $\mu_i$
übt auf keinen Knoten eine Kraft aus, $C^T\mu_i = 0$. Seine Arbeit auf den wirklichen Verformungen
der starren Stäbe ist deshalb null:
\[
  \mu_i^T (C\,u) = (C^T\mu_i)^T u = 0^T u = 0 .
\]
Multipliziert man das Elastizitätsgesetz $C\,u - g = \varepsilon\,W\,\lambda$ von links mit
$\mu_i^T$, steht links genau diese Arbeit, und es bleibt
\[
  \underbrace{\mu_i^T C\,u}_{=\,0} - \mu_i^T g = \varepsilon\,\mu_i^T W\,\lambda
  \qquad\Longrightarrow\qquad
  \varepsilon\,\mu_i^T W\,\lambda + \mu_i^T g = 0 .
\]
Die Temperatur $g$ bleibt dabei stehen. Ohne Temperatur fällt der zweite Term weg, das $\varepsilon$ kürzt sich heraus, und mit
$\lambda = \lambda_p + \sum_j t_j\,\mu_j$ bleibt
\[
  \sum_j \underbrace{\mu_i^T W\,\mu_j}_{\delta_{ij}}\;t_j
  = -\underbrace{\mu_i^T W\,\lambda_p}_{\delta_{i0}} .
\]
Das sind die Elastizitätsgleichungen des KGV (hier ist $\delta_{ij}$ wieder eine $\delta$-Zahl und
kein Kronecker-Delta). Weil $\varepsilon$ herausfällt, hängt das Ergebnis nicht davon ab,
\emph{wie} starr die Stäbe sind, nur davon, dass sie alle gleich starr sind. Deshalb existiert der
Grenzwert. Passt die Temperatur dagegen nicht zu den abhängigen Zeilen ($\mu_i^T g \ne 0$), muss
$\varepsilon\,\mu_i^T W \lambda$ einen festen Wert annehmen, und für $\varepsilon \to 0$ wird
$\lambda$ unendlich groß. Das ist die unendliche Zwangskraft, zu der das Programm einen Hinweis
zeigt.

\textbf{Wo die Ergänzungsenergie hineinspielt.} Die Gleichungen oben sind genau die Bedingung dafür,
dass
\[
  \Pi^*(t) = \tfrac12\,\lambda^T W \lambda , \qquad \lambda = \lambda_p + \sum_i t_i\,\mu_i ,
\]
minimal wird, denn $\partial\Pi^*/\partial t_i = \mu_i^T W \lambda$. Das ist der Satz von Menabrea:
Unter allen Kraftzuständen im Gleichgewicht stellt sich der mit der kleinsten Ergänzungsenergie
ein. \enquote{Kleinste Ergänzungsenergie} ist dabei nur eine andere Schreibweise für
\enquote{verträglich}, das KGV in \textbf{Statik} rechnet mit dem Prinzip der virtuellen Kräfte
genau dasselbe. Die Ergänzungsenergie ist also kein zusätzliches Prinzip, sondern die
Verträglichkeit in Energieform. Genau genommen ist die Ergänzungsenergie der fast starren Stäbe
$\varepsilon\,\Pi^*$, denn ihre Nachgiebigkeit ist $\varepsilon\,W$ und nicht $W$. Ein fester
Faktor ändert aber nichts daran, wo das Minimum liegt, deshalb steht $\varepsilon$ oben gar nicht
erst in $\Pi^*$.

\textbf{$W$ ausrechnen} geht wie eine $\delta$-Zahl mit der Koppeltafel, mit $EA = EI = 1$ und für
jeden Stab für sich (verschiedene Stäbe koppeln nicht):
\begin{itemize}
  \item Längszeile: $N = 1$ konstant, $\int N^2\,dx = L$.
  \item Drehzeile: $M = 1$ konstant, $\int M^2\,dx = L$.
  \item Querzeile: $M$ linear von $L$ auf $0$, $\int M^2\,dx = L^3/3$.
  \item Dreh- mal Querzeile: $\int 1 \cdot (L - x)\,dx = L^2/2$. Längs- mal Biegezeile: $0$.
\end{itemize}
Die Querkraftverformung fehlt, wie überall im Programm (Euler-Bernoulli). Dass $EA$ und $EI$ beide
$1$ sind, ist eine Annahme und keine Physik. Sie spielt keine Rolle, solange in einem
Eigenspannungszustand nur dehnstarre oder nur biegestarre Zeilen vorkommen, denn dann kürzt sich der
gemeinsame Faktor. Mischen sich beide, hängt die Aufteilung vom angenommenen Verhältnis $EA : EI$
ab.

So teilen sich zwei dehnstarre Stäbe nebeneinander eine Last je zur Hälfte, und ein starrer
Zweifeldträger unter Gleichlast hat die bekannten Lagerkräfte $\frac38\,ql$, $\frac{10}8\,ql$,
$\frac38\,ql$, genau wie mit dem KGV und dem Stützmoment als $X_1$. Eine gelbe Box weist darauf
hin, dass diese Aufteilung angenommen ist, aber nur, wenn sie das Ergebnis auch ändert. Ist
$\lambda_p$ auf allen Zeilen null, die über $W$ mit einem $\mu_i$ gekoppelt sind, dann ist
$\delta_{i0} = \mu_i^T W \lambda_p = 0$ für \emph{jede} Wahl der Nachgiebigkeiten, also $t_i = 0$, egal wie
starr die Stäbe im Verhältnis zueinander sind. Ein Beispiel ist eine starre Platte, die oben und
unten auf Loslagern steht: Die Längskraft zwischen den beiden Lagern ist unbestimmt, aber ohne Last
in Plattenrichtung immer null.

\textbf{Das Gesamtbild.} Beide Aufteilungen in partikulär und homogen gehören zusammen. Die
Kinematik $C\,u = g$ und das Gleichgewicht $C^T\lambda = r$ benutzen dieselbe Matrix, einmal so und
einmal transponiert:

\begin{center}
\begin{tabular}{@{}p{2.9cm}p{5.9cm}p{5.9cm}@{}}
  \toprule
   & \textbf{Kinematik} $C\,u = g$ & \textbf{Gleichgewicht} $C^T\lambda = r$ \\
  \midrule
  partikulär       & $u_0$ (Temperatur) & $\lambda_p$ (Hauptsystem) \\
  homogen          & $T\,u_r$: Bewegungen, die die starren Stäbe erlauben
                   & $\sum_i t_i\,\mu_i$: Eigenspannungszustände \\
  festgelegt durch & Gleichgewicht (Schritt~3) & Verträglichkeit \\
  als Energie      & kleinste potentielle Energie & kleinste Ergänzungsenergie \\
  wie im           & WGV & KGV \\
  \bottomrule
\end{tabular}
\end{center}

Der Kern von $C$ sind die Bewegungen, die die starren Teile frei lassen. Der Kern von $C^T$ besteht
aus den Eigenspannungszuständen $\mu_i$, seine Dimension ist der Grad der statischen
Unbestimmtheit der starren Teile untereinander, also die Anzahl der $t_i$ (im KGV der $X_i$).

Die rote Box \emph{Starrmodus \& statisch unbestimmt} kommt übrigens schon dann, wenn das System
statisch unbestimmt ist und alle Stäbe $EA, EI \ge 10^7$ haben, also auch mit den
Standardwerten.

%---------------------------------------------------------------------
\section{Das Rand-LGS: die Biegelinie wie von Hand}\label{sec:mathe-randlgs}

Der Export \emph{Neue Biegelinie} rechnet unabhängig vom Weggrößenverfahren genau das, was ihr in
\textbf{TM2} von Hand macht, nur für jedes System und mit allen Bedingungen auf einmal.

\subsection{Ansatz je Stab}

Mit $\varphi = -w'$, $Q' = -q$, $M' = Q - m$ und $N' = -n$ setzt das Programm für jeden Stab an

\[
  EI\,w(x) = w_p(x) + c_1\frac{x^3}{6} + c_2\frac{x^2}{2} + c_3 x + c_4, \qquad
  EA\,u(x) = u_p(x) + C_1 x + C_2 .
\]

Das sind sechs Unbekannte je Stab. Die Partikulärlösungen baut es nicht aus einer Tabelle,
sondern als Integral, das der CAS-Kern exakt auswertet:

\begin{align*}
  w_p(x) &= \int_0^x q(t)\,\frac{(x-t)^3}{6}\,dt + \int_0^x m(t)\,\frac{(x-t)^2}{2}\,dt
            - EI\,\kappa_T\,\frac{x^2}{2}, \\
  u_p(x) &= -\int_0^x n(t)\,(x-t)\,dt + EA\,\varepsilon_T\,x .
\end{align*}

Die Idee dahinter: Man denkt sich die Streckenlast in viele kleine Einzellasten $q(t)\,dt$ an den
Stellen $t$ zerlegt. Bei einer Einzellast $1$ an der Stelle $t$ springt die Querkraft dort um $-1$
(wegen $Q' = -q$), und ab $t$ kommt der Anteil $EI\,w = (x-t)^3/6$ dazu. Das ist $c_1 x^3/6$ mit
$c_1 = 1$, nur um $t$ verschoben, davor ist der Anteil null. Weil alles linear ist, addieren sich
die Beiträge aller kleinen Einzellasten, und aus der Summe wird das Integral. Genauso gibt ein
Einzelmoment $1$ bei $t$ den Anteil $(x-t)^2/2$ in $EI\,w$ und eine Längskraft $1$ den Anteil
$-(x-t)$ in $EA\,u$. In der Mathematik heißt so ein
Integral \emph{Faltung}. Deshalb funktioniert es für jede Lastfunktion, für Trapezlasten und für
die globalen Linienlasten (die vorher auf Quer- und Längsrichtung des Stabes zerlegt werden). Und deshalb heißt die Integrationsvariable \code{t}, und
ihr könnt \code{t} nicht als Symbol verwenden.

Aus dem Ansatz folgen die Größen an beiden Stabenden als \emph{lineare} Ausdrücke in den
Konstanten, am Anfang z.\,B.\ $w = c_4/EI$, $\varphi = -c_3/EI$, $M = -c_2$, $Q = -c_1$,
$u = C_2/EA$, $N = C_1$.

\subsection{Bedingungen an den Knoten}

An jedem Knoten schreibt das Programm die Bedingungen, die ihr auch aufschreiben würdet, nur
systematisch:

\begin{itemize}
  \item \textbf{Kinematik:} Alle Stabenden ohne N- bzw.\ Q-Gelenk haben dieselbe
        Knotenverschiebung $\vec U$, also $u_e = \vec U\cdot\vec t_e$ und
        $w_e = \vec U\cdot\vec n_e$, alle Enden ohne M-Gelenk dieselbe Verdrehung. $\vec U$ selbst
        wird durch das erste passende Stabende ausgedrückt, damit keine zusätzlichen Unbekannten
        entstehen.
  \item \textbf{Gelenke:} Ein M-Gelenk gibt $M_e = 0$, ein N-Gelenk $N_e = 0$, ein Q-Gelenk
        $Q_e = 0$.
  \item \textbf{Lager:} In einer gehaltenen Richtung $\vec U\cdot\vec e = 0$, in einer freien
        Richtung das Gleichgewicht der Stabendkräfte, Federn und Lasten in dieser Richtung. Die
        Richtungen sind $\vec\xi = (\cos\alpha, -\sin\alpha)$ und $\vec\eta = (\sin\alpha, \cos\alpha)$,
        so gehen gedrehte Lager ohne Sonderfall.
  \item \textbf{Drehung:} Einspannung $\varphi = 0$, sonst $\sum M_e - c_m\,\varphi + M = 0$.
\end{itemize}

Am Ende muss die Anzahl der Gleichungen genau $6 \cdot (\text{Stäbe})$ sein (plus zwei je Knoten,
an dem alle Stäbe über N- oder Q-Gelenke hängen, siehe unten). Sonst kommt \emph{nicht quadratisch}.

\textbf{Warum solche Knoten eigene Verschiebungen brauchen.} Normalerweise wird $\vec U$ durch ein
Stabende ohne N- und Q-Gelenk ausgedrückt: Dessen $u_e$ und $w_e$ sind die beiden Komponenten von
$\vec U$ in Stab- und Querrichtung, damit ist $\vec U$ ganz bekannt. Ein Stabende mit N-Gelenk darf
sich aber längs gegen den Knoten verschieben, es liefert nur noch $w_e = \vec U\cdot\vec n_e$. Eines
mit Q-Gelenk liefert nur $u_e = \vec U\cdot\vec t_e$. Hat jedes Stabende am Knoten so ein Gelenk,
gibt es kein Ende, das beide Komponenten festlegt. Dann bekommt $\vec U$ die eigenen Unbekannten
$u_x(K)$, $u_y(K)$, und jedes Ende schreibt seine Bedingung gegen $\vec U$. Normalerweise
\enquote{verbraucht} das Bezugsende seine zwei Komponenten, um $\vec U$ festzulegen, und liefert
selbst keine Gleichung. Hier verbraucht keines etwas, daher die zwei Gleichungen mehr. Beispiel:
zwei Stäbe, die sich an einem Knoten mit je einem Q-Gelenk treffen. Jeder sagt nur, wie weit sich
der Knoten in seiner Längsrichtung bewegt, erst beide zusammen bestimmen $\vec U$.

\subsection{Lösen: erst wie von Hand, dann das CAS}

Ein LGS mit 30 symbolischen Unbekannten direkt in \code{simult} zu stecken, ist auf dem Rechner
langsam und liefert unübersichtliche Ausdrücke. Das Programm macht deshalb vorher, was ihr auch
machen würdet:

\begin{enumerate}
  \item Jede Gleichung mit nur \emph{einer} Unbekannten wird sofort aufgelöst und der Wert überall
        eingesetzt (z.\,B.\ $M = 0$ am Gelenk liefert direkt $c_2$). Das wiederholt sich, solange
        etwas passiert.
  \item Eine Unbekannte, die nur noch in einer einzigen Gleichung vorkommt, wird mit dieser
        Gleichung beiseitegelegt und am Schluss aus ihr berechnet.
  \item Gleichungen der Form $a\,x_j + b\,x_m = 0$ mit Zahlen $a$, $b$ werden zu $x_j = f\,x_m$
        und eingesetzt.
\end{enumerate}

Was übrig bleibt, zerfällt oft in unabhängige Blöcke (Unbekannte, die nie zusammen in einer
Gleichung stehen). Jeder Block wird einzeln mit \code{exact(simult(\dots))} gelöst, danach wird
rückwärts eingesetzt. Aus demselben Protokoll entsteht auch \code{randbed}: Eine Größe steht dort
nur dann als Wert, wenn ihre Konstanten in einem \emph{früheren} Schritt schon bekannt waren,
also genau in der Reihenfolge, in der ihr es auf dem Papier bekommen würdet.

\textbf{Die Aufteilung in \enquote{nur Biegelinie} und \enquote{nur Längslinie}} kommt aus
demselben Protokoll: Das Programm verfolgt für jede Konstante der Biegelinie zurück, aus welchen
Gleichungen sie bestimmt wurde. Tauchen darin keine Längslinien-Konstanten auf, die nicht schon
vorher feststanden, reichen die Bedingungen der Biegelinie allein.

Bei Rahmen mit dehnsteifen Stäben reicht das nicht. Dort bestimmt das LGS die Biegelinie zusammen
mit $C_1$ und $C_2$ der dehnsteifen Stäbe, weil z.\,B.\ die Seitenverschiebung $u$ des Riegels in
die Biegelinie der Stiele eingeht. Bei der direkten Methode ist $u = C_2$ und $N = C_1$ aber an
beiden Stabenden dieselbe Unbekannte (ohne Längslast und Temperatur). Das Programm nimmt deshalb
eine Bedingung, in der nur eine dieser Unbekannten steht, löst sie danach auf und setzt sie in die
anderen Bedingungen ein, genau wie von Hand. Mathematisch ist das ein Gauß-Schritt mit dieser
Unbekannten als Pivot, das LGS bleibt dabei regulär. Bleiben danach genau $4\cdot(\text{Stäbe})$
Bedingungen übrig, in denen nur noch Konstanten der Biegelinie stehen, genügen sie für $w$. Bei
dehnbaren Stäben ist $u(L) = (C_1 L + C_2)/EA$, dort hängt die Längsverschiebung am Stabende von
beiden Konstanten ab, und das Durchreichen geht nicht.

\subsection{Starre Stäbe im Rand-LGS: direkt oder Grenzwert}

Für einen starren Stab gibt es kein $EI\,w'' = -M$ mehr, trotzdem braucht das Programm sechs
Gleichungen. Zwei Wege:

\begin{itemize}
  \item \textbf{direkt} (Standard): Der biegesteife Stab bewegt sich als Starrkörper,
        $w(x) = c_4 + c_3 x - \kappa_T x^2/2$ und beim dehnsteifen $u(x) = C_2 + \varepsilon_T x$.
        $c_1$, $c_2$ und $C_1$ bleiben als reine Kraftgrößen stehen und werden aus dem
        Gleichgewicht bestimmt.

        Warum das geht: Auch beim elastischen Stab sind $c_1$, $c_2$ und $C_1$ eigentlich
        Kraftgrößen, nämlich die Schnittgrößen am Stabanfang, $Q(0) = -c_1$, $M(0) = -c_2$,
        $N = C_1$. In der Biegelinie stehen sie nur deshalb, weil die Krümmung aus dem Moment kommt:
        Sie gehen als $c_1 x^3/(6EI)$ und $c_2 x^2/(2EI)$ in $w$ ein, also immer geteilt durch
        $EI$. $c_3$ und $c_4$ dagegen sind echte Weggrößen ($c_3/EI = w'(0)$, $c_4/EI = w(0)$).
        Wird der Stab starr, $EI \to \infty$, verschwindet bei endlichen Kräften der Beitrag von
        $c_1$ und $c_2$ aus der Biegelinie, und $c_3$, $c_4$ schreibt man gleich ohne $EI$
        ($c_4 = w(0)$, $c_3 = w'(0)$). Die Schnittgrößen bleiben aber da: $Q(x)$, $M(x)$ und $N(x)$
        folgen aus $c_1$, $c_2$, $C_1$ und den Lasten allein über das Gleichgewicht am Stabstück
        ($Q' = -q$, $M' = Q$, $N' = -n$), ganz ohne Steifigkeit.

        Die Bedingungen an den Knoten sind dieselben wie beim elastischen Stab. Jede enthält aber
        entweder nur Kraftgrößen ($c_1$, $c_2$, $C_1$) oder nur Weggrößen ($c_3$, $c_4$, $C_2$) des
        starren Stabes: Kinematik und Lager mit gehaltener Richtung enthalten Wege, Gelenk-,
        Gleichgewichts- und Drehbedingungen Kräfte. Beispiel: starrer Kragarm, eingespannt am
        Anfang, Kraft $F$ (in $+w$) am freien Ende. Die Einspannung gibt $w(0) = c_4 = 0$ und
        $w'(0) = c_3 = 0$, also bewegt sich nichts. Das freie Ende gibt $Q(L) = -c_1 = F$ und
        $M(L) = -c_2 - c_1 L = 0$, also $c_1 = -F$, $c_2 = F L$ und $M(0) = -F L$. Das ist genau die
        Rechnung des statisch bestimmten Kragarms von Hand. In \code{randbed} heißen die
        Kraftgrößen eines starren Stabes deshalb $Q_i$, $M_i$, $N_i$ statt $-EI\,w_i'''$ usw.
  \item \textbf{Grenzwert}: Jeder starre Stab bekommt $EI = 1/\varepsilon$ bzw.\
        $EA = 1/\varepsilon$ (im CAS heißt $\varepsilon$ \code{rleps}), das System wird symbolisch in
        $\varepsilon$ gelöst und danach \code{limit(\dots, rleps, 0)} gebildet.
\end{itemize}

Warum überhaupt der Grenzwert? Das direkte Verfahren wird singulär, wenn sich starre Teile
\emph{gegenseitig} statisch unbestimmt halten. Beispiel: zwei dehnstarre Stäbe hintereinander
zwischen zwei Festlagern, Horizontallast in der Mitte. Beide Stäbe ändern ihre Länge nicht, also
ist $u = 0$ überall. Aber wie sich die Last auf die beiden Stäbe verteilt, sagt euch das
Gleichgewicht nicht, nur $N_1 - N_2 = F$. Mit $EA = 1/\varepsilon$ ist das System regulär, und
der Grenzwert $\varepsilon \to 0$ liefert die Aufteilung, die zwei \enquote{gleich starren} Stäben
entspricht. Genau in diesem Fall wechselt das Programm von selbst auf den Grenzwert. In
\code{randinfo} steht dann \emph{direkt singulär, starre Teile gegenseitig statisch unbestimmt}.

\subsection{Knotenverschiebungen aus den Randwerten}\label{sec:mathe-knotenverschiebung}

Für die Knotenverschiebung (Knotenmenü Seite~4) muss nichts Neues gelöst werden. Die
Kinematik-Bedingungen drücken $\vec U$ ja schon durch ein Referenz-Stabende aus, und zwar durch das erste
Stabende am Knoten ohne N- und Q-Gelenk. Dessen Längs- und Querverschiebung stehen nach dem Lösen
fest, $u(0)$ und $w(0)$ am Stabanfang bzw.\ $u(L)$ und $w(L)$ am Stabende. Zurück ins globale
System:
\[
  U_x = u\,t_x + w\,n_x, \qquad U_y = u\,t_y + w\,n_y, \qquad
  \vec t = \frac{(\Delta x, \Delta y)}{L}, \quad \vec n = \frac{(-\Delta y, \Delta x)}{L} .
\]
Das sind dieselben $\vec t$ und $\vec n$ wie in den Bedingungen, deshalb zählt $U_y$ wie die
$y$-Achse des Programms nach unten. Die Verdrehung kommt vom ersten Stabende ohne M-Gelenk,
\[
  \varphi = -w' = -\frac{w_p'(x) + c_1 x^2/2 + c_2 x + c_3}{EI},
\]
beim biegesteifen Stab $\varphi = \kappa_T x - c_3$. Hängen alle Stäbe über N- oder Q-Gelenke am
Knoten, hat das LGS für ihn ohnehin die eigenen Unbekannten $u_x(K)$, $u_y(K)$, die werden direkt
genommen.

Damit das beim zweiten Knoten nicht wieder ein ganzes LGS kostet, hebt das Programm die Linien des
letzten Rand-LGS mit einem Schlüssel auf: dem Fingerabdruck aller Eingaben (derselbe, an dem die
Hinweisbox erkennt, ob sich etwas geändert hat. Er enthält auch die Knoten- und Stabnummern) und
den Schaltern, die das Rand-LGS sonst beeinflussen (Methode für starre Stäbe, symbolischer Modus,
Zahlenformat). Passt der Schlüssel, werden nur die Randwerte eingesetzt, das sind drei kurze
CAS-Aufrufe statt eines LGS. Gleicher Fingerabdruck heißt wirklich gleiches System, auch wenn ein
Knoten zwischendurch gelöscht und genau so wieder eingebaut wurde. Werden dabei Nummern
vertauscht, ändert sich der Fingerabdruck, und das ist auch richtig so: Dann gehören die Namen
\code{wv1}, \code{wv2}, \dots\ zu anderen Stäben.

%---------------------------------------------------------------------
\section{Symbolischer Modus: rechnen mit Platzhaltern}\label{sec:mathe-symbolisch}

Die Finite-Elemente-Rechnung für die Diagramme ist eine reine Zahlenrechnung. Symbolisch
invertiert der Rechner keine Steifigkeitsmatrix in vernünftiger Zeit. Der Trick:

\begin{enumerate}
  \item Jedes Symbol bekommt einen \textbf{Platzhalterwert}, standardmäßig~1. Die Länge $l$
        bekommt, wenn die Geometrie kein eigenes Symbol hat, den Rasterabstand.
  \item \textbf{Superposition:} Alle Lasten werden auf null gesetzt und dann \emph{einzeln} wieder
        eingeschaltet. Für jede Last rechnet das Programm die Stabendkräfte und teilt sie durch den
        Platzhalter der Last und die passende Potenz von $l$. Übrig bleibt ein reiner
        Zahlenfaktor~$c$:
        \[
          N = c\,F, \quad M = c\,F\,l \quad\text{(Einzellast)}, \qquad
          N = c\,q\,l, \quad M = c\,q\,l^2 \quad\text{(Streckenlast)} .
        \]
        Die Steifigkeitsmatrix wird dafür nur einmal zerlegt, jede Last kostet dann nur noch ein
        Vorwärts- und Rückwärtseinsetzen.
  \item Der Zahlenfaktor wird \textbf{zurück in einen exakten Ausdruck} verwandelt (siehe unten)
        und alle Anteile werden addiert.
\end{enumerate}

Das funktioniert, weil die Rechnung linear ist und alle Längen Vielfache von $l$ sind. Genau hier
kommen die Einschränkungen her: Eine Temperatur passt nicht in das Schema \enquote{Last mal
Länge hoch $k$} (deshalb fehlt sie in den Diagrammwerten; die Exporte \emph{Neue Biegelinie} und
\emph{Neue Längslinie}, also \code{wneu}, \code{uneu}, \code{wv}, \code{uv}, haben sie, weil das
Rand-LGS gar nicht mit Superposition rechnet), und ein Zahlen-$EI$
neben einem symbolischen $EI$ macht die Steifigkeitsverhältnisse vom Platzhalterwert abhängig
(deshalb der gelbe Hinweis). Die Stelle des Maximums in einem Feld sucht das Programm numerisch
und rastet sie auf typische Werte wie $1/\sqrt 3$ oder einfache Brüche ein.

\textbf{Aus einer Zahl wird wieder ein Bruch.} Aus $0{,}708333\ldots$ soll wieder
$\frac{17}{24}$ werden. Dafür entwickelt das Programm die Zahl in einen \emph{Kettenbruch}
$a_0 + \cfrac{1}{a_1 + \cfrac{1}{a_2 + \dots}}$ und nimmt den ersten Näherungsbruch, der bis
auf $10^{-12}$ stimmt (Nenner höchstens 10\,000). Klappt das nicht und steht das Zahlenformat auf
\emph{Wurzel/Bruch}, versucht es dasselbe mit dem \emph{Quadrat} der Zahl und zieht aus Zähler und
Nenner die Quadratfaktoren heraus. So wird aus $0{,}7071\ldots$ wieder $\frac{\sqrt 2}{2}$ und
aus $86{,}60\ldots$ ein $50\sqrt 3$. Scheitert beides, sucht es zuletzt noch einen Bruch mit
kleinem Nenner (höchstens 100), der nur bis auf die eingestellte Toleranz $10^{-x}$ stimmen muss
(Obermenü Seite~3, \emph{CAS Bruch-Tol.}). Das ist die einzige Stelle, an der diese Einstellung
wirkt, und gerechnet wird dabei in Lua, nicht im CAS. Erst wenn auch das nicht klappt, bleibt eine
Dezimalzahl stehen. Im Format \emph{Dezimal} fallen alle drei Versuche weg.

%---------------------------------------------------------------------
\section{Kraftgrößenverfahren: Hauptsystem und \texorpdfstring{$\delta$}{delta}-Zahlen}\label{sec:mathe-kgv}

\textbf{Das Hauptsystem} sucht das Programm \enquote{gierig}: Es löst der Reihe nach
Einspannmomente, Lagerkräfte, Momente an biegesteifen Stabanschlüssen und zuletzt unbelastete
Pendelstäbe. Es behält jede Lösung nur, wenn das System danach noch tragfähig ist (die
Steifigkeitsmatrix lässt sich noch zerlegen). Sobald nichts mehr geht, ist das System statisch
bestimmt, und die Anzahl der gelösten Bindungen ist der Grad $n$.

\textbf{Die $\delta$-Zahlen} integriert das Programm nicht über $\int M_i M_j/(EI)\,dx$ (so rechnet
ihr sie in \textbf{Statik} mit dem Prinzip der virtuellen Kräfte), sondern es
rechnet für jeden Zustand ($X_j = 1$ bzw.\ Lastzustand) die Verformung des Hauptsystems mit dem
Weggrößenverfahren und liest die zu $X_i$ gehörende Klaffung ab: Verschiebung am gelösten Lager,
gegenseitige Verdrehung am eingefügten Gelenk, Längenänderung am gelösten Pendelstab. Nach
Maxwell--Betti ist das dasselbe wie das Arbeitsintegral, inklusive der Anteile aus Normalkraft,
Federn und endlichem $EA$, die man von Hand gern vergisst. Ein gelöster Pendelstab ist im
Hauptsystem ganz entfernt. Seine eigene Längenänderung $L/(EA)$ fehlt deshalb in den
Knotenverschiebungen und wird zu $\delta_{jj}$ addiert. Bei gedrehten Lagern liegen die
Knotenfreiheitsgrade schon in Lagerrichtung ($u_\xi$, $u_\eta$), die Klaffung eines gelösten Lagers
ist genau dieser Freiheitsgrad, und die Einheitskräfte eines gelösten Pendelstabes werden in
diese Richtungen zerlegt.

%---------------------------------------------------------------------
\section{Polplan, Verschiebungszustand und Arbeitssatz}\label{sec:mathe-polplan}

\subsection{Scheiben finden}

Zuerst muss das Programm wissen, welche Stäbe zusammen eine starre Scheibe bilden. Das passiert
in zwei Durchgängen:

\begin{enumerate}
  \item \textbf{Ohne Lager:} An jedem Knoten ohne Vollgelenk fasst das Programm alle Stäbe
        zusammen, deren Enden dort kein Gelenk haben (\emph{Union-Find}: jede Gruppe hat einen
        Vertreter, zwei Gruppen werden vereinigt, indem man einen Vertreter unter den anderen
        hängt). Mit diesen Scheiben und nur den inneren Gelenken läuft der Polplan einmal durch.
        Findet er zwei Scheiben, die gegeneinander unverschieblich sind, z.\,B.\ die drei Stäbe
        eines Dreigelenk-Dreiecks, werden sie zu einer Scheibe zusammengefasst.
  \item \textbf{Mit Lagern:} Stäbe, die biegesteif an einer Einspannung hängen (oder an einem
        Festlager mit Drehfeder), gehören zur Scheibe~0, dem Boden. Mit diesen zusammengefassten
        Scheiben und den Lagern entsteht der Polplan, den ihr seht.
\end{enumerate}

\subsection{Pole und Polstrahlen}

An jedem Knoten schaut das Programm, welche Scheiben sich dort treffen und welche Relativbewegung
zwischen ihnen möglich ist:

\begin{longtable}{p{0.34\linewidth}p{0.60\linewidth}}
  \toprule
  \textbf{Verbindung} & \textbf{Folge im Polplan} \\
  \midrule \endhead
  M-Gelenk, Festlager & Pol im Knoten (Drehung möglich, keine Verschiebung) \\
  Rollenlager & Polstrahl durch den Knoten senkrecht zur Rollrichtung, also entlang der
    Lagerkraft \\
  N- oder Q-Gelenk & Pol im Unendlichen (reine Verschiebung längs bzw.\ quer zum Stab) \\
  Einspannung, biegesteifer Anschluss & kein Pol: die Scheiben sind starr verbunden \\
  \bottomrule
\end{longtable}

Danach wiederholt das Programm (höchstens zehnmal) zwei Schritte, bis nichts Neues mehr dazukommt:

\begin{enumerate}
  \item \textbf{Satz von den drei Polen:} Für je drei Scheiben $i$, $j$, $k$ liegen die Pole
        $(i,j)$, $(j,k)$ und $(i,k)$ auf einer Geraden. Sind zwei davon bekannt, entsteht ein
        Polstrahl für den dritten. Liegt einer der beiden im Unendlichen, geht der Strahl durch
        den anderen in dessen Richtung.
  \item \textbf{Strahlen schneiden:} Zwei Polstrahlen derselben Scheibenpaarung schneiden sich im
        Pol. Sind sie parallel, liegt der Pol im Unendlichen in dieser Richtung.
\end{enumerate}

Findet es für eine Paarung zwei verschiedene Pole, oder liegt ein Pol nicht auf einem Strahl, auf
dem er liegen müsste, sind die beiden Scheiben gegeneinander unverschieblich (im Polplan rot). Im
Durchgang ohne Lager werden solche Paare zu einer Scheibe zusammengefasst (siehe oben).

\subsection{Der Verschiebungszustand}

Eine Scheibe, die sich um ihren Hauptpol $(c_x, c_y)$ um den kleinen Winkel $\varphi$ dreht,
verschiebt jeden Punkt um $\vec r_\perp \varphi$, also senkrecht zum Polstrahl vom Pol zu diesem
Punkt. Das Programm gibt einer Startscheibe $\varphi = 1$ (bevorzugt eine, deren Hauptpol auf einem
Knoten liegt) und reicht die Bewegung über die Nebenpole weiter: Der Nebenpol $(A,B)$ ist der Punkt,
der auf beiden Scheiben dieselbe Verschiebung hat. Kennt man die Verschiebung dieses Punktes aus
Scheibe~$A$ und den Hauptpol von Scheibe~$B$, folgt $\varphi_B$. Hat eine Scheibe ihren Hauptpol im
Unendlichen, verschiebt sie sich nur.

\textbf{Verschobene Scheiben unterwegs.} Die Kette reißt an einer Scheibe, die sich nur verschiebt,
nicht ab. Bei ihr bewegt sich jeder Punkt um denselben Vektor $\vec v$, also auch der Nebenpol zur
nächsten Scheibe~$B$. Damit ist die Verschiebung des Nebenpols bekannt, und es geht weiter wie bei
einer drehenden Scheibe: Liegt der Hauptpol von $B$ im Endlichen, folgt $\varphi_B$ aus $\vec v$,
liegt er im Unendlichen, verschiebt sich auch $B$ nur. Liegt der Nebenpol $(A,B)$ selbst im
Unendlichen, drehen beide Scheiben gleich, $\varphi_B = \varphi_A = 0$: $B$ verschiebt sich dann auch
nur, oder sie steht still, wenn ihr Hauptpol im Endlichen liegt. Eine reine Verschiebung ist nur der
Grenzfall einer Drehung mit dem Pol im Unendlichen.

\textbf{Nebenpol gleich Hauptpol.} Der Hauptpol einer Scheibe ist der Punkt, der sich auf ihr nicht
bewegt. Fällt der Nebenpol $(A,B)$ mit dem Hauptpol von $B$ zusammen, bewegt sich dieser Punkt also
auf $B$ nicht, und weil er auf $A$ dieselbe Verschiebung haben muss, auch auf $A$ nicht. Kommt aus
der Startscheibe $A$ trotzdem eine Verschiebung ungleich null an, passt das nicht zusammen: Der Punkt
kann nicht zugleich stehen und sich bewegen. Dann kann sich $A$ gar nicht bewegen, die Startscheibe
war falsch gewählt, und das Programm nimmt die nächste.

Ein Beispiel ist euer Dreigelenktragwerk mit dem Gelenk $G$ senkrecht über dem linken Lager $A$,
bei dem ihr das linke Festlager zum Loslager macht (verschieblich in $x$). Das rechte Festlager $B$
ist der Hauptpol $(\mathrm{II})$ der rechten Scheibe, das Gelenk der Nebenpol $(\mathrm{I,II})$. Der
Hauptpol $(\mathrm{I})$ der linken Scheibe liegt auf der Senkrechten durch $A$ (Loslager) und nach
dem Satz von den drei Polen auf der Geraden durch $(\mathrm{II})$ und $(\mathrm{I,II})$, also auf
$BG$. Beide Geraden schneiden sich in $G$, damit ist $G$ zugleich $(\mathrm{I})$ und
$(\mathrm{I,II})$. Startet das Programm mit der rechten Scheibe, dreht sie um $B$ und schiebt $G$
seitlich weg, aber auf der linken Scheibe steht $G$ als Hauptpol still: Widerspruch, die rechte
Scheibe kann sich nicht bewegen. Mit der linken Scheibe als Start geht es: Sie dreht um $G$, ihr
Lager $A$ wandert waagerecht (senkrecht zum Polstrahl $GA$, und das lässt das Loslager zu). $G$
selbst bewegt sich dabei nicht. Die rechte Scheibe hat damit zwei Punkte, die stehen bleiben ($B$
und $G$), und bewegt sich deshalb gar nicht, $\varphi_\mathrm{II} = 0$. Das ist kein Fehler,
sondern die Kinematik dieses Systems: Nur die linke Scheibe kippt um das Gelenk.

\subsection{Der Arbeitssatz}\label{sec:mathe-arbeitssatz}

Nach dem Lösen einer Bindung ist das System eine kinematische Kette mit genau einem Freiheitsgrad.
Für jede virtuelle Bewegung, die die übrigen Bindungen zulassen, gilt
\[
  \delta W = \sum_k \vec F_k \cdot \delta\vec u_k + \sum_m M_m\,\delta\varphi_m = 0 ,
\]
und weil es nur einen Freiheitsgrad gibt, sind alle Wege Vielfache eines einzigen Parameters. Der
Export (Abschnitt~\ref{sec:arbeitssatz}) macht daraus drei Schritte.

\textbf{Weg einer Kraft (\code{pvv\_kin}, Spalte~2).} Eine Kraft mit Richtung
$\vec e = (e_x, e_y)$ ($y$ nach unten) greift im Punkt $P$ auf Scheibe~$i$ an. Ihr Weg ist die
Verschiebung von $P$ in Richtung $\vec e$:
\begin{itemize}
  \item Scheibe~$i$ dreht um ihren Hauptpol $(c_x, c_y)$ um $\varphi_i$ (gegen den Uhrzeigersinn).
        Jeder Punkt bewegt sich senkrecht zum Polstrahl vom Pol zu ihm, projiziert auf $\vec e$ ist das
        \[
          \delta_e = \varphi_i\,\bigl[e_x\,(P_y - c_y) - e_y\,(P_x - c_x)\bigr] .
        \]
        Die eckige Klammer ist der Hebel der Kraft um den Pol.
  \item Scheibe~$i$ verschiebt sich nur (Hauptpol im Unendlichen) um $\delta_i$ in Richtung
        $\vec t_i$: $\delta_e = \delta_i\,(\vec e\cdot\vec t_i)$.
  \item Ein Moment auf Scheibe~$i$ arbeitet am Drehwinkel: $\delta\varphi = \varphi_i$.
\end{itemize}
Verschieben \emph{und} drehen zugleich gibt es nicht als eigenen Fall: Eine kleine ebene
Starrkörperbewegung ist immer eine Drehung um einen Punkt, eben den Pol. Eine Verschiebung plus
Drehung ist nur eine Drehung um einen anderen Pol, die reine Verschiebung der Grenzfall mit dem Pol
im Unendlichen.
Streckenlasten dürfen exakt durch ihre Resultierende $R$ ersetzt werden. Auf einem geraden Stab
der Länge $L$, der zu einer starren Scheibe gehört, ändert sich die virtuelle Verschiebung linear
mit der Stabkoordinate $x$, denn ein Starrkörper verschiebt jeden Punkt linear mit seinem Ort:
$\delta\vec u(x) = \delta\vec u_A + \frac{x}{L}\,(\delta\vec u_B - \delta\vec u_A)$. Damit ist die
Arbeit der Streckenlast
\[
  \int_0^L q(x)\,\vec e\cdot\delta\vec u(x)\,dx
  = \vec e\cdot\delta\vec u_A \underbrace{\int_0^L q\,dx}_{R}
    + \vec e\cdot\frac{\delta\vec u_B - \delta\vec u_A}{L} \underbrace{\int_0^L q\,x\,dx}_{S = R\,x_R}
  = R\;\vec e\cdot\delta\vec u(x_R) ,
\]
also die Arbeit der Resultierenden $R$ an ihrem Angriffspunkt $x_R = S/R$. Auf einem Bogen ist die
Verschiebung nicht linear in der Bogenlänge, das Ergebnis gilt aber genauso: Dort wird die Last
durch ihre Resultierende (Kraft und Richtung) auf ihrer Wirkungslinie ersetzt. Die freigesetzte
Größe ist eine äußere Last in der gezeichneten Richtung, bei einem gelösten Gelenk arbeitet sie an
der gegenseitigen Verdrehung bzw.\ Verschiebung der beiden Seiten.

\textbf{Verhältnis der Scheiben (\code{pvv\_pol}).} Der Nebenpol $P_{AB}$ ist der Punkt, der auf
Scheibe~$A$ und Scheibe~$B$ dieselbe Verschiebung hat. Schreibt man diese Verschiebung mit der
Formel oben einmal über $A$ und einmal über $B$ hin, entsteht eine lineare Beziehung
$\varphi_B = \kappa_{AB}\,\varphi_A$. Die Verschiebung ist ein Vektor mit $x$- und $y$-Komponente,
das sind eigentlich zwei Gleichungen. Beide liefern aber dasselbe $\kappa_{AB}$: Über $A$ steht die
Verschiebung senkrecht auf dem Polstrahl von $(A)$ zum Nebenpol, über $B$ senkrecht auf dem von
$(B)$, und weil $(A)$, $(A,B)$ und $(B)$ auf einer Geraden liegen, sind beide Richtungen gleich.
Deshalb reicht eine Komponente. Das Programm nimmt die betragsmäßig größere, damit es nicht durch
eine Zahl nahe null teilt. Liegt der Nebenpol zweier drehender Scheiben im Unendlichen, drehen beide
gleich, $\varphi_A = \varphi_B$.

Von der Bezugsscheibe aus (der Startscheibe des Verschiebungszustands, ihr Winkel heißt
$\varphi_0$) laufen diese Beziehungen über die Nebenpole durch alle Scheiben und liefern
$\varphi_i = \kappa_i\,\varphi_0$. Das sind die unteren Zeilen von \code{pvv\_kin}. Gibt es mehr
Nebenpole als nötig, sind die übrigen Proben. Als Bezug nimmt das Programm eine Scheibe, die sich
dreht, wenn es eine gibt. Eine Scheibe, die stehen bleibt, bekommt einfach $\kappa_i = 0$, ihre
Lasten leisten keine Arbeit. Verschieben sich alle bewegten Scheiben nur (alle ihre Hauptpole im
Unendlichen), gibt es keinen Drehwinkel als Bezug. Dann ist die Bezugsgröße die Verschiebung $\delta_0$ der Startscheibe, und
alle anderen Bewegungen werden als Vielfache von $\delta_0$ geschrieben (siehe unten).

\textbf{Einsetzen (\code{w\_eq2}).} Mit $\varphi_i = \kappa_i\,\varphi_0$ wird jeder Weg zu
$h_k\,\kappa_{i(k)}\,\varphi_0$ (Spalte~3 von \code{pvv\_kin}). Dabei ist $i(k)$ die Scheibe, an der
die Kraft $F_k$ angreift, und $h_k$ ihr Hebel um den Pol dieser Scheibe, also die eckige Klammer
von oben, $h_k = e_x\,(P_y - c_y) - e_y\,(P_x - c_x)$. Das ist der senkrechte Abstand der
Wirkungslinie vom Pol mit Vorzeichen, nicht der Abstand des Angriffspunktes. Verschiebt sich die
Scheibe nur, tritt an die Stelle von $h_k\,\varphi_i$ der Weg $\delta_i\,(\vec e\cdot\vec t_i)$ aus
der Liste oben, $h_k$ ist dann der Richtungskosinus $\vec e\cdot\vec t_i$. Der Arbeitssatz zerfällt in
\[
  \delta W = \varphi_0 \Bigl(\sum_k F_k\,h_k\,\kappa_{i(k)} + \sum_m M_m\,\kappa_{i(m)}\Bigr) = 0 .
\]
Da $\varphi_0$ beliebig ist, muss die Klammer verschwinden. Sie ist linear in der freigesetzten
Größe, also eine Gleichung mit einer Unbekannten. Verschiebt sich die Bezugsscheibe nur, steht
$\delta_0$ statt $\varphi_0$ da, und die Faktoren der Drehwinkel haben die Einheit $1/$Länge.
Intern zählt das Programm Drehwinkel im Uhrzeigersinn. Im Export ist das Vorzeichen umgedreht,
damit alles zur Konvention der Vorlesung (gegen den Uhrzeigersinn positiv) passt.

%---------------------------------------------------------------------
\section{Gleichgewicht der Teilsysteme (LGS-Matrix)}\label{sec:mathe-gleichgewicht}

Für den Export \emph{LGS-Matrix} schneidet das Programm das System an allen Gelenken auf, genau
so, wie Seite~2 der Explosionszeichnung es zeigt. Beide benutzen dieselben drei Regeln im Code:
\code{gelenkKnoten} (welche Stäbe hängen mit Gelenk, welche biegesteif am Knoten, wie viele
Teilsysteme treffen sich), \code{gelenkRichtung} (positive Richtung der Gelenkgrößen) und
\code{lastStabKnoten} (an welchem Teilsystem Lasten und Lager eines Gelenkknotens hängen). Jedes
Teilsystem liefert drei Gleichungen: Kräfte nach rechts, nach oben und Momente gegen den
Uhrzeigersinn. Unbekannt sind die Lagerreaktionen und die Größen, die ein Gelenk überträgt: beim
M-Gelenk die beiden Kraftkomponenten $Gx$, $Gy$, beim N-Gelenk $Gq$ und $Gm$, beim Q-Gelenk $Gn$
und $Gm$. Treffen sich mehr als zwei Teilsysteme in einem Gelenk, bekommt jeder Gelenkstab eigene
Gelenkkräfte, gibt es dort keinen biegesteifen Rest, sagen zwei zusätzliche Gleichungen, dass der
Gelenkbolzen selbst im Gleichgewicht ist. Bei nur zwei Teilsystemen braucht der Bolzen keine
eigenen Gleichungen: Die Gelenkkraft steht in beiden mit entgegengesetztem Vorzeichen (actio =
reactio), damit ist sein Gleichgewicht schon erfüllt. Hängt am Gelenk ein biegesteifer Rest, gehört
der Bolzen zu ihm. Ein Beispiel ist eine Stütze, die durch den Knoten durchläuft, und zwei Riegel,
die gelenkig an ihr hängen. Die Gelenkkräfte der Riegel greifen dann direkt an der Stütze an und
stehen in deren drei Gleichungen. Der Bolzen ist kein eigener Körper, also gibt es für ihn keine
eigenen Gleichungen, sein Gleichgewicht steckt im Gleichgewicht der Stütze. Erst wenn alle Stäbe
gelenkig am Knoten hängen, gehört der Bolzen zu keinem Teilsystem und braucht seine zwei
Gleichungen selbst. Im symbolischen Modus stehen Lasten mit ihren Namen und
Hebelarme als Vielfache von $l$ in der Matrix.

Den \textbf{Momentenbezugspunkt} wählt das Programm so, wie ihr es in \textbf{TM1} lernt: Es
sucht den Punkt, durch den möglichst viele unbekannte Kräfte gehen, also die senkrechte Linie mit den
meisten unbekannten Vertikalkräften und die waagerechte mit den meisten unbekannten
Horizontalkräften, am liebsten ein Knoten, der auf beiden liegt. Die Momentengleichung wird dadurch
so kurz wie möglich.

Die Werte auf Seite~1 der Explosionszeichnung kommen aus dem Weggrößenverfahren. Löst ihr das LGS
(mit den oben beschriebenen Befehlen oder von Hand), müssen dieselben Zahlen herauskommen. Das ist
eure Kontrolle.

\textbf{Lösungsweg (\code{lgs\_weg}):} Gesucht ist eine Reihenfolge, in der jede Gleichung möglichst
nur eine neue Unbekannte enthält. Dafür braucht das Programm mehr Gleichungen als die in \code{lgs}.
Alle sind Linearkombinationen der Zeilen aus \code{lgs}, also nichts Neues, nur anders
zusammengefasst.

\textbf{Anderer Momentenpunkt.} Eine Kraft mit den Komponenten $F_\rightarrow$ (nach rechts) und
$F_\uparrow$ (nach oben) im Punkt $(x, y)$ hat um $P$ das Moment
$M_P = (x - x_P)\,F_\uparrow - (y - y_P)\,F_\rightarrow$ (gegen den Uhrzeigersinn positiv, hier $y$
nach oben). Die Pfeile sind dieselben wie in den Namen der Kraftzeilen von \code{lgs}. Um einen
anderen Punkt $Q$ ist es
\[
  M_Q = (x - x_Q)\,F_\uparrow - (y - y_Q)\,F_\rightarrow
      = M_P + (x_P - x_Q)\,F_\uparrow - (y_P - y_Q)\,F_\rightarrow .
\]
Über alle Kräfte eines Teilsystems summiert (reine Momente hängen nicht vom Punkt ab) heißt das
\[
  \textstyle\sum M_Q = \sum M_P + (x_P - x_Q) \sum F_\uparrow - (y_P - y_Q) \sum F_\rightarrow .
\]
Die Momentengleichung um $Q$ ist also die Momentenzeile plus Vielfache der beiden Kraftzeilen
desselben Teilsystems, die Faktoren sind Hebelarme. Im Programm zeigt $y$ nach unten, dort steht
deshalb $+(y_P - y_Q)$.

\textbf{Teilsysteme zusammenfassen.} Für eine Gruppe werden die Kraftzeilen addiert und die
Momentenzeilen, jede vorher auf denselben Punkt $Q$ umgerechnet. Eine Gelenkkraft zwischen zwei
Teilsystemen der Gruppe steht in beiden mit entgegengesetztem Vorzeichen und greift am selben Punkt
an. Sie fällt deshalb aus den Kräften \emph{und} aus den Momenten heraus. Alle Teilsysteme zusammen
sind das Gesamtsystem.

\textbf{Kräfte quer zu einer Richtung.} Zeigt eine Unbekannte in Richtung $(d_x, d_y)$ (wieder mit $y$
nach unten), dann enthält $d_y \sum F_\rightarrow + d_x \sum F_\uparrow$ (normiert) diese Unbekannte
nicht mehr.

\textbf{Warum höchstens drei Gleichungen je Scheibe?} Jede Gleichgewichtsbedingung einer Scheibe ist
entweder eine Kraftsumme in eine Richtung $\vec e$, also $e_x \sum F_\rightarrow + e_y \sum
F_\uparrow$, oder eine Momentensumme um einen Punkt, nach der Formel oben $\sum M_P$ plus Vielfache
der beiden Kraftsummen. Alle liegen also in dem Raum, den $\sum F_\rightarrow$, $\sum F_\uparrow$ und
$\sum M_P$ aufspannen, und der hat höchstens die Dimension~3. Mehr als drei linear unabhängige
Gleichungen bekommt ihr deshalb nie, egal um wie viele Punkte ihr die Momente bildet. Drei
Momentengleichungen um $A$, $B$ und $C$ sind genau dann unabhängig, wenn die drei Punkte nicht auf
einer Geraden liegen. Liegen sie auf einer Geraden, ist die dritte eine Kombination der beiden
anderen. Ein Durchlaufträger auf drei Lagern bleibt deshalb einfach statisch unbestimmt, auch wenn
man um jedes Lager die Momente bildet. Anschaulich (PvV): Jede Gleichung gehört zu einer virtuellen
Starrkörperbewegung, einer Verschiebung in Richtung $\vec e$ bzw.\ einer Drehung um $P$, und eine
Scheibe hat in der Ebene nur drei Freiheitsgrade.

\textbf{Welche Gruppen und Punkte?} Das Programm sucht die Gleichungen nicht gezielt, sondern stellt
alle aus einem festen Vorrat auf und wählt danach aus:

\begin{itemize}
  \item \textbf{Gruppen:} jedes Teilsystem einzeln, das Gesamtsystem und bei Rahmen alle
        zusammenhängenden Gruppen. Zwei Teilsysteme hängen zusammen, wenn sie eine
        Gelenkunbekannte teilen. Ohne gemeinsame Gelenkkraft hebt sich beim Zusammenfassen nichts
        auf. Die Gruppen entstehen, indem an jede Gruppe ein Nachbar angehängt wird, bis die
        Zahl unter \emph{Max. Teilgruppen} erreicht ist (ab Werk 30). Im Fachwerk gibt es nur die Knoten und das
        Gesamtsystem, denn eine Gruppe von Knoten wäre ein Ritterschnitt.
  \item \textbf{Momentenpunkte $Q$:} Knoten und Schnittpunkte zweier Wirkungslinien von
        Unbekannten bzw.\ Pendelachsen (innerhalb der doppelten Systemgröße), aber je Gruppe nur
        die, an denen eine \emph{äußere} Kraft der Gruppe herausfällt. Äußere Kräfte sind die, die
        nach dem Zusammenfassen noch in der Kraftsumme stehen. Gelenkkräfte zwischen zwei
        Teilsystemen der Gruppe heben sich auf und zählen nicht. Ein Knoten zählt, wenn die
        Wirkungslinie einer äußeren Kraft durch ihn geht, ein Schnittpunkt, wenn beide Linien zu
        äußeren Kräften gehören, und er wird nach diesem Paar benannt. Für \code{TS 1+2} ist das
        Gelenk zwischen TS~1 und TS~2 deshalb nur dann ein Kandidat, wenn noch eine andere äußere
        Kraft durch es hindurchgeht. $Q$ wird also nicht vorher ausgesucht: Jeder dieser Punkte
        liefert eine Gleichung, und die Auswahl unten nimmt die, in der nur noch eine Unbekannte
        übrig ist.
  \item \textbf{Kräfte quer} zu jeder schrägen äußeren Unbekannten und jeder schrägen Pendelachse.
\end{itemize}

Gleichungen, die sich nur um einen Faktor unterscheiden, bleiben einmal übrig, und zwar in der
einfachsten Form. Weil dabei nur mit Zahlen gerechnet wird, dauert das auch bei einigen hundert
Kandidaten kaum Zeit.

\textbf{Benennung.} Der Name nennt die Art, den Punkt und die Gruppe: \code{$\Sigma$M P5 TS 1+2}
sind die Momente um Knoten~5 an den Teilsystemen 1 und 2 zusammen, \code{$\Sigma$M S(A\_h,B\_v)
Ges} die Momente des Gesamtsystems (\code{Ges}) um den Schnittpunkt der Wirkungslinien von $A_h$ und
$B_v$, \code{$\Sigma$F quer zu B\_h TS 2} die Kräfte an Teilsystem~2 quer zu $B_h$. Neue Gleichungen
bekommen die nächste freie Nummer nach denen aus \code{lgs}.

\textbf{Pendelstäbe.} Ein Teilsystem eines Rahmens ohne Lasten, an dem nur an zwei Punkten $U$ und
$V$ Anschlusskräfte angreifen (Lager- bzw.\ Gelenkkräfte, keine Momente), nennen wir hier
Pendelstab. Ist es ein gerader Stab, ist das der Pendelstab aus \textbf{TM1} mit reiner
Normalkraft. Ist es geknickt (z.\,B.\ ein unbelasteter Winkel zwischen zwei Gelenken), hat es zwar
Momente und Querkräfte, für das Gleichgewicht verhält es sich aber genauso: Die Anschlusskraft liegt
auf der Verbindungslinie $UV$. Nur darum geht es hier. Sein
Momentengleichgewicht um $U$ enthält nur die beiden Kraftkomponenten bei $V$ und sagt: Die Kraft bei
$V$ liegt auf der Achse $UV$. Liegt der Punkt $Q$ einer anderen Momentengleichung auf dieser Achse,
stehen die beiden Komponenten bei $V$ dort im selben Verhältnis wie im Pendelgleichgewicht, nur mit
einem Faktor $\lambda$. Zieht man $\lambda$-mal das Pendelgleichgewicht ab, fallen beide heraus. Das
Programm prüft das für jede Gleichung und jeden Pendelstab mit den Zahlen und nimmt die Achse als
weitere Wirkungslinie für die Schnittpunkte dazu. So findet sich auch beim Dreigelenkrahmen mit
verschieden hohen Lagern und einer unbelasteten Hälfte ein Punkt, der direkt eine Lagerkraft
liefert. Ist die Achse waagerecht oder senkrecht, braucht es den Umweg nicht, dann hat das
Pendelgleichgewicht selbst schon nur eine Unbekannte. Dieselbe Gleichung entsteht oft auf zwei
Wegen, etwa am Gesamtsystem (Lagerkräfte über das Pendelgleichgewicht um das Gelenkende
herausgerechnet) und an der anderen Scheibe allein (Gelenkkräfte über das Pendelgleichgewicht um
das Lager herausgerechnet). Das Programm behält dann den Namen mit dem Gesamtsystem, weil das der
Vorstellung entspricht, dass die Lagerkraft auf der Pendelachse wirkt.

\textbf{Auswahl.} Dann geht es schrittweise: Gibt es eine Gleichung mit genau einer noch unbekannten
Größe, ist das der nächste Schritt. Unter mehreren gewinnt die einfachste: erst die aus \code{lgs},
dann andere Momentenpunkte am selben Teilsystem, dann das Gesamtsystem, dann Zusammenfassungen und
Schnittpunkte, zuletzt Gleichungen, in denen ein Pendelgleichgewicht steckt. Bei Gleichstand gewinnt
die mit weniger Unbekannten, danach die, die zuerst aufgestellt wurde (Knoten vor Schnittpunkten,
nach Nummern). Lagergrößen (Lager und Federn) haben dabei Vorrang, aber nur innerhalb derselben
Stufe: Ein Einzelschritt für eine Lagergröße geht vor jedem anderen Einzelschritt. Gibt es keinen,
nimmt das Programm einen Einzelschritt, nach dem eine Lagergröße als Einzelschritt folgt (einen
Schritt vorausgeschaut), sonst den einfachsten. Ein Block kommt erst, wenn es gar keinen
Einzelschritt mehr gibt. Gibt es keinen, sucht das Programm zwei Gleichungen mit denselben zwei neuen
Unbekannten, die voneinander unabhängig sind (Rang~2), danach drei mit drei. Unter gleich großen
Blöcken gewinnt der mit mehr Lagergrößen, dann der aus den einfacheren Gleichungen. Dafür hat jede
Gleichung eine Stufe (\emph{Klasse}) nach der Reihenfolge oben: $0$ für eine Zeile aus \code{lgs},
höher für andere Momentenpunkte, das Gesamtsystem, Zusammenfassungen, Schnittpunkte und zuletzt
Gleichungen mit Pendelgleichgewicht. Verglichen wird die Summe der Stufen der Gleichungen im Block,
die kleinere gewinnt. Erst wenn auch das nicht geht, löst es alle übrigen Unbekannten als einen
Block aus Gleichungen von \code{lgs}. Ein solcher Block wird durch Einsetzen gelöst: Die Gleichung
mit den wenigsten Unbekannten wird nach einer davon aufgelöst und in die übrigen eingesetzt, das
wiederholt sich, bis die letzte Gleichung nur noch eine Unbekannte hat. Dann wird von hinten
zurückeingesetzt. Bei gleich vielen Unbekannten nimmt das Programm am liebsten eine Unbekannte, die
in ihrer Gleichung den Faktor $+1$ oder $-1$ hat, denn dann entsteht beim Auflösen kein Bruch.

Die Auswahl läuft mit Zahlen, denn die Faktoren vor den Unbekannten hängen nur von der Geometrie ab,
auch im symbolischen Modus. Die Texte und Ergebnisse in \code{lgs\_weg} rechnet danach das CAS
exakt aus, sie sind also dieselben wie die Lösung des ganzen LGS.

\textbf{Option LGS wie Lösungsweg.} Der Lösungsweg benutzt genau $n$ Gleichungen für $n$
Unbekannte. Ordnet man sie nach den Schritten, hat ihre Matrix Blockdreiecksform: Jeder Schritt
enthält außer schon bekannten Größen nur seine neuen Unbekannten, und deren Block ist regulär
(Einzelschritt: Faktor ungleich null, Block: voller Rang). Damit ist die ganze Matrix regulär. Das
Programm lässt die Gleichungen aus \code{lgs} stehen, die der Weg benutzt, und setzt jede
abgeleitete Gleichung auf den Platz einer Gleichung, die er nicht benutzt, am liebsten auf eine,
aus der sie mit entsteht. Danach prüft es den Rang des neuen Systems noch einmal mit Zahlen. Allgemein
gilt: Ersetzt man Zeile $r$ durch $\sum_i f_i\, r_i$, bleibt die Lösung eindeutig genau dann, wenn
$f_r \neq 0$ ist (die Determinante ändert sich um den Faktor $f_r$).

\textbf{Gegenprüfung mit der FEM.} Für jede Unbekannte bestimmt das Programm aus den FEM-Ergebnissen
den Wert mit demselben Namen und derselben Richtung wie auf Seite~2: Lagerreaktionen direkt,
Gelenkgrößen aus den Stabendkräften des angeschlossenen Stabes (ohne Knotenlast und Lagerkraft, wenn
die zu diesem Stab gehören), Stabkräfte im Fachwerk aus der Normalkraft. Das vergleicht es mit dem
Wert aus \code{lgs\_weg}, im symbolischen Modus mit denselben Ersatzwerten für die Symbole, mit denen
die FEM gerechnet hat. Weicht etwas um mehr als $10^{-4}$ des größten Wertes ab, kommt die Warnung.

\textbf{Resultierende von Streckenlasten:} $R = \int_0^L q\,dx$ und das statische Moment
$S = \int_0^L q\,x\,dx$ um den Stabanfang, die Wirkungslinie liegt bei $x_R = S/R$. Für Trapezlasten
sind das die bekannten Formeln, für Lastfunktionen integriert das Programm mit der Simpsonregel
über elf Stützstellen. Die ist für $q\,x$ bis zu einem quadratischen $q$ exakt.

%---------------------------------------------------------------------
\section{Querschnittswerte}\label{sec:mathe-qswerte}

\textbf{Vielecke} (Rechteck, Dreieck, Trapez, Polygon, dicke Linie) rechnet das Programm alle mit
derselben Formel, der Gaußschen Trapezformel. Mit den Eckpunkten $(y_i, z_i)$ gegen den
Uhrzeigersinn und $c_i = y_i z_{i+1} - y_{i+1} z_i$ ist
\[
  A = \tfrac12 \sum c_i, \qquad
  y_S = \frac{\sum (y_i + y_{i+1})\,c_i}{6A}, \qquad
  \int z^2\,dA = \tfrac1{12} \sum \left(z_i^2 + z_i z_{i+1} + z_{i+1}^2\right) c_i ,
\]
entsprechend für $z_S$, $\int y^2\,dA$ und $\int yz\,dA$. Das ist exakt, es ist der Gaußsche
Integralsatz, angewandt auf die Randkanten. Deshalb ist es auch egal, wie kompliziert das Vieleck
ist, solange es sich nicht selbst schneidet.

\textbf{Kreisausschnitt und -abschnitt} mit halbem Öffnungswinkel $\alpha$: Der Schwerpunkt liegt
auf der Winkelhalbierenden im Abstand $\frac{2R\sin\alpha}{3\alpha}$ (Ausschnitt) bzw.\
$\frac{4R\sin^3\alpha}{3(2\alpha - \sin 2\alpha)}$ (Abschnitt). Die Trägheitsmomente werden erst um
die Winkelhalbierende bestimmt (beim Abschnitt wird das Dreieck zum Mittelpunkt abgezogen) und dann
ins KOS gedreht.

\textbf{Dünnwandige Linien} sind Mittellinie mal Dicke: $A = L\,t$,
$I_y = t L^3 \sin^2\beta/12$ usw. Den Eigenanteil $L t^3/12$ lässt das Programm weg, genau wie
die Vorlesung, weil er von höherer Ordnung in $t$ ist. \textbf{Dünnwandige Kreise und Bögen}
ebenso: $A = R\,\Delta\varphi\,t$, $I_y = R^3 t \int \sin^2\varphi\,d\varphi$ um den Mittelpunkt
(Vollkreis $\pi R^3 t$), Schwerpunkt des Bogens im Abstand $R\,\frac{\sin(\Delta\varphi/2)}{\Delta\varphi/2}$
auf der Winkelhalbierenden.

\textbf{Löcher} gehen mit negativer Fläche in die Summen ein, danach setzt der Satz von Steiner alles
zusammen. \textbf{Numerisch addierte} Elemente legt das Programm auf ein Raster von $40 \times 40$
Zellen über ihr umschreibendes Rechteck und zählt nur die Zellen, die nicht schon in einem anderen
Element liegen. Deshalb die Warnung, das ist eine Näherung.

\textbf{Hauptachsen:} Mit $m = \frac{I_y + I_z}2$ und
$d = \sqrt{\left(\frac{I_y - I_z}2\right)^2 + I_{yz}^2}$ ist $I_{1,2} = m \pm d$ und
\[
  \alpha = \tfrac12 \operatorname{atan2}(2 I_{yz},\ I_y - I_z) .
\]
Das \code{atan2} (der Arkustangens, der das Vorzeichen von Zähler und Nenner kennt) sorgt dafür,
dass $\alpha$ immer zur Achse~1 mit $I_1$ gehört. Mit dem einfachen $\tan 2\alpha$ aus der
Formelsammlung müsstet ihr das von Hand prüfen, dort kann auch die Achse mit $I_2$ herauskommen.

\textbf{Widerstandsmomente} $W = I/e$ nehmen als Randabstand $e$ den größten Abstand der Elemente
vom Schwerpunkt. Bei dünnwandigen Elementen ist das die Mittellinie, $t/2$ kommt nicht dazu.

%---------------------------------------------------------------------
\section{Normalspannung und neutrale Faser}\label{sec:mathe-sigma}

Mit $\Delta = I_y I_z - I_{yz}^2$ und $y$, $z$ vom Schwerpunkt aus rechnet das Programm immer mit
\[
  \sigma_x = \frac NA + \frac{M_y I_z - M_z I_{yz}}{\Delta}\,z
                      - \frac{M_z I_y - M_y I_{yz}}{\Delta}\,y .
\]
Für $I_{yz} = 0$ bleibt das bekannte $\frac NA + \frac{M_y}{I_y} z - \frac{M_z}{I_z} y$ übrig. Die
Formel kommt direkt aus $M_y = \int \sigma z\,dA$ und $M_z = -\int \sigma y\,dA$ mit dem
linearen Ansatz $\sigma = a + b\,y + c\,z$, also zwei Gleichungen für $b$ und $c$.

Weil $\sigma_x$ linear in $y$ und $z$ ist, speichert das Programm nur den Gradienten $(c_y, c_z)$ und
den Wert $c_0$ im Ursprung des gezeigten KOS. Daraus entstehen beide Gleichungen in der Tabelle: die
KOS-Form direkt, die HAS-Form, indem der Gradient um den Hauptachsenwinkel gedreht wird. Die
neutrale Faser ist die Gerade $c_0 + c_y y + c_z z = 0$.

Die \textbf{Extremwerte} sucht das Programm nur an den Eckpunkten des Querschnitts (Bögen alle
$10^\circ$ abgetastet). Eine lineare Funktion hat ihr Maximum immer in einer Ecke der konvexen
Hülle. Bei dünnwandigen Profilen sind das die Endpunkte der Mittellinien, also wieder die
Näherung aus der Vorlesung.

%---------------------------------------------------------------------
\section{Schubfluss in dünnwandigen Profilen}\label{sec:mathe-schub}

\subsection{Offene Profile}

Aus den dünnwandigen Linien baut das Programm einen \emph{Profilgraph}: Die Linien werden an
Kreuzungen, Berührpunkten und an den Schwerachsen geteilt, Punkte, die zusammenfallen, werden ein
Knoten. \textbf{Kreise und Bögen} kommen als Kette kurzer Sehnen dazu, geteilt an ihren Enden, an
Endpunkten anderer Elemente auf dem Bogen, an Kreuzungen mit geraden Linien und an den Schwerachsen,
dazwischen gleichmäßig mit höchstens $15^\circ$ je Sehne. Die Sehnen bestimmen nur die Knoten und
den Zusammenhang (Zellen, Symmetrie). Jede Sehne kennt ihren Bogen (Mittelpunkt, Radius, Winkel
$\varphi_1$ bis $\varphi_2$), und gerechnet wird auf dem Bogen: Länge $R\,|\varphi_2 - \varphi_1|$
statt der Sehne, Stützstellen auf dem Bogen. Dann läuft es los:

\begin{itemize}
  \item An jedem \textbf{freien Ende} ist der Schubfluss null, dort beginnt ein Pfad mit
        $S_y = S_z = 0$.
  \item Entlang eines Segments werden die statischen Momente aufintegriert,
        $S_z \mathrel{+}= (y - y_S)\,t\,\Delta s$ und $S_y \mathrel{+}= (z - z_S)\,t\,\Delta s$. Weil
        der Integrand auf einer geraden Linie linear ist, ist das mit der Mittelpunktsregel exakt.
        $S$ selbst ist quadratisch in $s$. Auf einem Bogen mit Mittelpunkt $(y_C, z_C)$ ist
        $ds = R\,d\varphi$, und das Integral ist geschlossen lösbar:
        \[
          \Delta S_z = t R \left[(y_C - y_S)\,\Delta\varphi + R\,(\sin\varphi_b - \sin\varphi_a)\right],
          \quad
          \Delta S_y = t R \left[(z_C - z_S)\,\Delta\varphi - R\,(\cos\varphi_b - \cos\varphi_a)\right]
        \]
        (für den Lauf gegen den Uhrzeigersinn, sonst mit umgekehrtem Vorzeichen).
  \item An einer \textbf{Verzweigung} gilt das Knotengleichgewicht: Kommen alle Äste bis auf einen an,
        startet der letzte mit der Summe der ankommenden statischen Momente.
\end{itemize}

Aus den statischen Momenten folgt der Schubfluss mit der gekoppelten Formel
\[
  \tau\,t = -\frac{Q_y\,(I_y S_z + I_{yz} S_y) + Q_z\,(I_z S_y + I_{yz} S_z)}{I_y I_z - I_{yz}^2} ,
\]
die für $I_{yz} = 0$ in das bekannte $\tau = -\frac{Q_z S_y}{I_y t}$ bzw.\ $-\frac{Q_y S_z}{I_z t}$
zerfällt. $\tau$ ist positiv in Laufrichtung $s$.

\subsection{Eine geschlossene Zelle: woher kommt \texorpdfstring{$q_0$}{q0}?}\label{sec:mathe-zelle}

Eine Zelle hat kein freies Ende, an dem der Schubfluss null sein müsste. Das Programm schneidet sie
deshalb an einer Stelle auf, rechnet wie beim offenen Profil und bekommt einen Verlauf, der sich
vom richtigen nur um einen \textbf{konstanten Umlaufanteil} $q_0$ unterscheidet. Diesen bestimmt
es auf zwei Wegen:

\textbf{(a) Symmetrie.} Liegt eine Symmetrieachse parallel zur Querkraft, ist der richtige Verlauf
antisymmetrisch und auf der Achse null. Das Programm schneidet die Zelle deshalb möglichst auf der
Achse auf. Klappt das nicht genau, verschiebt es den Umlaufanteil so, dass ein Punkt und sein
Spiegelpunkt entgegengesetzt gleiche Werte haben. Die Symmetrie selbst prüft es, indem es jede
Linie an 13 Stellen abtastet (jede Bogensehne an 3) und schaut, ob der gespiegelte Punkt wieder auf
einer Linie gleicher Dicke liegt. Weil Bögen auch an den Schwerachsen geteilt und die Stücke
gleichmäßig in Sehnen zerlegt werden, sind die Sehnen spiegelbildlicher Bögen selbst
spiegelbildlich.

\textbf{(b) Verträglichkeit.} Ohne passende Symmetrieachse kommt die Bedingung, dass sich die Zelle
unter einer Querkraft im Schubmittelpunkt nicht verdrillt:
\[
  \oint \frac{q}{G\,t}\,ds = 0 .
\]
Der Kniff ist, diese Bedingung nicht für die gerade eingegebene Querkraft zu lösen, sondern für
die statischen Momente selbst. Dazu sortiert man die Formel für den Schubfluss nach $S_z$ und $S_y$:
\[
  q = -\frac{a\,S_z + b\,S_y}{\Delta} , \qquad
  \begin{pmatrix} a \\ b \end{pmatrix}
  = \underbrace{\begin{pmatrix} I_y & I_{yz} \\ I_{yz} & I_z \end{pmatrix}}_{J}
    \begin{pmatrix} Q_y \\ Q_z \end{pmatrix} ,
\]
mit $\Delta = I_y I_z - I_{yz}^2 = \det J$. $a$, $b$, $G$ und $\Delta$ sind entlang der Zelle
konstant, also wird die Verträglichkeit zu
\[
  a \oint \frac{S_z}{t}\,ds + b \oint \frac{S_y}{t}\,ds = 0 .
\]
\emph{Warum $J$ regulär sein muss.} Die statischen Momente hängen nur vom Querschnitt ab, nicht von
der Last. Sollen sie für \emph{jede} Querkraft den richtigen Schubfluss liefern, muss die Gleichung
oben für jedes $(Q_y, Q_z)$ gelten. Weil $J$ regulär ist, gibt es zu jedem gewünschten Paar
$(a, b)$ eine Querkraft, die es erzeugt, nämlich $(Q_y, Q_z)^T = J^{-1}(a, b)^T$. Man darf also
$(a, b) = (1, 0)$ und $(a, b) = (0, 1)$ einsetzen und bekommt die beiden Bedingungen einzeln:
\[
  \oint \frac{S_z}{t}\,ds = 0 \quad\text{und}\quad \oint \frac{S_y}{t}\,ds = 0 .
\]
Wäre $J$ singulär, lägen alle Paare $(a, b)$, die überhaupt vorkommen können, auf einer Geraden durch
den Ursprung. Dann wüsste man nur, dass eine bestimmte Kombination der beiden Umlaufintegrale
verschwindet, und könnte die beiden statischen Momente nicht einzeln festlegen.

Regulär ist $J$ bei jedem echten Querschnitt: Nach der Cauchy-Schwarz-Ungleichung aus \textbf{HM1}
ist $I_{yz}^2 = \bigl(\int y z\,dA\bigr)^2 \le \int y^2\,dA \int z^2\,dA = I_z I_y$, also
$\Delta \ge 0$. Gleichheit gibt es nur, wenn die ganze Fläche auf einer Geraden durch den Schwerpunkt
liegt, und dann versagt schon die Biegeformel, in der $\Delta$ im Nenner steht.

Streng genommen braucht man die Regularität nur für diese Richtung des Schlusses. Dass zwei
einzeln verschobene statische Momente für jede Querkraft den richtigen Schubfluss liefern, folgt
schon daraus, dass $q_0 = -\oint (q_{\text{offen}}/t)\,ds \big/ \oint ds/t$ linear in $a$ und $b$
ist, also in einen Teil zu $S_z$ und einen zu $S_y$ zerfällt. Die Regularität sagt zusätzlich, dass
es keine andere Wahl der beiden Konstanten gibt. Deshalb sind auch die angezeigten $S$-Verläufe
eindeutig.

Das Programm verschiebt also jedes statische Moment in der Zelle um seine eigene Konstante $s_0$.
Aus $\oint (S + s_0)/t\,ds = 0$ folgt
\[
  s_0 = -\,\frac{\oint S/t\;ds}{\oint ds/t}
\]
(beide Integrale im Umlaufsinn gegen den Uhrzeigersinn). Damit stimmen nicht nur $\tau$, sondern
auch die angezeigten $S$-Verläufe. Integriert wird mit der Simpsonregel, die für das quadratische
$S/t$ exakt ist. Auf einem Bogen ist $S$ eine glatte Winkelfunktion, mit 16 Teilen je Sehne
von höchstens $15^\circ$ ist Simpson dort auf etwa $10^{-10}$ genau. Beim symmetrischen Profil kommt
so übrigens genau die antisymmetrische Lösung aus (a) heraus.

\subsection{Schubmittelpunkt}

Der Schubmittelpunkt ist der Punkt, durch den die Resultierende des Schubflusses geht. Das
Programm rechnet den Schubfluss für die Einheitslasten $Q_y = 1$ und $Q_z = 1$ (bei einer Zelle
inklusive $q_0$) und dessen Moment um den Schwerpunkt,
\[
  M_q = \int \left[(y - y_S)\,t_z - (z - z_S)\,t_y\right] q\,ds ,
\]
mit der Tangentenrichtung $(t_y, t_z)$ der Linie. Weil dasselbe Moment aus der Querkraft im
Schubmittelpunkt kommen muss, $M_q = (y_M - y_S)\,Q_z - (z_M - z_S)\,Q_y$, ist
\[
  y_M = y_S + M_q(Q_z = 1), \qquad z_M = z_S - M_q(Q_y = 1) .
\]
Der Integrand ist kubisch in $s$, die Simpsonregel also wieder exakt (auf Bögen, mit der
Bogentangente, wieder auf etwa $10^{-10}$). Die Momententabelle zeigt genau diese Summe, aufgeteilt
nach Elementen. Gegen die Handrechnung: offener Halbkreis $e = 4R/\pi$, D-Profil aus Halbkreis und
Platte und Kreisrohr stimmen auf etwa $10^{-9}$.

\subsection{Massive Querschnitte}

Für $\tau = \frac{Q\,S}{I\,b}$ teilt das Programm den Querschnitt in etwa 160 waagerechte (bzw.\
senkrechte) Streifen. Die Breite $b(z)$ kommt exakt aus allen Elementen: Löcher abgezogen, Polygone
über die Schnittpunkte der Kanten, Kreise über $2\sqrt{R^2 - \Delta^2}$, Kreisausschnitt und
-abschnitt als Kreis geschnitten mit den Halbebenen der Randstrahlen bzw.\ der Sehne. Die
Streifengrenzen liegen auf allen Knickstellen von $b$ (Ecken, Rechteckkanten, Kreisscheitel,
Lochränder) und auf der Schwerachse, dazwischen gleich weit. In jedem Streifen ist $b$ damit glatt,
und $S(z)$ wird mit drei Gauß-Punkten je Streifen von oben aufsummiert (an Knickstellen fein
unterteilt, wegen der Wurzel am Kreisscheitel).

$\tau$ wertet das Programm auf den Streifengrenzen aus, also auch genau am oberen und unteren Rand.
Dort ist $S = 0$, damit $\tau = 0$. Der kleine Rest, den die Quadratur am zweiten Rand übrig lässt,
wird anteilig nach der Fläche abgezogen. An Spitzen und am Kreisscheitel geht $b \to 0$, aber $S$
schneller, der Grenzwert $\tau \to 0$ wird direkt eingesetzt. An einem Breitensprung (Flansch und
Steg, Lochrand) gibt es zwei Stützstellen, eine mit der Breite darunter und eine mit der darüber,
der Verlauf zeigt also den Sprung. Früher wurde $\tau$ in der Streifenmitte ausgewertet, aber mit
$S$ bis zur Streifenunterkante: Der oberste Wert war deshalb nie null (beim Rechteck etwa
$2{,}5\,\%$ von $\tau_{max}$), und der Rand selbst kam gar nicht vor.

Wirken $Q_y$ und $Q_z$ gleichzeitig, ist das angezeigte $\tau_{max}$
$\sqrt{\tau_{y,max}^2 + \tau_{z,max}^2}$, also eine Abschätzung nach oben, denn die beiden Maxima
liegen meist nicht an derselben Stelle.

%---------------------------------------------------------------------
\section{Torsion und Verwölbung}\label{sec:mathe-torsion}

\textbf{Torsion:} Beim offenen Profil ist $I_T = \sum \frac{\xi}{3} L\,t^3$ und
$\tau = M_T\,t/I_T$. Bei einer Zelle kommt die Fläche $A_m$ der Mittellinie aus der Gaußschen
Trapezformel und $\oint ds/t = \sum L_i/t_i$, damit $I_T = 4A_m^2/\oint ds/t$ und
$\tau = M_T/(2A_m t)$. Läuft die Zelle über einen Bogen, nimmt die Trapezformel die Sehnen, und je
Sehne mit dem Winkel $\Delta\varphi$ kommt das Kreissegment
\[
  \Delta A = \frac{R^2}{2}\left(\Delta\varphi - \sin\Delta\varphi\right)
\]
zwischen Sehne und Bogen dazu (mit dem Vorzeichen des Umlaufs). $L_i$ ist dort die Bogenlänge
$R\,\Delta\varphi$. Beides ist exakt, die Sehnen sind also nur eine Hilfe für die Topologie.

\textbf{Zelle mit offenen Ästen} (z.\,B.\ ein Kasten mit überstehenden Gurten): Die Zelle trägt
nach Bredt das ganze $M_T$, und $I_T$ ist ihr Bredt-Wert. Daraus folgt die Verdrillung
$\vartheta = M_T/(G\,I_T)$, und die offenen Äste müssen sich genauso verdrillen wie die Zelle. Ein
Ast der Dicke $t$ hat dann am Rand $\tau = G\,\vartheta\,t = M_T\,t/I_T$. Streng genommen nehmen die
Äste damit zusätzlich das Moment $G\,\vartheta\,I_{T,\text{offen}}$ auf, und alle Schubspannungen
zusammen ergeben $M_T\,(1 + I_{T,\text{offen}}/I_{T,\text{Zelle}})$ statt $M_T$. Das ist eine
Näherung: $I_{T,\text{offen}}$ wächst mit $t^3$, $I_{T,\text{Zelle}}$ mit $A_m^2$, und der
Unterschied liegt bei üblichen Profilen im Promillebereich. Bei einem Kasten $200 \times 100$ mit
$t = 10$ und zwei $50$ langen Gurten derselben Dicke sind es etwa $0{,}1\,\%$. Genauso lässt Bredt
den kleinen offenen Anteil $\frac13 L\,t^3$ der Zellwände selbst weg. Exakt im Gleichgewicht wäre
$I_T = I_{T,\text{Zelle}} + I_{T,\text{offen}}$ mit einer Aufteilung von $M_T$ im Verhältnis der
beiden Steifigkeiten.

Ein
dünnwandiges Kreisrohr ist die Zelle mit $A_m = \pi R^2$ und $\oint ds/t = 2\pi R/t$, also
$I_T = 2\pi R^3 t$. Das Programm erkennt das Rohr aus einem einzelnen Kreis gesondert. Über die
Sehnenkette käme dasselbe heraus, ein Rohr aus zwei Halbkreisen rechnet so. Das Torsionsmoment
aus platzierten Kräften ist das Moment um den \emph{Schubmittelpunkt}, denn um den dreht sich der
Querschnitt.

\textbf{Verwölbung:} Auf einem geraden Segment ist der Fahrstrahlabstand $r_\perp$ linear, damit
ist die Verwölbung nach Formelsammlung geschlossen integrierbar:
\[
  u_x(s) = u_{x,0} + \frac{q_T}{G\,t}\,s - \vartheta\left[r_1 s + (r_2 - r_1)\frac{s^2}{2L}\right],
  \qquad q_T = \frac{M_T}{2A_m}, \quad \vartheta = \frac{M_T}{G\,I_T} ,
\]
wobei der Bredt-Anteil $q_T/(Gt)$ nur in der Zelle auftritt (mit Vorzeichen je nach Umlaufsinn).
Auf einem Bogen ist $r_\perp$ nicht linear, das Integral aber ebenfalls geschlossen lösbar. Mit
dem Pol $(y_M, z_M)$ und dem Winkel $\varphi$ ab dem Sehnenanfang $\varphi_1$ ist
\[
  \int r_\perp\,ds = R\left[(y_C - y_M)(\sin\varphi - \sin\varphi_1)
    - (z_C - z_M)(\cos\varphi - \cos\varphi_1) + R\,(\varphi - \varphi_1)\right] .
\]
Ein Kreisrohr verwölbt sich damit nicht: Mit dem Pol im Mittelpunkt ist $r_\perp = R$, und der
Bredt-Anteil hebt $R\,\vartheta$ genau auf.
Das Programm läuft von einem Knoten auf einer Symmetrieachse aus (dort ist $u_x$ bei symmetrischen
Profilen null) über den Profilgraph und setzt die Werte Segment für Segment fort. Beim Schließen der
Zelle muss es wieder beim Startwert ankommen, weil $\oint r_\perp\,ds = 2A_m$ ist. Das ist die
eingebaute Probe. Zum Schluss wird die Konstante so gewählt, dass $\int u_x\,t\,ds = 0$ ist, also
keine mittlere Längsverschiebung übrig bleibt.

%---------------------------------------------------------------------
\section{Kernfläche}\label{sec:mathe-kern}

Die Kernfläche ist die Menge aller Lastpunkte, deren neutrale Faser den Querschnitt nicht
schneidet. Gerechnet wird rückwärts:

\begin{enumerate}
  \item Das Programm sammelt die Eckpunkte des Querschnitts (Bögen abgetastet) und bildet ihre
        \textbf{konvexe Hülle}. Einbuchtungen spielen keine Rolle, denn eine Gerade, die die Hülle
        berührt, berührt auch den Querschnitt nicht innen.
  \item Jede Kante der Hülle ist eine mögliche \textbf{neutrale Faser}. Mit dem nach außen
        zeigenden Einheitsvektor $\vec n$ der Kante und ihrem Abstand $d$ vom Schwerpunkt gehört
        dazu der Lastpunkt
        \[
          \vec e = -\frac{1}{A\,d}
          \begin{pmatrix} I_z & -I_{yz} \\ -I_{yz} & I_y \end{pmatrix} \vec n .
        \]
        Für ein Rechteck und die obere Kante ist $\vec n = (0,1)$, $d = h/2$, und es kommt
        $e_z = -h/6$ heraus, das bekannte Kerndrittel.
  \item Die Kernpunkte aller Kanten, der Reihe nach verbunden, sind die Kernfläche. An Bögen legt
        das Programm alle $5^\circ$ eine Tangente an, beim Vollkreis 36 Stück.
\end{enumerate}

\textbf{Woher die Matrix kommt.} Die Formel für $\vec e$ ist die Formel für $\sigma_x$ aus
Abschnitt~\ref{sec:mathe-sigma}, rückwärts gelesen. Eine Normalkraft $N$ im Punkt
$\vec e = (e_y, e_z)$, gemessen vom Schwerpunkt, erzeugt $M_y = N e_z$ und $M_z = -N e_y$. In
$\sigma_x$ eingesetzt und nach $y$ und $z$ sortiert:
\[
  \sigma_x = \frac NA + \frac N\Delta\Bigl[\bigl(I_y e_y + I_{yz} e_z\bigr)\,y
                                       + \bigl(I_{yz} e_y + I_z e_z\bigr)\,z\Bigr]
           = \frac NA\Bigl(1 + \frac A\Delta\,\vec x^T J\,\vec e\Bigr) ,
\]
mit $\vec x = (y, z)$ und der Trägheitsmatrix $J = \begin{pmatrix} I_y & I_{yz} \\ I_{yz} & I_z
\end{pmatrix}$, $\Delta = \det J$. Die neutrale Faser ist $\sigma_x = 0$, also die Gerade
\[
  \vec x^T \Bigl(\frac A\Delta\,J\,\vec e\Bigr) = -1 .
\]
Die Kante mit dem nach außen zeigenden Einheitsvektor $\vec n$ und dem Abstand $d$ vom Schwerpunkt
ist die Gerade $\vec x^T \vec n = d$, oder genauso $\vec x^T(-\vec n/d) = -1$. Eine Gerade, die
nicht durch den Schwerpunkt geht, lässt sich nur auf eine Weise in der Form $\vec x^T \vec g = -1$
schreiben. Die neutrale Faser ist also genau dann die Kante, wenn die beiden Vektoren gleich sind:
\[
  \frac A\Delta\,J\,\vec e = -\frac{\vec n}{d}
  \qquad\Longrightarrow\qquad
  \vec e = -\frac{1}{A\,d}\,\bigl(\Delta\,J^{-1}\bigr)\,\vec n ,
  \qquad
  \Delta\,J^{-1} = \begin{pmatrix} I_z & -I_{yz} \\ -I_{yz} & I_y \end{pmatrix} .
\]
Die Matrix in der Formel ist also die Inverse der Trägheitsmatrix ohne den Vorfaktor $1/\det J$
(Diagonale vertauscht, Nebendiagonale mit anderem Vorzeichen). Auch hier braucht man, dass $J$
regulär ist, siehe Abschnitt~\ref{sec:mathe-zelle}.

\textbf{Warum man die Kernpunkte geradlinig verbinden darf.} Dreht man die neutrale Faser um eine
Ecke $\vec P$ der Hülle von einer Kante zur nächsten, gehen alle diese Geraden durch $\vec P$, also
gilt $\vec P^T(\vec n/d) = 1$. Der Vektor $\vec n/d$ wandert damit auf einer Geraden, und weil
$\vec e$ linear von $\vec n/d$ abhängt, wandert auch der Lastpunkt auf einer Geraden. Jeder Ecke der
Hülle entspricht so eine gerade Kante der Kernfläche, jeder Kante der Hülle eine Ecke.

Im Hauptachsensystem fällt die Kopplung weg ($I_{\eta\zeta} = 0$). Schreibt man die Kante mit
ihren Achsenabschnitten als $\eta/a_\eta + \zeta/a_\zeta = 1$, ist $\vec n/d = (1/a_\eta,
1/a_\zeta)$, und aus der Formel wird die bekannte Form $\eta_K = -i_\zeta^2/a_\eta$,
$\zeta_K = -i_\eta^2/a_\zeta$ mit $i^2 = I/A$. Deshalb beschriftet das
Programm die umhüllenden Geraden mit ihren Abschnitten auf den Hauptachsen, gemessen vom
Schwerpunkt, und nicht mit Koordinaten im gezeichneten KOS.

%---------------------------------------------------------------------
\section{Spannungskreis: vom Schnitt zum Zustand}\label{sec:mathe-kreis}

\textbf{Lineare Gleichungen.} Mit $\sigma_m$, $R_c$ und $R_s$ als Unbekannten ist jede Angabe mit
bekanntem Winkel eine Zeile
\[
  \sigma = \sigma_m + R_c\cos 2\alpha + R_s\sin 2\alpha, \qquad
  \tau = -R_c\sin 2\alpha + R_s\cos 2\alpha .
\]
\textbf{Punkte ohne Winkel} erfüllen $(\sigma - \sigma_m)^2 + \tau^2 = R_c^2 + R_s^2$. Zieht man
diese Gleichung für zwei solche Punkte voneinander ab, fallen alle quadratischen Terme weg und es
bleibt eine lineare Gleichung für $\sigma_m$ allein:
$2(\sigma_2 - \sigma_1)\,\sigma_m = \sigma_2^2 + \tau_2^2 - \sigma_1^2 - \tau_1^2$. Bei mehreren
Punkten ohne Winkel zieht das Programm die Gleichung des ersten Punktes von jeder anderen ab. Jede
Differenz ist linear und kommt zu den linearen Zeilen. Quadratisch bleibt damit nur noch eine
Gleichung, die des ersten Punktes selbst: Die Information der anderen Punkte steckt vollständig in
den Differenzen.

\textbf{Rang.} Die linearen Zeilen löst das Programm mit Gauß-Jordan und Spaltenpivotsuche und
zählt dabei, wie viele Pivotelemente es findet (Toleranz relativ zum größten Eintrag). Bleibt in
einer Zeile ohne Pivot rechts etwas übrig, widersprechen sich die Angaben.

\begin{itemize}
  \item \textbf{Rang 3:} Der Zustand ist eindeutig. Ein Punkt ohne Winkel muss dann nur noch auf
        dem Kreis liegen (Probe).
  \item \textbf{Rang 2 und ein Punkt ohne Winkel:} Die Lösungen der linearen Zeilen bilden eine
        Gerade $\vec x = \vec x_0 + t\,\vec d$ im $(\sigma_m, R_c, R_s)$-Raum. In die quadratische
        Gleichung eingesetzt, wird daraus $a t^2 + b t + c = 0$ mit keiner, einer oder zwei Lösungen.
        Genau daher kommt \emph{zwei Zustände möglich}.
  \item \textbf{Weniger:} Es fehlen Angaben.
\end{itemize}

Schnitte ohne Winkel mit nur $\sigma$ oder nur $\tau$ geben keine Gleichung, aber eine
Ungleichung: $|\sigma - \sigma_m| \le R$ bzw.\ $|\tau| \le R$. Damit sortiert das Programm von zwei
möglichen Zuständen den aus, auf dessen Kreis der Wert nicht passt.

\textbf{Die Spannungsscheibe} ist das Vieleck, dessen Kanten alle einen Kreis mit Radius $r$
berühren. Jeder Schnitt $i$ mit dem Normaleneinheitsvektor $\vec n_i$ gibt eine Kante: die Gerade
$\vec n_i \cdot \vec x = r$, senkrecht zu $\vec n_i$ im Abstand $r$ vom Mittelpunkt. Zur Scheibe
gehört alles auf der Seite des Mittelpunktes, also die Halbebene $\vec n_i \cdot \vec x \le r$. Die
Scheibe ist die Menge der Punkte, die in allen diesen Halbebenen zugleich liegen. Liegen zwei
benachbarte Normalen $160^\circ$ oder mehr auseinander, wäre die Scheibe
offen (oder eine sehr spitze Nadel). Dann ergänzt das Programm den $x$- oder $y$-Schnitt, dessen
Normale am nächsten an der Mitte der größten Lücke liegt, und das so oft, bis die Lücke kleiner ist.
